\begin{titelseite}
\quellseite{12}
\werktitel{Der Hauptgottesdienst\\
der\\
Evangelisch-Lutherischen Kirche.}

Zur Erhaltung des liturgischen Erbtheils\\
und zur\\
Beförderung des liturgischen Studiums in der amerikanisch-lutherischen Kirche erläutert und mit altkirchlichen Singweisen versehen

von\\
Friedrich Lochner, Pastor.

\begin{quote}
Lasset alles ehrlich und ordentlich zugehen. — Was wahrhaftig, was ehrbar, was gerecht, was keusch, was lieblich, was wohl lautet, ist etwa eine Tugend, ist etwa ein Lob, dem denket nach.\\
1 Cor. 14, 40. Phil. 4, 8.
\end{quote}

St. Louis, Mo.\\
CONCORDIA PUBLISHING HOUSE.\\
1895.
\end{titelseite}

\quellseite{14}
\chapter{Vorwort}


Seit meiner im Jahre 1876 erfolgten Berufung an die Trinitatisgemeinde in Springfield, Ills., bis zu meinem im Jahre 1887 erfolgten Abzug habe ich an der im Jahre 1875 von St. Louis nach Springfield verpflanzten praktisch-theologischen Lehranstalt der Missourisynode wie im Gemeindegesang, so auch in der Liturgie unterrichtet. Dadurch war ich veranlaßt, meine früheren, wie hymnologischen, so auch liturgischen Studien wieder aufzunehmen, zu welch letzteren mich vor mehr denn 40 Jahren während meiner Vorbereitung für den Dienst der amerikanisch-lutherischen Kirche der selige Pfarrer Löhe in Franken und sein noch lebender, damals bei einem benachbarten Gerichte angestellter Freund Hommel, nachmaliger Herausgeber einer musikalischen „Liturgie lutherischer Gemeindegottesdienste“, anregten und anleiteten. So entstanden denn in Folge noch eingehenderer Studien und behufs Unterrichts insbesondere in der Form und Weise des Hauptgottesdienstes und mit Zugrundelegung des 24. Artikels der Augsburgischen Confession nach und nach allerlei Aufzeichnungen, die auf Wunsch früherer und späterer Schüler möglichst eingehend und vervollständigt nun an die Oeffentlichkeit treten.

Unter den Vorarbeiten zu einem mir im Jahre 1861 für die damalige nördliche Distriktsynode aufgetragenen und auf deren Beschluß dann im Juni-, Juli- und Augustheft des achten Bandes von „Lehre und Wehre“ abgedruckten „Referat über die rechte Mitte der lutherischen Liturgie“ drängte sich mir schon damals der Wunsch auf, es möchte zum Frommen der auf Grund der reinen Lutherlehre sich erbauenden und unter den Segnungen der völligen Gewissensfreiheit sich neu und frisch gestaltenden Einwandererkirche dieses Landes einmal eine tüchtige Hand die durch Luther gereinigte und rechte „Messe“ des 24. Artikels der Augustana darstellen und erklären, wie sie dessen liturgische Schriften und die nicht geringe Anzahl der älteren Kirchenordnungen samt
\quellseite{15}
den liturgisch-musikalischen Schätzen der reformatorischen und nachreformatorischen Zeit aufzeigen und wie sie selbst in der Zeit des vulgären Rationalismus wenigstens den Hauptformen nach da und dort im alten Vaterlande sich noch behauptete, trotzdem daß dieser auch auf dem liturgischen Gebiet nach Vandalenart hauste. Wenn ich nun, da inzwischen eine tüchtige Hand nicht zum Werk gegriffen hat, es wage, mit meinen Aufzeichnungen an die Oeffentlichkeit zu treten, so bewegt mich dazu nicht nur das Verlangen, meinen früheren Schülern in dem mir zugewiesenen praktischen Fach des theologischen Unterrichts noch entsprechender zu dienen, als es durch einzelne kurze Dictate geschehen konnte, sondern auch die Befürchtung einer zukünftigen Gefahr für die nunmehr wieder in Gang gekommene alte Weise des Gottesdienstes, welcher Gefahr vorzubeugen ich eben durch dieses liturgische Werk meinerseits einen Beitrag thun möchte.

Als ich vor bald einem halben Jahrhundert mit andern Brüdern zum Dienst der Kirche in dies Land kam, kannte und übte man die schöne alte Weise des lutherischen Hauptgottesdienstes nur in den etlichen Gemeinden der zu Ende der dreißiger Jahre um des Glaubens willen eingewanderten sächsischen und preußischen Lutheraner, und bald hernach auch in den entstehenden fränkischen Colonien Michigans. Obwohl nun damals gerade für die hiesige lutherische Kirche Löhe seine Agende und hernach Hommel seine Altarweisen und übrigen liturgischen Gesänge herausgegeben hatte, so mußten und wollten wir, die wir die zerstreuten, meist aus rationalistischen oder unionistischen Gemeinden herübergekommenen Glaubensgenossen zu Gemeinden erst sammelten, zunächst gehörig Grund legen mit der reinen Lehre und uns daher mehr oder weniger mit der einfachsten Weise des Gottesdienstes behelfen. Verschrie man uns doch damals schon allein um unserer Lehre willen seitens der Secten, der Unirten und der fälschlich lutherisch sich nennenden General-Synode als Krypto-Papisten! Als nun aber hinreichend Grund in der Lehre gelegt war, kam es auch von selbst mehr und mehr in den Gemeinden zu liturgischer Gestaltung. Ueberall fand in unsern Gemeinden die im Jahre 1856 erschienene „aus den alten Sächsischen Kirchenagenden“ zusammengestellte „Kirchen-Agende“ der Synode von Missouri etc. Aufnahme, ja auch Eingang außerhalb ihres Gebietes, und in den schon etwas mehr geordneten Gemeinden konnte man den Hauptgottesdienst in der vollen Form sehen, wie sie diese Agende vorschreibt.

\quellseite{16}
Selbst da, wo man sich zwar der Agende bediente, aber noch vor dem liturgischen Singen sich scheute, als wäre es etwas „Katholisches“, schwand nach und nach das Vorurtheil auch gegen dieses.

Dennoch drängt sich nicht nur mir, sondern auch Andern auf Grund verschiedener Wahrnehmungen die Befürchtung auf, es möchte das, was als Frucht der reinen Lehre nun auch in der Liturgie uns zu Theil geworden ist, in der Folgezeit mit der kahlen puritanischen Weise, wie wir sie einst als die herrschende vorfanden, allmählich wieder vertauscht werden. Wir leben ja im Lande des wuchernden reformirten Sectenthums, also in einem im Grunde reformirten Lande, und bedarf es bei unserem heranwachsenden Geschlecht je länger je mehr des Belehrens, daß es sich von dem englisch-puritanischen Wesen in kirchlicher Beziehung nicht beeinflussen lasse. Dazu leben wir in einer Zeit, in der man je länger je weniger Geduld und Ausdauer für geistliche Uebungen hat, daher auch einen schon etwas abgekürzten liturgischen Gottesdienst da und dort bereits zu lang findet und den Prediger drängt, noch mehr zu kürzen. Dazu fehlt es, wie es die Verhältnisse mit sich bringen, vielfach noch so sehr an einem näheren liturgischen Verständniß, ohne welches man nur um so weniger das erst wieder erlangte liturgische Erbe festzuhalten sucht. Und ach, wenn schon in der Stellung zur gottentfremdeten Welt die theure Lehre von der christlichen Freiheit nicht allewege mehr recht angewendet wird, sondern der antinomistische Geist es uns zu guter letzt noch anthun möchte, so wird man in der Folgezeit nur um so mehr die alten schönen Ceremonien unbedenklich fahren zu lassen geneigt sein, da dieselben als Mitteldinge ganz und gar in das Gebiet der christlichen Freiheit gehören.

Zum Studium der Liturgie anzuregen und anzuleiten, insbesondere aber die Schönheit und Bedeutsamkeit des lutherischen Hauptgottesdienstes und den ihn durchdringenden Geist der Einfalt und des Bekenntnisses zu zeigen, in das Verständniß desselben einzuführen, und so beizutragen, daß das bereits empfangene liturgische Erbe uns auch erhalten bleibe und wo möglich noch mehr Gemeingut werde, ist daher die Absicht meiner Darlegung der Entstehung und Bedeutung des lutherischen Hauptgottesdienstes und seiner einzelnen Theile und meiner sie begleitenden Mittheilung aus dem reichen Schatze des alten liturgischen Gesanges unserer auch hier so reich begabten Kirche, sowie der
\quellseite{17}
gelegentlichen praktischen Winke, bei denen ich natürlich nur unser hiesiges Bedürfniß im Auge habe.

Damit ich jedoch in dem, was ich beabsichtige, richtig verstanden werde, erlaube ich mir noch einige Bemerkungen.

Bekanntlich ist in neuerer Zeit, besonders in der alten Heimath, das Interesse für altlutherische und altkirchliche Liturgie und darum auch für die alte Kirchenmusik überaus rege geworden. Mit einem großen Eifer wird das liturgische Studium betrieben und Schatz um Schatz gehoben und vorgezeigt. Ueberall ist man bestrebt, den Gottesdiensten wieder ihre alte schöne Gliederung zurückzugeben, kirchliche Sitten und Gebräuche in Bezug auf Altarbekleidung, Betglockenläuten etc. wieder zu pflegen u. dgl. Und so stark ist diese liturgische Strömung unter denen, die zum echten Lutherthum zurückzukehren meinen, daß sie sich auch hier zu Lande geltend macht und selbst die schon seit 1845 offenkundig vom Lutherthum abgefallene, aber immer noch lutherisch sein wollende Generalsynode sich ihr nicht entziehen kann. Ob dieser Bestrebungen könnte man sich ja nur freuen und sie als Kundgebung eines erwachten frischen und fröhlichen Lebens begrüßen, wenn man sie nur in zweiter Linie fände, Gottes Wort und Luthers Lehr aber in erster Linie. Allein, das ist, leider! nicht der Fall. Man greift nur nach der Schale und läßt den Kern liegen, ja nimmt die Schale für den Kern und meint, mit ihr der lutherischen Kirche der Gegenwart Luthers Sinn und Geist wieder einzuhauchen. Während der 7. Artikel der Augustana die Einträchtigkeit in der reinen Lehre und im rechten Brauch der Sacramente mit seinem „Es ist genug“ so gewaltig betont und die Gleichförmigkeit in den doch nur von Menschen eingesetzten Ceremonien für zur Einigkeit nicht wesentlich nothwendig bezeichnet, legt man kein Gewicht auf die reine Lehre, läßt die Uneinigkeit in der Lehre so groß sein, als sie will, und uniformirt durch Verfassung und altlutherische Liturgie. Während die Apologie im Artikel von der Messe bezeugt: „Es ist kein Ding, das die Leute mehr bei der Kirche behält, denn die gute Predigt“, meint man, durch Wiederherstellung der äußerlichen Schönheit und Mannigfaltigkeit der alten lutherischen Weise des Gottesdienstes die Leute wieder zur Kirche zu bringen und bei derselben zu behalten, statt vor allem der im Sinne der Apologie „guten“ Predigt sich zu befleißigen und vorerst die Wahrheit dieses symbolischen Ausspruchs zu erfahren. Und während man seinen zunehmenden Abfall von Gottes Wort
\quellseite{18}
und Luthers Lehre erkennen und über demselben in Sack und in der Asche Buße thun sollte, dankt man Gott über den in der Kirche der Gegenwart sich zeigenden „Sinn“ für liturgische Einrichtungen als für ein ihr verliehenes Charisma. Gott bewahre unsere americanisch-lutherische Kirche vor solcher Werthschätzung der Liturgie!

„Zwei thun dasselbe und ist doch nicht dasselbe.“ Wo „einträchtiglich nach reinem Verstande das Evangelium gepredigt und die Sacramente dem göttlichen Wort gemäß gereicht werden“ und wo man daher die Kirche nur mit dem Wort bauen und regieren will, da hat es eine weit andere Meinung mit den liturgischen Bestrebungen. Und der Art wollen auch die meinigen und der Dienst sein, den ich mit dieser Arbeit zu leisten suche. Weil es mir bei derselben auch nur um Erhaltung dessen zu thun ist, was von der alten Liturgie lediglich in Folge der Rückkehr zur alten Lutherlehre uns von selbst in den Schooß gefallen ist, so habe ich auch jenen Bestrebungen nicht das Wort reden mögen, bei denen man durch allerlei „Zusätze“ aus dem kirchlichen Alterthum und selbst auch aus den römischen Ritualien die alte lutherische Liturgie noch weiter auszugestalten und so über sie hinaus zu gehen sucht. Wenn ich daher auch bei meinem Unterricht in dem genannten theologischen Seminar meinen Schülern die Schönheit, den Reichthum und die Vorzüge der altlutherischen Liturgie zeigte, so machte ich sie doch immer wieder darauf aufmerksam, wie alles am reinen Wort und Sacrament liegt und alle Ceremonien nur dazu dienen sollen und nur dazu dienen dürfen, daß diese recht im Schwange gehen, wie aber auch diese Liturgie allewege diesen Dienst thue und es daher zu beklagen wäre, wenn sie da, wo sie eingeführt ist, in der Folgezeit durch die englisch-puritanische Kahlheit wieder verdrängt würde. Vor allem aber unterließ ich nicht, dieselben vor Ueberschätzung der Liturgie und vor unzeitigem Einführen derselben ernstlich zu warnen. Das konnten sie gerade an ihrem Orte auch in praxi sehen. Nach 26jährigem Dienst an einer unserer alten, nicht nur in der Erkenntniß geförderten, sondern auch liturgisch reich ausgestatteten Gemeinde war ich an eine Gemeinde gestellt, die für die alte Liturgie noch nicht völlig reif war, so daß auch das Singen am Altar zur Zeit noch unterbleiben mußte. Konnte ich nun auch meinen Schülern im Gottesdienst der Gemeinde nicht alles das vor Augen führen, was ich ihnen im liturgischen Unterricht gab, so konnte ich ihnen doch an meinem Exempel ein wenig zeigen, daß und wie ein lutherischer Prediger
\quellseite{19}
allen Werth auf eine „gute Predigt“ zu legen, sucht und der Zeit erharren kann, wo es auch in liturgischer Beziehung zu einer rechten Gestalt kommt, und daß man daher bei aller Liebe zur alten Liturgie und Begeisterung für dieselbe auch in dieser Beziehung durch Gottes Gnade mit einem Paulus sagen kann: „Ich kann beides: übrig haben und Mangel leiden — durch den, der mich mächtig macht, Christus.“ —

Ich erlaube mir schließlich noch, die Schriften anzugeben, aus denen ich mehr oder weniger schöpfte oder bei denen ich mir allerlei Rath holte.

Vor allem schöpfte ich aus Luthers liturgischen Schriften und ging namentlich immer wieder auf seine „Formula Missae“ oder seine „Deutsche Messe und Ordnung des Gottesdienstes“, 1523, zurück. An diese reihten sich die alten Kirchenordnungen — so viel mir nämlich zu Gebote standen. Es waren das: Die Brandenburg-Nürnberger K.-O., 1533; das Agendbüchlein von Veit Dietrich; die Agende Herzog Augusts von Braunschweig-Lüneburg, 1657; die Telser Agende, 1664; die Schwarzburger Agende, 1675; die Agende Herzog Heinrichs, Leipziger Ausgabe, 1681, und die Magdeburger K.-O., 1685. Namentlich aber suchte ich hierbei den aus drei umfangreichen Bänden bestehenden „Schatz des liturgischen Chor- und Gemeindegesangs nebst den Altarweisen“ von Prof. Dr. L. Schöberlein (Göttingen 1865) zu verwerten, indem dieses ausgezeichnete Werk eine gefüllte Rüstkammer der alten musikalischen Liturgie ist und am gründlichsten in das Verständniß des Hauptgottesdienstes und seiner einzelnen Theile einführt. Außerdem holte ich mir Rath in folgenden Schriften: 1. Dr. Freitag: Lutherus musico-liturgus, das ist, Luthers „deutsche Messe und Ordnung des Gottesdienstes“, 1871; 2. J. Lyra: „Die liturgischen Altarweisen des lutherischen Hauptgottesdienstes“ und dessen „Andreas Ornithoparcus von den Kirchenaccenten“, 1877; 3. Löhe: Sammlung liturgischer Formulare, und dessen Agende von 1844, 1853 und 1884; 4. F. Hommel: Liturgie lutherischer Gemeindegottesdienste, 1851; 5. Dr. F. Layriz: Kern des deutschen Kirchengesangs, IV. Liturgische Weisen, 1855; 6. Dr. L. Kraußold: Historisch-musikalisches Handbuch für den Kirchen- und Choralgesang, 1855; 7. Vielsecker: Die Lehre vom römischen Choralgesang, 1842; 8. Seminar-Director Joh. Zahn: Handbüchlein für evangelische Cantoren und Organisten; 9. Rudelbach: Sacramentworte;
\quellseite{20}
10. D. Kade: Der Luthercoder; 11. J. L. König: Die Hauptliturgien der alten Kirche in wortgetreuer Uebersetzung; 12. Naumann: Illustrirte Musikgeschichte etc., und noch etlichen Andern.

Noch erlaube ich mir, zu bemerken, daß, als die Arbeit bereits vollendet war, mir Dr. Herzogs vortrefflicher revidirter und ergänzter „Musikalischer Anhang zur Agende für die evangelisch-lutherische Kirche in Bayern. 1883“ zukam und es mir nicht geringe Freude gewährte, mich mit dessen Anschauungen allewege einig zu wissen. Seinem Werke entnahm ich noch den Satz für die schon früher durch Hommel bekannt gewordene Melodie des Nicänums.

Da im Lande des Puritanerthums und der Einwanderung, in welchem ja auch in Betreff des Kirchengebäudes zum Theil noch primitive Zustände vorhanden sind, unsere hier studirenden Prediger ja wenig Normales, wohl aber viel Abnormes und Unschönes zur Zeit noch vor Augen haben und es daher um so mehr noch vielfach um Verständniß dessen, was betreffs Kanzel, Altar und Taufstein kirchliche Form und Weise ist, fehlt, so hielt ich es für angezeigt, in einem besonderen Anhang über die Stätten für die Liturgie mich auszusprechen, dabei jedoch lediglich unser hiesiges Bedürfniß berücksichtigend. Hierbei benützte ich zum Theil ein in diesem Jahre erst erschienenes instructives Schriftchen, betitelt: „Das evangelische Kirchengebäude. Ein Rathgeber für Geistliche und Freunde kirchlicher Kunst, herausgegeben in Verbindung mit Baurath Dr. Mothes in Leipzig und Architekt Prüfer in Berlin von Victor Schulze, Prof. der Theologie. Leipzig. Georg Böhme. 1886.“

Gott erhalte uns bei seinem reinen Wort und Sacrament und lasse uns also in seinem Worte leben, daß auch in Bezug auf den öffentlichen Gottesdienst alles, was sich an Ceremonien und Weisen durch Menschenhand um die Gnadenmittel rankt, denselben diene und dabei nicht als ein gemachter, sondern als ein natürlicher und wirklicher Schmuck beim gemeinsamen Brauch der Gnadenmittel und im gemeinsamen Bekenntniß und Lob des Allerhöchsten erscheine.

\begin{unterschrift}
F. L.
\end{unterschrift}

\tableofcontents
\quellseite{24}
\mainmatter

\part[Erster Theil]{Erster Theil Die Entstehung und Gestalt des lutherischen Hauptgottesdienstes}

\quellseite{25}
Man leget den Unsern mit Unrecht auf, daß sie die Messe sollen abgethan haben, denn das ist öffentlich, daß die Messe, ohne Ruhm zu reden, bei uns mit größerer Andacht und Ernst gehalten wird, denn bei den Widersachern. So werden auch die Leute mit höchstem Fleiß zum öftermal unterrichtet vom heiligen Sacrament, wozu es eingesetzt und wie es zu gebrauchen sei, als nämlich die erschrockenen Gewissen damit zu trösten, dadurch das Volk zur Communion oder Meß gezogen wird. Dabei geschieht auch Unterricht wider andere unrechte Lehre vom Sacrament. So ist auch in den öffentlichen Ceremonien der Messe keine merkliche Aenderung geschehen, denn daß an etlichen Orten deutsche Gesänge (das Volk damit zu lehren und zu üben) neben lateinischem Gesang gesungen werden, sintemal alle Ceremonien fürnehmlich dazu dienen sollen, daß das Volk daran lerne, was ihm zu wissen von Christo noth ist.

Augsp. Conf. Art. 24.

\quellseite{26}
\chapter[§ 1. Der neutestamentliche Gottesdienst]{§ 1. Der neutestamentliche Gottesdienst. Sein Wesen und seine Form}


Der äußerliche gemeinschaftliche Gottesdienst der Christen beruht nicht wie der alttestamentliche Gottesdienst auf einem statuarischen (ceremonialgesetzlichen) Gebot des HEerrn und eben darum auch nicht auf irgend einem Gebote seiner heiligen Apostel, sondern er ist das Product, der Ausdruck, die nothwendige Offenbarung, Uebung und Bethätigung des inneren Lebens der Gemeinde der Heiligen, die da zwar unsichtbar ist nach ihrem Wesen, deren Dasein aber erkennbar wird am Wort und Sacrament, dadurch es geschaffen ist. Die innerliche Gemeinschaft der Heiligen besteht in dem Einen lebendigen Glauben aller einzelnen Glieder, die äußerliche Gemeinschaft in der Versammlung derselben zu gemeinsamem Gebrauch des Wortes und Sacramentes, wie zur Anbetung und Anrufung, zum Danken und zum Loben. Eph. 4, 3. Act. 2, 42. Hebr. 10, 25.

Da aber im gemeinsamen Gebrauch des Worts und der Sacramente nebst der Uebung im Gebet und Lob der öffentliche Gottesdienst der Christen besteht, so charakterisirt sich derselbe als eine öffentliche Bethätigung der darreichenden und segnenden göttlichen und der empfangenden und erwidernden menschlichen Liebe. Gott dient dem Menschen, indem er sich immer wieder mit seinem Heil ihm darzeigt und segnend zu ihm kommt im Wort und Sacrament, und der Mensch dient Gott, indem er ihn durch gläubige Annahme des dargereichten Heiles im Wort und Sacrament ehrt und in Gebet, Dank und Lob, wie im Opfer seiner selbst und im Gebrauch der ihm verliehenen Gaben des Geistes und des Leibes zu Gottes Ehre und zu des Nächsten Nutz die göttliche Liebe erwidert. 2 Mos. 20, 24. Ps. 100. Matth. 18, 20. Röm. 12, 1. ff. Hebr. 13, 15. 16., vgl. 1 Cor. 16, 2. „Christo dienen und Gott dienen“, sagt Luther, „heißt vornehmlich bei St. Paulo, ein Amt führen, das Christus ihm befohlen hat, nämlich das Predigen. Es ist ein Dienst, der von Christo, nicht zu Christo gehet, und der nicht von uns, sondern zu uns kommet.“ (Erl. Ausg. 7, 79.) Und zu den Worten Christi: „Solches thut zu meinem Gedächtnis“ schreibt er u. a.: „Ferner sein gedenken, das ist (wie gesagt) predigen, preisen, loben, zuhören und danken für die Gnade
\quellseite{27}
in Christo erzeigt. Thust du das, siehe, so bekennest du mit Herzen und Munde, mit Ohren und Augen, mit Leib und Seele, daß du Gott nichts gegeben habest noch mögest, sondern alles und alles von ihm habest und nehmest, sonderlich das ewige Leben und unendliche Gerechtigkeit in Christo. Wo aber das geschieht, so hast du ihn dir zum rechten Gott gemacht und mit solchem Bekenntniß seine göttliche Ehre erhalten. Denn das heißt ein rechter Gott, der da gibt und nicht nimmt, der da hilft und nicht ihm helfen läßt, der da lehret und regieret und sich nicht lehren und regieren lässet. Summa, der alles thut und gibt und er niemands bedarf und thut solches alles umsonst aus lauter Gnaden, ohn Verdienst der Unwürdigen und Unverdienten, ja den Verdammten und Verlornen. Solch Gedächtniß, Bekenntniß und Ehre will er haben. Siehe, dieser Gottesdienst gehet daher wohl ohne alle Pracht und füllet die Augen nicht nach dem Fleisch; er füllet aber das Herz, welches doch sonst weder Himmel und Erden füllen mag. Wenn aber das Herz voll ist, muß auch alsdann beide, Augen und Ohren, Mund und Nasen, Leib und Seele und alle Glieder voll sein. Denn wie sich das Herz hält, darnach halten und stellen sich alle Glieder und ist alles und alles eitel Zungen voll Lobens und Dankens gegen Gott.“ (Ibid. 23, 174 ff.) Eben darum sagt im Artikel von der Messe die Apologie: „Sacramentum ist eine Ceremonia oder äußerlich Zeichen, oder ein Werk, dadurch uns Gott gibt dasjenige, so die göttliche Verheißung, welche derselbigen Ceremonien angeheftet ist, anbietet. Als, die Taufe ist eine Ceremonie und ein Werk, nicht, das wir Gott geben oder anbieten, sondern in welchem uns Gott tauft oder der Diener an Gottes Statt. Da beut uns Gott an und gibt uns Vergebung der Sünde nach seiner Verheißung: Wer da glaubet und getauft wird, der soll selig werden. Wiederum, Sacrificium oder Opfer ist eine Ceremonia oder ein Werk, das wir Gott geben, damit wir ihn ehren.“ (Art. 12. Von der Messe.) Im Nachfolgenden führt die Apologie aus, daß der Christen Opfer nicht Sühnopfer, sondern eitel Dank- und Lobopfer für das vollgültige Versöhnopfer Christi seien.

Ist nun der öffentliche Gottesdienst der Christen ein naturgemäßer und nöthiger Ausdruck des Glaubenslebens der christlichen Gemeinschaft, so muß sich auch der Charakter desselben in Form und Weise abspiegeln. Mit letzterem aber hat es die Liturgie zu thun, sofern sie aufzeigt, was bei der Spendung und Empfangung der Gnadenmittel einerseits und bei der Darbringung des Gebets-, Lob- und Dankopfers anderseits geredet und gehandelt wird.

Anmerkung. Nach seiner wörtlichen und kirchlichen Bedeutung bezeichnet Liturgie die Verrichtung im Dienst der christlichen Gemeinschaft. In dem
\quellseite{28}
oben angeführten Artikel von der Messe definirt daher die Apologie: „So heißet Liturgia griechisch eigentlich ein Amt, darin man der Gemeinde dient. Das schicket sich wohl auf unsere Lehre, daß der Priester da als ein gemeiner Diener denjenigen, so communiciren wollen, dienet und das heilige Sacrament reichet.“ In den ersten Jahrhunderten hat die Kirche ihre Liturgie nicht schriftlich fixirt. Dies sagt deutlich Basilius der Große, † 379: „Die Worte der Anrufung bei der Weihung des Brodes und der Danksagung und des Kelches der Segnung, welcher der Heiligen hat schriftlich sie uns hinterlassen? Denn nicht begnügen wir uns ja mit demjenigen, was der Apostel oder das Evangelium erwähnt, sondern wir sagen sowohl vorher als nachher Anderes, als große Kraft habend, in Beziehung auf das Geheimniß, indem wir aus der nicht geschriebenen Lehre es entnommen haben.“ (König, Die Hauptliturgien der alten Kirche.) Durch die später nöthig gewordene Fixirung entstanden die liturgischen Bücher, die Ritualien, die Agenden etc.

\chapter[§ 2. Der Communiongottesdienst]{§ 2. Der Communiongottesdienst der Hauptgottesdienst unter den mancherlei Gottesdiensten}


Es ist ein Reichthum und eine liebliche Mannigfaltigkeit in den Formen und Weisen des öffentlichen Gottesdienstes der rechtgläubigen Kirche, in denen das Glaubens- und Gemeinschaftsleben seinen freien Ausdruck gefunden hat. Sie hat einen Communion-Gottesdienst, da auf die Predigt des Worts die Handlung des Abendmahls folgt, zum Wort also das Siegel und Pfand kommt; einen Predigt-Gottesdienst, in welchem namentlich die Schrift ausgelegt und angewendet oder für specielle Fälle eine Predigt (Casualpredigt) gehalten wird; einen Katechismus-Gottesdienst zur Gründung und Förderung in der Erkenntniß der Hauptstücke der christlichen Lehre (Katechismuspredigt und Katechizmusexamen), und einen Bet- und Lob-Gottesdienst in ihren Metten und Vespern, den täglichen kurzen Morgen- und Abendgottesdiensten, in welchen zwar nach Luthers wohlbegründeter Forderung Gottes Wort nicht nur gelesen, sondern auch aufs Kürzeste ausgelegt wird,\footnote{In der „Ordnung des Gottesdienstes in der Gemeinde 1523“ schreibt nämlich Luther: „Nun diese Mißbrauch abthun, ist aufs Erste zu wissen, daß die christliche Gemeine nimmer soll zusammenkommen, es werde denn daselbst Gottes Wort gepredigt und gebetet, es sei † auch aufs Kürzeste, wie Ps. 102 (23.): Wenn die Könige und das Volk zusammenkommen, Gott zu dienen, sollen sie Gottes Namen und Lob verkündigen. Und Paulus 1 Cor. 14 (31.) spricht: daß in der Gemeinde soll geweissagt, gelehrt und ermahnt werden. Darum wo nicht Gottes Wort gepredigt wird, ist's besser, daß man weder singe, noch lese, noch zusammenkomme. „Also ist's aber zugegangen unter den Christen zur Zeit der Apostel und sollt auch noch so zugehen, daß man täglich des Morgens eine Stunde früh um vier oder fünf zusammenkäme und daselbst lesen ließe, es seien Schüler oder Priester oder wer es sei, gleichwie man jetzt noch die Lection in der Metten liest. Das sollen thun einer oder zween, oder einer um den andern; wie das am besten gefällt. Darnach soll der Prediger, oder welchem es befohlen wird, herfürtreten und dieselbe Lection ein Stück auslegen, daß es die Andern alle verstehen, lernen und ermahnt werden. Das erste Werk heißt St. Paulus 1 Cor. 14 (26.) mit Zungen reden; das ander auslegen oder weissagen und mit dem Sinn und Verstand reden. Und wo dies nicht geschieht, so ist die Gemeine der Lection nichts gebessert, wie bisher in Klöstern und Stiften geschehen, da sie nur die Wände haben angebetet.“ (Erl. Ausg. 22, 154.) Damit nun in den Metten und Vespern auch Gottes Wort ausgelegt werde, aber „aufs Kürzeste“, so entstanden Veit Dietrichs Summarien. Vgl. die Vorreden Dietrichs und Vierlings in dem Altenburger Neuen Testament, amerikanische Ausgabe. In des Letzteren Vorrede findet sich auch die damalige Metten- und Vesperordnung.}  in denen aber doch Beten und Loben das herrschende Element bildet.

\quellseite{29}
Auf Grund von Act. 2, 42. und 1 Cor. 11. und nach dem Vorbild der alten Kirche betrachtet jedoch die lutherische Kirche den Communion-Gottesdienst als den herrlichsten und wichtigsten aller öffentlichen Gottesdienste und hat auch denselben am reichsten und sinnigsten liturgisch ausgestattet. Sie scheidet daher zwischen Haupt- und Nebengottesdienst. Nicht durch die Bedeutung des Sonn- oder Festtags, noch durch die Jahreszeit, noch durch liturgischen Reichthum wird ein Gottesdienst zum Hauptgottesdienst, sondern, wie es das schriftgemäße Verhältniß von Wort und Sacrament gibt, dadurch, daß auf die Verkündigung des Wortes des Evangeliums die Handlung des Sacraments des Leibes und Blutes Christi unmittelbar folgt und so diese als das Siegel des Wortes das Ziel und den Schlußstein des Gottesdienstes bildet. Alle übrigen Gottesdienste, in denen die Handlung des Sacraments von vornherein nicht beabsichtigt ist, werden zu Nebengottesdiensten, wie reich liturgisch manche derselben die Vorzeit auch gestaltet hat. Im Hinblick auf die Zusammengehörigkeit von Wort und Sacrament, wie auf das christliche Alterthum, da noch zu Augustins Zeiten (400 n. Chr.) das Abendmahl wenigstens allsonntäglich von der ganzen Gemeinde empfangen wurde, und in der Voraussetzung vorhandenen Verlangens, wenigstens in volkreicheren Gemeinden, war in der Reformationszeit und noch lange nach derselben der auf die Frühmette folgende
\quellseite{30}
Vormittagsgottesdienst regelmäßig Communiongottesdienst, also Hauptgottesdienst. Gegenüber den papistischen Winkelmessen, in welchen der fungirende Priester allein das Sacrament nimmt, hebt die Apologie Art. 8 hervor: „Bei uns aber braucht das Volk des heiligen Sacraments willig, ungedrungen alle Sonntage“; desgl. Art. 12: „Alle Sonntage und Feste werden in unsern Kirchen Messen gehalten, dabei das Sacrament gereicht wird denjenigen, die es begehren, doch also, daß sie erst verhört und absolviret werden.“ So kam es denn, daß, wo wegen Mangels an Communicanten das Abendmahl in einem solchen Gottesdienst nicht gehalten werden konnte, man ihm doch in allem Uebrigen die Form des Hauptgottesdienstes zu belassen suchte. An manchen Orten jedoch wurde nach der Predigt eine vorgeschriebene Vermahnung gelesen (Pommersche K.-O. 1563, Liegnitzer K.-O. 1594) oder eine mehr freie Vermahnung gehalten (Wittenb. K.-O. 1559, 1565; Mecklenb. K.-O. 1540, 1552), in welcher ohne gesetzliches Drängen der Mangel an Communicanten beklagt und zum oftmaligen Genuß des hochwürdigen Abendmahls vermahnt und gereizt wurde.\footnote{Anmerkung 1. Bekanntlich findet sich die älteste Schilderung des Communiongottesdienstes in der nachapostolischen Zeit von Justin, dem Märtyrer, † 166. Die Stelle, wie sie Dr. Uhlhorn in seiner Schrift: „Der Kampf des Christenthums mit dem Heidenthum“ in deutscher Uebersetzung gibt, lautet: „Am Sonntage geschieht eine Versammlung aller, die in den Städten oder auf dem Lande wohnen, und es werden darin die Denkschriften der Apostel oder die Bücher der Propheten vorgelesen, so lange wir Zeit dazu haben. Darnach, wenn der Vorleser geendigt hat, gibt der Vorsteher in einer Rede Erinnerung und Mahnung, jenen herrlichen Vorbildern nachzueifern. Alsdann stehen wir Alle mit einander auf und senden unsere Gebete empor. Und nachdem wir unser Gebet gethan haben, bringt man Brod und Wein und Wasser herbei und der Vorsteher verrichtet Gebet und Danksagungen, so viel er vermag. Die Gemeinde antwortet mit ihrem Amen und es geschieht die Austheilung der geweihten Dinge, welche jeder Anwesende empfängt, während sie den Abwesenden durch die Diakonen hingetragen werden. Die Wohlhabenden aber und die willig dazu sind, geben ein jeglicher nach seinem Gefallen, und die gesammelten Gaben werden vor dem Vorsteher niedergelegt, welcher damit den Wittwen und Waisen zu Hülfe kommt, auch der durch Krankheit oder sonstwie Heimgesuchten, der Gefangenen, der Fremdlinge, kurz, aller derer, die in Bedrängniß sind, sich annimmt.“}

\quellseite{31}
Zu Anfang des 4. Jahrhunderts zerfällt der Hauptgottesdienst in eine Missa catechumenorum (Predigtgottesdienst) und in eine Missa fidelium (eigentlich Gottesdienst der Gläubigen, Abendmahlsgottesdienst). Wenn die Predigt gethan war, wandte sich ein Kirchendiener an die anwesenden Nichtchristen und Katechumenen mit den Worten: „Ite, missa est!“ — geht, ihr seid entlassen. (Daher wohl das deutsche Wort Messe.) Namentlich die nun beginnende Missa fidelium zeigt im Vergleich zu obigem sehr einfachen Acte der Communion eine größere Ausdehnung und Feierlichkeit der Abendmahlsliturgie. Nach dem 8. Capitel der sogenannten „Apostolischen Constitutionen“ begann die Missa fidelium mit einem stillen Gebete der Zurückgebliebenen, dem dann das durch einen Diakon gesprochene allgemeine und ausführliche Kirchengebet folgte, bei dessen einzelnen Absätzen die Versammlung mit einem Kyrie eleison litaneiartig einfiel. Nach einem Collectengebet des Bischofs sammelten die Diakonen die dargebrachten Gaben an Brod und Wein und richteten sie zum Abendmahlsgebrauch zu, und nachdem die Männer unter sich und die Frauen unter sich einander den heiligen Kuß gegeben hatten, trat nach Anlegung eines besonderen festlichen Gewandes und umgeben von Kirchendienern der höheren Weihen der Bischof an den Altar, sprach den apostolischen Segenswunsch 2 Cor. 13, 13. über die Gemeinde, welche denselben mit den Worten: „Und mit deinem Geiste“ erwiderte. Darauf antiphonatim: B. Empor den Sinn! G. Wir haben ihn erhoben zum HEerrn. B. Lasset uns danken dem HEerrn. G. Würdig und gerecht ist es. B. Würdig fürwahr und gerecht ist es, vor allem hoch dich zu preisen, den wahrhaft seienden Gott“ etc. Dieses freilich sehr, sehr lange, doch überaus erhabene Dankgebet schließt mit dem von der ganzen Versammlung gesungenen Sanctus: „Heilig, heilig, heilig ist Gott, der HEerr Zebaoth. Voll sind Himmel“ etc. Hierauf folgt die Consecration, ein nochmaliges, aber durch den Bischof gesprochenes allgemeines Kirchengebet mit Fürbitten, der Friedensgruß und dann die Austheilung. Letztere, wie der Schluß, wird so beschrieben: „Und nachdem Alle (auf den Friedensgruß) gesprochen haben ‚Amen‘, sage der Diakonus: ‚Lasset uns aufmerken!‘ und der Bischof rede das Volk also an: ‚Das Heilige den Heiligen!‘ und das Volk antworte: ‚Einer ist heilig, Einer der HEerr, Ein Jesus Christus zur Verherrlichung Gottes, des Vaters, gepriesen in die Ewigkeiten. Amen. Ehre in der Höhe sei Gott, und Friede auf Erden und Wohlgefallen an den Menschen; Hosianna dem Sohne Davids, gelobt sei, der da kommt im Namen des HEerrn; Gott, der HEerr, erschien uns auch, Hosianna in der Höhe!‘ Und hierauf nehme der Bischof,
\quellseite{32}
darauf die Aeltesten und die Diakonen und Unterdiakonen und die Vorleser und die Sänger und die Asketen, und unter den Frauen die Diakonissen, und die Jungfrauen und Wittwen, dann die Kinder und dann das ganze Volk nach der Ordnung mit Scheu und Vorsicht ohne Lärm. Und zwar gebe der Bischof das Opfer, indem er sagt: „Der Leib Christi“, und der Empfangende sage: „Amen“, der Diakonus aber halte den Kelch und sage (ihn) reichend: „Das Blut Christi, der Kelch des Lebens“, und der Trinkende sage: „Amen“. Der 33. (34.) Psalm aber werde gesprochen (gesungen), während alle Uebrigen empfangen, und wann alle (Männer) und alle (Frauen) empfangen haben, sollen die Diakonen das Uebriggebliebene in die Nebengemächer bringen. Und der Diakonus spreche, nachdem der Sänger aufgehört hat: „Nachdem wir Theil genommen an dem theuren Leibe und dem theuren Blute Christi, lasset uns danksagen dem, der uns gewürdigt hat, Theil zu nehmen an seinen heiligen Geheimnissen, und lasset uns beten“ etc. Das durch den Bischof gesprochene Dankgebet bildet den Beschluß.“ (S. König: Die Hauptliturgien der alten Kirche.)

Anmerkung 2. Wiewohl die Brandenburg-Nürnb. Kirchenordnung 1533 den communionlosen Gottesdienst des Sonntags nach der Weise des Nebengottesdienstes zu formen sucht, so wurde es doch allgemeiner Gebrauch, demselben bis zur Predigt die Form des Hauptgottesdienstes zu belassen. Es folgt dann die Litanei oder das Te deum oder beides im Wechsel als „Gemeindegebet“ und „Gemeindedanksagung“, oder man beschränkte sich auf ein Loblied der Gemeinde, verbunden mit einer „christlichen Motette“ des Chors, die jenes einleitet, beides gleichsam das Deo gratias auf das gehörte Gotteswort. Die Braunschweig-Lüneburgische Kirchenordnung jedoch bestimmt: „Nach der Predigt soll gesungen werden die Präfation, Sanctus, deutsch Paternoster, Christe, du Lamm Gottes, eine deutsche Sonntagscollecte und die letzte Segnung.“

\chapter[§ 3. Die Messe durch Luther]{§ 3. Die Messe durch Luther in ihrer evangelischen Gestalt wiederhergestellt}


Wie die Augsburgische Confession, deren Apologie und Luthers liturgische Schriften zeigen, beließ man dem Hauptgottesdienst seinen herkömmlichen Namen und nannte ihn auch in der rechtgläubigen Kirche Messe. Es ist dies ein Ausdruck, welcher den heiligen Dienst, den Cult, bezeichnet und dem
\quellseite{33}
griechischen liturgia entspricht, wie die Apologie Art. 12 gegenüber der papistischen Herleitung desselben von dem hebräischen Misbeach, Opferaltar, hervorhebt, indem sie sagt: „Missa und Liturgia heißen nicht Opfer; Missa heißt hebräisch ein zusammengetragen Steuer, denn also ist die Weise gewesen, daß die Christen Speis und Trank zu gut den Armen in die Versammlung gebracht haben. Und solche Weise ist von Jüden herkommen, die auf ihre Feste mußten solche Steuer bringen, die nenneten sie Missa.“ Während in der skandinavisch-lutherischen Kirche dieser Name noch immer zur Bezeichnung des Hauptgottesdienstes gebraucht wird, ist er in der deutsch-lutherischen Kirche später ganz abgekommen, ohne daß man im Nürnbergischen noch zur Zeit unserer Jugend die Communion das „Amt“ und die Predigt im Hauptgottesdienst die „Amtspredigt“ zu nennen pflegte.

Schon aber die Beibehaltung des Namens „Messe“ deutet an, daß Luther im äußerlichen Gottesdienst nichts Neues anzurichten, sondern allein zur alten rechten Weise des Gottesdienstes zurück zu führen beabsichtigte. Die ganze Reformation Luthers ging ja überhaupt nicht auf Verneinung und Zerstörung, sondern auf Wiederherstellung. Sie wollte nur abthun, was in Lehre und Leben wider Gottes Wort war und sich als verderblicher Sauerteig im Lauf der Zeit aufgehäuft hatte. Als Luther zuerst, nachdem er die Reformation begonnen hatte, nun auch seinen Blick auf die Verwaltung der Sacramente in der römischen Kirche warf, rief er aus: „Die Worte des Heils und des Lebens sind verschlossen, wie die Arche des HEerrn in dem Götzentempel neben dem Dagon war, aber die Arche muß sich selbst den Weg bahnen!“ (Rudelbach, Sacramentsworte S. 9.) Seinem Scharfblick konnte es ja nicht entgehen, daß das Wort des HEerrn hier unter einer Menge von Ceremonien gleichsam vergraben sei, die schon durch ihre Masse betäubten und die Aufmerksamkeit der Gläubigen von dem Einen Nothwendigen ablenkten. Aber wie ließ er der Arche sich selbst den Weg bahnen? Nicht in der Weise eines Zwingli und der Seinen, die hier radical mit der Vergangenheit brachen und namentlich in äußerlichen Dingen abthaten, was nicht eine ausdrückliche Begründung im Schriftworte fand, sondern erst nach der Apostel Zeit in Aufnahme gekommen war. Luther verfuhr hier vielmehr mit großer Behutsamkeit und Schonung nach Grundsätzen, die aus dem innersten Lebensborn der Schrift und der rechten evangelischen Praxis geschöpft sind. Gemäß denselben setzte er eines theils unerbittlich alles aus, was eigentlicher papistischer Sauerteig war und was insbesondere dem eigentlichen Meßopfergreuel
\quellseite{34}
angehörte oder auch nur nach demselben „roch“; anderntheils aber sonderte er den bloß kirchlich gefärbten Wust von dem, was der reinen Lehre gemäß war und als Gemeingut der Kirche sich erwies, insonderheit aber dem Wort und Sacrament förderlich und dienstlich war. „Wir bekennen“, sagt Luther in seiner „Weise, christlich Meß zu halten“ (Formula Missae), „daß wir nie gedacht, allen äußerlichen Gottesdienst abzuthun, sondern den, so bisher im Brauch ist, aber mit vielen Zusätzen verderbt, wieder zu fegen und anzuzeigen, welches der rechte christliche Brauch ist.“ Ebenso erklärt er sich in seiner „Ordnung des Gottesdienstes in der Gemeinde 1523“: „Der Gottesdienst, der jetzt allenthalben gehet, hat eine feine christliche Abkunft, gleichwie auch das Predigtamt. Aber gleichwie das Predigtamt verderbt ist durch die geistlichen Tyrannen, also ist auch der Gottesdienst verderbt durch die Heuchler. Wie wir nun das Predigtamt nicht abthun, sondern wieder in seinen rechten Schwang begehren zu bringen, so ist auch nicht unsere Meinung, den Gottesdienst aufzuheben, sondern wieder in rechten Schwang zu bringen.“ (Erl. Ausg. 22, 153.)


Wenn Luther sagt, man habe den Gottesdienst, „so bisher im Brauch ist“, nur fegen und anzeigen wollen, „welches der rechte christliche Brauch ist“, dabei aber rühmt, daß „der Gottesdienst, der jetzt allenthalben im Brauch gehet, eine feine christliche Abkunft“ habe, so bezieht er sich auf die Gliederung der Messe in ihren Hauptstücken, wie sie die abendländische Kirche aus der ältesten morgenländischen Kirche überkommen hat, im Lauf der Zeit bis zur Reformation jedoch „verderbt mit vielen Zusätzen“ besaß und mit diesen seit 1570 bleibend und unverrücklich behält. (Ein Bild von der ursprünglichen Gestalt der Messe haben wir in jenem § 2, Anmerk. 1 angeführten Zeugniß Justins, und ihrer Erweiterung in der dort gleichfalls citirten Stelle der Apostolischen Constitutionen gesehen. Noch reicht auch in der letzteren nichts nach einem sühnenden Opfer; denn ob auch in dem mit der Consecration verbundenen längeren Gebete dieser Clementinischen Liturgie Brod und Wein als Opfergaben bezeichnet werden, so findet doch nicht eine meritorielle, sondern allein eine materielle Darbringung der Elemente statt, indem es daselbst ausdrücklich heißt: „... und bitten dich, daß du wohlwollend blickest auf diese vor dir daliegenden Gaben; du bedürfnißloser Gott, und Wohlgefallen an ihnen habest zur Ehre deines Christus, und herabsendest auf dieses Opfer deinen Heiligen Geist, den Zeugen der Leiden des HEerrn Jesus, daß er dieses Brod mache zu dem Leibe deines Christus, und diesen Kelch zum Blute deines Christus, damit die daran Theilnehmenden
\quellseite{35}
befestigt werden in der Gottseligkeit, die Vergebung ihrer Sünden erlangen, von dem Teufel und seinem Betrug befreit werden, vom Heiligen Geiste erfüllt, deines Christus würdig werden, das ewige Leben erlangen, nachdem du mit ihnen versöhnt bist.“ (König, Die Hauptliturgien etc. S. 114.) Als nun aber die seit dem 2. Jahrhundert übliche „Arcandisciplin“, das ist, die Geheimhaltung der Mysterien — Taufe, Abendmahl, apostolisches Symbolum, Vater unser — allmählich verschwunden war, so bildete sich seit dem 4. Jahrhundert die Ordnung der römischen Messe aus. Wohl ist dieselbe in ihrer aus der ältesten morgenländischen Kirche überkommenen Gliederung an Klarheit, Wirkung und Schönheit erheblich verbessert, aber auch — durch den in den Mittelpunkt gestellten Meßopfergreuel, nebst Transsubstantiation, Kelchraub, Heiligenanrufung, Winkel- und Seelmessen andererseits wieder sehr verderbt.

Diesen Gottesdienst nun, der eine „feine Abkunft“ hatte, aber hernach durch den Opfergreuel und andere „Zusätze“ im Papstthum so schrecklich verderbt worden ist, „fegte und reinigte“ Luther, wie man ein altes Kunstgemälde von Staub und den Flecken einer fremden Hand reinigt und es wieder auffrischt, oder wie man ein edles Bauwerk von theils entstellenden, theils überflüssigen Beiwerk späterer Zeit befreit, damit es in der Schönheit seiner wahren Gestalt wieder erscheine. In ihrer ursprünglichen und schönen, in ihrer dem Evangelium gemäßen Gestalt gab Luther die Messe dem Volke Gottes wieder!

\chapter[§ 4. Grundsätze und Verfahren Luthers]{§ 4. Grundsätze und Verfahren Luthers bei Wiederherstellung der evangelischen Messe}


Welche Gedanken es insbesondere waren, die Luther bei der Reformation des Gottesdienstes überhaupt leiteten, spricht er bereits in der „Ordnung des Gottesdienstes in der Gemeinde 1523“ also aus: „Drei große Mißbrauch sind in dem Gottesdienst gefallen. Der erst, daß man Gottes Wort geschwiegen hat und alleine gelesen und gesungen in den Kirchen; das ist der ärgste Mißbrauch. Der ander, da Gottes Wort geschwiegen gewesen ist, sind neben einkommen so viel unchristlicher Fabeln und Lügen, beide in Legenden, Gesängen und Predigten, daß greulich ist zu sehen. Der dritte, daß man solchen Gottesdienst als ein Werk gethan hat, damit Gottes Gnade und Seligkeit zu erwerben. Da ist der Glaube untergegangen
\quellseite{36}
und hat Jedermann zu Kirchen und Stiften gegeben, Pfaff, Mönch und Nonne werden wollen.“ (Erl. Ausg. 22, 153. St. Louiser Ausg. 10, 221.) Es sind eben die auch hier von Luther in Anwendung gebrachten Principien keine andern, als das Doppelprincip der lutherischen Kirche überhaupt, das sogenannte formale und materiale Grundprincip, das „Allein die Schrift“ und das „Allein durch den Glauben“.

Demgemäß setzte Luther bei Wiederherstellung der alten rechten Ordnung der Messe vor allem die Predigt des göttlichen Wortes in ihr volles Recht ein. Das erschien ihm ja als der ärgste Mißbrauch, daß man Gottes Wort „geschwiegen“ hat, daß die Predigt nach und nach aus der Messe verdrängt worden war, wie sie denn heute noch im römischen Cultus nicht einmal als ein organisches Glied, sondern mehr als eine freie Zugabe, als ein Einschiebsel, erscheint, bei dessen Wegfall die Messe nichts an ihrer papistischen Herrlichkeit verliert. Damit aber gab Luther der Messe einen bestimmten Halt. Als das Evangelium von der freien Gnade Gottes in Christo ist sie in den Mittelpunkt gestellt und bildet so den Grundgedanken des ganzen Hauptgottesdienstes, so daß alle die der Predigt vorausgehenden liturgischen Glieder ihre vorbereitende Stellung und Bedeutung und die ihr nachfolgenden ihre bestimmte Beziehung und Bedeutung für die Aneignung und Versiegelung dieser Gnade im Sacrament des Altars haben. Von selbst mußte daher Alles weichen und fallen, was diesem Evangelio von der freien Gnade Gottes und dem Sacrament als seinem Siegel zuwider war, was irgendwie an die Lehre vom Verdienst der Werke oder verdienstlicher Verrichtung des Gottesdienstes oder der Messe selbst sich anschloß, aus ihr floß, an sie erinnerte. Während daher die anglicanische Kirche jenes vorhin mitgetheilte, recht verstanden nicht schriftwidrige materielle Opfergebet der „Apostolischen Constitutionen“ fast wörtlich als “Oblation” und “Invocation” ins “Common-Prayer” aufgenommen hat, strich Luther auch dieses, auf daß nichts überall auch nur im entferntesten nach einem Opfer „rieche“ und von falschen Geistern irgendwie gemißbraucht werden möchte — und wir wollen keineswegs dies ihm mit manchen Neueren dahin deuten, als hätte er im reformatorischen Eifer des Guten zu viel gethan, sondern wollen es ihm vielmehr danken.

Vergegenwärtigen wir uns nun diese seine Arbeit im Prüfen der römischen Meßordnung, da er das Böse zu „fegen“ und das, „so das Beste ist“, allein „zu behalten“ suchte.

\quellseite{37}
\section[Die römisch-papistische Messe]{Die römisch-papistische Messe gliedert sich folgendermaßen:}

I. Initium missae solennis. Nachdem der Priester, im Meßgewand und hinter den Kerzen tragenden Meßdienern, zum Altar gegangen und den bedeckten Kelch auf denselben gestellt hat, steigt er wieder herab, stellt sich aufrecht am Fuß des Altars, bekreuzt sich und betet dann aus Ps. 43.: Introibo ad altare Dei etc. Es soll dies gleichsam die Vorbereitung des Priesters zum Eintritt in den (Opfer-) Altar sein.

II. Das Confiteor oder die Beichte, die Weihe des Priesters und seiner Ministranten zum Meßdienst. Der Priester spricht es vorwärts gebeugt am Fuß des Altars stehend. Bei den letzten Worten des damit verbundenen kurzen Gebetes besteigt er den Altar, küßt, beräuchert denselben, läßt sich selbst dann beräuchern und spricht hierauf auf der Epistel- oder rechten Seite des Altars:

III. Den Introitus oder die Eingangspsalmodie, die nach Zeit und Umständen verschieden ist.

IV. Das Kyrie in der Mitte des Altars neunmal und mit den Altardienern abwechselnd gebetet.

V. Das Gloria in excelsis mit dem Et in terra etc. und dem Laudamus te etc.

VI. Die Collecta(e). mit vorausgehendem Dominus vobiscum, nach Zeit, Fest und sonstigem Bedürfniß eine oder etliche.

VII. Epistola. Lectio derselben im liturgischen Ton.

VIII. Graduale (Stufengesänge) mit Halleluja oder Tractus (Gesang beim Fortziehen oder Gehen zur Evangelienseite).

IX. Evangelium, vom Diakon auf der linken Seite des Altars im liturgischen Ton gelesen, nachdem er zuvor sich hierzu vom Priester den Segen erbeten und diesem für denselben die Hand geküßt hat.

X. Das Credo Nicaenum, das der Priester, in der Mitte des Altars stehend, intonirt. Während es der Chor hinaussingt, betet ersterer.

XI. Das Offertorium oder Opfergebet (unter Ergießung des Weines mit etwas Wasser in den Kelch und Darbringung der Hostie nebst allerlei Beräucherung), woran sich sofort

XII. Die Secreta, die stillen oder heimlichen Gebete, reihen. Alles, was von der Intonation des Credo an bisher vorgenommen ist, spricht und handelt der Priester mit seinen Gehilfen in der Stille, während der Chor das Credo vollendet. Dann aber singt er laut den Uebergang: Dominus vobiscum etc. zur
\quellseite{38}
XIII. Praefatio, welcher der Chor mit dem Sanctus antwortet, während denn aber der Priester stille den

XIV. Canon missae, die Stillmesse, spricht. Diese besteht a. aus einem Gebet für die Kirche überhaupt, b. einer Fürbitte für einzelne Kirchglieder (commemoratio pro vivis), c. dem Erinnerungsgebet an die verstorbenen Heiligen, Maria an der Spitze, und ihrer verdienstlichen Fürbitte (Diptycha sanctorum), d. der Recitation der Einsetzungsworte oder Consecration, nach und unter eingetretener tiefer, allgemeiner Stille vom Priester leise gesprochen, e. dem Gebet um gnädige Aufnahme des jetzt vollendeten Opfers (canon missae post consecrationem), und f. dem Gebet für die Verstorbenen (Memento pro defunctis), denen das Meßopfer mit zu gute kommen soll. Hier schweigt der Chorgesang und der Priester singt laut die

XV. Praeparatio ad Communionem, bestehend aus dem Gesang des Vater Unsers und eines an die 7. Bitte sich anschließenden Gebetes, wobei ein Stückchen der Hostie gebrochen und in den Kelch geworfen wird, und worauf dann der Friedensgruß (Pax Domini vobiscum) folgt. Hierauf

XVI. Preces ante Communionem, bestehend aus dem Agnus Dei (Christe, du Lamm Gottes) mit etlichen Collecten;

XVII. Die Sumtio, Genießung des Brodes und Weines seitens des Priesters und Austheilung des Brodes an Communicanten, wenn solche vorhanden sind;

XVIII. Communio, das ist, ein nach Zeit und Umständen verschiedener biblischer Spruch, die Antiphone von dem Psalm, der ehedem während der Communion gesungen wurde.

XIX. Postcommunio, eine Collecte mit dem Ite, missa est;

XX. Finis missae, stilles Gebet des Priesters;

XXI. Benedictio, Segnung des Volkes;

XXII. Evangelium Joh. Cap. 1, v. 1–14., das der Priester am Altar noch liest oder singt, worauf das Deo gratias des Meßdieners schließt.

In seiner Formula missae 1523 nun bezeichnet Luther von dem ersten bis zum Credo incl. reichenden Theil der Messe als beizubehalten Nr. III bis X, also:

I. Introitus, wofür ihm jedoch eigentliche Psalmen lieber wären;

II. Kyrie, jedoch nicht das neun-, sondern das dreimalige;

III. Gloria mit dem Et in terra und dem Laudamus te:

IV. Collecta, jedoch nur eine;

\quellseite{39}
V. Epistola;

VI. Gradual von 2 Versen mit dem Halleluja, oder auch nur eines von beiden;

VII. Evangelium;

VIII. Credo oder Symbolum Nicaenum — worauf die Predigt folgt, wenn sie nicht dem Introitus bereits vorausgegangen.

Hiernach hat Luther gestrichen:

1. Initium missae sammt dem Confiteor. Er betrachtete somit den Introitus als den eigentlichen Anfang der Messe, wie er es auch ursprünglich war. Initium und Confiteor erschienen ihm wohl als bloße Sache des Priesters für sich, auch enthält das Confiteor allerlei papistische Zutaten.\footnote{Das Confiteor der römischen Messe lautet: „Confiteor Deo omnipotenti, beatae Mariae semper virgini, beato Michaeli archangelo, beato Joanni Baptistae, sanctis apostolis Petro et Paulo, omnibus sanctis et vobis, Fratres: quia peccavi nimis cogitatione, verbo et opere (percutit sibi pectus ter, dicens): mea culpa, mea culpa, mea maxima culpa! Ideo precor beatam Mariam semper virginem, beatum Michaelem archangelum, beatum Joannem Baptistam, sanctos apostolos Petrum et Paulum, omnes sanctos et vos, Fratres, orare pro me ad Dominum, Deum nostrum.“}

2. Tractus, Prosen (Reime) und Sequenzen. Nicht, als ob diese etwas Schriftwidriges enthalten hätten, sondern um dem Volke den Gottesdienst durch überflüssige Länge nicht zu verleiden, weshalb er denn auch das Gradual selbst auf 2 Verse beschränkte. Und in der wiederhergestellten Messe sollte ja als die Hauptsache das Wort und die Predigt recht hervortreten und ihr alles Vorausgehende dienen! Für sie mußte der nöthige Raum gesichert werden — wenn auch mit Drangeben der einen oder andern wirklichen Zierat.

Ganz anders aber gab es Durchstriche im zweiten Theil der Messe, der ja mit dem eigentlichen Opfergreuel untermengt ist. Hier sah Luther eigentlich „die Worte des Heils und des Lebens verschlossen, wie die Arche des HErrn in dem Götzentempel neben dem Dagon“. Gleich von vornherein fiel daher das Offertorium, denn von da aus „flingt und riecht nach Opfer fast Alles“. (An seine Statt trat später und behauptete bis heute seine Stelle u. A. das „Schaffe in mir, Gott“, wovon an seinem Orte das Nähere.) Dem Offertorium folgten unerbittlich die Secreta und der ganze Canon missae. Vom letzteren schreibt Luther: „Summa, was böse in der Messe ist, ist vom Opfer und Werk, das hat Gott wunderlich geschickt, daß fast alles der Priester heimlich liest,
\quellseite{40}
und heißt die Stillmesse.\footnote{Dieses „Leise und heimlich“ lesen hängt mit der papistischen Opferidee zusammen, denn wie Bellarmin offen erklärt, sei Opfern „nicht ein Reden, sondern ein Thun, und wenn es dennoch in gewissem Sinne auch zugleich ein Reden sei, so sei es ein Reden nicht in oder zu der Gemeinde (ecclesia), sondern zu Gott; wer Gott nämlich ein Opfer bringe, handle mit Gott, nicht mit Menschen, wenn schon sein Handeln keinen Privat-, sondern einen öffentlichen Charakter trage, weil er nicht bloß für sich, sondern für die ganze Kirche handle“. (S. Herzog, Real-Encycl. Art. Messe IX, 395.)} Was aber öffentlich durch den Chor und unter dem Haufen gesungen wird, ist fast eitel gut Ding und Lobgesang, als sollte Gott mit der That sagen, er wolle seine Christen mit der Stillmesse schonen, daß ihre Ohren solchen Greuel nicht müssen hören, und also die Geistlichen mit ihrem eigenen Greuel plagen lassen.“ (Jen. A. 5, 196 b.) Bekanntlich hat Luther fast gleichzeitig auch seine Schrift: „Von dem Greuel der Stillmesse, so man den Canon nennet“ ausgegeben lassen, in welcher er denselben für das Volk verdeutscht und glossirt, denn er wollte dem Teufel die Decke abthun“ und ihn auch vor dem Volke „an das Licht stellen, daß Jedermann sehe, was für schrecklich, greulich Ding in der Stillmesse der leidige Teufel treibt täglich in aller Welt“. (E. A. 29, 118.)

Aus diesem Meßkanon nun nahm Luther in seiner ersten liturgischen Arbeit die Präfation und die Einsetzungsworte heraus, letztere als eigentliche Consecratio, welche aber gleich dem Vater Unser nicht leise, wie in der papistischen Messe,\footnote{„Von besonderer Wichtigkeit ist es, daß auch die Consecrationsworte, die ja angeblich ebensowohl die Wesensverwandlung bewirken, als das Opfer constituiren, leise gesprochen werden; sie gehen die Gemeinde nicht an, sondern gelten nur den Elementen, um an diesen in magischer Kraft und Wirkung die Veränderung hervorzubringen, durch die der Priester täglich Christi Leib macht (conficere ist der bezeichnende Ausdruck).“ (Herzog, Real-Encycl. ibid.)} sondern laut recitirt, gesprochen oder gesungen werden sollen, „damit sie von den Umstehenden vernommen werden können“. Nach der Consecration sollte das Sanctus folgen, dann das Vater Unser und kein weiteres Gebet, sondern sofort darauf folgend Pax Domini. Während der Sumtio (Austheilung) soll das Agnus gesungen werden. Den Communion-Spruch läßt Luther gelten (seine Stelle vertrat dann bald das Communion-Lied); statt der üblichen Complet (Postcommunio, Collecten) will er ein anderes Gebet, „weil sie ziemlich vom Opfer lauten“. Statt dem Finis missae, der Benedictionsformel und der Johannislection, welche ohnehin die Perikope des dritten Weihnachtfeiertages bildet, setzt Luther den einfachen Schluß mit Dominus vobiscum, Benedicamus (Der HEerr sei mit euch — Und mit deinem Geiste.
\quellseite{41}
Laßt uns benedeien den HEerrn — (Gott sei ewiglich Dank!) und dem Aaronischen Segen, 4 Mos. 6.

So in der Formula Missae 1523. Noch mehr vereinfacht erscheint diese Abendmahlsliturgie in der deutschen Messe 1526. In derselben ist nämlich an die Stelle der Präfation die Paraphrase des Vater Unser mit kurzer Vermahnung gesetzt, worauf sofort die Consecration folgt und die Austheilung unter passenden Gemeindegesängen vor sich geht. Den Schluß macht eine Danksagungscollecte.

Ausdrücklich jedoch bemerkt Luther, daß er seine Formula Missae von 1523 nicht „aufgehoben oder verändert“ haben wolle. Sie ist auch der durchgreifende Typus für die lutherische Kirche geblieben, während von der 1526 gegebenen Form nur die an das paraphrasirte Vater Unser angehängte Vermahnung hie und da in Agenden überging, die Danksagungscollecte am Schluß aber in allen ihre Stelle fand. Es sind das namentlich die sächsischen und niederdeutschen Agenden einerseits und die österreichische und nürnberg-brandenburgische Agende andererseits, welche sich möglichst eng an Luthers Vorgang in seinen liturgischen Schriften anschlossen, während die süd- und westdeutschen Agenden, die Badener, Andorffer, Hanauer, Wormser, Hohenloher und auch später die Straßburger, bei aller Entschiedenheit im lutherischen Bekenntniß dem Formulare der reformirten Gottesdienstweise mehr Rechnung tragen. Gleichwohl aber banden sich die Agenden der ersteren Art bei aller Bewahrung des streng lutherischen Typus und bei allem Streben nach Gleichförmigkeit in den Ceremonien an Einzelheiten nicht sklavisch und pedantisch, und gestattete man sich selbst in der Consecrations- und Distributionsformel Mannigfaltigkeit innerhalb der Einheit des Bekenntnisses.

Alles zusammengenommen finden wir daher für den zweiten Theil der gereinigten Messe folgende beibehaltenen Stücke:

I. Praefatio mit Sanctus — wofür die Brandenburg-Nürnberger K.-O. von 1533 die Vermahnung — es ist dies die in der St. Louiser Agende S. 233 ff. enthaltene Volprecht’sche — gesetzt hat, während andere Agenden, namentlich sächsische, Präfation und Vermahnung haben.

II. Consecratio mit Vater Unser — bald das eine nach, bald das andere voraus, nebst dem Friedensgruß (Pax), welchen letzteren jedoch nicht alle Agenden aufgenommen haben.

\quellseite{42}
III. Agnus Dei (Christe, du Lamm Gottes) und andere Lieder, namentlich Abendmahlslieder, während der eigentlichen Communio oder Sumptio.

IV. Collecta mit: Dominus vobiscum und — jedoch nicht überall — vorausgehendem Gesang des Nunc dimittis (HEerr, nun lässest du deinen Diener in Friede fahren).

V. Aaronischer Segen mit vielfach vorausgehendem Benedictamus, worauf dann hie und da noch ein Lied vorgeschrieben wird, als z. B. „Gott sei gelobet“ etc., „Erhalt uns, HEerr“ etc., namentlich da, wo das Nunc dimittis weggefallen war.

\chapter[§ 5. Beibehaltung von Amtskleidern]{§ 5. Beibehaltung von Amtskleidern und deren Gebrauch in der Folgezeit}


„Servantur usitatae Ceremoniae publicae ordo lectionum, orationum, vestitus, et alia similia“ erklärt die lateinische Apologie im Artikel de Missa. Selbst also hinsichtlich der bisherigen Amtskleider that man nicht alles ab, sondern behielt Manches bei. In der „Weise, christliche Meß zu halten“ (Formula Missae), kommt Luther schließlich auf die Amtskleider zu sprechen und äußert sich also: „Von Kleidern haben wir nicht gesagt; halten aber davon, wie von andern dergleichen äußerlichen Geberden. Wir lassen geschehen, daß man ihrer frei brauchen möge, doch daß das Gepränge und übriger Ueberfluß vermieden werde. Denn du bist vor Gott nicht angenehmer, so du in priesterlichen Kleidern, noch unangenehmer, so du ohne solche Kleider das Sacrament handelst; denn uns auch die Kleider vor Gott nicht fördern. Ich wollte auch gern, daß man sie weder weihete, noch segnete, als sollten sie hinfort heiliger sein, denn andere Kleider; es wäre denn, daß man einen allgemeinen Segen gebrauchen wollte, da durch das Wort und Gebet, wie die Schrift lehrt 1 Tim. 4, 4. 5., alle gute Gottes Creatur geheiliget wird; sonst ist’s lauter Aberglaube und gottlos Wesen, eingeführt durch die Baalsbischöfe des höchsten und letzten Greuels in der Kirche, wie anderes dergleichen mehr.“ (St. L. A. X, S. 2246.) Wie man es nun dieserhalb hiernach z. B. in der Kirche zu Wittenberg hielt, darüber schreibt Luther in seiner Schrift „Wider die himmlischen Propheten“, 1525: „Im Kloster haben wir Meß gehabt ohn Kasel, ohn Aufheben, schlecht aufs allereinfältigst, wie Carlstadt Christus Exempel rühmet. Wiederum in der Pfarr haben wir noch Kasel, Alben, Altar, heben auf, wie lange es uns gelüstet.“ (E. A. 29, 191.)

\quellseite{43}
Das gewöhnliche Amtskleid bestand in dem langen schwarzen, bis an die Knöchel reichenden Priesterrock, jetzt Chorrock oder Talar genannt, und einem bis etwa an die Kniee reichenden weißleinenen Ueberwurf, Alba oder Chorhemd genannt, letzteres von der alten Kirche herrührend, in welcher das weiße Kleid als Symbol des Friedens, der Reinheit und der Würde der gottesdienstlichen Personen getragen wurde. Bis in das erste Viertel unseres Jahrhunderts trugen bei der Verwaltung der Sacramente und der Verrichtung der Altarliturgie die Pastoren in den lutherischen Kirchen Deutschlands das Chorhemd, in einigen Gemeinden, wie z. B. in Thüringen und selbst in dem sonst so liturgisch armen Württemberg an ein paar Orten, ist es noch im Gebrauch, selbst hier zu Lande bediente man sich, so viel ich weiß, seiner in den paar in Texas eingewanderten und zur Missourisynode gehörenden Gemeinden noch immer, wie auch unter unsern hiesigen norwegischen Brüdern. In Sachsen, in Braunschweig, im Brandenburg-Nürnbergischen Gebiete und noch anderwärts aber erschien zur Administration des Abendmahls der Liturg in dem eigentlichen Meßkleide, zu welchem außer dem Chorrock und dem Chorhemd noch die Casula oder die Kasel gehörte, der, mit einem goldgestickten Kreuz gezierte und verschiedenfarbige, ärmellose Umhang. So heißt es z. B. in einer Beschreibung des Gottesdienstes von Alt-Dresden aus der Zeit des 200jährigen Reformationsjubiläums: „Nach diesem gehet die Communion an, zu dem Ende der Diaconus mit der Alba und Meßgewand angethan in Begleitung zweier dazu besonders ausgekleideten und zu Haltung der Communiontüchlein oder Tüchel verordneten Knaben sich zu dem ausgezierten Altar verfüget.“ (Etwas zur K.-Hist. in Alt-Dresden etc. durch M. Paul Hilscher, Pastor allda. 1721.) Ebenso heißt es in Beschreibung der Einweihung der neuen Friedrichstädter-Kirche in Dresden 1730, bei welcher Dr. B. Löscher die Festpredigt hielt: „Dann ward das Te deum laudamus bei Trompeten und Pauken angestimmt. . . Mittlerweile der Pastor ord. templi, Herr M. David Mehner, nach dem Altar im grünen Meßgewand (Kasel) sich verfüget und nebst dem Vaterunser die Einsetzungsworte zur heiligen Communion abgesungen.“ In Nürnberg, der Vaterstadt des Verfassers, war noch bis zum Jahre 1790 bei Verwaltung des Abendmahls die Kasel im Gebrauch; ja, bei der Einweihung einer norwegisch-lutherischen Kirche in Wisconsin, der vor circa 40 Jahren derselbe beiwohnte, war der damalige Synodalpräses mit der Kasel bekleidet. Es war der Nationalismus, der zur „Religionsverbesserung“ auch in Betreff des Amtskleides seine Wandelbarkeit gethan, wenigstens in der deutsch-lutherischen Kirche. Bei der Rückkehr zur altlutherischen Liturgie
\quellseite{44}
kann um des Friedens und anderer Ursachen willen von einer Wiedereinführung selbst auch nur des Chorhemdes nicht mehr die Rede sein und können wir uns nur um so mehr an dem schwarzen Amtskleid genügen lassen, da selbst dieses im Lande des reformirten Sectenthums bei Gründung unserer Synode bereits zur Seltenheit geworden war.

\chapter[§ 6. Der Gebrauch der Landessprache]{§ 6. Der Gebrauch der Landessprache und die Einführung des geistlichen Volksliedes in die Liturgie — die merklichen Aenderungen. Bedeutung des Chorgesangs und Gebrauch der Orgel}


Zu den doch merklichen Aenderungen in der Messe gehört unstreitig der Gebrauch der Landessprache im Gottesdienste und die Einführung des geistlichen Volksliedes in die Liturgie. Durch diese Aenderung ist zugleich das geistliche Priestertum aller Christen in seinem gottgeordneten Verhältnis zum öffentlichen Beruf und Dienst am Wort auch durch die Liturgie zu seinem Recht und Brauch gekommen, und das in einem Maße, wie es seit der Apostel Zeit nicht der Fall gewesen ist, so daß nächst dem reinen Wort und Sacrament um dieser Aenderung willen die Augsburgische Confession wohl bekennen kann, „daß die Meß ohne Ruhm zu reden bei uns mit größerer Andacht und Ernst gehalten wird, denn bei den Widersachern“. Gerade in der durch den Gemeindegesang geschaffenen Theilnahme der Gemeinde am Gottesdienst spricht sich der der ganzen Reformation zu Grunde liegende Gedanke des allgemeinen Priesterthums recht „merklich“ aus — und in ihm liegt ja der tiefste Grundunterschied in der Anschauung von der Kirche selbst.

In der papistischen Kirche wird allerorts die Messe, wie die ganze Liturgie im Gottesdienste überhaupt, ausschließlich in der lateinischen Sprache gehalten. Sie ist die eigentliche Kirchensprache, die Universalsprache im Reiche des Antichrists, mit der gleich dem heidnischen Rom das päbstliche Rom sein starkes Joch den Völkern aufgehalst hat und deren Gebrauch als einer dem Volke unverständlichen Sprache mit der ganzen Auffassung des Meßopfers zusammenhängt. Da der Gebrauch einer fremden, dem Volke ganz und gar unverständlichen Sprache in directem Widerspruche mit 1 Cor. 14, 1–20. steht, so mußte nothwendig die Kirche der Reformation zum ursprünglichen Gebrauch der Landessprache in der Messe, wie im ganzen Gottesdienst zurückkehren, „sintemal“,
\quellseite{47}
wie es Art. 24 der Augspurgischen Confession heißt, „alle Ceremonien fürnehmlich dazu dienen sollen, daß das Volk daran lerne, was ihm zu wissen von Christo noth ist“. Zwar bemerkt dieselbe, „daß an etlichen Orten deutsche Gesänge (das Volk damit zu lehren und zu üben) neben lateinischem Gesang gesungen werden“; ja, von seiner erst in lateinischer Sprache herausgegebenen Formula missae erklärt Luther in seiner Schrift: „Deutsche Messe und Ordnung des Gottesdienstes 1526“: „Diese will ich hiermit nicht aufgehoben oder verändert haben; sondern wie wir sie bisher bei uns gehalten haben, so soll sie noch frei sein, derselbigen zu gebrauchen, wo und wenn es uns gefällt, oder Ursach bewegt. Denn ich in keinem Weg will die lateinische Sprache aus dem Gottesdienst lassen gar wegkommen.“ Aber warum? „Denn es ist mir alles um die Jugend zu thun“, fährt er fort, „und wenn ich’s vermöchte und die griechische und hebräische Sprache wäre uns so gemein als die lateinische und hätte so viel seiner Musika und Gesangs, als die lateinische hat; so sollte man einen Sonntag um den andern in allen vier Sprachen, Deutsch, Lateinisch, Griechisch und Ebräisch, Messe halten, singen und lesen.“ (Erl. A. 22, 229.) Also nur darum noch der Mitgebrauch der lateinischen Sprache in der Messe, damit die studirende Jugend in dieser Sprache geübt werde, und weil zu damaliger Zeit die lateinische Sprache mehr „gemein“ war und diese Sprache in der Messe „so viel seiner Musika und Gesangs“ hatte. Eben deshalb heißt es im Unterricht der Visitatoren“, 1528 und 1538: „Etliche singen deutsche, etliche lateinische Messen; welches wir lassen geschehen. Doch wird für nützlich und gut angesehen, wo das meiste Volk des Latein unverständlich, daß selbst deutsche Messe zu halten, damit das Volk den Gesang und anderes, was gelesen wird, desto besser verstehen möge: wie St. Paulus sagt 1 Cor. 14 (V. 6.): Wenn du aber segnest im Geist, wie soll der, so anstatt des Laien stehet, Amen sagen auf deine Danksagung, sintemal er nicht weiß, was du sagest? Du danksagest wohl fein, aber der andere wird nicht davon gebessert. Nun sagt ja Paulus auch an demselben Orte (V. 26.): lasset es alles geschehen zur Besserung.“ (Luther, St. Louiser Ausg. X, 1665.)

Schon im Jahre 1520 drängte sich Luthern der Wunsch auf: „Wollte Gott, daß wir Deutschen Meß zu deutsch läßen.“ Im Jahre 1523 erschien dann auf einem Quartbogen: „Von Ordnung des Gottesdienstes in der Gemeinde“, und bald darauf: „Formulae Missae et communionis pro Ecclesia Wittenbergensi“, nach welchen, da sie Paul Speratus alsbald ins Deutsche übersetzt
\quellseite{48}
hatte, bereits am Weihnachtsfest 1525 in der Pfarrkirche zu Wittenberg deutsche Messe gehalten wurde. Zur Fortsetzung und Vervollständigung gab Luther schließlich im Jahre 1526 auf fünf Quartbogen die Schrift: „Deutsche Messe und Ordnung des Gottesdienstes“ heraus, nach welcher schon an die Stelle der lateinischen Hymnen und Sequenzen deutsche Psalmen und Gesänge in bestimmter Reihenfolge treten sollten.

Während nun aber ein Zwingli nur von einem Singen deutscher Psalmen und Lieder etwas wissen wollte und in schwärmerischer Weise des liturgischen Singens spottete und es abthat, wollte Luther die Messe nach wie vor gesungen haben, aber nicht bloß in lateinischer, sondern ebenso auch in deutscher Sprache, und zwar nicht nur, weil das liturgische Singen bisher Brauch war und unverwerflich ist, sondern auch, weil das Singen im Gottesdienst überhaupt gut und schriftgemäß ist. Er, welcher ein eben solcher Liebhaber, als Kenner der Musik war und derselben ihre Stelle nächst der Theologie anwies, schreibt deshalb an dem schon angeführten Orte: „Das Gesänge in den Sonntagsmessen und Vesper laß man bleiben, denn sie sind fast gut und aus der Schrift gezogen, doch mag man’s wenigern oder mehren“; und ebenso in seiner „Weise, christlich Meß zu halten“: „Viel Gesang in der Meß ist fein und herrlich vom Danken und Loben gemacht und bisher geblieben.“ Unter Beihilfe der kurfürstlichen Kapellmeister Johann Walther und Konrad Rupff, die er sich für etliche Wochen zu einer „Cantorei“ in seinem Hause erbeten hatte, versah er daher die 1526 erschienene „Deutsche Messe“ mit den nöthigen Noten für den Liturgen und zugleich mit einer Anleitung zum liturgischen Singen. (S. Luther, St. Louiser Ausg., X, 235 ff.) Zwar wurde hier die herkömmliche gregorianische Weise zum Muster genommen, aber wie war er, der frühere talentvolle Kurrendschüler und Virtuose im Lautenschlagen, auch hier wieder der rechte Meister! „Wir werden behaupten können“, sagt C. Naumann, „Illustrierte Musikgeschichte“, 1883, „daß Dr. Martinus völlig dazu befähigt gewesen, jedem beliebigen Kirchentexte, selbst wenn dieser nicht metrisch gegliedert war (z. B. einem Introitus, Psalm, einer Collecte, Epistel, Präfation oder einer Evangelienstelle), schon vorhandene Melodien aus der Liturgie der älteren Kirche unterzulegen, und zwar mit allen in der deutschen Uebersetzung nothwendigen Veränderungen der Tonfolgen, sowie mit schärfster Berücksichtigung der Accentuation und Bedeutung der Worte hinsichtlich ihres besonderen Sinnes und Ausdrucks: Rücksichtnahmen, die eine große musikalisch-rhythmische und melodische Gestaltungs- und Umgestaltungsfähigkeit bei unserm
\quellseite{49}
Reformator voraussetzen. Als Zeuge für diese unsere Behauptung sei hier ... Johann Walther citirt, der sich hinsichtlich der berührten Dinge bei Michael Prätorius wie folgt über Luther äußert: „Er“ (Luther) „hat auch die Noten über die Episteln, Evangelia und über die Wort der Einsetzung des wahren Leibes und Blutes Christi selbst gemacht und mir vorgesungen und mein Bedenken darüber hören wollen.\footnote{„Der Tonsetzer Joh. Walther bezeichnet hier also seinen Freund Luther geradezu als Componisten. Denn Walther unterscheidet ... sehr deutlich den bloßen Bearbeiter einer Melodie von dem Erfinder einer solchen. ... Wenn sich nun auch, die Noten über die Episteln und Evangelia durchaus auf den lateinischen cantus gregorianus zurückführen lassen dürften, so ist es doch eben so gewiß, daß Luther bei einer solchen Arbeit vielfach mit selbständigen und in sein eigenes Ermessen gestellten Veränderungen und Umbildungen der von ihm zu benutzenden Tonfolgen vorgehen mußte — ein Verfahren, durch das sich ein Talent verräth, dem das Erfinden einer eigenen Melodie verhältnismäßig weit leichter werden mußte, als die in dem erwähnten Falle vorzunehmenden musikalischen Operationen.“ (Naumann.)} An einer andern Stelle desselben Schriftstückes sagt Walther: „Wie denn unter andern aus dem deutschen Sanctus zu ersehen, wie er alle Noten auf den Text nach dem rechten accent und concent\footnote{Sämtliche Gesänge des cantus gregorianus werden nämlich dem Gebrauche nach in den Concentus und Accentus eingetheilt. Der Concentus bezieht sich auf alle Responsorien- und Antiphonengesänge, die Psalmenmelodien und die eigentlichen Chorgesänge, der Accentus hingegen auf die, welche der Priester allein zu singen hat, wie die Collecten, Lectionen, Intonationen, Präfationen etc.} so meisterlich und wohl gerichtet hat, und ich auch die Zeit seine Ehrwürden zu fragen verursachet ward, woraus oder woher sie doch dies Stück oder Unterricht hätten: darauf der theuere Mann meiner Einfalt lachte, und sprach: „Der Poet Virgilius hat mir solches gelehret, der also seine Carmina und Wort auf die Geschichte, die er beschreibet, so künstlich appliciren kann: Also soll auch die Musica alle ihre Noten und Gesänge auf den Text richten.““ . . . Er kannte endlich auch ganz genau die Vortragsweise des gregorianischen cantus choralis und hat sich eingehend über die Natur, den besonderen Charakter und die besondere Ausdrucksfähigkeit eines jeden einzelnen der Kirchentöne, auf welche damals sowohl noch die gregorianischen Melodiengänge, als die gesammte Kunstmusik gegründet waren, unterrichtet. Walther sagt in dieser Beziehung von seinem großen Freunde: „Denn da er vor vierzig Jahren die deutsche Messe zu Wittenberg anrichten wollte, hat er durch seine Schrift an den Kurfürsten zu Sachsen, und Herzog Johannsen, hochlöbl. Gedächtnis, Seiner kurfürstlichen Gnaden die Zeit alten Sangmeister Ehre Conrad Rupff und mich gen Wittenberg erfordern lassen, dazumalen von den
\quellseite{50}
Choral Noten und Art der acht Ton Unterredung mit uns gehalten und beschließlich hat er von ihm selbst die Choral Noten octavi Toni der Epistel zugeeignet, und Sextum Tonum dem Evangelium geordnet, und sprach also: „Christus ist ein freundlicher HERR und seine Rede sind lieblich, darum wollen wir Sextum Tonum zum Evangelio nehmen; und weil St. Paulus ein ernster Apostel ist, wollen wir Octavum Tonum zur Epistel ordnen.“\footnote{Der 6. Kirchenton gleicht unserer heutigen C-dur-Scala; der 8. Kirchenton einer D-moll-Scala ohne Versetzungszeichen, daher ohne ein B oder Cis.} Somit waren Luthern „die einzelnen Kirchentonarten zu besonderen, von einander grundverschiedenen Gebieten musikalischen Ausdrucks geworden“.

Nicht minder aber, ja ganz besonders gehört zu der „merklichen“ Aenderung, daß das Volk in der Messe deutsche Gesänge singt. Zwar bemerkt in Verteidigung des 24. Artikels der Augustana die Apologie, daß dies „so neu nicht ist“, denn es „hat doch in allen Kirchen je etwas das Volk deutsch gesungen“. Es bezieht sich hier die Apologie auf die vorreformatorischen Lieder: „Ein Kindelein so löbelich“ — „Christ ist erstanden“ — „Nun bitten wir den Heiligen Geist“, z. B. — „Gott der Vater wohn uns bei“ — „Gott sei gelobet und gebenedeiet“, mit welch letzterem, wie Luther in seiner Schrift „Von der Winkelmesse und Pfaffenweihe“ sagt, zugleich „die Kirche öffentlich Zeter und Mordio über den Ehedienst und räubische Winkelpfaffen geschrieen“, die ihr den Kelch vorenthielten. Aber diese deutschen geistlichen Volksgesänge wurden nur bei Umzügen und Bittgängen oder bei dramatischen Aufführungen in der Kirche an hohen Festen, nicht aber als ein Bestandteil der Messe selbst gesungen. In dieser war allein der liturgische, lateinische gregorianische Gesang berechtigt. Der deutsche geistliche Volksgesang galt nur als ein „geduldeter, außerliturgischer Gesang“. Nur in den Gottesdiensten der böhmischen Brüder sang das Volk in böhmischer Sprache, also in der Landessprache, geistliche Lieder.

Im Gegensatz zu dieser geduldeten Stellung des deutschen Volksgesangs wies ihm daher Luther nur um so mehr eine gliedliche Stellung in der Messe an. „Ich wollte“, schrieb er, „daß wir viele deutsche Gesänge hätten, die das Volk unter der Messe sänge. Denn wer zweifelt daran, daß solche Gesänge, die nur der Chor allein singt oder antwortet, vorzeiten die ganze Kirche gesungen hat?“ Es kehrte aber hiermit Luther nicht nur zurück zum Brauch der alten Kirche, sondern er wies dem Volksgesang in der Messe, wie im ganzen Gottesdienst eine noch hervorragendene Stelle an, als die sie in der alten Kirche behauptete, wie am Tage ist. War doch nach langer, langer Nacht die
\quellseite{51}
Sonne des Evangeliums in einem Glanze wieder aufgegangen, wie er seit der Apostel Zeit nicht mehr vorhanden war! Hatte doch Gott sein Volk noch einmal heimgesucht! Und so hieß es denn wieder: „Singet dem HEerrn ein neues Lied; singet dem HEerrn alle Welt.“ Das eigentliche deutsche Kirchenlied ist erst mit der Reformation entstanden.

Zu dem Ende nahm Luther nicht nur die wenigen bereits vorhandenen deutschen geistlichen Volkslieder auf, jedoch zum Theil „gebessert und christlich corrigirt“ oder durch Hinzudichtung erweitert, und nicht nur verdeutschte er etliche der alten lateinischen Hymnen, sondern er dichtete auch selbst eine Anzahl Lieder und forderte die Befähigten unter seinen Freunden zu Beiträgen auf. So erschien denn schon im Jahre 1524 das erste lutherische Original-Gesangbuch, das „Enchiridion geistlicher Gesänge und Psalmen“, das 25 Lieder, darunter 18 von Luther, enthielt, im Auftrage Luthers von Justus Jonas herausgegeben, wie Wackernagel vermuthet. Etwas später, doch noch in demselben Jahre, kam das Walther’sche Chorgesangbüchlein zu Wittenberg heraus, eine gemeinschaftliche Arbeit Luthers und Walthers, bestimmt für den Gesangunterricht der Jugend, denn durch die sollte ja die Gemeinde im Gottesdienst singen lernen. Im Jahre 1529 erschien dann das Wittenberger Gemeindegesangbuch, in dessen Vorrede aber Luther nicht nur über die große Menge der seit der Herausgabe des Chorgesangbüchleins entstandenen untüchtigen Gesänge, sondern auch über die Fälschung seiner Lieder klagt und sich alle und jede unberufene Vermehrung oder Besserung des Gesangbuchs verbittet. Die letzte zu Luthers Lebzeiten im Jahre 1545 bei A. Bapst in Leipzig erschienene und mit einer neuen Vorrede Luthers versehene Ausgabe ist die innerlich vollendetere, wie äußerlich schön ausgestattete Gesangbuchsausgabe. Dieselbe enthält außer den schon 1542 besonders erschienenen Begräbnißgesängen auch eine Anzahl verdeutschter Collecten mit ihren Versikeln, und zwar stehen diese nicht als Anhang, sondern sind eingestreut, so daß bei den Weihnachtsliedern eine Weihnachtscollecte, bei den Osterliedern eine Ostercollecte etc. steht, was um so mehr den liturgischen Zweck vor Augen legt.

Diesen Liedern gab Luther mit seinen musikalischen Freunden die entsprechenden Melodien. In seiner Schrift „Wider die himmlischen Propheten“ schreibt er: „Ich wollte heute gern eine deutsche Messe haben, ich gehe auch damit um; aber ich wollte gern, daß sie eine rechte deutsche Art hätte. Denn daß man den lateinischen Text verdolmetscht und dazu lateinischen Ton oder Noten behält, lasse ich zwar geschehen; aber es lautet nicht
\quellseite{52}
artig noch rechtschaffen. Es müssen beide: Text und Noten, Accent, Weise und Geberde aus rechter Muttersprache und Stimme kommen, sonst ist alles ein Nachahmen, wie die Affen thun!“ (Erl. Ausg. 29, 203.) Damit nun für das deutsche geistliche Volkslied auch die Melodie so viel wie möglich „aus rechter deutscher Muttersprache und Stimme komme“ und so „eine rechte deutsche Art hätte“ griff Luther nicht nach dem cantus gregorianus, faßte also nicht den Most in alte Schläuche, indem er ihm seine und seiner Freunde deutschen Lieder schlechtweg unterzulegen suchte, sondern er schlug folgenden Weg ein.

Zunächst benützte er Vorhandenes. In der Vorrede zu den im Jahre 1542 herausgegebenen Begräbnißgesängen spricht er sich folgendermaßen aus: „Zudem haben wir auch zum guten Exempel die schönen Musica oder Gesänge, so im Pabstthum in Vigilien, Seelmessen und Begräbniß gebraucht sind, genommen ... doch andere Text darunter gesetzt ... Der Gesang und die Noten sind köstlich, Schade wäre es, daß sie sollten untergehen, aber unchristlich und unrein sind die Text oder Wort, die sollten untergehen. Gleichwie auch in allen andern Stücken thun sie (die Papisten) es uns weit zuvor, haben die schönsten Gottesdienste, schöne herrliche Stifte und Klöster. Aber das Predigen und Lehren, das sie drinnen üben, dienet das mehrer Theil dem Teufel und lästert Gott, denn er ist der Welt Fürst und Gott, darum muß er auch das Niedlichste, Beste und Schönste haben ... Also haben sie auch wahrlich viel treffliche schöne Musica oder Gesang, sonderlich in den Stiften und Pfarren, aber viel unflätiger abgöttischer Zier damit geziert. Darum wir solche abgöttische todte und tolle Zier entkleidet und ihnen die schöne Musica abgestreift und dem lebendigen heiligen Gotteswort angezogen, daselbst damit zu singen, zu loben und zu ehren, daß also solcher schöner Schmuck der Musica in rechtem Brauch ihrem lieben Schöpfer und seinen Schriften diene, daß er gelobt und geehrt, wir aber durch sein heiliges Wort mit süßem Gesang, ins Herz getrieben, gebeßert und gestärkt werden im Glauben.“ (Erl. Ausg. 56, 301. 302.) Das Vorhandene, was er nahm, waren zunächst Singweisen lateinischer Hymnen. Aber indem er den von ihm übersetzten guten Hymnen, wie „Nun komm, der Heiden Heiland“ u. a., ihre Melodie beließ und den abgöttischen, todten und tollen Texten „die schöne Musica abstreifte“, geschah solches doch zugleich durch Nachbildung oder Umbildung, damit auch die Melodie des lateinischen Textes für den deutschen Text eine „deutsche Art“ hätte. Ferner waren es die Melodien der aus der Vorreformatorischen
\quellseite{53}
Zeit von ihm aufgenommenen geistlichen Volkslieder, an denen er jedoch auch hie und da die bessernde Hand anlegte, wie z. B. „Christ ist erstanden“. Und endlich waren es schöne Melodien weltlicher Volkslieder, die nicht nur so manchem ehrbaren, sondern auch so manchem fleischlichen Text abgestreift und mit einem christlichen Text versehen wurden, wie z. B. für das Weihnachtslied: „Vom Himmel hoch, da komm ich her“ anfänglich die Melodie: „Aus fremden Landen komm ich her“, oder für das Lied von Elisabeth Creutziger: „HEerr Christ, der einig Gotts Sohn“ die Melodie: „Ich hört ein Fräulein klagen“ — eine Verwendung weltlicher Melodien auch in der nachfolgenden Zeit im Andenken an das Wort Luther’s: „Wie gehet es doch zu, daß wir in carnalibus so manch fein Poema und so manch schön Carmen haben, und in spiritualibus da haben wir aber so kalt, faul Ding!“

Sodann schuf Luther Originalmelodien, wie es noch feststeht, daß z. B. die Melodien zu „Ein feste Burg“ — „Wir glauben all“ — „Vater unser im Himmelreich“ — „Jesaja dem Propheten“ u. a. Luther’s Schöpfung sind. Hierdurch wurde er insbesondere der Vater unsers köstlichen und so volksthümlichen rhythmischen Kirchengesanges.

Hiernach entstammen unsere Kirchenmelodien einer vierfachen Quelle: 1. durch Nachbildung oder Umbildung dem lateinischen Hymnengesang, 2. dem religiösen, aber außerkirchlichen Volksliede der vorreformatorischen Zeit, 3. dem weltlichen Volksliede und 4. der eigenen Composition.

Ja freilich eine merkliche Aenderung, daß in der Messe und sonst bei uns „das Volk deutsche Gesänge singt“ und dadurch sonderlich „die Meß bei uns mit großer Andacht gehalten wird, denn bei den Widersachern“. Unsere lutherische Kirche heißt nun deswillen die singende Kirche. Auf welch eine Schaar von Psalmenfängern aus allen Ständen, darunter auch manche Mirjam, können wir hinweisen! Wie ist gerade durch das Kirchenlied jener Wunsch Mosis erfüllt: „Wollte Gott, daß alles Volk des HEerrn weissagte!“ Aber der Führer im Reigen dieser großen Sängerschaar nach Wort und Weise ist und bleibt doch Luther. Er war geistlicher Dichter und Musiker von Gottes Gnaden, also daß ihm selbst ein Paul Gerhardt nicht gleichkommt und in Betreff der Melodien kein Krüger und kein Ebeling; denn keiner war im Erfassen des Volkstones nach Wort und Weise so glücklich, keiner so der Mund der Kirche, der im Liede und seiner Melodie wiedergab, was ihre große Seele bewegte. Stehe daher hier jenes schöne Zeugniß von Cyriacus Spangenberg, welcher in der Vorrede zu seiner Cithara Lutheri 1569 S. 2
\quellseite{54}
sagt: „So muß man’s ja wahr sein lassen, daß Lutherus unter allen Meistersängern seit der Apostel Zeit der beste und kunstreichste gewesen, in dessen Liedern und Gesängen man kein vergeblich noch unnöthiges Wörtlein findet. Es fleußt und fället ihm Alles aufs lieblichste und artlichste voller Geists und Lehre, daß schier ein jedes Wort eine eigene Predigt oder doch zum wenigsten eine sonderliche Erinnerung gibt. Da ist nichts Gezwungenes, nichts Eingeflicktes, nichts Verbrochenes. Die Reimen sind leicht und gut, die Wort artlich und auserlesen, die Meinung klar und verständlich, die Melodeien und Ton lieblich und herzlich und in Summa alles herrlich und köstlich, daß es Safft und Krafft hat, herzet und tröstet. Und ist fürwahr seines Gleichen nicht, viel weniger seines Meisters zu finden, wie das alle fromme Herzen, denen anders Lutheri Gesangbüchlein rechtschaffen bekannt, mit uns bekennen müssen und sagen, daß uns Gott durch ihn an seinem Gesangbüchlein etwas Hohes, Wunderbarliches und Sonderliches geschenckt hat, dafür wir ihm in alle Ewigkeit nicht genugsam danken können.“


Wir dürfen jedoch diesen Abschnitt nicht schließen, ohne in Absicht auf den Hauptgottesdienst der lutherischen Kirche auch einen Blick sowohl auf den Chorgesang, als auch auf den Gebrauch der Orgel zu werfen.

a. Daß der Chorgesang in der Messe, wie in den übrigen Gottesdiensten der päbstlichen Kirche in sein sollender Vertretung der Gemeinde sowohl eine wichtige Rolle spielt, als auch mehr oder weniger Kunstgesang ist, ist bekannt. Wenn nun aber gleich die selbstthätige Mitwirkung der Gemeinde namentlich durch das Kirchenlied einen wichtigen Bestandtheil des lutherischen Hauptgottesdienstes bildet, so sollte doch damit keineswegs der Chorgesang durch dieselbe verdrängt werden. Es sollte auch die Kunstmusik in dem Dienst des Heiligen bleiben und zur Vermannigfaltigung und Bereicherung des liturgischen Theils des Gottesdienstes, zur Erhöhung der festtäglichen Feier und der Erbauung überhaupt gebraucht werden.

Zum Beweis, daß Luther noch etwas anderes für den Kirchengesang erstrebt habe, als dessen Beschränkung auf das Gemeindelied, stehe hier aus seinen Schriften noch folgende, auch von E. Naumann in seiner illustrirten Musikgeschichte citirte Stelle. Wir hören Luther sagen: „Wo aber die natürliche Musica durch die Kunst geschärft und polirt wird, da sieht und erkennt man erst die große Weisheit Gottes in seinem wunderbaren Werke der Musica, in welchem vor allem das seltsam und zu verwundern ist, daß einer eine schlichte Weise oder Tenor hersinget, neben welcher drei, vier oder fünf andere Stimmen
\quellseite{55}
auch gesungen werden, die um solche schlichte, einfältige Weise oder Tenor gleich als mit Jauchzen ringsherum spielen und springen und mit mancherlei Art und Klang dieselbige Weise wunderbarlich zieren und schmücken, gleichsam einen himmlischen Tanzreihen führen, freundlich einander begegnen und sich gleich herzen und lieblich umfangen, also daß diejenigen, so solches ein wenig verstehen und dadurch bewegt werden, sich des heftig verwundern müssen und meinen, daß nichts Seltsameres in der Welt sei, denn ein solcher Gesang mit viel Stimmen geschmückt. Die aber dadurch nicht gerühret werden, die sind fürwahr recht ungeschickte Esel.“ (Walch XIV, 411.) Das ihm hierbei vorschwebende Ideal fand er insonderheit wie in den Tonwerken des Niederländers Josquin, so in den Motetten des ihm befreundeten und von ihm hochgeachteten bayerischen Kapellmeisters Senfl verwirklicht, so daß er, als er an einem Decemberabend 1539 Sänger zu Gaste hatte und diese ihm nach Tisch „etliche seine und liebliche Motetten“ Senfls vorsangen, in die warmen Worte ausbrach: „Eine solche Motette vermöcht ich nicht zu machen, wenn ich mich auch zerreißen sollt.“\footnote{In der „Encyclopädie der gesammten musikalischen Wissenschaften“ von Dr. Gustav Schilling ist übrigens folgendes zu lesen: „Es ist bekannt, daß Luther außer manchen herrlichen Melodien, besonders zu seinen eigenen Kirchenliedern, auch verschiedene Motetten geschrieben hat. Eine Sammlung seiner Motetten liegt noch auf der Bibliothek zu München, mit dem Titel: Symphoniae jucundae 4 vocum seu Motettae. 52, cum praefatione Mart. Lutheri (Wittenb. 1535). Händel, der alle noch vorhandenen Werke in der That fleißig studirte, soll einmal geäußert haben, daß er diesem Studium sein Bestes verdanke. Luthers freier Sinn zerbrach zuerst die Fesseln der Canonik, die damals aus den Niederlanden hervorging und selbst in Italien die herrschende Tonweise war. Seine Freunde Walther, Senfl und Agricola standen ihm ehrlich und treulich darin bei. Der erste ästhetische Schimmer, der von ihm ausging, wirkte mächtig auf die Kirche und ging, sich später zur allerwärmenden gestaltend, auf das große weite Reich der Tonkunst über.“}

Wiewohl die Compositionen eines Gallus († 1591), Haßler († 1612) u. a. zeigen, daß in der lutherischen Kirche auch mitunter Männerchöre verwendet wurden, so bediente man sich doch am liebsten und häufigsten des gemischten Chors und ist daher für ihn auch bei weitem am meisten componirt worden. Da er die zwei Hauptgattungen der Stimme in sich vereinigt, auch den Stimmen eine viel freiere Bewegung gestattet, als der Männerchor, so ist der gemischte Chor unstreitig auch die vollkommenste Art von Chor. Man bediente sich aber ehedem für den Discant und Alt der Knabenstimmen durch die Lateinschulen. Erst am Weihnachtsfeste 1715 geschah es zum ersten Male, daß man bei dem gemischten Kirchenchor auch der Frauenstimmen sich bediente, indem der Hamburgische Tonmeister
\quellseite{56}
Mattheson bei seiner Aufführung drei Opernsängerinnen (!) in der Kirche auftreten ließ. Es ist daher auch heutzutage eine Streitfrage, ob der gemischte Kirchenchor nur aus Männern und Knaben bestehen soll, oder ob auch Frauen zugelassen werden dürfen. Vielleicht dürfte Palmer hier das Richtige getroffen haben, indem er in Herzogs Real-Encycl. Art. Gesang S. 106 und 107 sich also äußert: „Das Weib muß schweigen in der Gemeinde, wo es sich um liturgisches und homiletisches Auftreten handelt; aber auf musikalischem Gebiet ist auch dem Weib ein Charisma gegeben, dem seine ungeschmälerte Stelle im Cultus gebührt. — Jedoch muß zugestanden werden, daß, wenn der Chor im Raum der Kirche einen Standort hat, der ihn den Blicken Aller vorzugsweise aussetzt, dann aus anderweitigen Gründen allerdings die oberen Stimmen besser durch Knaben vertreten werden; diesem Gesichtspunkte, als dem ethischen, kann sich der ästhetische oder musikalische um so eher unterordnen, als bei sorgfältiger Auswahl und Pflege die Knabenstimme ihre eigenthümliche Schönheit besitzt und Treffliches leisten kann.“ Da der Raum unserer hiesigen Kirchen den entsprechenden Standpunkt für den Chor meistenteils nicht bietet, auch die Chorübungen gewöhnlich nur zur Nachtzeit vorgenommen werden können, so wäre es von dem ethischen Gesichtspunkte aus nur um so mühsamer, wenn in unsern gemischten Kirchenchören die oberen Stimmen wieder ausschließlich Knaben zugewiesen würden. Sollte es denn nicht möglich sein, namentlich in unsern Stadtschulen, eine Anzahl stimmbegabter Knaben auszulesen und als eigentliche Chorsänger nach und nach zuzurichten? Und da man mit Recht in den Gemeinden die Kinderstimmen so gern hört, dürfte es da wohl so schwer halten, die Gemeinde oder eine Anzahl ihrer Glieder zu bewegen, zur Ermunterung den „Chorschülern“ in irgend einer Form von Zeit zu Zeit eine kleine Belohnung zukommen zu lassen! Es wäre das wieder etwas von dem altkirchlichen Institut der Kurrende, dem einst auch ein Luther angehörte und ihm Lob spendete, das sonderlich in der lutherischen Kirche hochgehalten und zum Segen für Kirche und Schule gepflegt wurde und das erst der Rationalismus und seine Bestrebungen, Kirche und Schule zu emancipiren, mehr und mehr beseitigte. Ein um die Wiederaufnahme des Schülerchors und die Kurrende hochverdienter Pädagoge — Marquardt in Berlin — behauptet mit Recht: „Ein wohlorganisirtes und gepflegtes, der Kirche bei allen ihren Lebensäußerungen zur Seite stehendes Chorwesen ist ohne Benachtheiligung des Schulzweckes nicht nur möglich, sondern ein solches gibt einer ganzen Schule von unten herauf einen höheren Zug, der ihre erziehliche Kraft
\quellseite{57}
potenzirt.“ („Halleluja“, Organ für die geistliche Musik etc. 1884 Nr. 17, S. 195.)\footnote{Auch in dem hochkirchlichen Theil der Episcopalkirche dieses Landes hat man wieder angefangen, den Knabenstimmen im Kirchenchor den Vorzug zu geben, und darin schon Treffliches geleistet.}

Was nun den Zweck und die eigentliche Bedeutung des Chorgesangs im lutherischen Gottesdienste betrifft, so fehlt uns eine allgemein anerkannte Idee, der wir folgen könnten. Es herrschen daher mannigfache Auffassungen desselben. Da die Einführung des deutschen Gemeindegesangs durch die anfängliche Unbeholfenheit der Gemeinden große Schwierigkeiten hatte, so mußte der Chor die Kirchenlieder der Gemeinde allerdings erst vorsingen, auf daß sie durch Anhören und dann Nachsingen dieselben lerne, weshalb denn auch Luther im „Enchiridion“ von 1524, dem ersten lutherischen Gesangbüchlein, bemerkt: „... auf daß einmal der gemeine christliche Haufe mit der Zeit möge lernen verstehen, was man handle unter der Gemeine in Singen und Lesen.“ Allein wie bald sank doch bei der damaligen Begeisterung über den Wiederaufgang des Evangeliums das Volk sich ein, so daß es dann nur noch zu seiner Leitung und Unterstützung des Chors bedurfte. Weil nun aber der Chor durch den geistlichen Kunstgesang im Gottesdienst aus dem Haufen hervortritt, so ist es nach Schöberlein und andern neueren Liturgikern „die Kirche an sich, die ideale Gemeinde, die allgemeine Kirche im Gegensatz zur Localgemeinde, die der Chor vertritt“. Dies wird so begründet: „Was ist naturgemäßer, als daß die Localgemeinde, wenn sie zur Anbetung sich versammelt, sich als Glied der Einen allgemeinen Kirche fühle, auf welche ja ihre Gebete und Gesänge ausdrücklich Bezug nehmen, und daß sie deshalb Verlangen trage, auch die Stimme von dieser in ihre Gesänge hineintönen zu hören? Nun bildet aber offenbar dieselbe innerhalb der Versammlung der wirklichen Gemeinde ein ideales Element. Dies ist die Ursache, warum der Chor sich für seine Gesänge der idealen Formen der Kunst zu bedienen hat.“ Aber schon um dieser Auffassung von allgemeiner Kirche und „idealer“ Gemeinde willen, bleiben wir unsererseits nur um so mehr nach Hommel bei der praktischen Seite der Bedeutung des Chors, indem jener in ihm denjenigen Theil der Gemeinde erblickt, welcher mit seinem Charisma der Kirche dient. Hommel sagt nämlich: „Die praktische Auffassung der Sache ist wohl auch hier die beste. Auf keinen Fall darf er (der Chor) als Vertreter der Gemeinde oder als Vermittlung zwischen Geistlichem und Gemeinde genommen werden. Aber wenn nach den apostolischen Aussprüchen jeder mit seiner Gabe dem HEerrn dienen soll, so ist nicht abzusehen, warum nicht
\quellseite{58}
diejenigen, welche mit der Gabe des Gesanges vor andern begnadigt sind, diese auch in der Gemeinde, in den Gottesdiensten der Gemeinde ausüben sollen. Wie die christliche Gemeinde alles, was sie hat und was sie thut, dem HErrn opfert, so opfert sie ihm die Gaben des Gesanges, so in ihrer Mitte sind, und heiligt diese dadurch, und wie der Gedanke des Opfers in allen Stücken den Fleiß erweckt, mit den besten Gaben und Fähigkeiten dem HErrn zu nahen, so wird auch die Gemeinde im Opfer des Gesanges gern ihre schönsten Kräfte wirken lassen. Dies sind die schönen Gottesdienste des HErrn, schön, lieblich und mannigfaltig auch in der Form, nicht bloß nach dem Geiste, aus dem alles hervorgeht. Eben deshalb wird schon um der Lieblichkeit der Abwechslung willen ein Theil der Gemeinde hie und da aus ihr heraustreten und im öffentlichen Gottesdienste sein Opfer darbringen dürfen, unter lebendiger Theilnahme der Uebrigen im Geiste, so besonders zu größerer Verherrlichung der Feste.“ Gewiß ist dies auch die Anschauung von der Bedeutung des Chorgesangs, wenn der alte Dannhauer in seiner „Katechismus-Milch“ schreibt: „Mit Stimmen singen und also Gott loben, es geschehe nun choraliter oder figuraliter, mit schlechten oder mit Kunst-Stimmen, stehet uns nicht frei, sondern es lieget uns ob: dieweil es Gott befohlen, du sollt meinen Namen nicht mißbrauchen, ergo, sollt ihn recht brauchen, und consequenter auch die Stimmen, Discant, Alt, Tenor und Baß, mich dadurch zu ehren.“ Mit Recht bemerkt jedoch Hommel schließlich: „Nie aber sollte der Chor singen, ohne daß der Gemeinde der Text bekannt ist.“ Eben deshalb ist nöthig, daß der Chordirigent die Sänger an ein deutliches Aussprechen der Worte gewöhne, und nicht nöthig, daß immer wieder etwas Neues gesungen werde. Gute, schon mehrmals gesungene Stücke werden vielmehr vom Chor immer besser vorgetragen und der Gemeinde immer besser bekannt, durch beides aber der Erbauung desto mehr gedient werden.

Ist nun jenes Zweck und dieses, wenigstens nach der praktischen Seite, die Bedeutung des kirchlichen Chorgesangs, so muß beiden auch dessen Stellung in der Liturgie entsprechen.

Um der Noth willen dürfte er zunächst noch immer in seiner Stellung als Führer und Träger des Gemeindegesanges nicht überflüssig sein. Nicht nur ist erst die Zahl der der Gemeinde bekannten Melodien eine verhältnißmäßig kleine und bedarf es somit zur Vermehrung der Hülfe des Chors, und nicht nur muß in so mancher Gemeinde der rhythmische Gesang erst eingeführt
\quellseite{59}
werden und dazu bedarf es wieder der Mithilfe eines Chors, sondern er kann auch noch sehr gute Dienste leisten zur Erhaltung des rhythmischen Gesanges, denn bei der Unlust in geistlichen Dingen läßt man sich auch im rhythmischen Singen nur zu leicht gehen und kommt nur zu leicht wieder in den alten Schlendrian des geistlichen Singens, während man im Singen weltlicher Volkslieder allewege recht gut auf Takt und Rhythmus Acht hat.

Da jedoch nach der angegebenen Bedeutung der Chorgesang auch eine selbständige Stellung im Gottesdienste einzunehmen hat, so darf dieselbe aber gleichwohl keine isolirte Stellung sein, soll der Chor recht zur Erbauung der Gemeinde dienen und das Seinige zur wahren Verschönerung des Gottesdienstes beitragen. Die Gesänge, welche er vorträgt, müssen demnach in möglichstem Zusammenhang mit dem Handeln des Liturgen und der Gemeinde stehen, wenn nicht der Chorgesang mehr eine bloße Aufführung sein und den Charakter eines geistlichen Concerts annehmen soll. Und den soll er nicht haben. Nachdrücklich heben das auch gegenwärtig Sachverständige hervor. So berichtet z. B. der Musikdirector Otto Kade, der Herausgeber des „Luther-Codex“, über den von ihm gegenwärtig mit Meisterschaft und großem, auch liturgischem Verständnis geleiteten großherzoglichen Schloßchor in Schwerin: „Die rein äußerliche, fast nur auf den Vortrag geistlicher Concertstücke angewiesene Stellung, die mit der Liturgie selbst in gar keinem oder nur sehr losem Bezuge steht, ist in eine organisch mit dem Gange der geistlichen Handlung eng verknüpfte, wesentlich in dieselbe eingreifende verwandelt: die innigste Verbindung zwischen Liturg, Gemeinde und Chor erreicht.“ („Halleluja“, Organ für geistliche Musik etc. Jahrg. V, Nr. 8, S. 102.) So auch Dr. G. Herzog (Professor der Musik in Erlangen): „Seinen wahren Zweck wird der Chorgesang etc. nur dann erfüllen, wenn die zu wählenden Gesänge in innerem Zusammenhang mit dem Gottesdienste stehen, also organisch eingegliedert sind, echt kirchlich-liturgischen Charakter an sich tragen und allgemein verständlich sind.“ (Musik. Anhang zur bayerischen Agenda, Vorwort VII.) Dem sei auch folgendes beherzigenswerte Wort des Altdorfer Schullehrerseminar-Directors und gründlichen Kenners heiliger Musik, Johannes Zahn, beigefügt: „Der Chorgesang muß bezüglich seiner Ausführung für die Gemeinde ein Muster sein können; darum hat der Dirigent allen Fleiß anzuwenden, daß er die Mitglieder desselben befähigt, die vorzutragenden Chöre in möglichst vollkommener Weise auszuführen, soll darum vor allem nur solche Compositionen wählen, welche die Leistungsfähigkeit des Chors nicht übersteigen, und soll es für
\quellseite{60}
rühmlicher halten, eine leichte Composition gut auszuführen, als eine schwierige mangelhaft.“ (Handbüchlein für evangelische Cantoren und Organisten.)

Organisch eingegliedert dürfte im Hauptgottesdienst der Chor mit dem Liturgen und der Gemeinde etwa auf folgende Weise zusammenwirken:

1. Er übernimmt solche Theile der Liturgie, deren musikalische Ausführung der Gemeinde mehr oder weniger Schwierigkeit bieten dürfte. Dahin gehörten die Introiten im Psalmenton, in ihren ursprünglichen Weisen oder in der Figuralform; desgleichen die schönen Festkyries in der Liedform.

2. Wird an hohen Festen beim Gloria anstatt des „Allein Gott in der Höh sei Ehr“ die ursprüngliche Form gebraucht, so könnte dies in der Weise geschehen, daß, nachdem der Liturg intonirt: „Ehre sei Gott in der Höhe!“ die Gemeinde mit dem Chor gesungen hat: „Und auf Erden Fried und den Menschen ein Wohlgefallen“, der Chor alsdann weiter singt: „Wir loben dich, wir benedeien dich, wir beten dich an, wir preisen dich“ etc. und könnte die Gemeinde an dem „Amen“ sich betheiligen. Im II. Theil dieses Werkes ist das Gloria mit dem Et in terra und dem Laudamus te zu finden.

3. Bei antiphonischen Gesängen intonirt der Chor und die Gemeinde respondirt, wie z. B. bei dem gleichfalls im II. Theil mitgetheilten Doppel-Kyrie, nämlich: Chor: „Kyrie Eleison.“ Gemeinde: „HEerr, erbarme dich.“ Chor: „Kyrie Eleison.“ Gemeinde: „Christe, erbarme dich.“ Chor: „Kyrie Eleison.“ Gemeinde: „HEerr, erbarm dich über uns.“ Vornehmlich aber bei dem Te deum. Der Chor intonirte vierstimmig und ohne Orgelbegleitung: „HEerr Gott, dich loben wir“, die Gemeinde respondirte mit Orgelbegleitung: „HEerr Gott, wir danken dir.“ Der Chor: „Dich, Vater in Ewigkeit“; die Gemeinde: „Ehret die Welt weit und breit“ etc. (Will man jedoch auch den Chor mit der Orgel begleiten, so sind für denselben etwas schwächere Register zu nehmen.)

4. Zur lieblichen Abwechselung mit dem Gemeindegesang (Ps. 147, 7.) trägt nach Gelegenheit der Zeit und des Festes der Chor einen besonderen Kunstgesang an passender Stelle vor, worüber im II. Theile das Nähere. Ebenso könnte bei längeren Communionen auch dadurch eine liebliche Abwechselung zwischen Chor- und Gemeindegesang stattfinden, daß zwischen zwei Communionliedern, oder zwischen zwei Theilen eines längeren Communionliedes der Chor einen Lobgesang oder sonst zur Communion Passendes vorträgt. Das beförderte die Erbauung weit mehr, als wenn der Organist aus Furcht, es möchte das Communionlied während der Austheilung zu bald zu Ende sein, in
\quellseite{61}
Zwischenspielen von Vers zu Vers sich ergeht, die viel länger dauern, als ein Vers des Liedes und auf deren Ende die Gemeinde mit steigender Ungeduld wartet.

Wir haben bei diesen Andeutungen über die Stellung und Verwendung des Chorgesangs das kirchliche Vorbild vor Augen gehabt, wie wir es nicht nur in den rechtgläubigen Agenden erblickten, sondern auch in jenen gleichzeitigen Sammlungen, welche das Material für den Chorgesang enthielten, wie Johann Spangenbergs Kirchengesänge, 1545, die Psalmodia sacra von Lucas Lossius, 1553, die Wittenberger Kirchengesänge von Johann Reuchenthal, 1573, die Missodia etc. von Mich. Prätorius, 1611, u. A.

Zur Geschichte des Chorgesangs noch dieses. Mit besonderer Vorliebe setzte man in der nachreformatorischen Zeit für den besonderen Kunstgesang den Lutherischen Choral in motettenhafter Behandlung, worin der kurfürstliche Capellmeister Johann Eccard, † 1611, ein Schüler von Orlando di Lasso, das Größte leistete. Später wurden dann auch andere Texte für den Chor gewählt, zuerst Schriftsprüche und dann auch selbst eigene Dichtungen. Allmählich jedoch verließ man in den Formen den kirchlich-lutherischen Typus. Dies geschah durch Ariensolos, durch weitere Arten des Kunstgesangs im dramatischen Wechsel, durch die Instrumentalmusik, welche den Chorgesang begleiten sollte, aber immer selbständiger neben dem ersteren auftrat, zuletzt gab man sich immer rückhaltsloser den Einflüssen der in imponirender Weise in der Oper und in der Kammermusik sich entfaltenden weltlichen Musik hin, wenn schon ein Sebastian Bach durch seine ernste Pflege des Chorals und den heiligen Charakter seiner Cantaten und Passionen mit aller Macht die Verweltlichung der Kirchenmusik aufzuhalten suchte. Und selbst auch er, der bis jetzt in seinem Gebiet noch unübertroffene Meister, repräsentirt doch nicht mehr die volle Strenge des Kirchenstils.

Noch läßt uns unser hiesiger Chorgesang viel, viel zu wünschen übrig, wie sehr er auch an manchen Orten gepflegt wird, denn noch sind seine Vorträge mehr concertmäßige Aufführungen, und zwar Aufführungen vorherrschend moderner Erzeugnisse. Möchte man daher vor allem über Zweck, Bedeutung und Stellung des Chors sich klar zu werden suchen — dann fort mit den meisten Sammlungen moderner Chorliteratur und zurück nicht nur bis zu Bach, sondern bis zur classischen Periode der Kirchenmusik! „Die Kirche“, sagt treffend Dr. Schöberlein, „hat ihren eigenen Styl wie im Liede, so im Gesange. Kennt sie auch die Gegensätze heiliger Trauer und Freude, ja, des Jubels im höheren Chor, so sind doch bei ihr diese Gegensätze in den Schranken eines solchen Maßes gehalten,
\quellseite{62}
und über allen noch so lebendigen und tiefen Empfindungen der Buße und des Lobes waltet die selige Ruhe des Friedens Gottes. Auch die Harmonien sind frei von weichlicher Gefühligkeit und aufregenden Uebergängen; klar und rein, ruhig und erhaben schreitet ihr Gang fort und versetzt hiermit die Hörer aus dem Kreise subjectiver, weltlicher Empfindungen in das Reich des Heiligen. Die classische Periode für diese kirchliche Weise des Gesanges ist die des 16. bis in die Hälfte des 17. Jahrhunderts, indem sich eben in dieser Zeit Glaube, heiliger Gemeinsinn und strenge künstlerische Form verbanden und durchdrangen. Aus dieser Periode werden deshalb die Kirchengesänge vorzugsweise zu entnehmen sein. Und ob auch der Styl jener Zeit unserm musikalischen Gefühl und Bewußtsein zunächst ferner stehe, so ist dies doch kein anderes Verhältnis, als mit den Kirchenliedern jener Zeit, welche gleichfalls so viel mehr als die modernen Lieder das wahre Glaubensbedürfniß der Gemeinde befriedigen.“ Es geht einem hier wie mit den Melodien der reformatorischen Kirchenlieder oder auch den Schriften Luthers: anfangs wollen sie nicht munden und dann schmecken sie je länger, je besser. Möchten daher die Leiter unserer Kirchenchöre den Sinn für das Echtkirchliche zunächst durch ein eingehendes Studium des älteren Kirchengesangs zu bilden suchen. Wir machen hier auf das dreibändige Werk Dr. Schöberleins aufmerksam: „Schatz des liturgischen Chor- und Gemeindegesangs“ oder auf die „Musica sacra“, einem Auszug aus demselben, da ersteres für die Meisten doch zu kostspielig sein dürfte.\footnote{Ein Meisterwerk und eine reiche Schatzkammer muß das auf Anordnung des sel. Großherzogs von Mecklenburg-Schwerin herausgekommene „Cantionale für die Landeskirchen des Großherzogthums Mecklenburg-Schwerin“ sein. Die Zeitschrift „Halleluja“ berichtet folgendes: „Dasselbe sollte das ganze liturgische Material für die Sonn- und Feiertage wie für die Neben- und Wochengottesdienste, als: Mette, Vesper, Taufe, Hochzeit, Begräbniß etc. enthalten. Se. Königl. Hoheit der Großherzog setzte daher im Jahre 1862 eine Commission von sechs Fachmännern nieder, nämlich: Herrn Oberkirchenrath Dr. Kliefoth als Vorsitzenden, Kirchenrath Dr. Maßmann, Pastor in Wismar, Kirchenrath Dr. Wöhler, Pastor in Lichtenhagen, Musikdirector Pitschner am Seminar zu Neukloster, Musikdirector Dr. von Roda an der Universität Rostock ... und Musikdirector Kade, Dirigenten des Großherzoglichen Schloßchors in Schwerin. Letzterem wurde auch die Zusammenstellung des musikalischen Apparates und die künstlerische Ausarbeitung der mehrstimmigen Tonsätze übertragen.“ Da dieses köstliche Cantional bis noch vor ein paar Jahren im Buchhandel nicht zu haben war, so war es dem Herausgeber der „Siona“, Pfarrer Herold in Schwabach, gelungen, die Erlaubniß zu einem Auszug aus demselben zu erwirken. Derselbe ist bei C. Bertelsmann in Gütersloh unter dem Titel erschienen: „Der Hauptgottesdienst des Osterfestes und der österlichen Zeit. Liturgisch-musikalischer Auszug aus dem neuen Mecklenburgischen Cantionale. Herausgegeben von M. Herold.“} Ein anerkennenswerther Schritt in der Rückkehr zur classischen Kirchenmusik der
\quellseite{63}
lutherischen Kirche ist in dieser Beziehung auch durch das Chor- und Choralbuch von Endlich\footnote{Choralbuch mit Liturgie und Chorgesängen zum Kirchenbuch der Allgemeinen Kirchenversammlung. Bearbeitet von J. Endlich. Philadelphia, J. Kohler.“} geschehen, das den musikalischen Anhang zu dem „Kirchenbuch“ des General Council bildet. Auch innerhalb der Missourisynode ist in musikalischer Beziehung zur Rückkehr Mahnung und Anbahnung geschehen. Es sei hier u. a. auf die von Lehrer H. Ilse 1884 herausgegebene kleine Sammlung alter und neuer Weihnachtsgesänge\footnote{XIV alte und neue Weihnachtsgesänge für gemischte Kirchenchöre, bearbeitet und herausgegeben von Herm. Ilse. Brooklyn, 1884.} hingewiesen. Der Verfasser dieses liturgischen Werkes kann nicht umhin, aus der Vorrede zu dieser Sammlung folgende Worte zu den seinigen zu machen: „Freilich werden Viele, die sich ausschließlich mit moderner Musik beschäftigt haben, Zeit und Mühe brauchen, sich an den in diesen älteren Sätzen dargebotenen Stil zu gewöhnen. Man wird es zu diesem Ende von gutem Erfolg finden, wenn man sich vorher am Klavier so recht in die einzelnen Stimmen hineinsingt und dann erst das Ganze in seinem Zusammenhange hört. Während so manches Neuere gleich von allen Zuhörern verstanden und fast von keinem Sängerkreise ganz verdorben werden kann, werden die meisten älteren Sachen zu dem, was der Sänger aus ihnen zu machen weiß. Ungeachtet der großen Einfachheit, welche denselben eignet, wollen sie doch in großer Reinheit gesungen sein, und setzen, um wahrhaft verstanden und gewürdigt zu werden, ein tieferes Einleben in den Geist dieser Werke voraus, was immer den Ernst eines christlichen Gemüthes erfordern wird. Der Ausrede so Mancher, daß das Volk solche Musik nicht mehr verstehe, kann wenig Gewicht beigelegt werden. Es ist thatsächliche Erfahrung, daß die Gemeinden allmählig zu etwas Bessern herangezogen werden können, ja, daß auf der andern Seite, wenn das Volk den gesunden Geschmack verliert, niemand die Schuld trägt, als die Musiker selbst.“

b. Was nun die Orgel und deren Gebrauch im Gottesdienste betrifft, so fand Luther bekanntlich auch diesen vor, und es müßte befremden, wenn er sie aus dem Heiligthum verwiesen hätte, da er ja die Musik, wie die Kunst überhaupt, in den Dienst des Höchsten gestellt wissen wollte. Eignet sich doch kein Instrument so sehr, so ausschließlich zu gottesdienstlichem Gebrauche, wie die Orgel, diese Königin unter den musikalischen Instrumenten! Während man daher reformirterseits im bilderstürmerischen Eifer bekanntlich auch die Orgeln zerschlug, ja, während man sogar auf dem Tridentinischen Concil in Folge des
\quellseite{64}
eingerissenen Mißbrauchs der Orgel zu profanem und die Andacht störendem Spiele papistischerseits ernstlich auf die Abschaffung der Orgel dachte, blieb sie in der rechtgläubigen Kirche nicht nur im unbestrittenen Gebrauch, sondern sie erfuhr auch innerhalb derselben ihre großartige Ausbildung und fand in einem Sebastian Bach ihren bis jetzt noch unübertroffenen und kaum nunmehr erreichbaren Meister im kirchlichen Stil und Spiel. Wir finden daher, daß auch schon ältere Kirchenordnungen Vorschriften für ein würdiges Orgelspiel geben.

Im Gottesdienst der rechtgläubigen Kirche wird der Gemeindegesang von der Orgel begleitet und geleitet. Darin besteht wohl von Anfang an ihr wichtigster Gebrauch. Hinsichtlich der Harmonie aber halte man sich an den älteren Tonsatz, und nicht an den von Meister Bach, da dieser wegen seines Reichthums und seiner selbständigen Stimmenführung sich zwar zum kirchlichen Kunstgesang trefflich eignet und nach dieser Seite hin Unübertroffenes bietet, nicht aber für den Volksgesang. Ueber die Begleitung des liturgischen Gesangs durch die Orgel wird an seinem Orte das Nöthige gesagt werden. In Betreff der Frage aber, ob auch der Chorgesang von der Orgel begleitet werden solle, bemerkt Dr. J. G. Herzog:

\begin{quote}
„Ob der mit Orgel begleitete, oder der a capella-Gesang (Gesang ohne Instrumentalbegleitung) im evangelischen Cultus vorzuziehen sei, kann doch nur dahin beantwortet werden, daß beide Arten je nach Form des Gottesdienstes und Charakters der Tonstücke, sowie nach Beschaffenheit der ausführenden Kräfte, zulässig sind, wenn wir nicht auf ein gut Theil des vorhandenen Schatzes kirchlichen Gesanges verzichten wollen. Nicht die Begleitung der Orgel als einer Vermittlerin zwischen Chor und Gemeinde sanctionirt die Berechtigung desselben, wie Manche annehmen, sondern doch jedenfalls die innere Beziehung, in welcher er überhaupt zum Gottesdienst steht. Ich kann nicht verschweigen, daß ich im a capella-Gesang ein besonders geeignetes Mittel zur Erweckung religiöser \unsicher{Andacht} erblicke. In seiner reinen, unmittelbaren Wirkung auf das Gemüth, in dem klaren Heraustreten und Ineinandergreifen der verschiedenen Stimmencharaktere bildet er zugleich einen höchst abwechselungsreichen Gegensatz zum begleiteten Gemeindegesang. Nichts berührt den Menschen so innig, als ein rein vorgetragener mehrstimmiger Gesang ohne alle Begleitung. Wer einmal der Aufführung der Bach’schen Passionsmusik beigewohnt hat, wird die wunderbare Wirkung des Chorals: „Wenn ich einmal soll scheiden’ ohne Begleitung in seinem Leben nicht vergessen. In liturgischen Gottesdiensten nimmt er eine besonders bevorzugte Stellung ein.“ („Halleluja“, Organ für die geistliche Musik etc. Jahrg. 4, Nr. 9, S. 104.)
\end{quote}

\quellseite{65}
Es tritt die Orgel im Gottesdienst jedoch auch selbständig auf, indem sie durch ihr Vorspiel und Nachspiel ihn musikalisch einleitet und abschließt, durch ihr Zwischenspiel aber Cultusacte musikalisch verbindet. Und hier muß ein Jeder, der Urtheil und kirchlichen Geschmack hat, mit Dr. Herzog gewiß wünschen, „daß unser kirchliches Orgelspiel nach seiner selbständigen Seite, im Vor- und Nachspiel, wieder mehr und mehr Bachischen Geist in sich aufnehme, und jene beliebte weltlich süße Art, wie sich solche in so vielen neueren Orgelcompositionen, namentlich in den jetzt zur Mode gewordenen Orgelsonaten zu erkennen gibt, von sich ausscheide“.

Das einleitende Vorspiel hat dem Charakter des Sonn- oder Festtages Rechnung zu tragen. Wie die Glocken gleichsam rufen: „Dienet dem HEerrn mit Freuden, kommt vor sein Angesicht mit Frohlocken“ (Ps. 100, 2.), so soll dieses Vorspiel die nun sich sammelnde Gemeinde in jene Stimmung versetzen, da es heißt: „Kommet herzu, laßt uns dem HEerrn frohlocken und jauchzen dem Hort unsers Heils. Laßt uns mit Danken vor sein Angesicht kommen und mit Psalmen ihm jauchzen; denn der HEerr ist ein großer Gott und ein großer König über alle Götter. Kommt, laßt uns anbeten und knieen und niederfallen vor dem HEerrn, der uns gemacht hat. Denn er ist unser Gott und wir das Volk seiner Weide und Schafe seiner Hand.“ (Ps. 95, 1–3. 6. 7.) Ueber dies den Gottesdienst einleitende Vorspiel, das eigentliche Präludium, jedoch des Näheren im II. Theil bei Besprechung des Introitus.

Und wie das Vorspiel, so das Nachspiel, das doch ja nicht als ein bloß geduldetes Anhängsel oder eine Art Concert von dem Organisten behandelt und von der Gemeinde betrachtet werden soll, sondern dessen Bestimmung ist, als ein Nachklang der empfangenen Eindrücke und Empfindungen den Gottesdienst in würdiger Weise abzuschließen.

Das Zwischenspiel verbindet musikalisch die Cultusacte. Zu diesem aber rechnen wir nicht nur das zwischen den einzelnen Lied-Versen beim Gemeindegesang übliche Zwischenspiel, sowie das vom Gemeindegesang zum liturgischen Gesang überleitende, dem Liturgen den Ton angebende Spiel, sondern auch das Vorspiel zu den Zwischenliedern, zu den Liedern des Hauptgottesdienstes überhaupt und so auch zu dem sogenannten Hauptliede. Aber welcher liturgische Unverstand zeigt sich da oft hin und wieder — nichts zu sagen von jenem Leichtsinn, in dem mancher Organist das ihm zur Erbauung der Gemeinde anvertraute Amt ausrichtet. Wie lang sind da oft diese Zwischenspiele, selbst die von Vers zu Vers, sogar das vom Kirchenliede zum Schlußvers! Wie langweilt
\quellseite{66}
sich die Gemeinde, und harrt verdrießlich des Anfangs zum Singen, und wenn der Organist auch kein Stümper, sondern ein Meister ist, wie wird durch allerlei langes Zwischenspiel die liturgische Kette gleichsam zerrissen und wie verliert die ganze Liturgie, in der doch alles Schlag auf Schlag gehen soll, an Frische! Und wie verlängert es auch den Gottesdienst unnöthigerweise! Meint man an der Liturgie kürzen zu müssen, so sehe man nach, ob nicht vor allem eine Kürzung des Orgelspiels vorgenommen werden könnte oder müßte. Wie daher das gesammte Zwischenspiel würdig und entsprechend, so sei es auch kurz — kurz das sogenannte Vorspiel zum Liede, selbst zum Hauptliede, noch kürzer das Spiel von Vers zu Vers, am kürzesten aber das Spiel zur Intonation der Liturgie, vom Segen zum Schlußvers und, wo keine Intonation des Credo stattfindet, auch zu diesem.

\chapter[§ 7. Stellung des Worts und Sacraments]{§ 7. Stellung des Worts und Sacraments im genuin-lutherischen Hauptgottesdienst. Gliederung und Gang desselben. Christlich freie Abweichung rechtgläubiger Kirchen von der normalen Weise}


In seinem „Sermon von dem Neuen Testament, das ist, von der heiligen Messe, 1520“, sagt Luther: „In diesen Sprüchen (nämlich Luc. 22, 19. 1 Cor. 11, 26. Ps. 102, 22. 23. Ps. 111, 4. 5.) siehest du, wie die Meß eingesetzt ist, Christum zu predigen und zu loben, sein Leiden und alle seine Gnade und Wohlthat zu preisen, damit wir ihn zu lieben, hoffen, glauben bewegt, und also auf dieselben Wort oder Predigt auch ein leiblich Zeichen, das ist, das Sacrament empfahen, auf daß damit unser Glaube mit göttlichen Worten und Zeichen versorgt, befestigt, stark werde wider alle Sünd, Leiden, Tod und Hölle und Alles, was wider uns ist. Und wo die Predigt nicht hätt sollt sein, hätt er die Meß (Abendmahl) nimmermehr eingesetzt. Es ist ihm mehr am Wort denn am Zeichen gelegen.“ (Erl. Ausg. 27, 167 ff.) Die Messe soll sein „der Brauch und Nutz des Evangelii und Austheilung des Sacramentes“ (Weise, christl. Meß zu halten); es soll „alles um des Worts und Sacramenten willen unter den Christen geschehen im Gottesdienst“ (Deutsche Messe etc.). Somit ist die leitende Idee des Hauptgottesdienstes auch in der Gestalt, wie wir sie nun durch Luther aus dem christlichen Alterthum überkommen haben, das Wort der Gnade in Predigt und Sacrament und kann derselbe in seiner Zusammensetzung doch wohl kein zufälliges Conglomerat (Gehäufe) sein. Er muß ein geistlich-psychologisch geordnetes und gegliedertes Ganze sein.
\quellseite{67}
Bestimmen Wort und Sacrament die Ordnung des Heils, sollte dann diese nicht wieder dem sinnigen Alterthum die Gestaltung des Gottesdienstes bestimmt haben? Wenn z. B. der Anfang gerade mit dem Kyrie gemacht wird und diesem sofort das Gloria folgt, legt das nicht die Vermuthung nahe, daß in der Aufeinanderfolge dieser und der andern Bestandtheile ein Gedankengang aufwärts eingehalten sei? Im Anschluß an die von Löhe im Vorwort zur ersten Ausgabe seiner Agende gegebene Deutung hat Schreiber dieses bereits 1862 in einem „Referat über die rechte Mitte der lutherischen Liturgie“ versucht, diesen Gedankengang kurz darzulegen (siehe „Lehre und Wehre“ Jahrgang VIII, S. 205, 206), dabei jedoch betonend, daß im Abendmahl der Hauptgottesdienst nur insofern sein Ziel und seinen Schlußstein hat, insofern das Sacrament das Siegel des Worts ist. Da ihm seit dem keine andere geworden ist, so stehe auch hier die versuchte Deutung.

Nachdem der Introitus den Gottesdienst eingeleitet, den Charakter des Sonn- und Festtags bezeichnet hat, demüthigt sich zuvörderst die versammelte Gemeinde vor Gott im Kyrie, als in einem gemeinsamen Bekenntniß des ganzen Jammers, den Adams Fall uns gebracht hat, steigt aber in Erinnerung der in der heiligen Weihnacht geschehenen Geburt des verheißenen Heilandes, der Fleischwerdung des ewigen Wortes, auf zum Gloria der Engel („Ehre sei Gott in der Höhe“ etc. und „Allein Gott in der Höh’ sei Ehr’“) und bricht aus in einem erhabenen Preis des Dreieinigen („Wir loben dich, wir benedeien dich, wir beten dich an“ etc., oder „Wir loben, preisen, anbeten dich“ etc.). Vom Lobgesang schreitet die Gemeinde in der Collecte zur Bitte, in welcher das Verlangen je nach dem Charakter des Sonn- oder Festtags gleichsam in einem hellen Gedanken aufflammt. Und nun kommt zu ihr der HEerr im Worte. In der Epistel, deren Inhalt vornehmlich Lehre und Vermahnung ist, redet er durch apostolischen Mund, im Evangelio des Tages, das vornehmlich Verkündigung der betreffenden Heilsthatsache oder eines Wunderwerks des HEerrn ist, vernimmt man zugleich auch meist ipsissima verba. Auf erstere antwortet die Gemeinde mit einem Lobe des Höchsten im Hauptlied, mancherorts mit einem Halleluja vor demselben, auf letzteres (oder auf beides) aber mit ihrem freudigen und tapferen Glaubensbekenntniß, dem Credo („Wir glauben all an einen Gott“ oder auch an Festen das nicänische, apostolische Symbolum). Darauf sitzt sie mit Maria nieder zu des HEerrn Füßen und ruht, um zur Förderung ihres inwendigen Lebens die dem Glauben ähnliche Weissagung in der Predigt zu vernehmen, an die sich dann die öffentliche Kirchenbeichte
\quellseite{68}
und Absolution, wo dieselbe auf die Predigt folgt, als eine Uebung und Application des gehörten Gesetzes und Evangeliums unmittelbar anschließt. Reich gemacht durchs Wort und erfüllt mit Trost und Freudigkeit läßt sie im allgemeinen Kirchengebet, den speciellen Fürbitten und Danksagungen und dem heiligen Vaterunser ihr gemeinsames Anliegen in Bitte, Gebet und Fürbitte vor Gott kund werden, welchem dann auch hernach in der Präfation die Danksagung folgt.

So vorbereitet auf die Besiegelung der durch das Wort empfangenen Gnade und die Vereinigung mit dem HEerrn auch im Sacramente seines Leibes und Blutes naht man zum Altare, nachdem zuvor noch im evangelischen Opfertorium durch die Psalmodie: „Schaffe in mir, Gott, ein reines Herz“ aus Psalm 51 das Opfer eines bußfertigen Herzens und zugleich auch durch das Einsammeln der Collecte, wo dieselbe an dieser Stelle üblich ist, auch das Dankopfer der Liebe in freiwilliger Spende für den Gotteskasten, für Mission, für Lehranstalten u. dgl. dargebracht worden ist. Die nun beginnende Präfation ist ihrem Namen nach Einleitung zur Consecration, ihrem Inhalt nach aber die schon erwähnte Danksagung in erhabenster Form (Eucharistie) und wird dann im Sanctus, dem Dreimalheilig mit dem Benedictus, zur Anbetung des „der da kommt im Namen des HEerrn“ mit seinem Leib und Blut im hochwürdigen Sacramente. Im Vaterunser — da und dort mit vorausgehender Abendmahlsvermahnung — und den Testamentsworten vollzieht sich die Consecration, und nachdem im Agnus Dei („Christe, du Lamm Gottes“) die Gemeinde des HEerrn Tod einmütig noch besonders verkündigt, und im Friedensgruß, wo derselbe noch üblich, nochmals die Versicherung der Gnade Gottes und seines Wohlgefallens empfangen hat, erfolgt unter dem Lobgesang aller Versammelten, dem Communionlied, und mit dem Bekenntniß der wahrhaftigen Gegenwart des Leibes und Blutes Christi durch den Pastor (Distributionsformel) die Austheilung und Genießung des Sacraments. Trefflich schließt sich in manchen Liturgien an die geschehene Austheilung der Gesang Simeons, das Nunc dimittis: „HEerr, nun lässest du deinen Diener in Friede fahren“ oder: „Mit Fried und Freud ich fahr dahin“ an — denn was kann man im Jammerthal von Kräften der zukünftigen Welt mehr verlangen, als den Genuß des Leibes und Blutes Christi zur Vergebung der Sünden, zum Leben, zur Seligkeit, und wann könnte wohl die Sehnsucht, den HEerrn von Angesicht zu schauen und daheim bei ihm zu sein, größer sein, als da, wo er sich so nahe zu uns thut? Den Schluß des
\quellseite{69}
nun hierdurch zu seinem Höhepunkt gekommenen Gottesdienstes bildet die Danksagungscollecte, die zugleich Bitte um den Segen und die Frucht des Sacraments ist, und der Aaronische Segen „als Bekräftigung und Besiegelung aller im Wort und Sacrament empfangenen Gnadengaben durch eine letzte Gnadenmittheilung“ von dem Dreieinigen, der Eines gibt im Mannigfaltigen und in dem Einen Mannigfaltiges.

Wie man nun aber auch den Gang der Liturgie im Hauptgottesdienst sich deuten mag — so steht doch so viel fest, daß eine so schöne und liebliche, eine dem reinen Wort und Sacrament so entsprechende und dienende Form und Weise es wohl nicht gibt, wie die durch Luther gereinigte Messe. Da sie als aus dem christlichen Alterthum stammend eine so „feine Abkunft“ hat, so gilt auch von ihr, was Dr. Guerike in seiner Kirchengeschichte von der ausführlichen Liturgie des Hauptgottesdienstes der alten Kirche sagt, nämlich, daß „deren Zweck es war, die göttliche und gottmenschliche Lebensgemeinschaft der Gläubigen mit Christus und unter einander, vermöge der gemeinsamen reellen Theilnahme an Christi heiligem Leibe und Blute dem Bewußtsein der Abendmahlsgenossen recht anschaulich vorzuführen“ (5. Aufl. I. S. 313. 314). Je mehr man diese Form und Weise übt und noch mehr, je näher man sie kennen lernt, je mehr stimmt man Luthern bei, wenn er in seiner „Schrift von der Winkelmesse und Pfaffenweihe, 1533“ begeistert ausruft: „Gott sei gelobt, in unsern Kirchen können wir einem Christen eine recht christliche Messe zeigen nach Ordnung und Einsetzung Christi, auch nach der rechten Meinung Christi und der Kirche! Da tritt vor den Altar unser Pfarrherr, Bischof oder Diener im Pfarramt, recht und redlich und öffentlich berufen, zuvor aber in der Taufe geweihet, gesalbet und geboren zum Priester Christi, ungeachtet des Winkelchreiens. Der singet öffentlich und deutlich die Ordnung Christi, im Abendmahl eingesetzt, nimmt das Brod und Wein, dankt, theilet aus und gibt’s in Kraft der Worte Christi: Das ist mein Leib, das ist mein Blut; solches thut etc. uns Andern, die wir da sind und empfangen wollen. Und wir, sonderlich so das Sacrament nehmen wollen, knieen neben, hinter und um ihn her, Mann, Weib, Jung, Alt, Herr, Knecht, Frau, Magd, Eltern, Kinder, wie uns Gott allda zusammenbringt, allesammt rechte heilige Mitpriester, durch Christi Blut geheiliget und durch den Heiligen Geist gesalbet und geweihet in der Taufe. Und in solcher unserer angebornen, erblichen priesterlichen Ehre und Schmuck sind wir da, haben (wie Apokalypse am 4. Cap. abgebildet ist)
\quellseite{70}
unsere güldenen Kronen auf den Häuptern, Harfen in der Hand und gülden Rauchfässer, und lassen unseren Pfarrherrn nicht für sich, als für seine Person, die Ordnung Christi sprechen, sondern er ist unser aller Mund und wir Alle sprechen sie mit ihm von Herzen, und mit aufgerichtetem Glauben zu dem Lamm Gottes, das da für und bei uns ist und seiner Ordnung nach uns speiset mit seinem Leibe und Blut. Das ist unsere Messe und die rechte Messe, die uns nicht fehlet. — Denn hie geht’s erstlich Alles nach der Ordnung und Befehl Christi, daß es der Kirchen auch in beider Gestalt gereicht und gegeben wird durch die Worte Christi: Nehmet hin und esset, das ist mein Leib etc. Solches thut zu meinem Gedächtnis. Der Pfarrherr empfängt’s nicht allein für sich, wie der Pabstgreuel thut, er opfert’s auch nicht Gott für unsere Sünde und allerlei Noth, wie der Pabstgreuel thut, er theilt’s uns nicht mit, verkauft’s uns nicht für ein gut Werk, Gott zu versöhnen, wie der Pabstgreuel thut und solchen lästerlichen Jahrmarkt daraus gemacht hat: sondern reicht’s uns zu Trost und Stärkung unseres Glaubens. Hier verkündigt und predigt man von Christo. Sie kann kein Geiz, noch Abgötterei sein. Sie haben wir die Meinung Christi und der Kirchen gewiß. Sie dürfen wir nicht sorgen, ob der Pfarrherr die Wort heimlich spreche, oder ob er auch wandle, oder ob er auch glaube; denn wir hören die Worte der Einsetzung öffentlich und sprechen sie von Herzen mit ihm, und die Einsetzung Christi (nicht unser Thun, noch Chresem) wandelt und gibt uns den Leib und Blut Christi. Glaubt der Pfarrherr nicht oder zweifelt, so glauben wir: strauchelt er in den Worten oder wird irrig und vergißt, ob er die Wort gesprochen habe, so sind wir da, hören zu, halten fest und sind gewiß, daß sie gesprochen sind. Darum können wir nicht betrogen werden. Und weil die Ordnung und rechter Glaube da ist, muß es gewiß sein, daß wir den wahren Leib und Blut Christi empfangen. Und Gott sei Lob und Dank, daß ich die Zeit erlebt habe, zu sehen die rechte christliche Messe und den reinen christlichen Brauch des heiligen Sacraments. Ich sehe es mit Lust und Freude meines Herzens nach dem greulichen, schrecklichen Mißbrauch, so ich leider hab helfen treiben unter dem Pabstgreuel so viel Jahr.“ — (Erl. Ausg. 31, 370–372.)

Bekanntlich haben sich die lutherischen Gemeinden des südlichen und südwestlichen Deutschlands meistentheils weit weniger an Luthers Vorgang in seinen liturgischen Schriften angeschlossen. Sie gingen nicht nur ihren eigenthümlichen Weg und wichen in liturgischen Dingen von Luther und der herrschend werdenden Weise ziemlich ab, sondern gaben, bei allem Festhalten an der Einigkeit im Geiste und an der Lehre Luthers, im Formellen sogar der
\quellseite{71}
reformirten Weise in den Nachbarländern, der Schweiz und Frankreich, mehr oder weniger nach. So schaffte nicht nur Blarer bei der Reformation Württembergs aus der ersten Kirchenordnung vom Jahre 1536 mit Ausnahme der Perikopen Alles weg, was im Cultus an die Messe erinnerte, sondern auch C. Schnepf und Erasmus Alber gaben selber dem Gottesdienst eine weit mehr reformirte als lutherische Einfachheit, indem z. B. das Singen der Einsetzungsworte, der ganze Altardienst außer der Communion, alles Meßgewand, anfangs sogar der Gebrauch des Chorhemds, von diesen beiden so treuen Lutheranern abgeschafft wurde. Und ein Brenz, der mit Osiander die letzte Hand an die Nürnberger Kirchenordnung mit ihrer so reichen Liturgie gelegt hatte, gab in seiner von ihm verfaßten Württemberger Kirchenordnung vom Jahre 1553 dem Hauptgottesdienste eine Form, die wohl nicht kahler sein dürfte, nämlich: Lat. Introitus oder ein „deutsch geistlich Lied“, darauf Predigt und, wenn keine Communion — diese nur eine monatliche — Gesang und Segen. Oder wie einfach und kahl und zusammengewürfelt ist die Ordnung des Hauptgottesdienstes von Worms 1582: Litanei des Chors, auf der Kanzel 3 Collecten. Gesang der Gemeinde: „Nun bitten wir“ oder Te Deum, Predigt, Kirchengebet, Vaterunser, Vermahnung, offene Beichte und Absolution, wieder Ermahnung, Consecration, Austheilung, Gebet, Segen. Wollte doch auch Luther seine Gottesdienstordnung nicht als Gesetz für Andere irgendwie gehalten und aufgestellt wissen, und mit welchem Nachdruck hat die ganze rechtgläubige Kirche im 7. Artikel ihres Hauptbekenntnisses, der Augustana, erklärt: „Dieses ist genug zu wahrer Einigkeit der christlichen Kirchen, daß da einträchtiglich nach reinem Verstand das Evangelium gepredigt und die Sacramente dem göttlichen Worte gemäß gereicht werden. Und ist nicht noth zu wahrer Einigkeit der christlichen Kirchen, daß allenthalben gleichförmige Ceremonien, von den Menschen eingesetzt, gehalten werden, wie Paulus spricht, Eph. 4.: Ein Leib, Ein Geist, wie ihr berufen seid auf Einerlei Hoffnung eures Berufs, Ein HEerr, Ein Glaube, Eine Taufe.“

Wir haben es indeß hier mit dem genuin-lutherischen Hauptgottesdienst zu thun, bei welchem die von der alten Kirche her vorgefundenen guten oder auch einer späteren Zeit angehörenden, aber schriftgemäßen Formen und Weisen möglichst beibehalten worden sind und so „in den öffentlichen Ceremonien keine merkliche Aenderung geschehen“ ist. Und das ist die liturgisch-normale und in der lutherischen Kirche einst herrschend gewesene Ordnung und Weise des Hauptgottesdienstes.

\quellseite{72}
\chapter{§ 8. Vergleichende Tabelle}


Um „die evangelische Messe im Vergleich mit der römischen Messe“ übersichtlich vor Augen zu führen, gibt Löhe in seiner „Sammlung liturgischer Formulare“ etc., 1839, eine Tabelle, in welcher die Ordnung des Hauptgottesdienstes aus den namhaftesten Kirchenordnungen der reformatorischen und nachreformatorischen Zeit neben die der römischen Messe gestellt ist. Diesem Beispiele folgend schließen wir diesen ersten Theil mit einer ihr nachgebildeten Tabelle. Selbstverständlich stellen wir in derselben die beiden liturgischen Hauptschriften Luthers voran — seine „Weise, christlich Meß zu halten“, 1523 und 1524, und seine „deutsche Messe und Ordnung des Gottesdienstes“, 1526. An diese reiht sich die Weise der Brandenburg-Nürnberger Kirchenordnung von 1533, die Chyträus eine „fons et mater agendorum“ nennt und die nicht nur namentlich in Franken, sondern auch im Erzherzogthum Oesterreich u. d. E. eingeführt war, desgleichen auch die Weise der im Jahre 1657 revidirten Kirchenordnung von Braunschweig-Lüneburg, welche ursprünglich von Bugenhagen 1528 verfaßt worden ist und nach deren Gottesdienstordnung vornehmlich die Kirchen in Niedersachsen und Pommern sich richteten. Hierauf folgen ein paar der vornehmenden sächsischen Ordnungen, nämlich die von Justus Jonas, Spalatin, Creuziger und Myconius 1536 unterschriebene und nach Kurfürst Augusts Agenda 1581 gebesserte Ordnung des Herzogs Heinrich zu Sachsen, sowie die Agenda Schwartzburgica, 1675. Die letzte Rubrik füllt die Ordnung des Hauptgottesdienstes der „Kirchen-Agenda für Ev.-Luth. Gemeinden u. A. C. Zusammengestellt aus den alten rechtgläubigen Sächsischen Kirchenagenden und herausgegeben von der Allgemeinen deutschen Evangelisch-Lutherischen Synode von Missouri, Ohio u. a. St. St. Louis, Mo. 1866“. Letztere Rubrik soll zeigen, daß und wie diese Agenda der genuin-lutherischen Weise des Hauptgottesdienstes getreu geblieben ist.

Wir möchten unserer Tabelle ein paar Bemerkungen nun noch vorausschicken.

1. Nicht nur soll die Tabelle zeigen, welche Stücke als papistischer Unrath oder überflüssiger Zierrath in der „evangelischen Messe“ abgethan sind und doch gleichwohl „in den öffentlichen Ceremonien der Messe keine merkliche Aenderung“ geschehen ist, sondern auch, wie bei mancherlei Verschiedenheit,
\quellseite{73}
die man sich in christlicher Freiheit gestattete, doch eine sehr merkliche Gleichförmigkeit betreffs der Bestandtheile des Hauptgottesdienstes und deren Aufeinanderfolge vorherrscht.

2. Auch die große Stetigkeit im lutherischen Hauptgottesdienst springt beim Blick auf die Tabelle in die Augen. Kyrie, Gloria, Collecte, Epistel, Gemeindegesang (Hauptlied), Evangelium, Credo und Predigt und die sich ihr anschließenden Acte bilden allüberall den ersten Theil des Hauptgottesdienstes; und wenn nun auch der zweite Theil einige Verschiedenheit aufzeigt, so betrifft dieselbe außer der Aufeinanderfolge von Vaterunser und den Testamentworten doch hauptsächlich nur den Gebrauch der altkirchlichen Präfation, indem an deren Stelle in Luthers deutscher Messe, 1526, und in der Brandenburg-Nürnberger K.-O., 1533, die Vermahnung tritt.

\section[Vergleichende Tabelle]{Vergleichende Tabelle der Ordnung der römischen Messe und des lutherischen Hauptgottesdienstes}

\begin{vergleichstabelle}
\hline
 & & \multicolumn{2}{c}{Luther.} & \\
\cline{3-4}
 & Römische Messe. & Weise, christlich Meß zu halten, lateinisch 1523, deutsch 1524. & Deutsche Messe etc. 1526. & Brandenburg-Nürnberg. 1533. \\
\hline
\endhead
1. & \textbf{Initium missae.} &  &  &  \\ \hline
2. & \textbf{Confiteor.} &  &  & Confiteor, „oder was den Priester seine Andacht erinnert“. \\ \hline
3. & \textbf{Introitus.} & Introitus. „Obwohl die Psalmen uns dafür lieber wären, aus denen sie genommen sind.“ & „Ein geistlich Lied oder deutscher Psalm.“ & Introitus. Auf den Dörfern etwa deutsche Gesänge. \\ \hline
4. & \textbf{Kyrie} neun Mal. & Kyrie. „In mancherlei Melodie oder Weise nach Unterschied der Zeit.“ & Kyrie. 3 Mal, nicht 9 Mal. So auch in allen andern Kirchen-Ordnungen. & Kyrie. \\ \hline
5. & \textbf{Gloria in excel. Et in terra. Laudamus te.} & Gloria etc. &  & Gloria etc. \\ \hline
6. & \textbf{Collecta (e).} & Eine Collecte. & Collecte. & Der HErr sei mit euch. Eine oder mehr Collecten. \\ \hline
\end{vergleichstabelle}

\begin{vergleichstabelle}
\hline
 & & \multicolumn{2}{c}{Sachsen.} & \\
\cline{3-4}
 & Braunschweig-Lüneburg. 1657. & Herzog Heinrich. K. O. Leipzig 1681. & Agenda Schwartzburg. 1675. & Kirchen-Agende der Synode von Missouri, Ohio u. a. St. St. Louis 1866. \\
\hline
\endhead
1. &  &  &  &  \\ \hline
2. &  &  &  &  \\ \hline
3. &  & Introitus „von der Dominica oder Festen“. & „Komm, Heiliger Geist, erfüll“, sammt einer Collecte um den Heiligen Geist. &  \\ \hline
4. & Kyrie. „Der Cantor, nebenst Schulcollegen u. Schülern singen: O Vater, allmächtiger Gott“ etc. & Kyrie eleison. & Kyrie de tempore. & Kyrie, Gott Vater etc. St. Louiser Gesangbuch Nr. 7. \\ \hline
5. & Gloria in excel. Deo. „Von der ganzen Gem.: Allein Gott i. d. Höh’“ etc. Der HErr sei etc. & Gloria in excel. Deo und Et in terra lateinisch. & Ehre sei Gott etc. Mel. nach Zeit und Fest. Gem.: „All Ehr und Lob soll“ etc.; oder: „Allein Gott in der Höh“; oder des etwas. & P.: Ehre sei Gott etc. G.: „Allein Gott in der Höh“, Ges. Nr. 1. \\ \hline
6. & Collecte „de tempore oder Festo, wie dieselbe sich auf das Evangelium schicket“. & „Die Collecten“, Deutsch oder Lateinisch. & „Collecta de Tempore.“ & Der HErr sei etc. Versikel. Collecte. \\ \hline
\end{vergleichstabelle}

\begin{vergleichstabelle}
\hline
 & & \multicolumn{2}{c}{Luther.} & \\
\cline{3-4}
 & Römische Messe. & Weise, christlich Meß zu halten, lateinisch 1523, deutsch 1524. & Deutsche Messe etc. 1526. & Brandenburg-Nürnberg. 1533. \\
\hline
\endhead
7. & \textbf{Epistola.} & Epistel. Nebenbei Wunsch, auch andere epistol. Lectionen zu gebrauchen. & Epistel. & Ein Kapitel aus den Episteln der Apostel Pauli, Petri und Johannis. Besondere Eingangs- und Schlußformel. \\ \hline
8. & \textbf{Graduale} mit Alleluja oder \emph{Tractus.} (\emph{Prosa, Sequentia.}) & „Gradual, etwa mit 2 Versen sammt dem Halleluja.“ Doch nicht die langen Festgradualia! Sequenz und Pros nur an Weihnachten: grates nunc omnes; und Pfingsten: Spiritus Sancti etc. Veni sancte etc. & „Ein deutsch Lied: ‚Nun bitten wir den Heiligen Geist‘, oder sonst eins, und das mit dem ganzen Chor.“ & Nach der Epistel mag er lesen ein Alleluja mit seinem Vers lateinisch oder ein Gradual aus der heil. Schrift genommen; dasselbe mögen auch die Schüler lateinisch singen. \\ \hline
9. & \textbf{Evangelium.} & Evangelium. Lichter und Rauchwerk beliebig. & Evangelium. & „Kapitel aus dem Evangelium oder Geschichten der Apostel.“ Anm. Man blieb bei den Perikopen. \\ \hline
10. & \textbf{Credo.} & Credo. Predigt hier, oder lieber vor der Messe, denn die Predigt ruft zu Gottes Abendmahl; die Messe ist für die, welche den Ruf angenommen haben, für die Gläubigen. & \emph{Credo:} „Wir glauben all“ von der Gemeinde zu singen. Predigt. Gemein Gebet, ob auf der Kanzel oder am Altar zwar freigelassen; doch: „Es sieht, als haben’s die Alten bisher auf der Kanzel gethan.“ & \emph{Credo,} „das sollen die Schüler lateinisch singen, wie der Gebrauch ist, oder aber das Volk soll den Glauben deutsch singen.“ Darauf soll dann am Feiertag die gewöhnliche Predigt folgen. „Daß zwischen die Predigt oder das heilige Mahl das gemeine Gebet trat \dots{} ist bekannt.“ (Löhe.) \\ \hline
\end{vergleichstabelle}

\begin{vergleichstabelle}
\hline
 & & \multicolumn{2}{c}{Sachsen.} & \\
\cline{3-4}
 & Braunschweig-Lüneburg. 1657. & Herzog Heinrich. K. O. Leipzig 1681. & Agenda Schwartzburg. 1675. & Kirchen-Agende der Synode von Missouri, Ohio u. a. St. St. Louis 1866. \\
\hline
\endhead
7. & Epistel „fein laut, deutlich und langsam“. & Epistel „gegen dem Volk deutsch“. & Epistel etc. & Epistel. \\ \hline
8. & „Nach Verlesung der Epistel singe man aus den gemeinen Gesangbüchern einen deutschen Psalm oder Gesang \emph{de Tempore}. Und kann der Organist auf der Orgel die Gesänge fein langsam in Contrapuncto, wie es die Musici nennen, mit musiciren.“ & Nichts erwähnt. & „Ein deutsch Choral-Lied, welches der Prediger anzuordnen.“ & Hauptlied. \\ \hline
9. & Evangelium „aus unserer Ev. Kirchen-Harmonie, ganz unverstümmelt, hell und deutlich der Gemeine vorgelesen“. & Evangelium „von der Dominica oder vom Fest, auch gegen dem Volk deutsch gelesen“. & Evangelium. & Evangelium. \\ \hline
10. & \emph{Credo} in unum Deum von dem Presbyter gesungen und darauf von der Gemeinde: „Wir gläuben all an Einen Gott.“ Predigt mit vorausgehendem, eingeleiteten Vater Unser und Kanzellied. Beichte und Absolution mit Retention. Allgemeines Gebet. Fürbitten. Etwaige Excommunication etc. etc. Vater Unser. Votum. & \emph{Credo} in unum Deum und das lateinische Patrem. Darnach den Glauben deutsch: „Wir gläuben all.“ Predigt. Beichte mit Absolution ohne Retention. Allgemeines Gebet. & \emph{Credo} in unum Deum „Auf welches ein Stück zu musiciren und der Glaube, wo es die Zeit leidet, zu singen etc. Sodann wird die Predigt gehalten und bei derselben Anfang vor dem Vater Unser ein Lied de Tempore \dots{} gesungen.“ Gemeine Beichte. Allgemeines Gebet. Fürbitten etc. & Glaube: „Wir gläuben all.“ Predigt mit vorausgehendem stillem Gebet (Vater Unser) und Kanzelvers vor Verlesung des Textes. Beichte mit Absolution ohne Retention. Kirchengebet. Fürbitten und Danksagungen. Vater Unser. Votum. \\ \hline
\end{vergleichstabelle}

\begin{vergleichstabelle}
\hline
 & & \multicolumn{2}{c}{Luther.} & \\
\cline{3-4}
 & Römische Messe. & Weise, christlich Meß zu halten, lateinisch 1523, deutsch 1524. & Deutsche Messe etc. 1526. & Brandenburg-Nürnberg. 1533. \\
\hline
\endhead
11. & \textbf{Offertorium.} & Unter dem Gesang trägt man Brod und Wein hervor. &  & „Gesang oder Offertorium.“ \\ \hline
12. & \textbf{Secreta.} &  &  &  \\ \hline
13. & \textbf{Praefatio mit Sanctus.} & Einleitung zur Präfation. \emph{Praefatio} bis: „per Christum Dominum nostrum“, worauf eine kleine Stille erfolgt und dann die & Paraphrase des Vater Unsers mit Abendmahlsvermahnung. & Abendmahlsvermahnung (von Joh. Volprecht). Später mit angehängter allgemeiner Beichte und nochmaliger (allgemeiner) Absolution. \\ \hline
14. & \textbf{Canon Missae:} a. ante consecrationem. b. commemoratio pro vivis. c. diptycha sanctorum. &  &  &  \\ \hline
 & d. \textbf{Consecratio;} e. post consecrationem; f. memento pro defunctis. & \emph{Consecratio} (laut oder leise) und \emph{Sanctus} folgt. Unter dem Benedictus des Sanctus Elevation zugegeben. & \emph{Consecratio} und zwar wird das gesegnete Brod sogleich nach der Segnung und Elevation ausgetheilt. Man singt indeß das deutsche Sanctus: „Jesaia, dem Propheten“, desgl. „Gott sei gelobet“, „JEsus Christus, unser Heiland“. Dann wird der Kelch gesegnet und ausgetheilt. Man singt dabei die angefangenen Lieder zu Ende und zum Schluß das Agnus Dei. & \emph{Consecratio.} Nach derselben das \emph{Sanctus.} \\ \hline
15. & \textbf{Praeparatio ad Communionem.} a. Pater noster. b. Oratio. & Vater Unser. &  & Vater Unser. Auch wohl lateinisch mit den Einleitungsworten aus der Messe: „Praeceptis salutaribus moniti.“ \\ \hline
 & c. \textbf{Pax.} & Pax. &  & Pax. Antwort: Amen. \\ \hline
\end{vergleichstabelle}

\begin{vergleichstabelle}
\hline
 & & \multicolumn{2}{c}{Sachsen.} & \\
\cline{3-4}
 & Braunschweig-Lüneburg. 1657. & Herzog Heinrich. K. O. Leipzig 1681. & Agenda Schwartzburg. 1675. & Kirchen-Agende der Synode von Missouri, Ohio u. a. St. St. Louis 1866. \\
\hline
\endhead
11. & „Weil der Priester von der Kanzel tritt und nach dem Altar gehet, sollen sich die Communicanten auf den (Altar-) Chor verfügen, und ein andächtiger Psalm \emph{de tempore} gesungen werden, dadurch die Gemeine zu mehrerer Andacht bewogen werde.“ &  &  & Gesang: „Schaffe in mir, Gott.“ \\ \hline
12. &  &  &  &  \\ \hline
13. & Präfation deutsch mit Einleitung und dem Sanctus, in welchem „die fürnehmste Partikul der Danksagung bestehet“. Kurze Vermahnung und Gebet. & Präfation lateinisch mit Luthers „Paraphrasie des Vater Unsers“ und der Vermahnung zum Sacrament für dem Altar. & „Deutsche \emph{Praefationes}, welche, wo sie eingeführet, auf die hohen Feste zu gebrauchen.“ \emph{Sanctus} v. Chor. Luthers kurze Abendmahlsvermahnung. & Präfation mit dem Heilig. \\ \hline
14. &  &  &  &  \\ \hline
 & Vater Unser deutsch gesungen. & Vater Unser deutsch gesungen. & Vater Unser deutsch gesungen. & Vater unser gesungen. \\ \hline
15. & Consecratio. & Consecratio. & Consecratio. & Consecration. \\ \hline
 &  &  &  &  \\ \hline
\end{vergleichstabelle}

\begin{vergleichstabelle}
\hline
 & & \multicolumn{2}{c}{Luther.} & \\
\cline{3-4}
 & Römische Messe. & Weise, christlich Meß zu halten, lateinisch 1523, deutsch 1524. & Deutsche Messe etc. 1526. & Brandenburg-Nürnberg. 1533. \\
\hline
\endhead
16. & Preces ante communionem: \textbf{Agnus} und Vorbereitungsgebet. & Agnus. Auch das Gebet aus der römischen Messe: Domine Jesu Christe, Fili Dei. &  & Agnus Dei. Bei großer Menge der Communicanten „soll man nicht allein eine Communio (aus der Schrift genommen) singen, sondern mag u. soll auch etwas mehr singen \dots{} als: das Responsorium \emph{Discubuit}“. \\ \hline
17. & \textbf{Sumtio.} (Nehmen des Sacraments seitens des Meßhaltenden Priesters.) & Während des Agnus Sumtio oder Communio. Dabei aus der römischen Messe: Corpus Domini etc. Sanguis Domini etc. &  & Indeß Communion. \\ \hline
18. & \textbf{Commun.} & Commun., wenn man will. &  &  \\ \hline
19. & \textbf{Post communio.} & Statt der letzten Collecte, die meist aufs Opfer geht, das Gebet aus der römischen Messe: Quod ore sumsimus. & Dankcollecte. & Collecte. \\ \hline
20. & \textbf{Finis Missae.} & Benedicamus mit Halleluja. &  & Benedicamus etc. \\ \hline
21. & \textbf{Benedictio.} & Segen aus 4 Mose 6. & Segen etc. & Segen. (4 Formeln.) \\ \hline
22. & Joh. 1, 1–14. &  &  & Gesang: Erhalt uns, HErr, bei deinem etc. \\ \hline
\end{vergleichstabelle}

\begin{vergleichstabelle}
\hline
 & & \multicolumn{2}{c}{Sachsen.} & \\
\cline{3-4}
 & Braunschweig-Lüneburg. 1657. & Herzog Heinrich. K. O. Leipzig 1681. & Agenda Schwartzburg. 1675. & Kirchen-Agende der Synode von Missouri, Ohio u. a. St. St. Louis 1866. \\
\hline
\endhead
16. & O Lamm Gottes. & \emph{Agnus Dei} lateinisch. (Während des Gesangs beginnt schon die Austheilung.) &  & Christe, du Lamm Gottes, Nr. 69. \\ \hline
17. & Indeß Communion unter weiterem Gesang von Abendmahlsliedern etc. & Indeß Communion unter weiterem Gesang von: JEsus Christus, unser etc. Man mag auch den 111. Psalm singen, nachdem der Communicanten viel sind. & Communion. & Communion unter dem Gesang von Abendmahlsliedern. \\ \hline
18. &  &  &  &  \\ \hline
19. & Ps. 23 vorgelesen. Collecte. & Collecte. & Collecte. & Collecte. \\ \hline
20. &  &  &  &  \\ \hline
21. & Segen etc. & Segen. & Segen. & Segen. \\ \hline
22. & Darauf singe die ganze Gemeine choraliter: „HErr, nun lässest du deinen Diener in Friede fahren“ etc. &  &  & Gesang: Gott sei gelobet und gebenedeiet, Nr. 195, 1 etc., oder einen andern passenden Schlußvers. \\ \hline
\end{vergleichstabelle}


\quellseite{74}
Und doch bei aller Stetigkeit welche Mannigfaltigkeit und darum welcher Wechsel in den Worten und Weisen dieser feststehenden Bestandtheile! Da ist allsonn- und festtäglich der Introitus ein anderer, da ist die Form von Kyrie, Gloria und Credo in Wort und Melodie mancherlei, wie der zweite Theil insonderheit zeigen wird, da gibt es eine Präfation nicht nur für die gewöhnlichen Sonntage, sondern auch je eine für die Festtage und Festzeiten. Und welcher Wechsel im Gemeindegesang! Man denke sich nun eine mit diesen verschiedenen Formen und Weisen einigermaßen vertraute Gemeinde, und dazu einen seine Aufgabe erfassenden Kirchenchor, und es wird wohl schwerlich die Klage über ein Einerlei der Form des Gottesdienstes zu vernehmen sein oder das Gelüsten nach Aenderungen in der Liturgie bei Prediger und Gemeinde Berücksichtigung finden.

\quellseite{75}
\quellseite{76}
\quellseite{77}
\quellseite{78}
\quellseite{79}
\quellseite{80}
\quellseite{81}
\quellseite{82}
\part[Zweiter Theil]{Zweiter Theil Die einzelnen Bestandtheile des lutherischen Hauptgottesdienstes}
\quellseite{83}
\quellseite{84}
\chapter{Einleitung}


\begin{quote}
„Die Noten machen den Text lebendig.“ Luther.
\end{quote}

\section{§ 1}


Am Weihnachtsfeste des Jahres 1525 wurde die erste deutsche Messe in der Pfarrkirche zu Wittenberg gehalten. Sie wurde aber nicht gesprochen, sondern gesungen. Wir wissen bereits, daß und wie Luther in besonderer dreimonatlicher Arbeit mit dem Kapellmeister Johann Walther den musikalischen Theil der deutschen Messe schuf. Schrieb er doch schon zu Anfang des Jahres 1525 in seiner Schrift „wider die himmlischen Propheten“: „Ich wollte heute gerne eine deutsche Messe haben, ich gehe auch damit um, aber ich wollte ja gerne, daß sie eine rechte deutsche Art hätte. Denn daß man den lateinischen Text verdolmetscht und lateinischen Ton oder Noten behält, laß ich geschehen; aber es lautet nicht artig noch rechtschaffen. Es muß beide Text und Noten, Accent, Weise und Geberde aus rechter Muttersprach und Stimme kommen: sonst ist’s alles ein Nachahmen, wie die Affen thun.“ (Luther, Erl. Ausg. 29, S. 203.)

Was bewog ihn nun, auch in der Vortragsweise keine „merkliche Aenderung“ vorzunehmen, ohne daß man nicht in lateinischer, sondern in deutscher Sprache und auf deutsche Art die Messe sang? War’s lediglich Rücksicht auf die nun einmal herrschend gewordene liturgische Vortragsweise? Der Grund ist ein tieferer.

Was im Gottesdienst der Christen von Seiten des Pastors und der Gemeinde gehandelt wird, besteht nicht nur im Sprechen, sondern auch im Singen des Worts; denn „das in gehobener Gemüthsstimmung gesprochene Wort geht in Folge des dem Menschen angebornen Kunsttriebes in Gesang über, indem sich die Modulation der Stimme nach der in den Naturgesetzen begründeten Tonleiter fixirt“ (Zahn). Dieser Gesang ist aber nicht nur ein metrischer, sondern auch ein recitativischer. Und letzterer ist wohl der ursprüngliche. Sagt doch auch Prof. Max Müller, der anerkannte Sprachkenner, wie
\quellseite{85}
Lyra in seinen „Altarweisen“, aus dessen kleinerer Ausgabe heiliger Lieder indischer Brahmanen citirt: „Auch halte ich dafür, daß, was wir Accent nennen und was für uns durchaus nichts Musikalisches hat, ursprünglich in allen Sprachen rein musikalisch war. ... Ich weiß, daß diese Ansicht von dem ursprünglichen Charakter der Accente weiterer Begründung bedarf. ... Für das Sanskrit steht sie aber ziemlich fest“; desgleichen mit Bezug auf die älteste Prosodie besonders, aber auch auf die ursprüngliche Form der Recitation in allen Sprachen insgemein: „Die Cantilene (der Singeton) der Stimme stieg wirklich beim Hochton, fiel beim Tiefton und hatte im Uebrigen eine mittlere Haltung.“

Nach den Zeugnissen aus dem kirchlichen Alterthum ist für gewiß anzunehmen, daß nach dem Vorbild der jüdischen Kirche, welche die Psalmen antiphonatim nach dem Parallelismus sang (Ps. 147, 7.: „Singet um einander dem HEerrn“), frühzeitig in den Gottesdiensten der Christen ein Wechselgesang stattfand. Auf die clericalischen Intonationen und Gebete antwortete die Gemeinde, indem sie den von einem oder etlichen der Kirchendiener intonirten Psalm oder Abschnitt eines Psalms fortführte oder wiederholte oder beschloß oder mit dem Gesang eines bloßen „Amen“ besiegelte (cantus responsorius, hipophonicus). Diesem trat dann später im 4. Jahrhundert noch eine andere Art des Wechselgesangs zur Seite, der antiphonische Gemeindegesang selbst, bei welchem die Gemeinde sich in zwei oder etliche Chöre theilte — in Männer-, Frauen-, Jungfrauen- und Knabenchöre —, die „um“ oder „gegeneinander“ sangen (cantus antiphonicus).\footnote{Eusebius, Ephraim, vornehmlich aber Ambrosius, der Dichter und Sänger unsterblicher Hymnen der alten Kirche, gedenken auch dieser Wechselgesänge. Wir entnehmen der „Siona“ folgende Uebersetzung der schönen Worte aus dessen Hexaem. III, 9.: „Was ist jenes Tosen der Meereswogen anders, als ein Typus der Jubelgesänge (Hymnen) der Gemeinde? In der That — nur sehr passend wird dem Meere die gottesdienstliche Gemeinde verglichen, indem diese bei dem Hereinströmen der Kommenden gleichsam das erste sanfte Wellengeräufel zeiget, welches bei den vereinigten Gebeten der Versammlung wie in murmelnder Brandung sich lauter erhebt, und endlich bei den Responsorien der Psalmen, sowie in dem Wechselgesang der Männer, Frauen, Jungfrauen und Knaben, gleich dem concertirenden Widerhall der Wogen, mächtig abrauset.“ Daß Luthern bei Anrichtung des deutschen Gottesdienstes auch diese letztere Art Wechselgesänge vor Augen standen, beweist nicht nur die Verdeutschung des Te deum und der Litanei, sondern auch die bei dem Weihnachtsliede: „Vom Himmel hoch da komm ich her“ gemachte Anordnung, nach welcher die Worte des Engels V. 1–5 von Knaben allein gesungen werden sollen.}

Durch Ambrosius kam bekanntlich in der abendländischen Kirche ein besonderer Gemeindegesang noch auf, indem jener metrisch gebaute Hymnen
\quellseite{86}
dichtete und diesen eine modulirte Singweise gab, welcher geistliche Volksgesang ja durch Luther erst recht zu einem neuen und unvergänglichen Leben in der Kirche der Reformation erstanden ist. Als nun aber hernach dieser ambrosianische Hymnengesang in Verfall gerieth und dabei eine große Mannigfaltigkeit und Unsicherheit der Singweise in den Kirchen hin und her sich immer mehr geltend machte, so trat eine neue Epoche des Kirchengesangs ein. Wir meinen den römischen Choral\footnote{Choral heißt dieser Gesang, weil er im Chor, das ist, im Altarraum von den Clerikern vorgetragen wurde.} oder den Gregorianischen Gesang (cantus choralis Gregorianus). Papst Gregor I. (geb. 540, gest. 604) suchte dadurch eine Reform des Kirchengesanges herbeizuführen, daß er die vorhandenen guten, aber aus der Uebung gekommenen Kirchengesänge sammelte, verbesserte, vermehrte und ordnete, ihnen ihre Stellung in der Liturgie anwies, für deren Vortrag und Singweise ein musikalisch-theoretisches System feststellte und dabei die bisherigen, aus der ambrosianischen Zeit stammenden 4 Kirchentöne oder Tonarten auf 8 erhöhte, welche Tonarten sich auch durch das ganze Mittelalter erhalten haben und noch jetzt nach ihren besonderen Namen (ionisch, dorisch, phrygisch, lydisch etc.) in der Musik bekannt sind, wenn schon ein völliger Umschwung des ganzen Musiksystems eingetreten ist. Zur Fixirung des von ihm eingeführten Gesangs erfand Gregor auch eine besondere Tonschrift und zur Pflege derselben gründete er besondere Schulen, weshalb man ihm auch nachher als Schutzpatron der Schulen einen besonderen Tag zum Gedächtniß widmete, den St. Gregoriustag, und diesen sogar in der lutherischen Kirche noch lange Zeit hindurch als jährliches Schulfest feierte. Unter den von Gregor gemachten Anordnungen lautete eine: „Keiner kann zum Priesteramte gelangen, der nicht wohl erfahren ist im Gesang.“

„Der gregorianische Gesang“, sagt Jahn in seinem „Handbüchlein für evangelische Cantoren und Organisten“, ist ein Erzeugniß ernster christlicher Kunst und Frömmigkeit und hat einen unvergänglichen Werth für die ganze christliche Kirche. Es ist ihm besonders eine feierliche Erhabenheit und Würde eigenthümlich.“ Gregor ging nämlich von dem Grundsatz aus, daß aller Gesang in der Kirche Träger des göttlichen Wortes sein müsse. Damit keinerlei außer dem Worte Gottes liegende Factoren sich geltend machten, mußten daher nicht nur alle diejenigen Elemente im kirchlichen Gesange weichen, welche dazu angethan waren, überwiegend ästhetisches Interesse zu erwecken, sondern es mußten bei dem gregorianischen Gesang, dem cantus firmus, auch
\quellseite{87}
alle künstlichen Formen wegfallen. So entstand denn eine Singweise, bei welcher der Gesang einstimmig\footnote{„Der gregorianische Choralgesang war, da er vor Erfindung des Contrapunktes (Tonsatzes) entstanden ist, ursprünglich bloß für einstimmigen Vortrag bestimmt. Er wurde zwar später mehrstimmig gesetzt und auf die verschiedenste Weise von der Kunst des Contrapunkts benützt; er widerstrebt aber vielfach solcher Benützung.“ (Jahn.)} und in lauter einfachen Haupttönen einherschreitet, die nur bei gewissen Abschnitten des Textes durch Hebung und Senkung der Stimme nach bestimmten Regeln melodisch werden. Das Charakteristische dieser Singweise zeigt auch deren Terminologie. Wegen des ebenmäßigen Fortschreitens der Melodie, die sich meistens nur in Sekunden und Terzen bewegen, hieß er cantus planus (ebenmäßiger Gesang, englisch: plain chant). Weil die einmal festgesetzte Melodie von den Sängern nicht verändert werden darf, so redet man von einem „cantus firmus“ (canto firmo, feststehender Gesang), und weil dieser Gesang nicht bloß Einzelgesang, sondern zum Theil auch von Vielen oder Mehreren zugleich vorgetragen wurde, zum Theil also auch Chorgesang war, so nannte man ihn den „Cantum choralem“ (Allgemein- oder Vereintgesang), daher auch die Vorsänger, welche genau die feststehende Singweise, den canon, kennen mußten, ursprünglich canonici und die von ihnen geleiteten Chorsänger chorales hießen.

Man denke nun aber nicht, daß der gregorianische Gesang schleppend und monoton langweilig ist, wie Unkundige meinen, weil man dessen Töne durch Noten von gleichem Zeitwerthe darzustellen pflegt. Auch er hat Rhythmus, indem beim Vortrag die Töne durch die Länge und Kürze vielmehr ihr Gewicht erhalten und er daher nur insofern taktlos ist, als seine Töne nicht in einem genau abgemessenen Zeitmaße abgesungen werden, wie dies z. B. bei unsern Gesangbuchsliedern geschieht. Dies ist auch ersichtlich aus den in der alten Tonschrift oft genug angebrachten Zeichen, welche bald eine schnellere, bald eine langsamere Vortragsweise andeuten, und zwar nicht für einzelne Töne, sondern für ganze Wort- und Tonphrasen, z. B. t. (tarditas = Zögerung), m. (mediocritas = Mittelbewegung), c. (celeritas = Schnelligkeit). Ebenso bezeichnen folgende Buchstaben die Tonstärke: b (bene) bedeutet eine Verstärkung; b t (bene teneatur) ein längeres Anhalten mit erhobener Stimme, x (expectare) eine Ruhepause.

Weil der Gregorianus zum Theil Einzelgesang, zum Theil Vereintgesang ist, so theilen sich seine Gesänge in zwei Gattungen: den Accentus und den Concentus. Zu dem Accentus gehören alle Gesänge, welche von einer
\quellseite{88}
Person allein anzustimmen oder abzusingen sind, also die Gesangesvorträge des Liturgen: die Intonationen, Collecten, Lectionen, Präfationen etc., welche Vorträge mehr den Leseton, den Recitationston, haben, was man in späterer Zeit choraliter lesen (modus choraliter legendi) nannte. Wir kommen darauf noch einmal zurück. Zum Concentus gehören alle von mehreren Personen zugleich anzustimmende Gesänge, also alle Responsorien- und Antiphonengesänge, die Psalmodien und die eigentlichen Chorgesänge. Für sie waren bestimmte Formen vorgeschrieben und bezog deren Verschiedenheit sich auf Anfang, Mitte und Ende.

Wohl ist es nun wahr, daß mit dem gregorianischen Gesang die Theilnahme des Volkes am Kirchengesang förmlich und völlig aufhörte, und zwar theils durch die Fertigkeit, welche er erfordert, theils durch den ausschließlichen Gebrauch der lateinischen Sprache, auch nachdem dieselbe bereits zur todten Sprache geworden war. Wie schon an seinem Orte bemerkt, ging der Gesang gänzlich in die Hände des hohen und niederen Clerus über, so daß dem Volke im Gottesdienst nichts als das bloße Zusehen und Zuhören verblieb. Doch es kam ja die gesegnete Zeit, in welcher es auch hier wieder hieß: „Alles ist euer!“ — also auch Ambrosius und Gregorius. Luther führte das Christenvolk nicht nur zur Betheiligung am Kirchengesang wieder zurück, und zwar zur reichlichen Theilnahme durch die Einführung des geistlichen Volksliedes als integrirenden Theil des Gottesdienstes, sondern auch zum Mitbesitz und köstlichen Mitgenuß des cantus Gregorianus durch die Verwendung, welche letzterer von ihm und den ihm befreundeten Tonmeistern fand. Und dazu war er der Mann. Von seiner Schulzeit her lebte ja der wie sprachlich und poetisch, so auch musikalisch hochbegabte Luther im Tonheiligthum des cantus choralis Gregorianus. Mit dessen Tongrammatik war er vertraut, wie mit der Sprachgrammatik. Im Hinblick auf dieses ihm so bekannte und von ihm verstandene gregorianische Tonheiligthum fiel sein Urtheil über die Musik aus, also daß er bekanntlich ihr unter sämmtlichen Künsten die nächste Stelle nach der Theologie einräumte, und es nur um so mehr fein und herrlich fand, daß „viel Gesang in der Meß“ sei. Wie er daher in den wichtigsten Bestandtheilen und der ganzen Composition des bisherigen Hauptgottesdienstes keine merkliche Aenderung machte, weil derselbe „eine christliche feine Abkunft“ hatte, ohne daß er ihn von den Zusätzen fegte, die ihn verderbt hatten: so behielt er auch das Singen der Messe und die bisherige Weise dieses Singens bei, weil er in dem cantus Gregorianus die den heiligen Texten entsprechendste musikalische Form vor sich
\quellseite{89}
hatte und dieses Kunstgebilde in den rechten Dienst dessen gestellt wissen wollte, der aller Künste Urheber und Geber ist. Doch „jede Sprache hat ihre andere Musik, ein anderes Gesetz ihrer Klangordnung“. Niemand aber mußte das besser, als Luther, der darüber an verschiedenen Orten sich ausläßt, und Niemand verstand es auch besser, den lateinischen Meßgesang auf deutsche Art zu gestalten, als er, der Urheber unserer schönen deutschen Kirchen- und Schriftsprache, der deutsche Uebersetzer der Bibel. Und wie ihm auch in musikalischer Beziehung sein Werk gelungen ist, das haben wir aus dem Munde seines musikalischen Mitarbeiters, Johann Walther, vernommen. Letzterer aber macht sich nebst andern Tonmeistern dabei auch noch durch die Verbindung des cantus Gregorianus mit dem mehrstimmigen Tonsatze um den kirchlichen Kunstgesang in der lutherischen Kirche verdient.

Durch Luther haben wir somit neben dem geistlichen Volkslied auch den gregorianischen Kirchengesang in unserm öffentlichen Gottesdienst nach Bedarf verwerthet und zwar für deutsche Texte auch „auf deutsche Art“. In den beibehaltenen Bestandtheilen der Messe, im Ambrosianischen Lobgesang und in der Litanei, im Lobgesang des Zacharias und der Maria, wie im Gesang Simeons, und endlich im Psalmengesang finden wir ihn verwerthet, und jene herrlichsten musikalischen Erzeugnisse auf dem Gebiete des geistlichen Volksliedes, wie „Ein feste Burg ist unser Gott“ — „Jesaja dem Propheten das geschah“ — „Allein Gott in der Höh sei Ehr“ verleugnen bei aller Originalität nicht ihre gregorianische Abstammung. — Werke, wie die „Kirchengesänge“ von J. Spangenberg, Nordhausen 1545, und von J. Keuchenthal, Wittenberg 1573 — die „Psalmodia sacra“ von Lucas Lossius, Lüneburg 1553 und 1579, ein auch in katholischen Kreisen hochgeschätztes Quellenwerk — die „Cantica sacra“ von Franz Eler in Hamburg 1588 — das „Missale, Vesperale und Matutinale“ von M. Ludecus in Havelberg 1589 — die „Missodia etc.“ von M. Prätorius 1611 u. a. entstammen der lutherischen Kirche. Auch in dieser Anerkennung und Verwerthung des cantus Gregorianus zeigt sich wieder die charakteristische Verschiedenheit der schweizerischen und der deutschen, der Lutherischen Reformation. Während die Schweizer von vornherein gegen denselben sich ablehnend verhielten, ja die zweite helvetische Confession ihn ausdrücklich als „höchst abirrend“ und dessen Fernhalten aus dem reformirten Gottesdienst als ein Verdienst bezeichnen, hegen die Lutheraner des mittleren und nördlichen Deutschlands gegen denselben die größte Hochachtung und bekunden ein treffliches Verständnis seiner Schönheit und seiner Verwendung für seinen Dienst
\quellseite{90}
zur Förderung des Evangeliums. Bei der Rückkehr zu der genuin-lutherischen Gottesdienstweise in der Gegenwart ist man denn auch auf die Bedeutung des cantus Gregorianus und seine Verwendung in der lutherischen Kirche wieder aufmerksam geworden und hat sich hier der Großherzogliche Musikdirector Otto Kade zu Schwerin, Mecklenburg, besondere Verdienste erworben, und zwar nicht nur durch seinen im Jahre 1870 erschienenen „Luther-Codex“, sondern auch durch ein von dem sel. Großherzog Friedrich Franz II. kräftig gefördertes Unternehmen. Es ist das „die Verarbeitung, Verwendung und Verwerthung der alten herrlichen Gesangstücke des Gregorianischen Melodienschatzes, die in der protestantischen Kirche leider längst in Vergessenheit gerathen und verklungen sind, in vorsichtiger Auswahl von neuem in wohlgefälliger ansprechender Gestalt zu vier Stimmen der Gemeinde wieder zuzuführen“.\footnote{Der erste, schon im Jahre 1867 auf großherzogliche Kosten gedruckte Theil dieses Werkes enthält die einstimmigen gregorianischen Tonweisen zu dem Hauptgottesdienst, sowohl für den Liturg, wie für den Chor. Der zweite Theil, der 1875 erschien, enthält die einstimmigen gregorianischen Melodienstücke für die Nebengottesdienste, Mette und Vesper, und der dritte Theil brachte 1880 die mehrstimmige Bearbeitung des ersten Theils durch D. Kade. Ob der vierte Theil, der die Tonsätze zum zweiten Bande bringen sollte, nunmehr erschienen ist, ist mir zur Zeit unbekannt. Der deutsche Text ist durch Dr. Klieboth besorgt.} Groß ist das Lob, das dieser Kenner des alten Kirchengesanges dem cantus Gregorianus spendet. „Nichts in der Welt“, ruft er u. a. im Vorwort zu seinem „Luther-Codex“ aus, „ersetzt den tiefen Werth dieser Charaktertypen und Gesangsformen, an deren Vollendung die Kirche tausend Jahre arbeitete. Keine Musik erreicht denselben an eindringlichen Motiven, es ist die geheimnißvolle Tonsprache, die es in der Welt gibt, es ist der köstlichste Besitz einer Gemeinde, die in dieser reichen Auswahl von Singweisen, womit sämmtliche liturgische Textstücke nicht nur einmal, sondern bisweilen doppelt, je nach der liturgischen Stellung, belegt sind, einen Mittelpunkt fand, in welchem sich Kirche und Kunst begegnen. Es ist die in Musik gesetzte Bibel.“\footnote{Nämlich „die in Musik gesetzte Vulgata! Aber so wie der, welcher die lateinische Grammatik fest gelernt hat, auch gleichzeitig für die deutsche Grammatik die sämmtlichen Generalregeln mitgelernt hat, so eben auch hat Luther aus dieser musikalischen Vulgata die Generalregeln gelernt, deren er sich bei Melodisirung der Lutherbibel zu bedienen hatte“. (Blätter f. Hymnol., Juni 1884.)} — Nie und unter keiner Bedingung hätte der Sinn für diese Gesangweisen aus dem Gedächtniß der Gemeinde schwinden, nie die Verwendung dieser herrlichen Charaktertypen in dem mehrstimmigen kirchlichen Tonsatze aufgegeben werden
\quellseite{91}
dürfen.“ In einem, im „Halleluja“, einer Zeitschrift für geistliche Musik, im Februar 1884 erschienenen Artikel: „Der großherzogliche Schloßchor in Schwerin“ wiederholt Kade fast wörtlich dieses Urtheil und fügt demselben noch hinzu: „Dieses Fundament aller Melodie . . . vermag allein den Begriff Kirchenmusik zu vermitteln, der ohne ihn schwer faßlich ist. Die freie Composition, das heißt, eine Composition, welche diesen alten ehrwürdigen Choral nicht zu Grunde legt, war dadurch vom Gottesdienste principiell ausgeschlossen, oder wenigstens auf ein Minimum beschränkt.“ —

Mit Beseitigung der alten rechtgläubigen Agenden hat der Rationalismus unser lutherisches Volk auch um das Erbe des liturgischen Gesanges größtentheils gebracht oder doch den Gebrauch desselben, wo er noch fortbestand, meist auf ein Minimum beschränkt. Insbesondere aber ist unserm hiesigen, im Lande der Schwärmer aufgewachsenen Volke so sehr der Sinn und das Verständnis für den liturgischen Gesang überhaupt abhanden gekommen, daß man trotz aller Belehrung und Reizung auch selbst in so mancher erkenntnißreicheren, namentlich großentheils aus Württembergern und andern Süddeutschen bestehenden Gemeinde von Wiedereinführung des liturgischen Singens zur Zeit noch Abstand nehmen muß, will man als rechtschaffener Lutheraner um einer äußerlichen, wenn auch noch so herrlichen und wünschenswerthen, aber doch immerhin menschlichen Ordnung willen nicht Trennung und Spaltung herbeiführen. Singendes Respondiren seitens der Gemeinde oder des Chors läßt man sich vielleicht noch da und dort gefallen; nur der Liturg soll am Altar in keinerwege sich des kirchlichen Lesetons bedienen, denn das ist „katholisch“. Und ob auch da und dort der liturgische Altargesang sich behauptet hat oder wieder in Aufnahme gekommen ist, man singt meist nur, wie man's von andern gehört; denn es ist nicht zu leugnen, „daß die positiven Maximen, die sehr einfachen Grundbegriffe der kirchlichen Accentlehre, leider den Unsrigen verloren gegangen sind in einer Zeit, worin . . . an die Stelle der kirchlichen Kunst die Willkür des subjectiven Geschmackes gesetzt und am Altar ein schaler Abhub von den Prunkstücken der Salons und Opernbühnen den Kirchgängern geboten wurde“, und worin die alten liturgischen Accente mit ihrer gesunden Kraft und Einfalt aus der Technik des Cultus verdrängt waren „durch die nämliche Tyrannei des Zeitgeistes, womit er in den kirchlichen Gebäuden die gothischen Säulenknäufe schonungslos unter den Spitzhammer brachte, um corinthische Capitälter von Gyps und allerlei gepinseltes Stuccaturwerk daraufzuleimen“ (J. Lyra). Und der Chorgesang? Der Opernstil hat bei Sängern
\quellseite{92}
und Gemeinde den Geschmack so gründlich verdorben, daß der, welcher die opernstilartigen Chorstücke in die Rumpelkammer verweisen und nach den Schätzen des altkirchlichen Chorgesangs greifen will, vorerst noch gegen den Strom zu schwimmen hat.

Indes ist ein guter Anfang zur Bekanntschaft mit den in Vergessenheit geratenen musikalischen Schätzen unserer Kirche und zum Verständnis der Entstehung unserer Altarliturgie, ihrer positiven Maximen und der einfachen Grundbegriffe der liturgischen Accente gemacht. Außer D. Kade haben F. Hommel in seiner „Liturgie des luth. Gemeindegottesdienstes“ (1851) und in seinem „Psalter“ etc. (1859); Dr. F. Layriz in seinem „Kern des deutschen Kirchengesanges“ 1855, dessen 4. Band 120 liturgische Weisen enthält; Ph. Herold und Prof. Dr. Krüger in ihrer Zeitschrift „Siona“, vor allen Dingen aber die Prof. Dr. Schöberlein und Riegel in ihrem dreibändigen und unvergleichlichen „Schatz des lit. Chor- und Gemeindegesangs nebst den Altarweisen“ etc. 1872 u. a. die Schätze gehoben, während Dr. L. Kraußold in seinem „Handbuch für den Kirchen- und Choralgesang“ 1855 und namentlich Pastor Justus W. Lyra in „Die Liturg. Altarweisen des luth. Hauptgottesdienstes nach ihrer Reinheit und Einheit in musikal. Beziehung untersucht und festgestellt etc. 1873“, und im Anschluß daran in seinem „Andreas Ornithoparchus aus Meiningen, der Zeitgenosse Luthers\footnote{Andreas Ornithoparchus („Vogelsang“) ist einer der ältesten und bedeutendsten deutschen Musikpädagogen. In demselben Jahre, in welchem Luther seine 95 Thesen anschlug, gab er zu Leipzig zuerst, dann hernach öfter sein „Musicae activae Micrologus“ heraus. Das dritte Buch dieser Schrift enthält die Lehre vom liturgischen Accent, die systematische Darstellung und methodische Weiterführung der Grundregeln für das choralmäßige Lesen. D. gilt als Originalschriftsteller für diese Lehre.} und dessen Lehre von den Kirchenaccenten“ 1877 uns in die Entstehung und die Grammatik des liturgischen Gesanges einführen. Und als einer auf amerikanischem Boden gereiften Frucht so mancher der genannten Arbeiten und ihres eingehenden Studiums sei hier nochmals auf das im Jahre 1879 zu Philadelphia bei Jg. Kohler, Nr. 911 Arch Straße, erschienene „Choralbuch mit Liturgie und Chorgesängen zum Kirchenbuch der Allgemeinen Kirchenversammlung. Bearbeitet von J. Endlich“ hingewiesen. Das Werk zerfällt in zwei Theile. Der erste enthält die Liturgie und Chorgesänge, der zweite die sogenannten Choräle oder die Weisen zu den Liedern des Gesangbuchs, und zwar durchweg in ihrer rhythmischen Gestalt und nach dem älteren harmonischen Satz. In
\quellseite{93}
Bezug auf den ersten Theil sei für unsern Zweck rühmend hervorgehoben, daß der Bearbeiter als Hauptquellen für liturgische Musik außer Schöberleins „Schatz“ etc. aus vergangener Zeit benützte: Spangenberg, Kirchengesänge 1545; Pfalzgraf Wolfgangs Kirchenordnung 1570; Pommersche Kirchenordnung 1563, und andere ältere und neuere Werke; desgleichen, daß dies Werk die Intonationen und Responsorien nur nach dem reinen und richtigen Accent enthält und daß es für den Liturg(en), für Chor und Gemeinde nur Bewährtes, nur Altkirchliches darbietet. Möge dies kirchliche Musikwerk auch in unsern Kreisen die ihm im Interesse insbesondere des liturgischen Gesanges gebührende Beachtung finden!

Nachtrag. Den vorgenannten Werken hat sich neuerdings noch hinzugesellt: „Der Vortrag des liturgischen Gesanges. Ein Handbuch für ev. Geistliche, bearbeitet von H. Franke. Leipzig. Verlag von Carl Klinner, Musikalienverlag und Sortiment. 1891.“ (S. dessen Beurtheilung in „Lehre und Wehre“, 39. Jahrg. März. S. 87 ff.)

\section[§ 2. Lehre vom Kirchenaccent]{§ 2. Zum richtigen Vortrag auch der lutherischen Altarweisen bedarf es einiger Kenntniß der Lehre vom Kirchenaccent}


Unter Kirchenaccent (accentus ecclesiasticus) versteht man die Tonweise, worin die Silben beliebiger Textworte nach Maßgabe ihrer natürlichen Aussprache mit regelrechter Hebung und Senkung der Stimme vorgetragen werden. Es ist das also nicht ein grammatischer, sondern ein musikalischer Accent, und zwar als solcher speciell in Bezug nicht auf den Rhythmus, sondern auf die Melodie. Die Lehre vom Kirchenaccent befaßt sich demnach mit der Aufstellung der Regeln für das Choraliterlesen (modus choraliter legendi). In der lutherischen Kirche jedoch befaßt sich die Accentlehre nur mit den elementarsten Formen des Gregorianischen Gesanges: der Psalmenmelodien nach den acht Kirchentönen, der Tonweise der Lectionen und der Testamentsworte, der Orationen (Collecte, Präfation, Vater Unser), der Versikel und der Benedictionen.

(Es ist an seinem Ort schon erinnert worden, daß Luther seinem Mitarbeiter am Musikalischen der deutschen Messe auf die Frage, wer ihm denn so meisterlich die Anwendung des Accentus und Concentus gelehrt habe, antwortete: „Der
\quellseite{94}
Poet Virgilius hat mir solches gelehrt, der also seine Carmina und Wort auf die Geschichte, die er beschreibt, so künstlich richten kann. Also soll auch die Musika alle ihre Noten und Gesänge auf den Text richten.“


Nach diesem Grundsatz verwerthete Luther nun die Accentlehre, wie sie kaum ein Jahrzehnt vorher von Meister Ornithoparchus in seinem Micrologus vorgetragen worden war. Es geschah dies von ihm in seiner „Deutschen Messe und Ordnung des Gottesdienstes 1526“. „Sie bildet alsdann ein stehendes Kapitel in den zahlreichen Anweisungen zum Gesange für die Schulen des XVI. Jahrhunderts. Der häufig mißverstandene Satz: accentus pater musices, ‚der Accent ist der Vater des Gesanges‘, ward dadurch in Fleisch und Blut der Zeitgenossen verwandelt, deren musikalische Gemeinbildung nicht wie die heutige sich auf den sinnlichen Klang musikalischer Instrumente, sondern auf den sinnvollen Ton der articulierenden Menschenstimme stützte.“ (Lyra.) Wie dies geschah, zeigen die Weisen zur deutschen Messe in der Musikbeilage zu Band 22 von Luthers Werken Erl. Ausg. wie in Band X. der St. Louiser Ausgabe, die auch andernwärts verschiedentlich transportirt sich vorfinden. Wir sehen aus denselben zugleich, in welchem Sinn bei Herstellung einer deutsch zu singenden Messe für die betreffenden Theile derselben Luther die gregorianischen Tonweisen formte und dabei über Ornithoparchus und die bisherige Tradition selbst hinausging, und durch den an der deutschen Bibelübersetzung erwachten Geist der deutschen Kirchensprache bewogen und getrieben, oft, wie z. B. im Evangelien- und Epistelton und den Einsetzungsworten schwungreichere, ausdrucksfähigere Formen zu setzen, die aber gleichwohl das Gepräge bestimmter gregorianischer Tonarten fest und deutlich an der Stirne tragen.

Wohl haben wir nun in den vorhandenen alten Agenden, wie in den heutigen gehobenen Schätzen genug Mustervorlagen für die liturgischen Altarweisen. Dennoch wollen wir durch Hülfe der genannten Schriften von Kraußold und Lyra\footnote{Und nun auch des im Nachtrag zum vorigen Paragraphen genannten Werkchens von Francke.} wenigstens für die einfachste Form des liturgischen Gesanges einen kurzen Leitfaden uns zu verschaffen suchen, um nicht bloß nach der Tradition, die heutzutage sehr unzuverlässig ist, noch viel weniger nach dem noch unzuverlässigeren subjectiven Ermessen, liturgisch singen zu lernen und so mit Nutz die vorhandenen Muster desto mehr studiren und gebrauchen zu können.

Für diese einfachste Form nun bezeichnet Lyra diejenigen Tonweise als die beste, welche
\quellseite{95}
1. der freien Rhythmik eines wohlarticulirten, würdig gemessenen affectlosen Redevortrags allen Spielraum läßt und\\
2. a. eine mäßige Hebung und Senkung der Stimme im Anschluß an die unwillkürliche Tonbewegung beim feierlichen Sprechen erlaubt, auch\\
b. innerhalb des beschränktesten Umfangs der Stimmintervalle eine edle Mannigfaltigkeit der melodischen Wendungen begünstigt.

Der Forderung sub 1. wird genügt durch die im lutherischen Gottesdienst noch üblich verbliebene Methode des choralmäßigen Lesens — choraliter legere —, welche nicht allein bei der Recitation der Collecten, sondern auch bei den responsorisch gesungenen Versikeln, sowie bei der Psalmodie zur Anwendung kommt, sofern sich hier die Stimme lange Zeit auf einem herrschenden Ton, der sogenannten Dominante, also in gleicher Tonhöhe rhythmisch fortbewegt,\footnote{Dies Lesen oder Singen, da die Stimme immer in ein und derselben Tonhöhe bleibt, ohne jedem melodischen Intervallwechsel, heißt nach dem Sprachgebrauch der alten Schule „in Unisono“ lesen. Wenn daher Luther in seiner „deutschen Messe und Ordnung des Gottesdienstes“ (Bd. 22, S. 237 f. A. und Bd. X, S. 239 St. L. A.) anordnet: „Darnach lieset der Priester eine Collecten in F faut, in Unisono, wie folget: „Allmächtiger Gott, der du bist ein Beschützer aller“ etc., so bedient er sich zur Angabe der Tonhöhe für den laufenden Ton, in welchem die Collecte gesungen werden soll, der damals überall gebräuchlichen, von Guido von Arezzo 1038 aufgebrachten Solmisation, nach welcher die verschiedene Höhe der Töne für Gesang durch die Silben ut, re, mi, fa, sol, la, si bezeichnet werden. (Es ist hier nicht der Ort, auf diese guidonische Methode einzugehen; es genüge, nach Lyra zu bemerken, daß das F faut oder richtig F fa ut in Unisono einfach heißt: der herrschende, der laufende Ton, der tonus currens für die Collecte soll das F auf der vierten Linie im Baßschlüssel sein. Damit will Luther jedoch nicht sagen, die Collecte dürfe nur in F gesungen werden. Er will nur den Durchschnittston für die Collecte beispielsweise angeben. Je nachdem die Stimme des Liturgen höher oder tiefer ist, kann der Durchschnittston für die Collecte auch G oder A oder B oder auch C sein. — Hiernach erklärt sich auch die andere Stelle bei Luther: „Darnach die Epistel in octavo Tono, daß er im Unisono der Collecten gleich bleibe.“ Wie an seinem Ort näher dargethan werden wird, wurden eben auch Epistel und Evangelium vom Altar aus in gregorianischer Weise gesungen, und hat eben hierfür Luther in der genannten Schrift für den Leseton der Perikopen nicht nur Regeln aufgestellt, sondern dieselben auch durch Notenbeispiele erläutert. Wenn er nun hier von einem octavo Tono redet, so meint er nicht den achten Ton der Tonleiter, sondern den achten der alten Kirchentöne oder Melodien. Nach dieser achten alten Kirchenweise solle die Epistel gesungen werden, wie das beigegebene Beispiel zeigt, die dabei einzuhaltende Höhe des laufenden Tones (Unisono) jedoch soll der Collecte gleich sein. Obwohl also nach zweierlei Weise gesungen, sollen Collecte und Epistel gleiche Tonhöhe haben.} mit sinnenentsprechender
\quellseite{96}
Verwendung der Athmungspausen nach alleiniger Maßgabe des Sprachgefühls die Gliederung des Redeflusses wirkt. Z. B.:

\notensatz{music/096-01}{Versikel: Danket dem HErrn: denn er ist freundlich etc.}

\begin{verse}
Collecte: Allmächtiger HErre Gott, der du bist etc.\\
Psalm: Wohl dem, der nicht wandelt etc.
\end{verse}

\subsubsection{Ausführung}

\notensatz{music/096-02}{Versikel: Danket dem HErrn: denn er ist freundlich etc.}

\notensatz{music/096-03}{Collecte: Allmächtiger HErre Gott, der du bist etc.}

\notensatz{music/096-04}{Psalm: Wohl dem: der nicht wandelt im Rath etc.}

Was nun die für das choralmäßige Lesen sub 2a und b aufgestellte Forderung betrifft, so kommen hier nun die besonderen Accentregeln der gregorianischen Schule in Betracht. In dem laufenden Ton, dem tonus currens, wird nämlich fortgesungen, bis durch Interpunction im Texte, durch die Distinctio oder unser Komma, durch das Semicolon und den Doppelpunkt, durch den Punkt und das Fragezeichen ein Einschnitt oder Abschnitt im Satz gemacht und dieser durch Hebung und Senkung der Stimme nach einer bestimmten Regel im liturgischen Vortrag zum Gehör gebracht wird.

1. Die Distinctio oder das, was wir jetzt Komma nennen, hat je nach der größeren oder geringeren Bedeutung des Satzes dreierlei Accente: accentus immutabilis, accentus acutus und accentus moderatus.

Beim accentus immutabilis findet weder ein Heben noch Senken der Stimme statt. Die Beobachtung des durch ein Komma gemachten Abschnitts geschieht bloß durch eine Dehnung der Endsilbe mit kurzer Pause, etwa des Athemholens:

\notensatz{music/096-05}{Accentus immutabilis: gleichbleibender Ton G mit Dehnung und Atempause}

\quellseite{97}
Beim accentus acutus kehrt die Stimme für die letzte oder vorletzte Silbe in den laufenden Ton wieder zurück, nachdem sie vorher zwei Stufen abwärts gefallen war:

\notensatz{music/097-01}{Accentus acutus: drei Formeln mit Rezitation auf G, Terzfall nach E und Rückkehr zu G.}

Beim accentus moderatus kehrt (nach Ornithoparch) bei der letzten Silbe des Einschnitts die Stimme nicht in den laufenden Ton zurück, wie beim acutus, sondern nimmt ihren Halt im Zwischenton:

\notensatz{music/097-02}{Accentus moderatus: Rezitation auf G, dann Fis–E–Fis, mit Brevis auf dem Schlußton.}

2. Das Semikolon und der Doppelpunkt (größerer Einschnitt) hat den accentus medius, indem die Stimme auf der oder den letzten Silben in die Terz herunterfiel:

\notensatz{music/097-03}{Accentus medius: Rezitation auf G, dann E–E, mit Brevis auf dem Schlußton.}

3. Das Punktum (colon) hat den accentus gravis, das heißt, der Ton fällt eine Quinte abwärts:

\notensatz{music/097-04}{Accentus gravis: viermal G, dann Quintfall nach C.}

4. Das Fragezeichen oder der Fragesatz hat den accentus interrogativus, bei welchem die Stimme auf der letzten Silbe um einen oder etliche Töne steigt oder auch um zwei Stufen vom laufenden Tone fällt und dann durch die Zwischenstufe zu demselben zurückkehrt:

\notensatz{music/097-05}{Accentus interrogativus: D–E–Fis–G–G.}

\subsubsection{Oder:}

\notensatz{music/097-06}{Alternative Frageformel: G–G–G–E–Fis–G, mit Brevis auf dem Schlußton.}

5. Schließlich rechnen Einige noch den Final-Accent, die clausula finalis, hinzu, der nach den Einen ein Hinaufgehen in die Terz mittelst der Secunde und Wiederherabsteigen, nach Andern ein allmähliches Herabsteigen in den letzten Silben bis zur Quart des laufenden Tones ist:

\notensatz{music/097-07}{Clausula finalis: G–A–H–A–G, mit Brevis auf dem Schlußton.}

\subsubsection{Oder:}

\notensatz{music/097-08}{Alternative Clausula finalis: G–Fis–E–D.}

\quellseite{98}
Nach diesen Accenten, von denen angenommen wird, daß sie mit den hebräischen Accenten in Verbindung stehen, konnten alle Stellen der heiligen Schrift gelesen werden. Da wir sie im Versikel- und Orationston noch haben, falls wir bei unserm liturgischen Singen uns nicht schlechtweg an die Tradition,\footnote{Traditionell ist z. B. namentlich die sehr verbreitete Formel g f e fis g. Lyra bezeichnet sie als incorrect und nennt sie „jene weinerliche Formel ..., welche thatsächlich noch mancherwärts beim liturgischen Recitiren von Versikeln, Collecten etc. das ‚Recept für alles‘ bildet“, und mit der als Weiterbildung von g f e g „der Proceß der willkürlichen Corruption unserer liturgischen Melodien sich vollendet“. Es sei auch in diesem Werke von ihr kein Gebrauch gemacht.} sondern an die älteren Vorbilder unserer singenden Kirche halten, so mögen folgende Beispiele die Anwendung jener Accente in etwas veranschaulichen:

\subsection{Acc. immutabilis}

\subsubsection{Salutation}

\notensatz{music/098-01}{Accentus immutabilis, Salutation: Der HErr sei mit euch! Fünfmal G, abschließend als Brevis.}

\subsubsection{Versikel}

\notensatz{music/098-02}{Accentus immutabilis, Versikel: Weise mir, HErr, deinen Weg, daß ich etc. Mit Brevis und Atempause nach Weg.}

\subsubsection{Collecte}

\notensatz{music/098-03}{Accentus immutabilis, Collecte: Allmächtiger HErr Gott, der du der Elenden Seufzen etc. Mit Brevis und Atempause nach Gott.}

\subsection{Acc. acutus}

\subsubsection{Salutation}

\notensatz{music/098-04}{Accentus acutus, Salutation und Antwort: Der HErr sei mit euch! Und mit deinem Geiste. Terzfall nach E, Rückkehr zu G.}

\subsubsection{Oder:}

\notensatz{music/098-05}{Alternative Salutation und Antwort: Der HErr sei mit euch! Und mit deinem Geiste. Abstieg über Fis nach E, Rückkehr zu G.}

\subsubsection{Versikel}

\notensatz{music/098-06}{Accentus acutus, Versikel: HErr, lehre mich thun nach deinem Wohlgefallen. Danket dem HErrn, denn er ist freundlich. Halleluja. Und seine Güte währet. Fortsetzung auf Scan 99.}

\quellseite{99}
\notensatz{music/099-00}{Fortsetzung des Versikels: Ewiglich. Halleluja. E–G–G, dann G–A–G–E, mit Breven auf den beiden Phrasenenden.}

\subsubsection{Collecte}

\notensatz{music/099-01}{Accentus acutus, Collecte: Allmächtiger, ewiger Gott, der du die ganze Christenheit heiligest und regierest etc. Breven mit Fermaten auf Gott und regierest, Fis auf und.}

\subsection{Acc. moderatus}

\subsubsection{Präfation}

\notensatz{music/099-02}{Accentus moderatus, Präfation: Wahrhaft würdig und recht, billig und heilsam ist’s, daß wir etc. Breven auf recht und ist’s.}

\subsection{Acc. medius}

\subsubsection{Versikel}

\notensatz{music/099-03}{Accentus medius, Versikel: Rufe mich an in der Noth, so will ich dich erretten, und du sollst mich preisen. Brevis und Atempause nach erretten.}

\subsubsection{Collectenschluß:}

\subsubsection{Clausula finalis}

\notensatz{music/099-04}{Clausula finalis, Collectenschluß: Durch JEsum Christum, deinen Sohn, unserm HErren. Um JEsu Christi, unsers HErren, willen. Deinen Sohn, unserm HErren. Mit Breven und Melismen auf wil und dem letzten HErr.}

Nach alledem noch einige Winke behufs der Weise des musikalischen Vortrags. Soll derselbe nämlich der Gemeinde zur Erbauung dienen, so ist noch folgendes zu beachten:

1. Geschrei und Geplärr, das kein Gesang ist, muß streng vermieden werden, namentlich wenn die Melodie in höhere Töne steigt. Stärker kann die Stimme in der Tiefe hervortreten, doch darf sie auch hier nicht in Brüllen ausarten. Ein ungezwungener, ohne besondere Anstrengung in gleichmäßiger Stärke hinfließender Gesang ist zur Erbauung der beste. Eben darum soll jeder Liturg den Umfang seiner Stimme durch Uebung genau kennen lernen und deshalb sich über Höhe oder Tiefe der Intonation zuvor mit dem Organisten
\quellseite{100}
verständigen, wo derselbe den Liturgen oder allein die responsorirende Gemeinde mit der Orgel zu begleiten hat. Nur eine natürliche Stimme, nur Bruststimme, und nicht eine solche, die wie aus einem zusammengeschnürten Halse gequetscht wird!

2. Die Buchstaben und Silben sollen rein ausgesprochen werden und ganz besonders bedarf die Aussprache der Vocale einer sorgfältigen Uebung. Die Tonleitern abwechselnd mit den fünf Vocalen durchzusingen, führt am sichersten zum Zwecke. Das a soll hell klingen, nicht wie o oder oa; das e nicht wie ae; das i nicht wie e oder ae; das o nicht wie a und das u nicht wie o. Also: Vater, nicht Voater; denn, nicht dänn; Christum, nicht Chrestum; durch, nicht dorch. Eben so wenig darf man am Anfang eines Wortes einen Buchstaben hinzusetzen, z. B. A-der-HEerr segne dich etc., oder die eine Silbe in zwei Silben spalten, wie gleichfalls beim aaronischen Segen manchmal zu vernehmen ist, nämlich: Der Her-herr segne dich etc.

3. Die Haltung des ganzen Körpers sowohl als einzelner Theile desselben sei anständig und würdevoll. Der Mund darf nicht zu weit aufgerissen oder seitwärts gezogen, sowie überhaupt das ganze Gesicht nicht verzerrt werden.

4. Wo immer möglich, bereite sich der Liturg nicht nur für die Auswahl der betreffenden liturgischen Stücke, der Versikel, Collecten etc. vor, sondern auch auf den Vortrag, so z. B. der Collecten betreffs der Accente.

\section[§ 3. Gebrauch der Orgel]{§ 3. Und nun noch ein Wort über den Gebrauch der Orgel beim liturgischen Gesang}


Wohl ist der gregorianische Choralgesang überhaupt und der priesterliche Altargesang insbesondere seinem Ursprung und Wesen nach nicht, wie bereits bemerkt, für die Harmonie, also auch nicht für harmonische Begleitung der Orgel berechnet und bestimmt, daher auch beim liturgischen Gesang in der That viel später, als beim kirchlichen Gesang im Allgemeinen der Gebrauch der Orgel eingeführt ist. Immerhin meinen Männer wie Kraußold, Lyra, Herzog, Schöberlein, daß derselbe unter Voraussetzung einer besonnenen Registerwahl und einer stilmäßigen Begleitung zugelassen und vertheidigt werden könne. Da am eingehendsten sich Kraußold hierüber ausspricht, so möge aus dessen historisch-musikalischem „Handbuch für den Kirchen- und Choralgesang“ der betreffende
\quellseite{101}
Paragraph, der zugleich den Schluß des ganzen werthvollen Werkchens bildet, fast unverkürzt hier stehen. Es heißt daselbst:

\begin{quote}
„Der Zweck des Altargesangs scheint der harmonischen Begleitung der Orgel zu widersprechen, indem in dem modus choraliter legendi eben das ausgedrückt werden sollte, daß das eigentlich musikalische Moment, das Melodiöse, mehr zurücktrete, wogegen beim Chorgesang, bei den Psalmen und Hymnen etc. das musikalische Moment des Sängers mehr hervortrat. Indessen hat sich zwischen dem strengen modus choraliter legendi und dem Psalmen- und Hymnengesang eine Mittelgattung einzelner Altargesänge ausgebildet, welche sowohl in der katholischen als protestantischen Kirche die Orgelbegleitung zuläßt. Das strenge Chorallesen nämlich blieb vorzugsweise den biblischen Lectionen, der historischen Erzählung, und den Collectengebeten eigen. Dagegen erhoben sich einzelne liturgische Stücke, in denen sich mehr ein bestimmtes Gefühl ausspricht, die mehr ein Erguß eines religiösen Gefühls sein sollten, aus dem bloßen Leseton zu bestimmten Modulationen, zu melodiösem Ausdruck. So die Intonation des Kyrie und des Gloria, dann das Vater Unser, die Einsetzungsworte, die Präfation. So wenig nun die eigentlichen Lectionen (Epistel- und Evangelienton) und Collecten ihrem Inhalt und Zweck nach zur Orgelbegleitung sich eigneten, so sehr luden dagegen die letztgenannten melodiösen Stücke dazu ein, und nur wo ein eigentlicher Wechselgesang zwischen dem Geistlichen und der Gemeinde eintrat, blieb die Orgelbegleitung bei ersterem weg. So hat sich denn die traditionelle Praxis gebildet, daß zum Vater Unser, Einsetzungsworten, Präfation etc., aber nicht zu den Responsorien, der Gesang des Geistlichen mit der Orgel begleitet wurde. Wenn nun die Frage gestellt wird, ob überhaupt eine Orgelbegleitung beim Altargesang stattfinden soll, oder nicht, so muß allerdings festgehalten werden, daß sie weder wesentlich noch nothwendig ist, daß sie eben so gut angewendet werden, als wegbleiben kann, und es wird lediglich darauf ankommen, die Gründe zu erwägen, die etwa einer solchen Begleitung der Orgel das Wort reden möchten. Und da wird gewöhnlich als erster Grund angeführt, daß die Orgelbegleitung den Zweck habe, den Geistlichen im Gesang zu unterstützen und zu halten. Wir wollen nicht leugnen, daß manche Geistliche in der Orgelbegleitung wirklich einen Halt für ihre Stimme finden mögen, aber eine vielfache Beobachtung und Erfahrung hat uns gelehrt, daß dies eine große Beschränkung erleidet. Als einzigen Grund können wir dies schon von vorn herein nicht gelten lassen, weil dann bei einem musikalisch gebildeten Geistlichen, der solche Unterstützung nicht bedarf, die Orgelbegleitung
\quellseite{102}
zwecklos wäre und wegfallen müßte, während im Gegentheil gerade bei einer guten musikalisch gebildeten Stimme die Orgelbegleitung recht zweckmäßig und wohlthätig erscheint. Wir setzen voraus und fordern es als Pflicht des Geistlichen, daß er seinen Altargesang fest inne habe, daß er nicht von der Orgel geführt werde, und gleichsam erst von den Orgeltönen seine Töne treffen lernen wolle. Er muß seines Gesanges aus vorausgegangener Einübung so sicher und gewiß sein, als seiner Predigt. Die Orgelbegleitung hätte aber bei diesem Zweck überhaupt nicht sowohl der Stimme als solcher zu Hülfe zu kommen, als vielmehr dem Gehör des Singenden. Ist nun der Geistliche mit einem feinen musikalischen Gehör begabt, so kann ihm allerdings, wie gesagt, die Orgelbegleitung gute Dienste leisten. Denken wir uns aber das Gegentheil — und das ist gewiß nicht selten der Fall — denken wir uns einen Geistlichen mit schlechtem musikalischen Gehör, so müßte die Orgel, wenn sie ihm etwas helfen sollte, so stark gespielt werden, daß sie den Gesang und die Worte übertönt — und das wäre dem Zwecke des Altargesangs (des choraliter Lesens) ganz entgegen, da das Wort bei dem Gesange die Hauptsache sein und bleiben muß —, oder der Gesang des Geistlichen und die Orgeltöne stehen in einem solchen Mißverhältniß, daß es nicht auszuhalten ist. Wir erinnern uns aus unserer frühesten Jugendzeit mehr als einmal, daß der Organist z. B. bei dem Consecrationsgesang des Geistlichen mitten drin die Tonart um einen halben Ton wechseln und transponieren mußte, um den Geistlichen — was er ja sollte — bei seinem Gesang begleiten zu können. Thut das der Organist nicht, sondern bleibt er in seiner Tonart, während der Geistliche um einen oder einen halben Ton detonirt oder hinaufgeht, meint der Organist, sein Spiel sei bestimmt, nicht sowohl den Geistlichen zu begleiten, als ihn zu leiten, im Ton zu halten und auf den rechten Ton zurückzuführen, so tritt eben jenes ohrenzerreißende und alles musikalische Gefühl empörende Mißverhältnis ein, das dem Hörer den unaufhörlichen stillen Wunsch auspreßt: O, daß doch die Orgel schwiege! Und wir müssen es offen aussprechen, daß, wo der Geistliche seinem Gehör nicht trauen darf, er besser thut, wenn er ja singen will, daß er ohne Orgelbegleitung singe. Es soll ja seinem Zweck nach immerhin nur ein singendes Lesen sein, und kann als solches weit eher mit einigen musikalischen Irregularitäten ohne Orgelbegleitung bestehen, als mit Orgelbegleitung.

„Offenbar aber hat die Orgelbegleitung noch einen andern Zweck, nämlich den, die Feierlichkeit zu erhöhen und die Andacht zu heben. Das liegt einmal in der unleugbaren wunderbaren Wirkung des Orgeltons überhaupt,
\quellseite{103}
und dann eben in der der Einzelstimme beigegebenen Harmonie, die für das menschliche Gemüth neben der Melodie das eigenthümliche beruhigende Moment ist, das der religiösen Stimmung des Gemüths so wohlthuend zusagt. Aber eben darin liegt es auch, daß Melodie und Harmonie, das ist, die Melodie führende Stimme des Geistlichen und die harmonische Begleitung der Orgel zusammenstimmen müssen, wenn nicht das Gegentheil von dem herauskommen soll, was mit der Orgelbegleitung bezweckt wird.

„Wenn der Geistliche seinen Gesang nicht fest und sicher inne hat, wenn er erst durch die Orgel geleitet und geführt werden will, so tritt ein anderer Uebelstand ein, der uns auf die Art und Weise der Orgelbegleitung beim Altargesang des Geistlichen führt. Diese muß sich nämlich von der Begleitung des Chor- und Gemeindegesangs wesentlich unterscheiden. Der Gesang des Geistlichen ist Einzelgesang, und zwar recitirender Einzelgesang. Das Singen soll eigentlich mehr ein musikalisches Sprechen für die Gemeinde sein. Die Gemeinde soll nicht bloß die Töne des Gesangs, sondern vor Allem und zumeist die Worte hören, welche gesungen werden. Wollte nun die Orgel jeden Ton des Geistlichen mit allen Accorden begleiten, oder gar auch mit stärkeren Registern spielen, so würde gerade, wie schon erwähnt, die Hauptsache verloren gehen, und es ist besser, ohne Orgelspiel zu singen, oder auch selbst nur zu sprechen, als wenn durch das Orgelspiel die eigentliche Absicht des Altarsingens verloren geht. Eine solche durchgehende Orgelbegleitung ist aber auch schon darum unthunlich, weil der Altargesang nicht mensurirt ist, nicht in einem bestimmten Tact, sondern nur nach dem rhetorischen Rhythmus sich bewegen kann. Die einzelnen Noten werden nicht nach ihrem bestimmten Maß und Gehalt vorgetragen, sondern nach der Länge und Kürze der Silben, nach dem hervorzuhebenden größern oder geringern Nachdrucke der einzelnen Worte. Darnach muß sich denn auch die Orgelbegleitung durchaus richten, wenn sie ihrem Zwecke entsprechen soll. Vier Punkte sind es, die hierbei vorzugsweise in Betracht kommen und hier schließlich zur genauen Beobachtung empfohlen werden.

„1. Der Organist nehme zur Begleitung des Geistlichen im liturgischen Gesang nie starke Register; je sanfter und lieblicher, desto besser.

„2. Der Organist spiele die einzelnen Accorde nie in ihrer vollkommenen Zeitdauer oder im strengen Tact, sondern genau nach dem recitirenden Gesang des Geistlichen. Wir wiederholen: er soll ihn nicht leiten, sondern begleiten. Daher muß der Organist den Geistlichen hören
\quellseite{104}
während des Spiels; und um ihn hören zu können, dazu dienet neben den obenerwähnten leichteren Registern insbesondere dies, daß er seine Accorde nicht voll austönen lasse, sondern abbreche und dem Gesang Raum lasse.

„3. Der Gesang des Geistlichen werde überhaupt nicht in continuo begleitet. Wo kein Wechsel der Harmonie ist, wo oft mehrere Töne der Melodie denselben Accord zur Grundlage haben, schweige die Orgel ganz. Nur der Schluß verlangt vollen harmonischen Gang.

„4. Endlich vermeide der Organist das Pedal. Im Gemeindegesang hat das Pedal seine Wirksamkeit, wo es gilt, eine große Masse Singender im rechten Ton und Tact zu halten und die Grundlage (Baß) des Gesanges recht bemerklich zu machen. Im Einzelgesang des Geistlichen ist das Pedal überflüssig und höchst lästig; außer etwa wieder am Schlusse. —

„Es ist unglaublich, wie durch Nichtbeachtung dieser Punkte der Altar-Gesang dem Geistlichen erschwert wird, und zwar, wie sich von selbst versteht, in demselben Maße, als eine größere Kirche dazu kommt. Das ist wahr, wenn der Organist so ungeschickt ist, mit starken Registern und vollen Accorden, ohne Unterbrechung mit des Pedalbasses Grundgewalt, den Gesang des Geistlichen zu begleiten — eine vierfache Plage, davon eine schon lästig genug ist! — da ist es nicht zum Aushalten im vollen Sinne des Wortes, nämlich weder von Seiten des Zuhörenden, noch von Seiten des Singenden. Und wer seine Brust lieb hat und nicht bei Zeiten nicht bloß für den Altargesang, sondern auch für sein Amt überhaupt untauglich werden will, der habe Acht auf diese vier Punkte. Ein Geistlicher, der die Altarliturgie singen soll unter dem Drucke dieser vierfachen Orgellast, hat seine physische Kraft für die Predigt schon halb verzehrt! Weg mit der Orgelbegleitung, wenn sie belästigend sein will! Lieber singe er ganz ohne Orgel! und man wird die Klage nicht mehr hören, daß die Liturgie zu anstrengend sei.“
\end{quote}

\chapter{I. Introitus}


\quellseite{105}
Die römische Messe beginnt mit dem Initium missae, dem Eingangsgebet des Priesters aus dem 43. Psalm, und dem Confiteor, der öffentlichen Beichte, worauf dann erst der Introitus, der Eingangs- oder Eröffnungsgesang, folgt. Mit Weglassung jenes Eingangsgebetes beginnt auch in der Brandenb.-Nürnb. K.-O. 1533, in der Märkischen K.-O. 1540 und in noch manch anderer der ersteren K.-O. sich anschließenden Agende der Hauptgottesdienst mit dem Confiteor, während die an Luthers deutsche Messe sich mehr oder weniger anschließenden sächsischen und meistentheils auch die späteren Agenden die Beichte nach der Predigt setzen.

Den Namen Introitus (Einzug, Eingang) hat dieser Gesang von der Weise des christlichen Alterthums, nach welcher man, wie es jetzt auch wieder ritualistischerseits in der Episcopalkirche geschieht, in Procession und mit Gesang in die Kirche einzog. Dieser Gesang bestand in einem Psalm, welchen die Kirchendiener mit den Choristen antiphonatim sangen, und einer Antiphonie oder einem Gegengesang der Gemeinde zwischen einzelnen Psalmversen, welche Antiphonie aber immer nach der Bedeutung des Sonn- oder Festtages sich richtete. Als nun später dieser feierliche Einzug abkam, behielt man die Antiphonie nebst dem ersten Vers des betreffenden Psalms bei und beließ diesem Gesang, weil Eröffnungsgesang, den Namen „Introitus“, weshalb auch Luther sagt: „Nachdem aber das Singen in der Kirche aufkommen, hat man die Psalmen verwandelt in den Introitum.“

Sowohl in seiner „Weise, christlich Meß zu halten, 1523“, als auch in seiner „deutschen Messe etc., 1526“, will Luther den Introitus als Eingangsgesang überhaupt beibehalten wissen. In ersterer heißt es: „Die Introite der Sonntage, und so man singt auf die Feste Christi, als Ostern, Pfingsten, Weihnachten, loben wir, halten sie auch; obwohl uns die Psalmen lieber wären, aus welchen sie genommen sind, wie ehemals gehalten ward; doch wollen wir hierin dem gemeinen Brauch weichen.“ (St. L. Ausg. X, 2238.) In der letzteren jedoch trägt er schon seinem Wunsche etwas Rechnung: „Zu Anfang aber singen wir ein geistlich Lied oder einen deutschen Psalmen in primo Tono, auf die Weise, wie folget.“ (Erl. Ausg. 22, 237; St. L. Ausg. X, 235 ff.) Es folgt dann der ganze 34. Psalm: „Ich will
\quellseite{106}
den HEerrn loben allezeit“, in Noten nach dem ersten der acht kirchlichen Psalmtöne. Luther ordnet also für den Introitus „ein geistlich Lied oder einen deutschen Psalm“. Nach Luthers Vorgang lassen denn auch die übrigen den Introitus dem Kyrie vorangehen. Jedoch in mancherlei Form und Weise. Viele der Agenden, besonders der nachreformatorischen, ordnen für den Inhalt „ein geistlich Lied“, namentlich „Komm, Heiliger Geist, erfüll“ etc. oder „Komm, Heiliger Geist, HEerre Gott“, während gleichwohl wieder die Braunschweig-Lüneburger Ordnung 1657 ordnet: „Die Meß oder Communion auf den Dörfern soll der Pfarrer nebenst dem Custode und den Schülern anfahen mit einem deutschen Psalm, darnach das deutsche Kyrie: O Vater, allmächtiger Gott“; desgleichen sollen nach der Leipziger Agende 1681 „die Schüler singen erstlich (zu Anfang) den Introitum von der Dominica (dem Sonntag) oder Festen, darauf das Kyrie eleison“ etc.

Der Introitus besteht, wie bereits angedeutet, aus der sogenannten Antiphona ad introitum, dem Psalmus, einem Vers oder ein paar Versen eines Psalms, und der kleinen Doxologie, dem Gloria patri, dem „Ehre sei dem Vater“ etc., jener schon aus den Kämpfen mit den Arianern stammenden doxologischen Bekenntnißformel. So z. B. der Introitus des ersten Weihnachtstages:

\begin{verse}
Antiphona: Uns ist ein Kind geboren, ein Sohn ist uns gegeben, welches Herrschaft ist auf seiner Schulter, und er heißet: Wunderbar, Rath, Kraft, Held, Ewigvater, Friedefürst.\\
Psalmus: Singet dem HEerrn ein neues Lied, denn er thut Wunder.\\
Gloria Patri: Ehre sei dem Vater und dem Sohn und dem Heiligen Geiste, wie es war im Anfang, jetzt und immerdar und von Ewigkeit zu Ewigkeit. Amen.
\end{verse}

Der Introitus soll der Herold des gottesdienstlichen Tages sein, welcher der versammelten Gemeinde gleich von vornherein die Bedeutung des Sonntages, des Festes, der Festzeit verkündet. Sein Schwerpunkt fällt daher in die Antiphone, indem dieselbe mit einem kurzen Schriftwort die Bedeutung des Tages angibt und die zu feiernde Heilsthatsache verkündet. Jeder Sonntag und Festtag hatte seinen besonderen Introitus und manche Sonntage haben ihre Namen von dessen erstem lateinischen Worte, wie Estomihi, Invocavit, Reminiscere oder Quasimodogeniti, Misericordias Domini, Jubilate etc. Gemeinden, die nicht dem heutigen Gelüsten nachgeben, den Gottesdienst immer mehr abzukürzen, sondern die für die
\quellseite{107}
Dauer eines vollen Hauptgottesdienstes etwas mehr Ausdauer von früherer Zeit her sich bewahrt haben, wäre die Wiederaufnahme des Introitus zu empfehlen. Dies könnte am besten nach kirchlichem Vorgang durch sanglustige und sangfähige Schüler in Ausführung gebracht werden. Können sie auch nicht mehrstimmig singen, so thut’s auch ein zweistimmiges, ja auch ein einstimmiges Singen, und klingt dasselbe durch die hellen Kinderstimmen gar lieblich; und reichen unsere Kräfte auch nicht mehr aus für die reichere Weise des psalmenartigen oder des hymnenartigen festlichen Introitusgesanges früherer Zeiten, so thäten’s für alle Introiten auch ein paar der Psalmentöne: einer für die Sonntage, einer für die Passionszeit und einer für die Festtage. Auf diese Weise, die sich von dem kirchlichen Vorgang nicht entfernte, brächte man zugleich wieder etlichermaßen den so abhanden gekommenen Psalmengesang in unsere Gottesdienste, wie dazu auch ein ähnlicher Anfang in Endlichs Choralbuch gemacht ist. Auch Löhe schreibt: „Noch ist zu bemerken, daß der Introitus am besten nach einem der 8 Psalmentöne gesungen wird, so lange die alten Melodien der Introiten unserer Kirche nicht wieder erstattet sind.“ Es geschähe dies sogar nach Luthers Vorgang. Derselbe nahm bekanntlich in seiner „deutschen Messe und Ordnung des Gottesdienstes“ für den Introitus den 34. Psalm, gab ihm den 1. Psalmenton und fügte die Noten bei (S. Luthers Werke, St. L. Ausg. X, S. 235 ff. oder die musikalische Beilage zu Band 22 der Erl. Ausg.)

Es seien daher die sämmtlichen Introiten der älteren Zeit hier mitgetheilt. Versuchsweise sind denselben die Singweisen des I., III., V. und VIII. Psalmentons beigegeben nebst einem zu der Orgelbegleitung passenden zweistimmigen Satze für einen Kinderchor.

Da es sich um die so unbekannt gewordene Weise des Psalmensingens handelt, so möge hier einiges zur Erklärung und Anweisung folgen, soweit es hierher gehört.

1. Der Anfang eines jeden Psalmentons, wie er an Sonn- oder Festtagen gesungen wird, heißt die Intonation; der Ton, in welchem darauf in gleichbleibender Höhe fortgesungen wird, die Dominante, der herrschende oder der laufende Ton und die Töne, welche in der ersten Hälfte des Verses auf die Dominante folgen, die Mediation, die aber, welche in der zweiten Hälfte folgen und die Melodie beschließen, heißen Termination oder Differenz. Hier eine Veranschaulichung beispielsweise nach dem I. Ton:

\quellseite{108}
\subsubsection[Intonation. Dominante. Mediante]{Intonation. Dominante. Mediante. Dominante. Termination oder Finale}

\notensatz{music/108-01}{Erster Psalmenton: Intonation, Dominante und Mediante, dann Dominante und Termination. Wohl denen, die ohne Wandel leben, die im Gesetz des HErren wandeln.}

2. Obwohl hier wie anderwärts beim gregorianischen Gesang die Noten von gleichem Zeitwerth sind, so dürfen es doch nicht die Töne nach demselben sein; es ist hier das Singen kein syllabirendes Singen, gleich dem Lesen eines Buchstabirschülers, sondern ein recitativisches Singen, ein accentuirtes, singendes Lesen. Der Rhythmus dieser Psalm-melodien (wie auch der Versikel, Collecten etc.) richtet sich nach dem Accent des Textes. Die allgemeine Regel ist hier, daß man die Intonation langsam anhebt, die Dominante bei öfterer Wiederholung etwas rascher, den Schluß beider Hälften, Mediation und Termination, wieder ruhiger und gedehnter singt. (Es ist nämlich die Schreibweise diese:

\subsubsection{Intonation. Dominante. Mediation}

\notensatz{music/108-02}{Schreibweise des ersten Psalmentons, ohne Textunterlegung: Intonation F–G–A, Dominante A und Mediation B–A–G–A.}

die Ausführung aber ohngefähr diese:

\notensatz{music/108-03}{Rhythmische Ausführung des ersten Psalmentons. 1. Ehre sei dem Vater und dem Sohne. 2. Wie es war im Anfang: jetzt und immerdar. Die Cadenz lautet B–A–G–G–G–A; beide Verse haben unterschiedliche Melismen.}

\subsubsection{Schreibweise:}

\subsubsection{Dominante. Termination}

\notensatz{music/108-04}{Schreibweise der zweiten Psalmenhälfte, ohne Textunterlegung: zehnmal A als Dominante, dann G–F–G–A–G als Termination.}

\subsubsection{Ausführung:}

\notensatz{music/108-05}{Rhythmische Ausführung der zweiten Psalmenhälfte, mit gedehntem Schluß. 1. Und dem Heiligen Geiste. 2. Und von Ewigkeit zu Ewigkeit. Amen. Sechs Rezitationsnoten auf A, dann G–G–G–F–G–A–G; die Verse teilen dreizehn Noten mit unterschiedlicher Silbenunterlegung.}

\quellseite{109}
3. Wie bei jedem Psalm, so wird auch bei dem nach dem Psalmenton gesungenen Introïtus nur für den Anfang des Introïtus und den des Gloria Patri die Intonation gebraucht, alle übrigen Zeilen aber sind mit der Dominante zu beginnen.

4. Der Introïtus wird von 2 Chören im Wechsel (antiphonatim) gesungen, wobei der zweite Chor etwas stärker als der erste zu besetzen ist.\footnote{In den einzelnen Psalmversen tritt nämlich nicht nur in der Regel eine Gliederung des Verses in zwei einander ergänzende Halbverse und ein Parallelismus der darin ausgesprochenen Gedanken zu Tage, sondern es steht auch gewöhnlich der zweite Halbvers zu seinem Vorgänger im Verhältnis der Steigerung. Z. B.:\\
Lobe den HEerrn, meine Seele,\\
Und was in mir ist seinen heiligen Namen!\\
Der erste Halbvers ist nur der Positiv der hier anklingenden Stimmung, der zweite dessen Comparativ, denn nun soll nicht mehr allein des Psalmisten Seele, sondern alles, was in ihm ist, seinen heiligen Namen loben. Diese Steigerung ist der Grund, daß der andere, den zweiten Halbvers singende Chor stärker zu besetzen ist, und daß auf der Orgel für denselben schon ein etwas stärkeres Register allemal gezogen werden kann.} Das Gloria Patri aber, falls man nicht auch hier den Wechsel durchführen will, kann von beiden Chören zusammen gesungen werden, was sich recht gut ausnimmt. In jedem Fall sind bei dem Gloria Patri allemal stärkere Register auf der Orgel zu ziehen. Dasselbe ist auch zur Reverenz stehend zu singen (Preis des Dreieinigen).

5. Die in den Texten der Introiten durch den Druck hervorgehobenen Silben der Mediation und Termination zeigen den Eintritt derselben an. Der Doppelpunkt (:) gibt an, wo im Singen der Dominante wie zum Athemholen einen Augenblick abzusetzen ist, was gerade beim Psalmensingen im Chor oft nöthig wird.\footnote{Siehe hierbei auch: „Kleines Psalterium. Etliche Psalmen zum Singen in Christenlehren und andern Nebengottesdiensten, wie auch in Jünglings- und Jungfrauenvereinen und im Familienkreise. Ein geringer Beitrag zur Wiederaufnahme und zur Beförderung des altkirchlichen Psalmen-Gesanges in den Gemeinden der lutherischen Kirche von Friedrich Lochner, luth. Pastor. (Im Selbstverlage.) Milwaukee, Wis. 1889.“}

Da man, wo eine Orgel vorhanden ist, den ganzen Gottesdienst durch ein Orgelpräludium zu eröffnen pflegt, so dürfte in Bezug darauf wohl schließlich eine Erinnerung am Orte sein.

1. Während alle Präludien zu den Zwischenliedern, wie z. B. das Gemeindelied vor und nach der Predigt, möglichst kurz sein müssen, da jene
\quellseite{110}
kein organisches Glied im Hauptgottesdienst bilden, sondern mehr nur als Ueberleitung zum Gemeindegesang nebst der Begleitung desselben ihren Dienst zu thun haben, so ist dieses Eröffnungspräludium eigentlich nicht als Vorspiel zum Eingangsgesang, sondern als musikalische Einleitung zum Hauptgottesdienst selbst zu betrachten. Als solches kann es bei weitem länger sein, als jene, nur wenige Takte zu umfassenden Zwischenpräludien; als solches soll es aber auch nicht so sehr dem Charakter des Eingangsliedes, als vielmehr dem Charakter des ganzen Haupt- (wie auch Neben-)gottesdienstes entsprechend sein. Von vornherein soll daher dies Präludium die versammelte Gemeinde durch seine Art und seine Registrirung freudig, festlich stimmen, und das nicht nur am hohen Festtag, sondern auch am gewöhnlichen Sonntag, dem wöchentlichen Gedächtnistag der Auferstehung des HEerrn, dem „Tag des HEerrn“, an welchem zum Ausdruck der Freude die alte Kirche ja nicht knieend, sondern stehend zu beten pflegte. Nur an den Sonntagen der Fastenzeit oder an einem besonderen Sonntag, wie z. B. dem 10. Sonntag nach Trinitatis, dem Gedächtnistag der Zerstörung Jerusalems, dürfte der Charakter dieses Präludiums ein mehr ernster sein.

2. Ist nun aber dies der Zweck des den ganzen Gottesdienst eröffnenden Präludiums, so ist um so mehr zu wünschen, daß der Organist, dem sowohl die contrapunktische Bildung, als auch die Gabe der Erfindung abgeht, sich allein an eine gute Präludiensammlung halte, vorzugsweise an eine solche, welche dem reinen Kirchenstil Rechnung trägt, wie z. B. außer den Werken von Herzog und dem „Präludienbuch von J. Zahn“ Prof. Riegels „Praxis organoedi in ecclesia“ (Kirchliches Orgelspiel), das Lyra, wie „Siona“ als ein „an klassischen Orgel-Präludien aus der Vergangenheit und Gegenwart ungemein reiches Werk“ allen Organisten aufs wärmste empfehlen, was vorliegende Proben dem Schreiber dieses bestätigen. (Auch auf die nunmehr vollständig zu St. Louis, Mo., erschienene Sammlung sei hier aufmerksam gemacht. Sie ist betitelt: „Vorspielbuch. Ein Magazin von Orgelvorspielen aus alter und neuer Zeit zu den gebräuchlichen Chorälen der ev.-luth. Kirche. Gesammelt und herausgegeben von Dietrich Meibohm.“)

\quellseite{111}
\section{Psalmentöne zum Introitus}

\subsection{Erster Ton}

\notensatz{music/111-01}{Erster Psalmenton, I. Chor, mit Orgelbegleitung. 1. Nach dir, HErr, verlanget mich. 2. Laß mich nicht zu Schanden werden. 3. Denn keiner wird zu Schanden. 4. HErr, zeige mir deine Wege.}

\notensatz{music/111-02}{Erster Psalmenton, II. Chor, mit Orgelbegleitung. 1. Mein Gott: ich hoffe auf dich. 2. Daß sich meine Feinde: nicht freuen über mich. 3. — — der dein harret. 4. Und lehre mich deine Steige.}

\notensatz{music/111-03}{Erster Psalmenton, beide Chöre, mit Orgelbegleitung. Ehre sei dem Vater und dem Sohne, und dem Heiligen Geiste. Wie es war im Anfang: jetzt u. immerdar, u. v. Ewigkeit zu Ewigkeit. Amen.}

\quellseite{112}
\subsection{Dritter Ton}

\notensatz{music/112-01}{Dritter Psalmenton, I. Chor, mit Orgelbegleitung. 1. Sei mir ein starker Fels und eine Burg. 2. Denn du bist mein Fels und meine Burg. 3. HErr: auf dich traue ich: laß mich nimmermehr zu Schanden werden.}

\notensatz{music/112-02}{Dritter Psalmenton, II. Chor, mit Orgelbegleitung. 1. Daß du mir helfest. 2. Um deines Namens willen: wollest du mich leiten und führen. 3. Errette mich: durch deine Gerechtigkeit.}

\notensatz{music/112-03}{Dritter Psalmenton, beide Chöre, mit Orgelbegleitung. Ehre sei: dem Vater: und dem Sohne: und dem Heiligen Geiste. Wie es war im Anfang: jetzt und immerdar: u. von Ewigkeit zu Ewigkeit. Amen.}

\quellseite{113}
\subsection{Fünfter Ton}

\notensatz{music/113-01}{Fünfter Psalmenton, I. Chor, mit Orgelbegleitung. 1. Uns ist ein Kind geboren, 2. Welches Herrschaft ist auf seiner Schulter. 3. Singet dem HErrn ein neues Lied,}

\notensatz{music/113-02}{Fünfter Psalmenton, II. Chor, mit Orgelbegleitung. 1. Ein Sohn ist uns gegeben. 2. Und er heißt Wunderbar: Rath: Kraft: Held: Ewigvater, Friedefürst. 3. Denn er thut Wunder.}

\notensatz{music/113-03}{Fünfter Psalmenton, beide Chöre, mit Orgelbegleitung. Ehre sei: dem Vater und dem Sohne: und dem Heiligen Geiste. Wie es war im Anfang: jetzt und immerdar: und von Ewigkeit zu Ewigkeit. Amen.}

\quellseite{114}
\subsection{Achter Ton}

\notensatz{music/114-01}{Achter Psalmenton, I. Chor, mit Orgelbegleitung. 1. Er ist auferstanden: was suchet ihr: den Lebendigen bei den Todten: Halleluja! 2. Des Menschen Sohn: muß gekreuzigt werden. 3. Mit Ehre und Schmuck: hast du ihn gekrönet.}

\notensatz{music/114-02}{Achter Psalmenton, II. Chor, mit Orgelbegleitung. 1. Gedenket daran: wie er euch sagte: Halleluja! 2. Und am dritten Tage auferstehen: Halleluja: Halleluja! 3. Und hast ihn zum HErrn gemacht: deiner Hände Werk.}

\notensatz{music/114-03}{Achter Psalmenton, beide Chöre, mit Orgelbegleitung. Ehre sei: dem Vater: und dem Sohne: und dem Heiligen Geiste. Wie es war im Anfang: jetzt und immerdar: und von Ewigkeit zu Ewigkeit. Amen.}

Anmerkung. Nach ein paar überleitenden Accorden folgt unverweilt auf den Introitus das Kyrie, sei es das liedweise oder antiphonische Kyrie.

\quellseite{115}
\section[Introiten für die Sonn- und Festtage]{Introiten für die Sonn- und Festtage des Kirchenjahres}

Anmerkung. Die erste Zeile singt der erste Chor, die zweite Zeile der zweite Chor. Zum Schluß singen die beiden Chöre zusammen die Doxologie, wie die Musikbeilage zeigt. Unnöthige Wiederholung zu vermeiden, ist daher hier nur bei dem ersten Introitus die Doxologie angegeben.

\section{Erster Adventssonntag}

Ps. 25, 1–4.

\subsubsection{I. Ton}
\begin{verse}
Nach dir, HErr: ver-lang-et mich.\\
Mein Gott: ich hof-fe auf dich.\\
Laß mich nicht zu Schan-den werden,\\
Daß sich meine Feinde: nicht freu-en über mich;\\
Denn keiner wird zu Schan-den\\
. . . . . der dein har-ret.\\
HErr, zeige mir dei-ne Wege,\\
Und lehre mich dei-ne Steige.
\end{verse}

\subsubsection{Beide Chöre}
\begin{verse}
Eh-re sei dem Vater: und dem Sohne\\
Und dem Hei-li-gen Gei-ste,\\
Wie es war im Anfang: jetzt und immer-dar,\\
Und von Ewig-keit: zu Ewig-keit. Amen.
\end{verse}

\section{Zweiter Adventssonntag}

Jes. 62, 11. 30, 29. 30. Ps. 80, 2.

\subsubsection{I. Ton}
\begin{verse}
Toch-ter Zion, siehe: der HErr wird kom-men,\\
Zu hel-fen den Völ-kern.\\
Er wird seine herrliche Stimme schal-len las-sen,\\
Und euer Herz: wird sich freu-en.\\
Du Hirte Isra-els, hö-re,\\
Der du Josephs hütest wie der Schafe.
\end{verse}

\section{Dritter Adventssonntag}

Phil. 4, 4. ff. Ps. 85, 2.

\subsubsection{I. Ton}
\begin{verse}
Freu-et euch in dem HErrn al-le-we-ge,\\
Und abermal sage ich euch: Freu-et euch.\\
\quellseite{116}
Eure Lindigkeit : lasset kund sein al-len Men-schen.\\
Der HEerr ist na-he.\\
Sorget nichts : sondern in al-len Din-gen\\
Lasset eure Bitte : mit Gebet und Flehen : vor Gott kund wer-den.\\
HEerr : der du vormals bist gnädig gewesen : dei-nem Lan-de\\
Und hast die Gefangenen Ja-kobs er-lö-set.
\end{verse}

\section{Vierter Adventssonntag}

Jes. 45, 8. Ps. 19, 2.

\subsubsection{I. Ton}
\begin{verse}
Träufelt, ihr Himmel, von o-ben,\\
Und die Wolken regnen die Ge-rech-tig-keit.\\
Die Erde thue sich auf —,\\
Und bringe Heil (f g a g).\\
Die Himmel : erzählen die Eh-re Got-tes,\\
Und die Veste : verkündiget sei-ner Hän-de Werk.
\end{verse}

\section[Weihnachten]{Weihnachten (erster und zweiter Feiertag)}

Jes. 9, 6. Ps. 98, 1.

\subsubsection{V. Ton}
\begin{verse}
Uns ist ein Kind ge-bo-ren,\\
Ein Sohn : ist uns ge-ge-ben.\\
Welches Herrschaft : ist auf seiner Schul-ter,\\
Und er heißt : Wunderbar, Rath, Kraft, Held : Ewig-va-ter, Friede-fürst.\\
Singet dem HEerrn : ein neues Lied,\\
Denn er thut Wun-der.
\end{verse}

\section{Sonntag nach Weihnachten}

Ps. 93, 5. 2. 1.

\subsubsection{V. Ton}
\begin{verse}
Dein Wort ist die rechte Lehre : Heiligkeit ist die Zierde deines Hauses e-wiglich.\\
Dein Stuhl stehet fest : Du bist e-wig.\\
Der HEerr ist König : und herrlich ge-schmü-cket,\\
Der HEerr ist geschmückt : und hat ein Reich an-ge-fan-gen.
\end{verse}

\quellseite{117}
\section{Tag der Beschneidung, Neujahr}

Ps. 8, 1. 5. Jes. 63, 16.

\subsubsection{V. Ton}
\begin{verse}
HErr, un-ser Herrscher : wie herrlich ist dein Name : in al-len Lan-den,\\
Da man dir dan-ket im Him-mel.\\
Was ist der Mensch : daß du sein ge-den-kest,\\
Und des Menschen Kind : daß du dich sein an-nimmst?\\
Du, HErr, bist unser Vater : und unser Er-lö-ser,\\
Von Alters her : ist das dein Na-me.
\end{verse}

\section{Sonntag nach Neujahr}

(Wie am Sonntag nach Weihnachten.)

\section{Epiphanias}

Mal. 3, 1. Ps. 72, 1.

\subsubsection{V. Ton}
\begin{verse}
Sie-he, nun kommt der HErr, der HEerr-scher,\\
Und in seiner Hand ist das Reich : und die Kraft : und die HEerr-lichkeit.\\
Gott, gib dein Gericht dem Kö-nig,\\
Und deine Gerechtigkeit des Kö-nigs Sohne.
\end{verse}

\section{Erster Sonntag nach Epiphanias}

Jes. 6, 1. Offenb. 4 u. 5. Ps. 100, 1.

\subsubsection{I. Ton}
\begin{verse}
Auf einem hohen und erhabenen Throne : sah ich den HErren sit-zen,\\
Und sein Saum : fül-le-te den Tem-pel.\\
Und es beteten ihn an : die Heerschaaren der Engel : und sprachen\\
mit gro-ßer Stim-me.\\
Seines Reiches Ehre währet in E-wig-keit.\\
Jauchzet dem HErrn : al-le Welt,\\
Dienet dem HErrn mit Freuden.
\end{verse}

\section{Zweiter Sonntag nach Epiphanias}

Ps. 66, 4. 1. u. 2.

\subsubsection{I. Ton}
\begin{verse}
Alles Land bete dich an : und lob-sin-ge dir —\\
Lobsinge deinem Namen, du Höch-ster.\\
Jauchzet Gott, al-le Lan-de,\\
Lobsinget zu Ehren seinem Namen : rühmt ihn herr-lich.
\end{verse}

\quellseite{118}
\section{Dritter Sonntag nach Epiphanias}

Ps. 97, 8. u. 1.

\subsubsection{I. Ton}
\begin{verse}
Be-tet ihn an : alle sei-ne En-gel!\\
Zion hört es und ist froh.\\
Die Töchter Ju-da sind fröh-lich,\\
HEerr, über dei-nem Re-gi-ment.\\
Der HEerr ist König : des freue sich das Erd-reich,\\
Und seien fröhlich die In-seln, so viel ih-rer ist.
\end{verse}

\section[Vierter und fünfter Sonntag]{Vierter und fünfter Sonntag nach Epiphanias}

(Wie am dritten Sonntag.)

\section{Sechster Sonntag nach Epiphanias}

Ps. 77, 19. 84, 1. 2.

\subsubsection{I. Ton}
\begin{verse}
Dei-ne Blitze leuchteten auf dem Erd-bo-den.\\
Das Erdreich regte sich und be-bete da-von.\\
Wie lieblich sind deine Wohnungen : HEerr Ze-ba-oth,\\
Meine Seele verlanget und sehnet sich : nach den Vor-hö-fen des HEerr-ren.
\end{verse}

\section{Sonntag Septuagesimä}

Ps. 18, 5–7. 2. u. 3.

\subsubsection{I. Ton}
\begin{verse}
Es umfingen mich des To-des Bande,\\
Und die Bäche Beli-als erschreck-ten mich.\\
Da mir angst war : rief ich den HEerrn an : und schrie zu meinem Gott,\\
Und er erhörte meine Stimme : von seinem hei-ligen Tempel.\\
Herzlich lieb habe ich dich : HEerr, mei-ne Stär-ke,\\
HEerr, mein Fels : meine Burg : und mein Er-retter.
\end{verse}

Oder: Jes. 43, 22. Jer. 31, 3.

\begin{verse}
So spricht der HEerr : Nicht daß du mich hättest ge-ru-fen, Ja-kob,\\
Oder daß du um mich gearbeitet hät-test, Is-ra-el.\\
Ich habe dich je und je ge-lie-bet,\\
Darum habe ich dich zu mir gezogen : aus lau-ter Gü-te.
\end{verse}

\quellseite{119}
\section{Sonntag Sexagesimä}

Ps. 44, 24–27. u. 2.

\subsubsection{III. Ton}
\begin{verse}
Erwecke dich, HErr: war-um schläfest du?\\
Wache auf: und ver-stoß uns nicht so gar!\\
Warum verbirgst du dein Ant-litz,\\
Vergissest unsers E-lends und Dranges?\\
Denn unsre Seele ist ge-bugt zur Erde,\\
Unser Bauch klebt am Erd-bo-den.\\
Mache dich auf, HErr: hilf uns: und er-lö-se uns!\\
Gott: wir haben es mit unsern Ohren gehört: unsre Väter haben es\\
uns er-zäh-let.
\end{verse}

\section{Sonntag Estomihi}

Ps. 31, 3. 4. u. 2.

\subsubsection{III. Ton}
\begin{verse}
Sei mir ein starker Fels: und eine Burg: daß du mir helfest.\\
Denn du bist mein Fels und meine Burg: und um deines Namens\\
willen: wollest du mich leiten und führen.\\
HErr, auf dich traue ich: laß mich nimmermehr zu Schan-den\\
wer-den.\\
Errette mich: durch deine Ge-rech-tig-keit.
\end{verse}

\section{Sonntag Invocavit}

Ps. 91, 15. 16. u. 1.

\subsubsection{III. Ton}
\begin{verse}
Er ruft mich an: so will ich ihn erhören: ich bin bei ihm in der Noth;\\
Ich will ihn herausreißen: und zu Eh-ren machen.\\
Ich will ihn sättigen: mit lan-gem Le-ben,\\
Und will ihm zei-gen mein Heil.\\
Wer unter dem Schirm des Höch-sten sitzet,\\
Wird unter dem Schatten des All-mäch-ti-gen blei-ben.
\end{verse}

\section{Sonntag Reminiscere}

Ps. 25, 6. 1. u. 2.

\subsubsection{VIII. Ton}
\begin{verse}
Gedenke, HErr: an deine Barmherzigkeit: und an deine Güte,\\
Die von der Welt her ge-we-sen ist,
\end{verse}

\quellseite{120}
\begin{verse}
Daß nicht unsere Feinde über uns herr-schen.\\
Gott Israels: erlöse uns aus al-ler Noth.\\
Nach dir, HEerr, verlanget mich: mein Gott, ich hoffe auf dich,\\
Laß mich nicht zu Schan-den werden.
\end{verse}

\section{Sonntag Oculi}

Ps. 25, 15. 16. u. 1.

\subsubsection{VIII. Ton}
\begin{verse}
Meine Augen sehen stets zu dem HEer-ren,\\
Denn er wird meinen Fuß aus dem Net-ze ziehen.\\
Wende dich zu mir: und sei mir gnä-dig,\\
Denn ich bin ein-sam und elend.\\
Nach dir, HEerr: ver-lan-get mich.\\
Mein Gott: ich hoffe auf dich: Laß mich nicht zu Schan-den\\
werden.
\end{verse}

\section{Sonntag Lätare}

Jes. 66, 10. 11. Ps. 122, 1.

\subsubsection{V. Ton}
\begin{verse}
Freuet euch mit Je-ru-salem\\
Und seid fröhlich über sie: alle, die ihr sie lieb habt.\\
Freuet euch mit ihr: alle, die ihr über sie traurig ge-we-sen seid.\\
Denn dafür sollt ihr saugen: und satt werden: von den Brüsten\\
ih-res Trostes.\\
Ich freue mich des, daß mir ge-redt ist,\\
Daß wir werden ins Haus des HEer-ren gehen.
\end{verse}

\section{Sonntag Judica}

Ps. 43, 1–3.

\subsubsection{VIII. Ton}
\begin{verse}
Richte mich, Gott: und führe meine Sache: wider das un-hei-li-ge Volk,\\
Und errette mich von den falschen und bösen Leuten; denn du bist\\
der Gott mei-ner Stär-ke.\\
Sende dein Licht: und deine Wahrheit: daß sie mich lei-ten,\\
Und bringen zu deinem heiligen Berge: und zu dei-ner Wohnung.
\end{verse}

\quellseite{121}
\section{Palmsonntag}

Ps. 22, 20. 22. 1.

\subsubsection{III. Ton}
\begin{verse}
HErr, sei nicht ferne : meine Stärke : eile, mir zu hel-fen.\\
Hilf mir aus dem Rachen des Löwen : und errette mich : von den Ein-hör-nern.\\
Mein Gott : warum hast du mich ver-las-sen?\\
Ich heule, aber meine Hilfe ist ferne.
\end{verse}

\section{Gründonnerstag}

Phil. 2, 10. 8. 11. Ps. 102, 2.

\subsubsection{III. Ton}
\begin{verse}
Im Namen JEsu : sollen sich beugen alle derer Kni-e,\\
Die im Himmel : und auf Erden : und unter der Er-de sind.\\
Denn er ist gehorsam worden bis zum Tode : ja zum Tod am Kreu-ze.\\
Darum ist JEsus Christus : der HErr zur Ehre Gottes des Va-ters.\\
HErr : hö-re mein Ge-bet,\\
Und laß mein Schreien zu dir kom-men.
\end{verse}

\section{Charfreitag}

Jes. 53, 4–6. Ps. 102, 2.

\subsubsection{III. Ton}
\begin{verse}
Fürwahr, er trug uns-re Krank-heit.\\
Und lud auf sich uns-re Schmer-zen.\\
Er ist um unserer Missethat willen ver-wun-det\\
Und um unserer Sünde willen zer-schla-gen.\\
Wir gingen alle in der Irre wie Scha-fe :\\
Aber der HErr warf unser aller Sünde auf ihn.\\
HErr : hö-re mein Gebet,\\
Und laß mein Schreien zu dir kommen.
\end{verse}

\section{Ostersonntag und -Montag}

Luc. 24, 6. 7. Ps. 8, 6. 7.

\subsubsection{VIII. Ton}
\begin{verse}
Er ist auf-erstanden, was suchet ihr den Lebendigen bei den Todten. Halle-lu-ja!\\
Gedenket daran : wie er sagte : Hal-le-lu-ja!
\end{verse}

\quellseite{122}
\begin{verse}
Des Menschen Sohn : muß gekreuzigt wer-den,\\
Und am dritten Tage auferstehen : Halleluja, Hal-le-lu-ja!\\
Mit Ehre und Schmuck hast du ihn ge-krö-net.\\
Und hast ihn zum HEerrn gemacht : dei-ner Hände Werk.
\end{verse}

\section{Sonntag Quasimodogeniti}

1 Petr. 2, 2. Ps. 81, 2.

\subsubsection{I. Ton}
\begin{verse}
Als die jetzt gebornen Kindlein, Hal-le-lu-ja!\\
Seid begierig nach der vernünftigen lautern Milch : Halleluja :\\
Halleluja : Hal-le-lu-ja!\\
Singet fröhlich Gott : der unsre Stär-ke ist;\\
Jauchzet dem Got-te Ja-kobs.
\end{verse}

\section{Sonntag Misericordias Domini}

Ps. 33, 5. 6, 1.

\subsubsection{I. Ton}
\begin{verse}
Die Er-de ist voll der Güte des HErrn, Hal-le-lu-ja!\\
Der Himmel ist durchs Wort des HErrn gemacht : Halleluja, Hal-le-lu-ja!\\
Freuet euch des HErrn, ihr Ge-rech-ten;\\
Die Frommen sollen ihn schön prei-sen.
\end{verse}

\section{Sonntag Jubilate}

Ps. 66, 1–3.

\subsubsection{V. Ton}
\begin{verse}
Jauchzet Gott, alle Lande, Halle-lu-ja!\\
Lobsinget seinem Namen, Hal-le-lu-ja!\\
Rühmet ihn herr-lich.\\
Halleluja : Halleluja : Hal-le-lu-ja!\\
Sprechet zu Gott : Wie wunderbarlich sind deine Wer-ke;\\
Es wird deinen Feinden fehlen : vor deiner gro-ßen Macht.
\end{verse}

\section{Sonntag Cantate}

Ps. 98, 1. 2.

\subsubsection{V. Ton}
\begin{verse}
Sin-get dem HErrn : ein neues Lied, Halle-lu-ja!\\
Denn der HErr thut Wunder, Hal-le-lu-ja!
\end{verse}

\quellseite{123}
\begin{verse}
Vor den Völkern: läßt er seine Gerechtigkeit offenbaren, Hal-le-lu-ja!\\
Er sieget mit seiner Rechten,\\
Und mit seinem heil-gen Arme.
\end{verse}

\section{Sonntag Vocem Jucunditatis oder Rogate}

Jes. 48, 20. Ps. 66, 1. 2.

\subsubsection{V. Ton}
\begin{verse}
Predigt es mit fröhlichem Schall: daß man’s höre: Hal-le-lu-ja!\\
Bringet es aus bis an der Welt Ende: Der HErr hat seinen Knecht\\
Jakob erlöset: Halleluja, Hal-le-lu-ja!\\
Jauchzet Gott, alle Lan-de.\\
Lobsinget seinem hei-li-gen Na-men.
\end{verse}

\section{Himmelfahrtsfest}

Apost. 1, 11. Ps. 47, 2.

\subsubsection{V. Ton}
\begin{verse}
Ihr Män-ner von Galiläa: was stehet ihr: und sehet gen Himmel, Hal-le-lu-ja!\\
Wie ihr ihn gesehen habt gen Himmel fahren: wird er kommen.\\
Halleluja: Halleluja: Hal-le-lu-ja!\\
Jauchzet Gott, alle Lan-de!\\
Lobsinget zu Ehren: seinem hei-li-gen Na-men.
\end{verse}

\section{Sonntag Exaudi}

Ps. 27, 7–9. 1.

\subsubsection{I. Ton}
\begin{verse}
HErr, hö-re meine Stimme: wenn ich rufe. Hal-le-lu-ja!\\
(Sei mir gnädig, und er-hö-re mich.)\\
Mein Herz hält dir vor dein Wort: Ihr sollt mein Ant-litz su-chen.\\
Darum suche ich auch, HErr, dein Ant-litz.\\
Verbirg dein Antlitz nicht vor mir, Halleluja: Halle-lu-ja!\\
Der HErr ist mein Licht und mein Heil: vor wem sollt ich mich\\
fürchten?
\end{verse}

\quellseite{124}
\section{Pfingstfest}

Weish. 1, 7. Ps. 68, 29. 30.

\subsubsection{VIII. Ton}
\begin{verse}
Der Weltkreis ist voll Geistes des HEerrn: Halle-lu-ja!\\
Und der die Rede kennet: ist allenthalben: Halleluja: Halle-luja: Hal-le-lu-ja!\\
Dasselbige dein Werk: wollest du, Gott, uns stär-ken,\\
Von deinem heiligen Tempel zu Je-ru-salem.
\end{verse}

(Pfingstmontag dasselbe.)

\section{Trinitatisfest}

Tob. 12, 7.

\subsubsection{VIII. Ton}
\begin{verse}
Gebene-deit sei die heilige Dreifaltigkeit: und unzertheilte Einigkeit.\\
Lasset uns loben und danken Gott vom Himmel bei jedermann: daß\\
er an uns seine Barmherzig-keit er-zeiget hat.\\
Lasset uns loben den Vater und den Sohn\\
Sammt dem Hei-li-gen Gei-ste.
\end{verse}

\section{Erster Sonntag nach Trinitatis}

Ps. 13, 6. 2.

\subsubsection{I. Ton}
\begin{verse}
HEerr, ich hof-fe darauf: daß du so gnädig bist: mein Herz freuet sich: daß du\\
so ger-ne hil-fest.\\
Ich will dem HEerrn singen: daß er so wohl an mir thut.\\
HEerr: wie lange willst du meiner so gar ver-gessen?\\
Wie lange verbirgest du dein Ant-litz vor mir?
\end{verse}

\section{Zweiter Sonntag nach Trinitatis}

Ps. 18, 19. 20. 2. 3.

\subsubsection{V. Ton}
\begin{verse}
Der HEerr ward meine Zu-versicht: und er führte mich aus in den Raum.\\
Er riß mich heraus: denn er hat-te Lust zu mir.\\
Herzlich lieb habe ich dich, HEerr: meine Stär-ke,\\
HEerr, mein Fels und mei-ne Burg.
\end{verse}

\quellseite{125}
\section{Dritter Sonntag nach Trinitatis}

Ps. 25, 16. 18. 1. 2.

\subsubsection{III. Ton}
\begin{verse}
Wende dich zu mir : und sei mir gnä-dig,\\
Denn ich bin ein-sam und e-lend.\\
Siehe an meinen Jammer und E-lend,\\
Und vergib mir alle meine Sün-de, mein Gott.\\
Nach dir, HErr, ver-lan-get mich,\\
Mein Gott : ich hoffe auf dich : laß mich nicht zu Schan-den werden.
\end{verse}

\section{Vierter Sonntag nach Trinitatis}

Ps. 27, 1–3.

\subsubsection{VIII. Ton}
\begin{verse}
Der HErr ist mein Licht und mein Heil : vor wem soll ich mich fürch-ten?\\
Der HErr ist meines Lebens Kraft, vor wem sollt mir grau-en?\\
Darum : so die Bö-sen : meine Widersacher und Feinde : an mich\\
wollen : mein Fleisch zu fres-sen :\\
Müssen sie an-laufen und fal-len.\\
Wenn sich schon ein Heer wider mich le-get,\\
So fürchtet sich den-noch mein Herz nicht.
\end{verse}

\section{Fünfter Sonntag nach Trinitatis}

Ps. 27, 7. 9. 1.

\subsubsection{VIII. Ton}
\begin{verse}
HErr, höre meine Stimme, wenn ich ru-fe.\\
Du bist mei-ne Hilfe.\\
Laß mich nicht : und thue nicht von mir die Hand ab, Gott, mein Heil.\\
Der HErr ist mein Licht und mein Heil : vor wem sollt ich mich\\
fürchten?
\end{verse}

\section{Sechster Sonntag nach Trinitatis}

Ps. 28, 8. 9. 1.

\subsubsection{I. Ton}
\begin{verse}
Der HErr ist seines Vol-kes Stär-ke.\\
Er ist die Stärke, die sei-nem Gesalb-ten hilft.
\end{verse}

\quellseite{126}
\begin{verse}
Hilf deinem Volk, HEerr: und seg-ne dein Er-be,\\
Und wei-de sie: und er-hö-he sie e-wig-lich.\\
Wenn ich rufe zu dir: HEerr, mein Hort: so schwei-ge mir nicht,\\
Auf daß nicht: wo du schweigest: ich gleich werde denen: die in\\
die Höl-le fah-ren.
\end{verse}

\section{Siebenter Sonntag nach Trinitatis}

Ps. 47, 2. 4.

\subsubsection{I. Ton}
\begin{verse}
Froh-locket mit Händen, al-le Völ-ker,\\
Und jauchzet Gott: mit fröh-lichem Schal-le.\\
Denn der HEerr: der Allerhöchste: ist er-schreck-lich,\\
Ein großer König: auf dem gan-zen Erd-bo-den.
\end{verse}

\section{Achter Sonntag nach Trinitatis}

Ps. 48, 10. 11.

\subsubsection{VIII. Ton}
\begin{verse}
Gott, wir war-ten deiner Gü-te\\
In dei-nem Tem-pel.\\
Gott: wie dein Name: so ist auch dein Ruhm: bis an der Welt\\
En-de.\\
Deine Rechte ist voll Ge-rech-tig-keit.\\
Groß ist der HEerr: und hoch-berühmt\\
In der Stadt unsers Gottes: auf seinem hei-li-gen Ber-ge.
\end{verse}

\section{Neunter Sonntag nach Trinitatis}

Ps. 54, 6. 7. 3.

\subsubsection{VIII. Ton}
\begin{verse}
Siehe, Gott steht mir bei.\\
Der HEerr erhält mei-ne See-le.\\
Er wird die Bosheit meinen Feinden be-zah-len.\\
Zerstöre sie durch deine Treue: HEerr, mei-ne Stär-ke.\\
Hilf mir: Gott: durch deinen Na-men,\\
Und errette mich durch dei-ne Stär-ke.
\end{verse}

\quellseite{127}
\section{Zehnter Sonntag nach Trinitatis}

Ps. 55, 17. 19. 20. 23. 2.

\subsubsection{VIII. Ton}
\begin{verse}
Ich will zu Gott ru-fen,\\
So wird er meine Stim-me hö-ren.\\
Er erlöset meine Seele: von denen, die an mich wollen: und schafft ihr Ru-he.\\
Gott wird hören: und sie demüthigen: der all-weg bleibt.\\
Wirf dein Anliegen auf den HErrn: der wird dich ver-sor-gen.\\
Gott: hö-re mein Gebet.\\
Und verbirg dich nicht vor mei-nem Fle-hen.
\end{verse}

\section{Elfter Sonntag nach Trinitatis}

Ps. 68, 6. 7. 36. 1.

\subsubsection{V. Ton}
\begin{verse}
Er ist ein Gott in seiner heiligen Woh-nung,\\
Ein Gott: der den Einsamen das Haus voll Kin-der gibt.\\
Der Gott Israels: wird dem Volke Macht und Kraft ge-ben.\\
Es stehe Gott auf: daß seine Feinde zerstreuet werden: und die ihn hassen, vor ihm fliehen.
\end{verse}

\section{Zwölfter Sonntag nach Trinitatis}

Ps. 70, 2. 3.

\subsubsection{VIII. Ton}
\begin{verse}
Ei-le, Gott: mich zu er-ret-ten,\\
HErr: mir zu hel-fen.\\
Es müssen sich schämen: und zu Schanden werden: die nach meiner Seele ste-hen.\\
Sie müssen zurück kehren: und gehöhnet werden: die mir Ue-bels wün-schen.
\end{verse}

\section{Dreizehnter Sonntag nach Trinitatis}

Ps. 74, 20. 21–23. 1.

\subsubsection{VIII. Ton}
\begin{verse}
Ge-den-ke, HErr, an dei-nen Bund.\\
Laß den Geringen nicht mit Schanden da-von ge-hen.\\
Mache dich auf, Gott: und führe aus deine Sa-che,\\
\quellseite{128}
Und vergiß nicht des Geschreies dei-ner Feinde.\\
Gott : warum verstößest du uns so gar?\\
Und bist so grimmig zornig über die Schafe dei-ner Wei-de?
\end{verse}

\section{Vierzehnter Sonntag nach Trinitatis}

Ps. 84, 10. 11. 2. 3.

\subsubsection{I. Ton}
\begin{verse}
Gott, un-ser Schild, schaue doch : siehe an das Reich dei-nes Gesalb-ten.\\
Denn ein Tag in deinen Vorhöfen : ist besser : denn sonst tausend.\\
Wie lieblich sind deine Wohnungen : HErr Ze-ba-oth,\\
Meine Seele verlanget und sehnet sich : nach den Vor-hö-fen des HErr-ren.
\end{verse}

\section{Fünfzehnter Sonntag nach Trinitatis}

Ps. 86, 1–4.

\subsubsection{III. Ton}
\begin{verse}
HErr, neige deine Ohren und er-hö-re mich.\\
Hilf du, Gott, mein Gott, deinem Knechte : der sich ver-läßt auf dich.\\
HErr : sei mir gnädig : denn ich ru-fe täg-lich zu dir.\\
Erfreue die Seele deines Knechtes : denn nach dir, HErr, verlangt mich.
\end{verse}

\section{Sechzehnter Sonntag nach Trinitatis}

Ps. 86, 3. 5. 1.

\subsubsection{III. Ton}
\begin{verse}
HErr, sei mir gnädig : denn ich ru-fe täg-lich zu dir;\\
Denn du, HErr : bist gut und gnädig : von großer Güte, allen, die dich an-ru-fen.\\
HErr : neige deine Ohren und er-hö-re mich.\\
Denn ich bin einsam und elend und arm.
\end{verse}

\section{Siebzehnter Sonntag nach Trinitatis}

Ps. 119, 137. 124. 1.

\subsubsection{I. Ton}
\begin{verse}
HErr, du bist gerecht : und dein Wort ist recht.\\
Handle mit deinem Knecht nach dei-ner Gna-de.\\
Wohl denen : die ohne Wan-del leben,\\
Die im Gesetz des HErren wan-deln.
\end{verse}

\quellseite{129}
\section{Achtzehnter Sonntag nach Trinitatis}

Sir. 36, 18. 19. Ps. 122, 1.

\subsubsection{V. Ton}
\begin{verse}
HErr, vergilt denen, so auf dich harren : daß deine Propheten wahrhaftig erfunden wer-den.\\
Erhöre das Gebet deines Knechts : und deines Vol-kes Is-rael.\\
Ich freue mich des, daß mir ge-redt ist,\\
Daß wir werden ins Haus des HErren gehen.
\end{verse}

\section{Neunzehnter Sonntag nach Trinitatis}

Ps. 35, 3. 34, 18. 78, 1.

\subsubsection{III. Ton}
\begin{verse}
Der HErr spricht zu seinem Volke : Ich bin dei-ne Hilfe.\\
Wenn sie schreien : so will ich hören : und sie erretten aus aller ihrer Noth : und ihr HErr sein für und für.\\
Höre, mein Volk, mein Ge-setz;\\
Neiget eure Ohren zu der Rede mei-nes Mun-des.
\end{verse}

\section{Zwanzigster Sonntag nach Trinitatis}

Dan. 9, 14. Ps. 119, 1.

\subsubsection{III. Ton}
\begin{verse}
Der HErr, unser Gott, ist gerecht in allen seinen Wer-ken, die er thut.\\
Denn wir gehorchten seiner Stim-me nicht.\\
Nun aber, HErr, gib Ehre dei-nem Na-men,\\
Und thue mit uns nach deiner großen Barm-her-zig-keit.\\
Wohl denen : die ohne Wan-del le-ben,\\
Die im Gesetz des HErren wan-deln.
\end{verse}

\section[Einundzwanzigster Sonntag]{Einundzwanzigster Sonntag nach Trinitatis}

Stücke in Esther 2, 1–3. Ps. 119, 1.

\subsubsection{VIII. Ton}
\begin{verse}
HErr Gott, du bist der allmächtige Kö-nig.\\
Es steht alles in deiner Macht : und deinem Willen kann niemand wi-der-ste-hen.
\end{verse}

\quellseite{130}
\begin{verse}
Denn du hast Himmel und Erde gemacht : und alles, was\\
dar-in-nen ist.\\
. . . Du bist al-ler HErr.\\
Wohl denen, die ohne Wandel le-ben,\\
Die im Gesetz des HErren wan-deln.
\end{verse}

\section[Zweiundzwanzigster Sonntag]{Zweiundzwanzigster Sonntag nach Trinitatis}

Ps. 130, 3. 4. 1. 2.

\subsubsection{III. Ton}
\begin{verse}
So du willst, HErr, Sünde zurechnen : HErr : wer wird be-ste-hen?\\
Denn bei dir ist die Ver-gebung, Gott Is-ra-el.\\
Aus der Tiefe ruf ich, HErr, zu dir.\\
HErr : höre mei-ne Stim-me.
\end{verse}

\section[Dreiundzwanzigster Sonntag]{Dreiundzwanzigster Sonntag nach Trinitatis}

Jer. 29, 11. 12. 14. Ps. 85, 2.

\subsubsection{I. Ton}
\begin{verse}
Ich weiß wohl, was ich für Gedanken über euch habe, spricht der HErr,\\
Nämlich Gedanken des Friedens : und nicht des Leides.\\
Und ihr werdet mich anrufen, und mich bitten;\\
Und ich will euch er-hö-ren.\\
Und ich will euer Ge-fäng-niß wenden,\\
Und euch sammeln aus allen Völkern und von allen Orten.\\
HErr, du bist vormals gnädig gewesen dei-nem Lan-de,\\
Und hast die Gefangenen Ja-kobs er-lö-set.
\end{verse}

\section[Vierundzwanzigster Sonntag]{Vierundzwanzigster Sonntag nach Trinitatis}

Ps. 95, 6. 7. 1.

\subsubsection{V. Ton}
\begin{verse}
Kommt, lasset uns anbeten : und knieen : und niederfallen vor dem HErrn :\\
der uns ge-macht hat.\\
Denn er ist un-ser Gott.\\
Kommt herzu : lasset uns dem HErrn froh-lo-cken,\\
Und jauchzen dem Hort un-sers Hei-les.
\end{verse}

\quellseite{131}
\section[Fünfundzwanzigster Sonntag]{Fünfundzwanzigster Sonntag nach Trinitatis}

Ps. 31, 10. 16. 18. 2.

\subsubsection{VIII. Ton}
\begin{verse}
HErr, sei mir gnädig, denn mir ist angst.\\
Errette mich von der Hand meiner Feinde : und von denen, die mich ver-fol-gen.\\
HErr, laß mich nicht zu Schanden werden : denn ich ruf dich an.\\
HErr, auf dich traue ich : Laß mich nimmermehr zu Schanden werden : errette mich durch deine Ge-rech-tig-keit.
\end{verse}

\section[Sechsundzwanzigster Sonntag]{Sechsundzwanzigster Sonntag nach Trinitatis}

Ps. 54, 3. 4. 7.

\subsubsection{VIII. Ton}
\begin{verse}
Hilf mir, Gott, durch deinen Na-men,\\
Und schaffe mir Recht durch dei-ne Gewalt.\\
Gott, er-hö-re mein Gebet,\\
Vernimm die Rede mei-nes Mun-des.\\
Er wird die Bosheit meinen Feinden be-zah-len.\\
Zerstöre sie durch dei-ne Treu-e.
\end{verse}

\section[Siebenundzwanzigster Sonntag]{Siebenundzwanzigster Sonntag nach Trinitatis}

(Wie am Trinitatisfest.)

\section{Tag der Reinigung Mariä}

(Wie am achten Sonntag nach Trinitatis.)

\section{Tag der Verkündigung Mariä}

(Wie am vierten Adventssonntag.)

\section{Tag St. Johannis des Täufers}

Jes. 49, 1. 2. Ps. 92, 1.

\subsubsection{V. Ton}
\begin{verse}
Der HErr hat mich von Mutterleibe an bei meinem Namen ge-ru-fen,\\
Und hat meinen Mund gemacht wie ein scharf Schwert.\\
Mit dem Schatten seiner Hand hat er mich be-de-cket,\\
Und hat mich zum rei-nen Pfeil gemacht.\\
Das ist ein köstlich Ding, dem HErrn dan-ken,\\
Und lobsingen deinem Na-men, du Höchster.
\end{verse}

\quellseite{132}
\section{Tag der Heimsuchung Mariä}

(Wie zu Mariä Verkündigung.)

\section{Tag Michaelis}

Ps. 103, 20. 21. 1.

\subsubsection{V. Ton}
\begin{verse}
Lobet den HErrn : ihr seine Engel,\\
Ihr starken Helden : die ihr seinen Befehl ausrichtet : daß man höre die Sti-mme sei-nes Worts.\\
Lobet den HErrn, alle seine Heer-scharen,\\
Seine Diener : die ihr sei-nen Willen thut.\\
Lobe den HErrn, meine See-le,\\
Und was in mir ist, seinen heili-gen Na-men.
\end{verse}

\section{Aposteltage}

2 Tim. 4, 8. 1, 12. Ps. 139, 1. 2.

\subsubsection{V. Ton}
\begin{verse}
Ich weiß, an welchen ich glaube,\\
Und bin gewiß, daß er mir meine Beilage kann bewahren bis an jenen Tag.\\
Hinfort ist mir beigelegt die Krone der Ge-rech-tig-keit,\\
Welche mir der HErr : der gerechte Rich-ter, geben wird.\\
HErr, du erforschest mich und ken-nest mich.\\
Ich sitze oder stehe auf, so weißt du es.
\end{verse}

\section{Fest der Reformation}

Ps. 46, 8. 3. 2.

\subsubsection{VIII. Ton}
\begin{verse}
Der HErr Zebaoth ist mit uns,\\
Der Gott Ja-kob ist unser Schutz.\\
Darum fürchten wir uns nicht : wenn gleich die Welt unter-gin-ge,\\
Und die Berge mitten ins Meer sän-ken.\\
Gott ist unsere Zuversicht und Stär-ke,\\
Eine Hilfe in den großen Nöthen : die uns trof-fen ha-ben.
\end{verse}

\quellseite{133}
\section{Buß- und Bettag}

Jes. 1, 2. 4. Ps. 130, 3.

\subsubsection{III. Ton}
\begin{verse}
Höret, ihr Himmel, und Erde, nimm zu Ohren : denn der HEerr re-det:\\
Ich habe Kinder auferzogen, und erhöhet : und sie sind von mir ab-ge-fal-len.\\
Sie haben den HEerrn ver-las-sen,\\
Den Heiligen in Israel gelästert : und wei-chen zu-rück.\\
So du, HEerr, Sünde zu-rech-nest,\\
HEerr, wer wird be-ste-hen?
\end{verse}

(Das Gloria Patri fällt an diesem Tage aus.)

\section{Erntefest und Sonntag}

Ps. 65, 12. 10. 1.

\subsubsection{V. Ton}
\begin{verse}
Du krönest das Jahr mit dei-nem Gut.\\
Deine Fußstapfen trie-fen von Fett.\\
Du suchest das Land heim : und wäs-serst es.\\
Gottes Brünnlein hat Was-sers die Fül-le.\\
Gott, man lobet dich in der Stil-le zu Zi-on,\\
Und dir bezahlet man Ge-lüb-de.
\end{verse}

\section{Kirchweihstag}

Ps. 100, 1. 2. 4. 5.

\subsubsection{V. Ton}
\begin{verse}
Jauchzet dem HEerrn, al-le Welt.\\
Dienet dem HEerrn mit Freuden.\\
Danket ihm : lobet seinen Na-men.\\
Denn der HEerr ist freundlich : und seine Güte wäh-ret e-wiglich.
\end{verse}

\section{Nachbemerkung}

Sollte es an manchem Orte zu schwer fallen, alle die vorstehenden Introiten zu singen, so kann man die Menge der einzuübenden Introiten dadurch beschränken, daß man nicht für jeden Sonntag, sondern für jede Kirchenjahreszeit nur je einen Introitus festsetzt, also z. B. den des ersten Advents für die Adventszeit überhaupt singen läßt, wie dies dermalen in der Mecklenburgischen Landeskirche geschieht.

\quellseite{134}
\chapter{II. Kyrie}


Das Kyrie eleison ist Schriftwort; denn Matth. 15, 22. stößt in griechischer Sprache das cananäische Weib diesen Ruf nach dem göttlichen Erbarmen aus, da sie bei dem menschgewordenen Gottessohne, dem verheißenen Davidssohn, Hilfe wider des Teufels Plage für ihre Tochter sucht. Um zu erinnern, daß durch die Mannigfaltigkeit der Zungen die Völker der ganzen Welt der Heilige Geist in Einigkeit des Glaubens versammelt hat, gab man neben dem hebräischen Amen, Halleluja und Hosianna auch dem griechischen Kyrie Bürgerrecht im Gottesdienst.

Der erste gottesdienstliche Gebrauch des Kyrie gehört dem 3. Jahrhundert an. Er bestand darin, daß bei dem Fürbittengebet, mit dem man ursprünglich die Messe zu beginnen schien, das Volk auf die einzelnen Bitten mit dem Kyrie antwortete, ähnlich wie jetzt noch in der Litanei mit dem Rufe: „Behüt uns, lieber HEerre Gott“, und: „Erhör uns, lieber HEerre Gott.“ Gregor der Große, wegen seines liturgischen Ausbaues der römischen Messe und des Gottesdienstes überhaupt auch Pater ceremonarium benannt, beließ das Fürbittengebet der Messe nur für die kirchliche Buß- und Trauerzeit, Adventszeit und Passionszeit, und löste für die übrige Zeit das Kyrie ab, machte es aber nicht nur zu einem dreimaligen Ruf durch Hinzufügung des: Christe eleison, sondern auch zu einem neunmaligen Rufe durch dreimalige Wiederholung eines jeden der drei Glieder, wie es denn noch jetzt in der römischen Messe als ein neunmaliger Ruf gebraucht wird. Seit dem 13. Jahrhundert erweiterte man das Kyrie in mannigfacher Weise. Es entstanden so unter anderm auch jene Festkyries, die zum Theil hernach Aufnahme in die lutherische Kirche fanden. Während nun aber Gregor der Große das Kyrie zum Wechselgesang zwischen Clerus und Chor machte und so dem Volke die Betheiligung entzog, ließen sich die Deutschen dies doch nicht gefallen, sondern fingen seit dem 10. Jahrhundert an, das Kyrie mit ein paar deutschen Festliedern dergestalt zu verbinden, daß es den Refrain eines jeden Verses bildete, woher denn diese Lieder den bekannten Namen „Leise“ erhielten. Sah man sich doch genöthigt, die verbreitetste dieser Leisen, das „Christ ist erstanden von der Marter alle“ bildungsweise in Missalien Deutschlands aufzunehmen.

In seiner „Weise, christlich Meß zu halten“ sagt Luther: „Darnach gefallen mir die wohl, so das Kyrie eleison hinzugesetzt haben; denn wir
\quellseite{135}
lesen, daß zu Basilii Zeiten, den man Magnum nennt, das Kyrie eleison in gemeinem Gebrauch sei gewesen allem Volk.“ (St. L. Ausg. X, 2235.) Als Schriftwort und als ursprünglich vom Volk allgemein gebraucht, ließ daher Luther das Kyrie stehen, beschränkte es jedoch auf einen dreimaligen Gebrauch. Den vorgefundenen Kyrie-Schatz bildete dann die lutherische Kirche nach dem Vorgang der böhmischen Brüder in deutscher Sprache noch weiter fort. Hierdurch entstanden zweierlei Kyries für die Sonn- und Festtage nicht nur dem Inhalte, sondern auch der Form nach. In Bezug auf die letztere sind es 1. antiphonische Kyries und 2. liedförmige Kyries. Im Nachfolgenden sind von beiden Arten Beispiele nach Form und Weise mitgetheilt.

Die Ausführung kann demnach auf verschiedene Weise geschehen. Wird nach Luther etc. das Kyrie antiphonisch gesungen, so kann der Wechselgesang stattfinden 1. zwischen Pastor und Gemeinde oder, was wir vorziehen würden, zwischen Cantor und Gemeinde, wobei dann gleichfalls schon beim Beginn des Kyrie der Pastor den Altar zu betreten hätte, um nach demselben sofort das: „Ehre sei Gott in der Höhe“ zu intoniren; 2. zwischen Chor und Gemeinde und 3. zwischen zwei Chören. Wird aber das liedförmige Kyrie gebraucht, wie z. B. das im St. Louiser Gesangbuch unter Nr. 7 abgedruckte: „Kyrie, Gott Vater in Ewigkeit“, so singt es die Gemeinde, und will sie sich an Festen eines Festkyries bedienen, so kann, falls die Melodie der Gemeinde zu schwer wäre, die Ausführung in antiphonischer Weise geschehen, z. B.:

\begin{verse}
Chor: Kyrie, Gott, aller Welt Schöpfer und Vater,\\
Gem.: Eleison.\\
Chor: Christe, wahrer Gott und Mensch geborn, der du für uns trägst\\
Gottes Zorn,\\
Gem.: Eleison.\\
Chor: Kyrie, Gott Heiliger Geist, mit Vater und Sohn Ein Gott,\\
Gem.: Eleison.
\end{verse}

Die Deutung des Kyrie in seiner Stellung ist verschieden. Manche neuere Liturgiker betrachten dasselbe als Sündenbekenntniß, mit dem sich gleich zu Anfang des Gottesdienstes die Gemeinde dem HEerrn naht, und daher das darauf folgende Gloria als die göttliche Absolution darauf. Da jedoch in der römischen Messe, wie in der hier sich anschließenden Kirchenordnung von Brandenburg-Nürnberg und ihr folgenden Agenden ein Confiteor dem Kyrie vorausgeht, so kann diese Deutung doch wohl nicht zutreffen. Im Hinblick auf den biblischen Ursprung des Kyrie als Ruf um Erbarmung unter
\quellseite{136}
des Teufels Plage und auf seine liturgische Verbindung mit dem engelischen Lobgesang „Ehre sei Gott“ etc. und dem sich anschließenden weiteren Preis des Dreieinigen („Wir loben, preisen, anbeten dich“ etc.) möchten wir lieber, wie schon angedeutet, mit Andern im Kyrie das einmüthige und demüthige Bekenntniß des ganzen Elends, der „tausend Plagen und großen Jammerlast, die kein Mund kann aussagen“, erkennen, um welches Elends willen Gottes eingeborner Sohn Mensch geboren und von dem Lobgesang der himmlischen Heerschaaren begrüßt wurde, und welches Elend die Erlöseten des HEerrn um der Sünde im Fleisch willen noch immer, wiewohl befreit von der Last des Fluches, drückt, bis sie einst auch davon völlig erlöst werden. Als Bekenntniß des mannigfaltigen Elends, das in Folge des Sündenfalls uns von Kind auf drückt, faßt auch der als Meister in liturgicis gleichfalls von Lutheranern anerkannte Cardinal Bona diesen Ruf um Erbarmung auf. —

Da mit dem Kyrie, sei es schon zu Anfang oder gegen das Ende, der Liturg den Altar betritt, um so nun die Abhaltung des ganzen Amtes der Communion — einschließlich der Predigt — zu beginnen, so möge hier die schöne Erinnerung der Braunschweig-Lüneburgischen Kirchenordnung aus dem Capitel „von dem Ornatu Ecclesiastico“ stehen: „So sollen die Pastores und Kirchendiener, so das h. Amt halten wollen, ... fein ehrerbietig mit großer Andacht und Anrufung des Sohnes Gottes vor den Altar treten und wohl bedenken, was für hohe Mysteria, ipsis etiam sanctis angelis tremenda, sie allda, wiewohl sie sündige Menschen sein, verrichten sollen. Einem jeden verständigen, gottesfürchtigen, gewissenhaften Priester wird seine Devotion selbst lehren, was und wie er andächtig beten soll, damit er solches hohe Amt wohl verrichten möge.“

\section{A. Antiphonisches Kyrie}

\subsection{Nr. 1}

\notensatz{music/136-01a}{Antiphonisches Kyrie Nr. 1, Melodie aus Luthers deutscher Messe von 1526. Pastor oder Cantor: Kyrie! Gemeinde mit oder ohne Orgel: Eleison! Pastor: Christe! Gemeinde: Eleison!}

\quellseite{137}
\notensatz{music/137-01b}{Antiphonisches Kyrie Nr. 1, Schluß. Pastor: Kyrie! Gemeinde: Eleison!}

P. Ehre sei Gott etc.

\subsection{Nr. 2}

Chor oder erster Chor.

Gem. oder zweiter Chor.

Harm.: F. Riegel.

\notensatz{music/137-02}{Antiphonisches Kyrie Nr. 2. Kyrie! Eleison! Christe! Eleison! Kyrie! Eleison!}

P. Ehre sei Gott etc.

\subsection{Nr. 3}

Mel.: Straßb. R. G. Harm.: F. Riegel.

\notensatz{music/137-03a}{Antiphonisches Kyrie Nr. 3. Chor oder Cantor: Kyrie! Eleison! Gemeinde: HErr, erbarme dich!}

\quellseite{138}
\notensatz{music/138-03b}{Antiphonisches Kyrie Nr. 3. Chor oder Cantor: Christe! Eleison! Gemeinde: Christe, erbarme dich!}

\notensatz{music/138-03c}{Antiphonisches Kyrie Nr. 3. Chor oder Cantor: Kyrie! Eleison! Gemeinde: HErr, erbarm dich über uns!}


\subsection[Mecklenburgischen Cantional]{Dasselbe aus dem Mecklenburgischen Cantional}

\notensatz{music/138-m1}{Mecklenburgisches Cantional, erstes Kyrie. Liturg: Kyrie! Eleison! Respons: HErr, erbarme dich!}

\notensatz{music/138-m2}{Mecklenburgisches Cantional. Liturg: Christe! Eleison! Respons: Christe, erbarme dich!}

\quellseite{139}
\notensatz{music/139-m3}{Mecklenburgisches Cantional, letztes Kyrie. Liturg: Kyrie! Eleison! Respons: HErr, erbarm dich über uns!}

Mel.: Deutsch Kirchenamt 1525. Zweibrückener K.-O. 1557 und 1570.

\section{B. Liedförmiges Kyrie}

\subsection{1. Für die Adventszeit}

Mel.: Böhm. Brüder. Harm.: F. Riegel.

\notensatz{music/139-advent}{Adventskyrie, Anfang: O milder Gott, allerhöchster Hort … der Kirche Haupt, Freud, Ehr und Kron und.}

\quellseite{140}
\notensatz{music/140-advent}{Adventskyrie, Schluß: frei offner Heilsbronn … Erneurer und Regierer, erbarm dich unser!}

\subsection{2. Von Weihnachten bis Ostern}

Mel.: Spangenberg, Lossius etc. Harm.: Mich. Prätorius.

\notensatz{music/140-kyrie2}{Kyrie von Weihnachten bis Ostern, Anfang: O Vater, allmächtiger Gott! Zu dir schreien wir in der Noth.}

\quellseite{141}
\notensatz{music/141-01}{Kyrie von Weihnachten bis Ostern: durch dein groß Barmherzigkeit, erbarm dich über uns!}

\notensatz{music/141-02}{Kyrie von Weihnachten bis Ostern: O Christe, wollst uns erhören, für uns bist}

\notensatz{music/141-03}{Kyrie von Weihnachten bis Ostern: du geboren von Maria, erbarme dich über uns!}

\notensatz{music/141-04}{Kyrie von Weihnachten bis Ostern: O Heilger Geist, wollst uns geben, dich allzeit herzlich zu lieben}

\quellseite{142}
\notensatz{music/142-01}{Kyrie von Weihnachten bis Ostern: und nach deim Willen streben, erbarme dich über uns!}

\subsection{3. Von Ostern bis Pfingsten}

Mel.: Spangenberg u. Harm.: F. Riegel.

\notensatz{music/142-02}{Kyrie von Ostern bis Pfingsten, Chor: Kyrie, Gott, aller Welt Schöpfer und Vater,}

\notensatz{music/142-03}{Kyrie von Ostern bis Pfingsten, Gemeinde: Eleison! Chor: Christe, wahrer Gott und}

\notensatz{music/142-04}{Kyrie von Ostern bis Pfingsten, Chor: Mensch geborn, der du für uns trugst Gottes Zorn. Gemeinde: Eleison!}

\quellseite{143}
\notensatz{music/143-01}{Kyrie von Ostern bis Pfingsten, Chor: Kyrie, Heiliger Geist, mit Vater und Sohn Ein}

\notensatz{music/143-02}{Kyrie von Ostern bis Pfingsten: Gott. Gemeinde: Eleison! Chor: Kyrie, hilf uns, daß}

\notensatz{music/143-03}{Kyrie von Ostern bis Pfingsten, Chor: wir in solchem Glauben rein dich anbeten allein und bleiben die}

\notensatz{music/143-04}{Kyrie von Ostern bis Pfingsten: Diener dein. Gemeinde: Eleison!}

\subsection{4. Für Trinitatis und die Trinitatiszeit}

das Kyrie summum: „Kyrie, Gott Vater in Ewigkeit“ (No. 7 des Gesangbuchs und No. 177 des neuen Choralbuchs von St. Louis, Mo. 1886.)

\quellseite{144}
\chapter{III. Gloria}


Wie das Kyrie, so ist auch nach Inhalt und Wortlaut das Gloria in excelsis biblischen Ursprungs, denn es ist bekanntlich der Lobgesang der himmlischen Heerschaaren bei der Geburt Christi: „Ehre sei Gott in der Höhe, Friede auf Erden und den Menschen ein Wohlgefallen“ (Luc. 2, 14.), daher auch Hymnus angelicus genannt.

Anfangs wurde dieser Hymnus nur in der Weihnachtspigilie gesungen. Von da ging er zu regelmäßigem Gebrauch in der Matutin der griechischen Kirche, dem Gottesdienst in früher Morgenstunde, über, „um in dem liturgisch-symbolischen Drama ihres Cultus auf die heilige Nacht hinzudeuten, welche der Menschheit die Erlösung aus der Noth der Sünde gebracht hat“.

Die Einführung des Gloria in die Messe soll durch den römischen Bischof Telesphorus um das Jahr 126 geschehen sein, wo es später nach dem Kyrie eingereiht wurde. In den apostolischen Constitutionen, deren Sammlung wohl dem 4. Jahrhundert angehört, findet sich schon die Erweiterung durch das prachtvolle, erhabene Laudamus, das ist, „Wir loben dich“ etc. Im Lateinischen und in seiner von der lutherischen Kirche gebrauchten Verdeutschung lautet der Hymnus also:

\begin{quote}
Gloria in excelsis Deo! Et in terra pax hominibus bonae voluntatis. Laudamus te. Benedicimus te. Adoramus te. Glorificamus te. Gratias agimus tibi propter magnam gloriam tuam. Domine Deus, Rex coelestis, Deus Pater omnipotens; Domine fili unigenite, Jesu Christe, Altissime, Domine Deus. Agnus Dei, filius Patris: qui tollis peccata mundi, miserere nostri, qui tollis peccata mundi, suscipe deprecationem nostram. Qui sedes ad dextram Patris, miserere nostri. Quoniam tu solus sanctus. Tu solus Dominus. Tu solus altissimus, Jesu Christe, cum Sancto Spiritu in gloria Dei Patris. Amen.
\end{quote}

\begin{quote}
Ehre sei Gott in der Höhe, und auf Erden Fried, den Menschen ein Wohlgefallen. Wir loben dich. Wir benedeien dich. Wir beten dich an. Wir preisen dich. Wir sagen dir Dank um deiner großen Ehre willen, HEerr Gott, himmlischer König! Gott, allmächtiger Vater! HEerr, eingeborner Sohn Jesu Christe, du Allerhöchster, und dir, du Heiliger Geist! HEerr Gott, Lamm Gottes, ein Sohn des Vaters, der du hinnimmst die Sünd der Welt, erbarm dich unser! Der du hinnimmst die Sünd der Welt, nimm auf unser Gebet! Der du sitzest zur Rechten des Vaters, erbarm dich unser! Denn du allein bist heilig, du bist
\quellseite{145}
allein der HEerr, du bist allein der Höchst, Jesu Christe, mit dem Heiligen Geist in der Herrlichkeit deines Vaters. Amen.
\end{quote}

Seit Gregor dem Großen wurde jedoch das Gloria mit dem Laudamus für jede Sonn- und Festtagsmesse verwendet.

„Viel Gesangs in der Meß“, sagt Luther, „ist fein und herrlich vom Danken und Loben gemacht und bisher blieben, als Gloria in excelsis et in terra, das Alleluja, das Patrem, die Präfation, das Sanctus, das Benedictus, das Agnus Dei. In diesen Stücken findest du nichts vom Opfer, sondern eitel Lob und Dank, darum wir sie auch in der Messe behalten.“ (Jen. Ausg. V, f. 196 b.) Während nun aber die anglicanische Kirche, wohl nach Vorgang der orientalischen Kirche, im Common Prayer-Book dem Gloria seine Stellung nach der Consecration, ja sogar nach der Austheilung anwies, schloß sich Luther auch hier der römischen Messe an und ließ es auf das Kyrie folgen.

Ueber die Bedeutung des Gloria mit seiner Erweiterung in dieser Stellung und in der Messe überhaupt haben wir uns bereits in der Deutung des Ganges des Hauptgottesdienstes etwas ausgesprochen. Wir haben ihn vornehmlich als Lobpreisung des „kündlich großen Geheimnisses: Gott ist geoffenbart im Fleisch“ aufgefaßt, in Folge dessen nach geschehener Erlösung uns nicht nur das Heil in Christo verkündigt, sondern auch darauf im Hauptgottesdienst Christus seinen wahrhaftigen Leib und sein wahrhaftiges Blut in sacramentlicher Vereinigung mit Brod und Wein zum mündlichen Genuß uns darreicht. Ob daher nicht überhaupt die Fleischwerdung des ewigen Wortes der rothe Faden sein soll, an dem sich die einzelnen Stücke der Messe wie Perlen reihen, da die römische Kirche nach Schluß derselben vom Priester noch Joh. 1, 1–14. vorlesen läßt?

Die Behandlung und Ausführung dieses Stückes der Liturgie war in der lutherischen Kirche eine verschiedene und mannigfaltige. Wie die Präfation sang der Chor es vielfach lateinisch, an manchen Orten bis ins 18. Jahrhundert, zum Theil aber auch gesellte sich der deutsche Gesang der Gemeinde in der Liedform hinzu. Aber auch schon frühzeitig bediente man sich einer Verdeutschung des lateinischen Textes, so z. B. der Nürnberger Kaplan Andreas Döber in seiner „Evangelischen Meß 1525“. Man zog jedoch bald den Gebrauch der Liedform dieses Hymnus vor, wie sie Decius in seinem: „Allein Gott in der Höh sei Ehr“ (St. Louiser Gesb. Nr. 1) und Luther in seinem: „All Ehr und Lob soll Gottes sein“ (ibid. Nr. 142) der Kirche geboten hatte.

\quellseite{146}
Was nun insbesondere die Ausführung betrifft, so intonirte der Liturg,\footnote{In der alten Kirche intonirte der Bischof das Gloria. In der lutherischen Kirche war es, wie der ganze Altardienst bei der Communion, dem Diakonus zugewiesen, wo ein solcher an einer Gemeinde als Hilfsprediger oder zweiter Pastor angestellt war, während dann der erste Pastor die sogenannte „Amtspredigt“ hielt und bei der Communion assistirte. Wir finden es passend, daß, wo mehr als ein Prediger an einer Gemeinde angestellt ist, im Hauptgottesdienst der eine als Liturg und der andere als Prediger fungirt.} gegen den Altar gekehrt das Gloria: „Ehre sei Gott in der Höhe!“ worauf der Chor mit dem Et in terra einfiel oder die Gemeinde bald mit dem „Allein Gott in der Höh“, bald mit dem „All Ehr und Lob soll Gottes sein“ respondirte. Beigefügt finden wir die Ermahnung, daß der Liturg „mit großer Reverenz“ das Gloria handeln (Preuß. K.-O. 1525 u. 1544) und „sein ehrlich und mit großer Andacht und Anrufung des Sohnes Gottes vor dem Altar stehen und nach dem dritten Kyrie das Gloria singen solle“. (Oldenb. K.-O. 1573.) Bemerkt sei beiläufig, daß der Liturg eben deswegen nicht erst beim dritten, sondern schon beim ersten Kyrie an den Altar trat.

Für die Intonation des Gloria finden sich mancherlei Singweisen—Singweisen auf die Sonntage und Singweisen auf die verschiedenen Feste, alle jedoch nach Einem Typus. Die Intonation des Sonntags ist überall einfacher, als die des Festtags.

Es wäre nun zu wünschen, daß, wie ehemals sowohl das Luther’sche „All Ehr und Lob soll Gottes sein“, als auch der Originaltext in seiner Verdeutschung neben dem „Allein Gott in der Höh sei Ehr“ wieder in allgemeinerem Brauch käme, etwa so, daß man die beiden Gloria-Lieder abwechselnd an den gewöhnlichen Sonntagen gebrauchte, die Originalfassung aber für die hohen Feste und sonstige Anlässe verwendete. Das gäbe eine liebliche Abwechselung, dabei das Gloria unverrückt seine Stellung in der Liturgie behielte und nur die Form desselben abwechselte, gleichwie das auch in den meisten unserer hiesigen Gemeinden mit dem sogenannten „großen“ und „kleinen“ Glauben (St. Louiser Synod. Nr. 183 und 184) geschieht.\footnote{Das Gesangbuch der Synode von Buffalo enthält dies dreierlei Gloria.}

Was nun insonderheit zunächst das Luther-Gloria betrifft, das namentlich eine getreue metrische Uebersetzung des Laudamus ist, so pflegte man es häufig der Art auszuführen, daß der Pastor intonirte: „Ehre sei Gott in der Höhe!“, der Chor einfiel: „All Ehr und Lob soll Gottes sein, er ist und bleibt der Höchst
\quellseite{147}
„allein“ und die Gemeinde dann von den Worten an: „Sein Zorn auf Erden hat ein End“ das Lied gar hinaussang. Für dieses Lied ist hier aus Layriz eine Melodie aufgenommen, die ursprünglich zwar für Luthers Vaterunser-Lied componirt war, ganz trefflich aber auch zu seinem Glorialied paßt. Es findet sich dieselbe auch in dem von Hölter herausgegebenen St. Louiser Choralbuch unter Nr. 16 und in dem Brauer’schen Choralbuch unter Nr. 9.

Den Originaltext in seiner Verdeutschung anlangend, so ist hier eine der ältesten lutherischen Singweisen aufgenommen, wie dieselbe außer Hommel auch bei Schöberlein, in dem musikalischen Theil des bayrischen Agenden-Kerns und im Mecklenburgischen Cantional und manch anderm neueren liturgischen Werke sich findet, und wie dieselbe ehedem auch in den hiesigen fränkischen Gemeinden gesungen wurde. Kann es nun nicht erreicht werden, daß die Gemeinde auch das Laudamus mitsingt, so übernimmt der Chor dasselbe, und die Ordnung wäre dann: Pastor: „Ehre sei Gott in der Höhe!“ Gem.: „Und auf Erden Fried und den Menschen ein Wohlgefallen.“ Chor: „Wir loben dich! Wir benedeien dich“ etc. Im letzteren Fall könnte sich mitunter der Chor auch der einfachen und schönen Composition des Laudamus von D. Bortniansky bedienen. Obschon der Neuzeit angehörend, trägt sie durchaus ein kirchliches Gepräge.

Schließlich sei noch erwähnt, daß nach alter Sitte an manchen Orten das Gloria in der Fastenzeit unterlassen wurde, so daß auf das Kyrie sogleich Salutation und Collecte folgte.

\section{I. Gloria-Intonationen}

\subsection{1. Für Sonntage (gloria dominicale)}

Magdeb. 1615.

\notensatz{music/147-01}{Gloria dominicale, Magdeburg 1615: Ehre sei Gott in der Höhe! Elf Noten mit Breven auf Gott und he; Fis am Anfang und vor der abschließenden F-Brevis, drei aufwärts gerichtete Bögen. Die gedruckten Doppelpunkte stehen bei Eh und Hö.}

Nürnb. Agendb. 1639.

\notensatz{music/147-02}{Gloria dominicale, Nürnberger Agendbüchlein 1639: Ehre sei Gott in der Höhe!}

Pommern 1563.

\notensatz{music/147-03}{Gloria dominicale, Pommern 1563: Ehre sei Gott in der Höhe!}

\quellseite{148}
Spangenberg 1545.

\notensatz{music/148-01}{Gloria dominicale, Spangenberg 1545: Ehre sei Gott in der Höhe!}

\subsection{2. Für die hohen Feste}

Nürnb. Agendb. 1639. Schwarzb. A. 1675 u. a. D.

\notensatz{music/148-02}{Gloria für die hohen Feste, Nürnberger Agendbüchlein 1639 und Schwarzburger Agende 1675: Ehre sei Gott in der Höhe!}

Bayreuth. Chorord. 1755.

\notensatz{music/148-03}{Gloria für die hohen Feste, Bayreuther Chorordnung 1755: Ehre sei Gott in der Höhe!}

\subsection{Ostern. (Gloria paschale.)}

M. Prätorius 1611. Liturgodia Sionia.

\notensatz{music/148-04}{Gloria paschale, M. Prätorius 1611, Liturgodia Sionia: Ehre sei Gott in der Höhe!}

Schwarzburg 1675.

\notensatz{music/148-05}{Gloria paschale, Schwarzburg 1675: Ehre sei Gott in der Höhe!}

\subsection{3. Für andere Festtage}

\subsection{a. Advent und Michaelis}

Nürnb. Agendb. 1639.

\notensatz{music/148-06}{Gloria für Advent und Michaelis, Nürnberger Agendbüchlein 1639: Ehre sei Gott in der Höhe!}

\subsection{b. Aposteltage}

Schwarzburg 1675.

\notensatz{music/148-07}{Gloria für Aposteltage, Schwarzburg 1675: Ehre sei Gott in der Höhe!}

\quellseite{149}
\section{II. Luthers Gloria-Lied}

\notensatz{music/149-01}{Pastor, Intonation zu Luthers Gloria-Lied. Ehre sei Gott in der Höhe!}

Wittenberg 1544. Ursprünglich: Vater unser im Himmelreich.

\notensatz{music/149-02}{Gemeinde: All Ehr und Lob soll Gottes sein. Vierstimmiger Satz.}

\quellseite{150}
\section[III. Gloria mit Et in terra und Laudamus]{III. Gloria mit Et in terra und Laudamus antiphonisch}

\subsubsection{Pastor}

Straßb. K.-A. 1525. — Pflz. K.-O. 1570. — Harm.: Bayr. Agende.

\notensatz{music/150-01}{Gloria mit Et in terra und Laudamus, Intonation des Pastors. Ehre sei Gott in der Höhe!}

\subsubsection{Gemeinde}

\notensatz{music/150-02}{Gemeinde: Und auf Erden Fried und den Menschen ein Wohlgefallen.}

\subsubsection{Gemeinde oder Chor}

\notensatz{music/150-03}{Gemeinde oder Chor: Wir loben dich … Gott, allmächtiger Vater!}

\quellseite{151}
\notensatz{music/151-01}{Laudamus: HErr, eingeborner Sohn … nimm auf unser Gebet!}

\notensatz{music/151-02}{Laudamus: Der du sitzest zur.}

\quellseite{152}
\notensatz{music/152-01}{Laudamus: Rechten des Vaters … mit dem Heiligen Geist in der Herrlichkeit deines Vaters.}

\notensatz{music/152-02}{Laudamus: Amen.}

\quellseite{153}
\section{IV. Laudamus}

D. Bortniansky (1752–1825).

\subsubsection{Chor}

\notensatz{music/153-01}{Bortniansky, Laudamus: Wir loben dich, wir benedeien dich, wir beten dich}

\notensatz{music/153-02}{Bortniansky, Laudamus: an, wir preisen dich, wir sagen dir Dank um deiner großen}

\notensatz{music/153-03}{Bortniansky, Laudamus: Herrlichkeit willen. HErr Gott, himmlischer König, allmächtiger}

\notensatz{music/153-04}{Bortniansky, Laudamus: Vater. HErr, du eingeborner Sohn, JEsus Christus.}

\quellseite{154}

\notensatz{music/154-01}{Bortniansky, Laudamus: HErr Gott, du Lamm Gottes, Sohn des Vaters, der du die}

\notensatz{music/154-02}{Bortniansky, Laudamus: Sünde der Welt trägst, erbarme dich unser! Der du die}

\notensatz{music/154-03}{Bortniansky, Laudamus: Sünde der Welt trägst, nimm auf unser Gebet! Der du sitzest zur}

\notensatz{music/154-04}{Bortniansky, Laudamus: Rechten des Vaters, erbarme dich unser! Denn du allein bist}

\quellseite{155}

\notensatz{music/155-01}{Bortniansky, Laudamus: heilig, du allein bist der HErr, du allein bist der Allerhöchste,}

\notensatz{music/155-02}{Bortniansky, Laudamus: JEsus Christus mit dem Heiligen Geiste in der Herrlichkeit}

\notensatz{music/155-03}{Bortniansky, Laudamus: Gottes des Vaters. Amen. Amen. Amen.}

\chapter{IV. Salutation}


Der Gruß, mit welchem der Liturg nach dem Gloria sich zu der Gemeinde wendet, und der Gegengruß, mit dem diese antwortet, ist Schriftwort und hat biblischen Vorgang. In einem älteren liturgischen Werke, (Miscellanea etc. 1565) heißt es: „Das ist ein freundlicher, heilsamer Gruß, so der Priester dem Volk entbeut. ... Ich acht, es komme aus dem alten Testament,
\quellseite{156}
da der Engel habe zu dem Gideon gesagt: Der HEerre sei mit Euch! (Der HEerr mit dir, du streitbarer Held! Richt. 6, 12.) Auch so habe Boas zu seinen Messoribus gesagt: Dominus vobiscum! (Ruth 2, 4.) Auch so sei die Wiederantwort: „Und mit deinem Geiste!“ aus der Epistel Pauli ad Timotheum. (2 Tim. 4, 22.: „Der HEerr Jesus sei mit deinem Geiste!“) Das bezeugt der Pabst Soter, der solches zu halten eingesetzt hat.“ (S. Löhe, Lit. Formulare 1842. 3. Heft S. 13.) Seiner ursprünglichen Form nach war er Friedensgruß und durfte zur Zeit der Arkan-Disciplin nur in der Missa fidelium gebraucht werden. Wie Calvör erinnert, bediente man sich nach Aufhören derselben der jetzt bräuchlichen Form.

Wo immer der Liturg als Wortführer der Gemeinde in Gebet und Lob, sei es im Haupt- oder im Nebengottesdienste, gegen Gott handeln soll, fordert er mit diesem Gruße die Gemeinde zum inneren Mitwandeln auf. Eben darum hat auch Gregor der Große den Gegengruß: „Und mit deinem Geiste“ der ganzen Gemeinde zugetheilt. Im 12. Jahrhundert entzog man jedoch der Gemeinde auch diese Betheiligung an der Liturgie und ließ sie hierin von den Meßdienern oder Ministranten vertreten, wie dies noch heutiges Tages in der römischen Kirche geschieht.

Indem die lutherische Kirche diesen Gruß in ihre Gottesdienste überhaupt aufnahm, machte sie von demselben gerade im Hauptgottesdienst noch mehr Gebrauch, als es jetzt der Fall ist. Sie stellte ihn nämlich nicht nur der Präfation voran, sondern auch der Lesung des Evangeliums und dem Kirchensegen. Die herrschende Weise von Gruß und Gegengruß war und blieb die des Singens. Aus Unkenntniß sang man schon im 17. Jahrhundert in manchen Gemeinden: „Und mit seinem Geist“, statt: „mit deinem Geist“.

Die liturgische Bedeutung des schönen, biblischen Grußes und Gegengrußes liegt auf der Hand. Mit ihm begrüßen sich Pastor und Gemeinde im Hause des HEerrn, mit ihm wünschen sie sich gegenseitig zum Gebet und Lob, wie zur Handlung und zum Brauch des Wortes und Sacramentes „die Gegenwärtigkeit des HEerrn“, der Mensch geworden ist und die Verheißung gegeben hat: „Wo zwei oder drei versammelt sind in meinem Namen, da bin ich mitten unter ihnen.“ Calvör erinnert an folgende schöne Auslegung des Cardinales Bona: „Durch die Priester grüßt uns Christus, indem er in uns Wohnung zu machen und uns mit sich zu vereinigen begehrt, ja uns durch die Empfangung des Sacramentes des Altars zu seinem Tempel zu bereiten sucht. Wie der Engel zur heiligen Jungfrau spricht: „Der HEerr mit dir!“ ehe der HEerr in ihrem
\quellseite{157}
Leibe Fleisch an sich nahm.“\footnote{Per sacerdotes salutat nos Christus, cupiens in nobis mansionem facere et nos secum unire, ac templum suum eucharistici sacramenti perceptione efficere. Sicut angelus St. virgini dixit: „Dominus tecum“, antequam in ejus utero carnem susciperet.} (Calvörs Rit. I, 468.) Und von dem Gegengruß theilt Bona folgende Erklärung von Florus mit: „die Kirche habe nichts Besseres dem Priester zu wünschen gewußt, als was der Priester der Kirche wünscht, nämlich, daß derselbe HErr, der die Kirche gewürdigt hat, bei ihr zu sein, auch den Prediger würdige, mit seinem Geiste zu sein.“

Aus verschiedenen Werken sind hier etliche Melodien zu beliebigem Gebrauche mitgetheilt, mit dem Bemerkten jedoch, daß nur bei der Präfation die Melodie für das Dominus vobiscum eine andere ist, als die bei der Collecte. Hat man sich aus den hier mitgetheilten Melodien eine erwählt, so bleibe man bei derselben, so oft der Gruß der Collecte vorausgeht. Die in D-dur gesetzten Melodien sollen beispielsweise nur andeuten, daß für eine Bariton- und Tenorstimme diese Tonart die entsprechendere ist, wie für den Gruß, so auch für Versikel und Collecte.

Selbstverständlich singt (oder spricht) der Liturg den Gruß gegen die Gemeinde und kehrt sich dann während des Gegengrußes wieder dem Altar zu, um nach Beendigung desselben sofort den Versikel zu intoniren. Die für den Gruß des Liturgen entsprechende Geberde ist: die Hände wie zum Segen ein wenig auseinander thun und dann falten.

\subsubsection{Salutation}
\notensatz{music/157-01}{Salutation mit Orgelintonation, C-Dur. Pastor: Der HErr sei mit euch!}
\notensatz{music/157-02}{Hannoversche Kirchenordnung, vierstimmige Antwort der Gemeinde. Und mit deinem Geiste.}
\notensatz{music/157-03}{Salutation mit Orgelintonation, D-Dur. Der HErr sei mit euch!}

\notensatz{music/157-04}{Vierstimmige Antwort in D-Dur. Und mit deinem Geiste.}

\quellseite{158}
\notensatz{music/158-01}{Salutation mit Orgelintonation, C-Dur. Der HErr sei mit euch!}
\notensatz{music/158-02}{Traditionelle vierstimmige Antwort. Und mit deinem Geiste.}

\notensatz{music/158-03}{Salutation mit Orgelintonation, D-Dur. Der HErr sei mit euch!}

\notensatz{music/158-04}{Bayreuther Chorordnung 1755, vierstimmige Antwort. Und mit deinem Geiste.}

Anmerkung. Die römische Weise, der offenbar die vorstehenden Beispiele nachgebildet sind, ist folgende:

\notensatz{music/158-05}{Römischer Gruß und Antwort. Dominus vobiscum. Et cum spiritu tuo.}

\chapter{V. Versikel}


Unter Versikel (versiculus), hin und wieder auch Antiphon und Intonation genannt, versteht man einen vor der Collecte antiphonisch gesungenen Bibelvers oder Theil eines Bibelverses. Wie es scheint, hat der Versikel seinen Ursprung im kirchlichen Psalmengesang, nach welchem die zwei parallelen Hälften je eines Verses in melodischem Zuruf und Gegenruf gesungen werden. Es ist ein Singen „um einander“ (Ps. 147, 7.) oder „gegen einander“ (Vgl. Gesangb. 343, 2.).

Die römische Kirche bedient sich des Versikels vornehmlich in den Nebengottesdiensten, zum Theil aber auch in der Messe, doch dann mit verschiedener Melodie. Die lutherische Kirche aber gab dem Versikel nicht nur im Nebengottesdienst,
\quellseite{159}
sondern vielfach und zuletzt allgemein auch für die Anfangs- und Schlußcollecte der Messe seine Stelle. So heißt es in der Magdeburger Kirchenordnung: „Folgendes mit der ganzen Gemeine der Lobgesang: Allein Gott in der Höh sei Ehr, und darauf von dem Diacono nach geschehenem Gruß: Der HEerr sei mit euch! eine deutsche Collecte in den Festtagen mit dem Versikel, außer denselben aber an Sonntagen ohne Versikel gesungen“ etc. Während aber diese den Versikel nur für die Festtagsmesse gebraucht wissen will, bestimmen ihn die meisten Kirchenordnungen Thüringens auch für die Sonntagsmesse, und folgte man dann diesem Brauch fast überall. An Festtagen wurden zur Auszeichnung alsdann jedoch zwei Versikel gesungen. So ordnet es auch die St. Louiser Agenda.

Bei der Wahl des Versikels richtet man sich nach der Festzeit, außer derselben nach dem Inhalt der betreffenden Collecte. In der Passionszeit pflegt man das Halleluja eines Versikels wegzulassen.

Anfangs sang nur der Chor den Gegenruf des Versikels. Damit aber auch die Gemeinde denselben mitsingen könne, gab man später dem Gesangbuch eine Versikel-Sammlung bei, welchem Vorgang denn auch allhier die beiden Gesangbücher der Synodalconferenz, sowie das „Kirchenbuch“ des Generalconcils (General Council) gefolgt sind. Da nun aber von der Zeit des herrschenden Rationalismus her und absonderlich hier im Lande des reformirten Sectenthums Verständnis und Geschmack in liturgischen Dingen erst wieder gewonnen werden müssen, so ist es nöthig, auf einige Ungehörigkeiten im Singen des Versikels beim Respondiren hinzuweisen.

Gewöhnlich singt man den Gegenruf fast ohne allen Rhythmus, das heißt, man singt eine Silbe so lang gedehnt wie die andere. Dadurch wird bei längeren Versikeln das Respondiren ungemein schleppend und ermüdend und gleicht eher dem Lesen des ABC-Schützen in seiner Fibel, während es doch ebenso wie die Intonation und das Collectensingen des Liturgen mehr ein singendes Lesen (choraliter legendi) sein soll, also ein Singen unter Beobachtung des Wortaccentes und der Interpunction. Eine andere Ungehörigkeit betrifft die Orgelbegleitung, deren man sich beim Singen des Gegenrufs zu bedienen pflegt: Für diejenigen Versikel namentlich, bei denen die Gemeinde etwas länger in dem laufenden Ton zu singen hat, sucht man sich nämlich mit einem möglichst großen Wechsel der Accorde zu helfen. Je geringer aber die Kenntniß der Harmonisirung seitens des Organisten ist, desto ungehöriger fällt dieser Accordenwechsel aus und es entzieht dann ein um so widerwärtigeres
\quellseite{160}
Auf und Nieder, ein Hin- und Hertaumeln in der Begleitung. Ebenso ungehörig ist es, bei der Begleitung den laufenden Ton als Tonica behandeln, während derselbe gleichwie meistens in der Collecte immer Quinte, der 5. Ton, ist und als solcher den Dominanten-Accord fordert. Ist also z. B. der laufende Ton g, so darf nicht der G-Accord genommen werden, sondern es muß der C-Accord sein: c e g. Gleichermaßen hat sich im Gebrauch der Modulation bei einem Einschnitt, Absatz, Ruhepunkt eine fehlerhafte Weise unter uns verbreitet. Man singt nämlich:

\notensatz{music/160-01}{Dein Wort ist die Wahrheit: Beispiel der beanstandeten Singweise, mit Melisma auf Wahr}

Es soll aber so gesungen werden:

\notensatz{music/160-02}{Dein Wort ist die Wahrheit: berichtigte Singweise, ohne Melisma}

Auf die auch bei dem Versikel stattfindende recitativische Vortragsweise kommen wir bei der Collecte nochmals zurück. Daher hier nur noch ein paar praktische Bemerkungen zu den hier mitgetheilten musikalischen Beispielen.\footnote{In der St. Louiser Agende ist nicht nur für den Schluß eines jeden Versikels mit oder ohne Halleluja, sondern auch für den Gegengruß der Salutation, der Responsorien bei der Präfation und für sämmtliche Amen folgende Weise vorgeschrieben:
\notensatz{music/160-03}{St. Louiser Agende, gemeinsame Weise mit fünf Unterlagen: Dein Wort ist die Wahrheit. Halleluja! Dasselbe ist unsers Herzens Freude und Trost. Erheben wir zum HErrn. Und mit deinem Geiste. Amen. Neun G, dann A–H–C–D–D; alle Noten als Ganze, Bogen von A nach C.}
Der Verfasser konnte diese, wie es scheint, aus der Zeit des Verfalls stammende Singweise in keiner der ihm zu Gebote stehenden verschiedenen Liturgien, Schöberleins reichen Schatz etc. mit eingeschlossen, finden. Bloß in der von Schenk 1851 herausgegebenen musikalischen Agende steht sie unter etlichen Beispielen, und auch da nicht für den Versikel, für sonstige Responsorien, noch weniger für die Amen, sondern allein für die Salutation, auch glaubt er sie an dieser Stelle auch einmal in Rußwurms musikalischer Agende 1826 gesehen zu haben. Möge sie mehr und mehr der alten weit schöneren Weise weichen, nachdem diese neben jener in das neue St. Louiser Choralbuch von Hölter aufgenommen ist.}

\quellseite{161}
1. Diese verschiedenen Beispiele sollen nicht nur zeigen, daß und wie man an verschiedenen Orten je eine besondere Weise des Modulirens im Versikel hatte, sondern auch aufmerksam machen, daß und wie man zu beachten hat, ob der Versikel eingliederig oder zweigliederig ist, ob die Zeile mit einem einsilbigen oder mehrsilbigen Worte, mit oder ohne Halleluja endet. Man mag sie daher als ein Schema gebrauchen, um nach ihr sich eine stehende Weise zu wählen. Mögen aber Liturg und Vorsänger behufs Vortrags den betreffenden Versikel zuvor sich ein wenig ansehen.

2. Da für den singenden Lector beim Respondiren der Gemeinde im Accentuiren der Worte und Silben schwerlich das gewünschte Ziel sich wird erreichen lassen, so ist bei diesen Beispielen durch den Gebrauch von Noten verschiedener Werthe auf Grund gemachter Erfahrung versucht, zu zeigen, wie auf einfache Weise etwas Rhythmus in dieses Respondiren zu bringen wäre. Singt man nämlich bei Absätzen und besonders die Worte und Silben am Schluß der Zeile gedehnt, alles Uebrige, im gleichen Ton gehende kürzer, rascher, und spricht der Vorsänger Wort und Silbe recht deutlich aus, und spielt, wo die Orgel begleitet, der Organist dabei im Pedal staccato, so läßt sich annähernd schon etwas erreichen. Es verlohnte sich daher schon der Mühe, wenn der Cantor oder Organist nach der hier angegebenen Weise sämmtliche Festtagsversikel, sowie die sonst am häufigsten vorkommenden Versikel alphabetisch geordnet in Noten setzte.

\section{1. Ohne Halleluja}
\subsection{a. Eingliederig}

\notensatz{music/161-01}{Versikel des Pastors nach Wollersheim 1858. HErr, erhalte uns dein Wort.}

\notensatz{music/161-02}{Vierstimmige Gemeindeantwort. Dasselbe ist unsers Herzens Freude und Trost.}

\quellseite{162}
\notensatz{music/162-01}{Versikel des Pastors. HErr, lehre mich thun nach deinem Wohlgefallen.}

\notensatz{music/162-02}{Bayreuther Chorordnung 1755, alternative Intonation. HErr, lehre mich etc.}

\notensatz{music/162-04}{Vierstimmige Gemeindeantwort: Dein guter Geist führe mich auf ebner Bahn. Mit alternativem Anfang.}

\subsubsection{Traditionell}
\notensatz{music/162-03}{Traditioneller Versikel des Pastors. Weise mir, HErr, deinen Weg, daß ich wandle in deiner Wahrheit.}

\notensatz{music/162-07}{Vierstimmige Gemeindeantwort: Erhalte mein Herz bei dem Einigen, daß ich deinen Namen fürchte.}

\section{2. Mit Halleluja}
\subsubsection{Traditionell}
\notensatz{music/162-05}{Traditioneller Versikel des Pastors. Bereitet den Weg dem HErrn, Halleluja!}

\notensatz{music/162-06}{Alternative Halleluja-Kadenz des Pastors.}

\notensatz{music/162-08}{Vierstimmige Gemeindeantwort: Machet seine Steige richtig. Zwei Halleluja-Kadenzen.}

\quellseite{163}
\notensatz{music/163-01}{Versikel des Pastors: Hosianna dem Sohne David. Zweite Unterlage: Heilige uns, HErr, in deiner Wahrheit.}

\notensatz{music/163-02}{Alternative Halleluja-Kadenz.}

\notensatz{music/163-03}{Vierstimmige Gemeindeantwort: Hosianna in der Höhe. Zweite Unterlage: Dein Wort ist die Wahrheit. Zwei Halleluja-Kadenzen.}

\subsection{b. Zweigliederig}

\notensatz{music/163-04}{Pastor: Du Hirte Israels, höre, der du Joseph hütest wie der Schafe. Zwei Halleluja-Kadenzen.}

\notensatz{music/163-05}{Vierstimmige Gemeindeantwort: Erscheine, der du sitzest über Cherubim, und komme uns zu Hilfe. Zwei Halleluja-Kadenzen.}

\quellseite{164}
\chapter{VI. Collecte}


Der Ursprung dieser Art von Kirchengebeten gehört fast dem apostolischen Zeitalter an, denn schon Tertullian und Justin gedenken ihrer. Woher nun aber der Name für diese Gebete kommt — ob daher, daß es ein Gebet der Versammlung, oder eine Summa, eine Zusammenfassung der Bitten dieser Versammlung ist, oder daß bei den altkirchlichen Oblationen, der Einsammlung der Opfergaben, etwa dergleichen kurze Gebete gebraucht wurden, wird wohl unentschieden bleiben. Luther leitet ihn von jenen Oblationen her, indem er in seiner Schrift wider die himmlischen Propheten bemerkt: „So haben nu die Apostel und erste Christen, da noch Ebräische Sprache unter ihnen gemein war, das Brod und Wein, so sie zusammentrugen zum Sacrament, Missa ebräisch genennet, auf Jüdische Weise, daß darnach eins Theils gesegnet zum Sacrament, das ander unter die Diener der Gemeinde und Armen ausgetheilet ward, welches sie hernach lange Zeit auch Collecten hießen, von demselbigen Zusammentragen, wie die Historia Tripartita zeuget, davon noch das Wort Collecta in der Papisten Messe blieben ist.“ (Erl. Ausg. 29, 181.) Unter Anwendung eines in der Uebersetzung nicht wiederzugebenden Wortspiels äußert sich Calvör in seinem Rituale so: „Die Collecten sind ohne Zweifel kurze Gebete der Priester, welche sie nicht für sich, sondern mit der ganzen Versammlung (Collecta) und Kirche verrichten und welche daher das versammelte Volk mitliest (collectus col-legit) oder zugleich liest und mit dem Priester für das gemeine Wohl und die gemeine Nothdurft bittet.“ (I. C. Cap. XIII.) Aehnlich die „Cöllnische Reformation“ von 1544, nach welcher sie ist „ein gemein Gebet der ganzen Gemein, darum es auch ein Collect genannt ist, ein gesammelt Gebet von dem Herzen und dem Leben aller Gegenwärtigen oder ein Gebet der gesammleten Gemein“. Genug, der Sache nach ist die Collecte ein Gebet der versammelten Gemeinde, in welchem sie alles, worinnen sie nach Zeit und Gelegenheit zu bitten eins ist, aufs kürzeste zusammenfaßt und aufs bestimmteste ausdrückt. „Eine einzige Bitte — gerichtet zu Gott dem Vater — im Namen seines eingebornen Sohnes, wie es der Sohn selbst gebeut: ein einziger, vereinter Seufzer der ganzen Gemeinde, ein Wort aus ewige Vaterherz, über welchem alle eins geworden sind, — das ist die Collecte, und ist es um so mehr, je mehr sie ist, was sie soll. Je weniger sie das ist, je weniger ist sie den herrlichen Collecten der uralten Zeit der Christen zu vergleichen, je mehr ist sie aus der Art geschlagen.“ (Löhe.)

\quellseite{165}
Die Reformation fand einen Reichthum antiker Collecten vor, und unter ihnen nicht wenige, die noch rein waren. Sie nahm daher eine Anzahl mit herüber, verdeutschte sie, brauchte sie aber auch nebenbei im Latein. Außerdem aber gingen aus der lutherischen Kirche selbst eine Anzahl Collecten hervor. Die vorzüglichsten sind wohl die von Veit Dietrich 1541 und Johann Mathesius 1568. Die Collecten von Veit Dietrich sind Evangeliencollecten, also Textcollecten. Dieselben sollten jedoch allein in den Schulen und Häusern\footnote{Zum Hausgebrauch sind sie daher auch in den in St. Louis erschienenen (großen) „Gebetsschatz“ S. 48–70 aufgenommen.} von der Jugend gelesen und auswendig gelernt werden, und obwohl durch die von M. A. Pancratius 1572 herausgegebenen „Kurzen Summarien und Gebetlein über Sonn- und Festtagsepisteln und Evangelien“ nebst den von Major und Pancratius verfaßten Epistelcollecten in die Kirchen gekommen, wollten sie doch keineswegs die Collecten vor der Epistel verdrängen, denn sie sollten nebst den Summarien nach den Texten verlesen werden. Die des originellen, „viel guten alten HEerrn“ Mathesius sind „Andächtige und christliche gemeine Gebetlein, für allerlei Noth der Christenheit“ und haben besonders in der an Collecten reichen „christlichen Kirchen-Agenda von Oestreich u. E.“ eine besondere Berücksichtigung gefunden. Die St. Louiser Agenda enthält für den Altardienst allein solche Collecten, welche Gemeingut der lutherischen Kirche geworden sind.

Obwohl nach der Weise der römischen Messe manche der frühesten Kirchenordnungen vor der Epistel mehr als eine Collecte ordnen — zuerst eine Collecte de festo oder de tempore und sodann zwei Collecten für gemeine Noth —, so wurde doch der Brauch von nur Einer Collecte herrschend: für die Feste eine Festtagscollecte, für die Sonntage eine gemeine, dem Charakter des Tages möglichst entsprechende Collecte.

Der Collecte geht auch in der lutherischen Kirche das altkirchliche „Oremus“ oder „Laßt uns beten!“ voraus:

\subsubsection{(Ursprünglich.)}

\notensatz{music/165-01}{Ursprüngliches Oremus: Oremus! Laßt uns beten! G–F–G–G und G–G–F–G–G, mit Breven am Phrasenschluß und Melismen auf re und be.}

\subsubsection{(Später.)}

\notensatz{music/165-02}{Späteres Oremus: dieselben beiden Formeln mit Fis statt F; im Original ohne eigene Textunterlegung.}

\quellseite{166}
Dieser Aufruf zeigt der Gemeinde an, daß der Liturg das Gebet im Namen der Gemeinde zu Gott thue, weshalb diese durch ihr Amen zu solchem Gebete sich bekennt und es bekräftigt. Diese u. a. nach Nehem. 8, 6. und 1 Cor. 14, 16. alte biblische und apostolische Sitte, die gottesdienstlichen Gebete durch das Amen der Gemeinde bekräftigen zu lassen, wurde in der Kirche der ersten Jahrhunderte allgemein beobachtet und erhielt sich in der römischen Kirche bis ins 12. Jahrhundert und theilweise noch drüber hinaus. Wie sollte daher Luther diese Sitte nicht wieder zurückführen, da sich nicht nur durch das Oremus des Liturgen, sondern auch und vornehmlich durch das Amen der Gemeinde die Collecte als Gemeingebet erweisen soll. Wo daher die Collecte selbst auch nur vom Liturgen gesprochen werden kann, sollte von der Gemeinde das Amen gesungen werden, am besten dann in der einfachsten Weise, wie:

\notensatz{music/166-01}{Amen der Collecte: zweimal G, mit Fermate auf der zweiten Silbe}

Die Collecte hat ihre bestimmte Schlußformel. Ist jene an den Vater gerichtet — und das ist sie meistentheils —, so lautet diese entweder: „Durch Jesum Christum, unsern HErrn“, oder: „durch Jesum Christum, deinen Sohn, unsern HErrn“, oder: „um Jesu Christi, deines Sohnes, unsers HErrn willen“, manchmal mit der doxologischen Erweiterung: „Der mit dir in Einigkeit des Heiligen Geistes (oder: und dem Heiligen Geiste) lebet und regieret (oder: herrschet) von Ewigkeit zu Ewigkeit“ (oder: „immer und ewiglich“). Geht die Collecte an den Sohn, so schließt sie: „Der du lebest und regierest mit Gott, dem Vater und dem Heiligen Geiste von Ewigkeit zu Ewigkeit“ oder: „immer und ewiglich.“ Geschieht Erwähnung des Sohnes in der Collecte, so ist der Schluß: „Durch denselben deinen Sohn“ etc., und geschieht Erwähnung des Heiligen Geistes, was jedoch selten vorkommt, so schließt man: „In der Einigkeit desselben Heiligen Geistes, der mit dir und dem Sohne lebet“ etc.

Nach altkirchlichem Gebrauch geschieht solch „kurz Gemeingebet“ dem Altar zugekehrt, und nicht, was gegen allen liturgischen Anstand und Verstand wäre, gegen die Gemeinde. Die meisten älteren Kirchenordnungen schreiben ausdrücklich diese Stellung des Liturgen vor, und die es nicht thun, setzen sie als selbstverständlich und bekannt voraus. So heißt es in einer derselben, der Braunschweig-Lüneburgischen: „Darauf (nach dem Gruß) wende
\quellseite{167}
sich der Priester wieder gegen den Altar und singe eine Collectam de Tempore oder Festo, wie dieselbe sich auf das Evangelium schicket.“ Wie sollte die Stellung eine andere sein? Einem bayerischen Pfarrer, der beim Erscheinen der Löhe’schen Agende einen Feuerruf erhob (später jedoch widerrief) und der auch in dieser Stellung Papisterei sah, antwortete Bucherer: „Wir haben in unsern lutherischen Kirchen Altäre, wir haben nicht bloß Abendmahlstische. Vom Begriff des Altars ist aber der eines darauf darzubringenden Opfers nicht zu scheiden. Von einem unblutigen Opfer im Abendmahl wissen wir zwar nichts und wollen in Ewigkeit nichts davon wissen, aber von einem Rauchopfer des Gebets (Offenb. 8, 3. 4.), von einem Opfer bußfertiger Herzen (Ps. 51, 19.), von einem Opfer des Dankes (Ps. 50, 14. 23.) wissen wir; und solche geistliche Opfer bringen wir Gott auf unsern Altären dar. Darum steht auch die ganze Gemeinde als ein priesterlich Volk (1 Petr. 2, 9.) bei solchen Opfern dem Altare zugekehrt, und der Pfarrer, der . . . die Gebete im Namen der Gemeinde spricht, sollte dabei dem Altar den Rücken zukehren?“ Zur Befestigung dieser Antwort beruft sich dann Bucherer auf den Liturgiker Dr. Höfling und führt aus dessen Schrift „von der Composition der christlichen Gemeindegottesdienste“ folgende Stelle an: „Wie Kanzel und Lesepult in der Kirche die Richtung auf die Gemeinde repräsentiren, so ist der Altar umgekehrt der Ort, welcher die unmittelbare Richtung auf Gott symbolisch darstellt. . . . Und zwar liegt es weiter in der Natur der Sache, daß der Liturg, welcher das Gebet der Gemeinde als ihr Wortführer und Stellvertreter vor Gott bringt, also mit einem Auftrage von der Gemeinde herkommt, diese nicht vor, sondern hinter sich hat, nicht ihr, sondern dem Altar zugewendet spricht. Wenn Gegentheiliges in der kirchlichen Praxis vorkommt, so können wir es immer nur als eine liturgische Unschicklichkeit betrachten.“

Wir kommen nun noch insbesondere zu der Weise des liturgischen Vortrags der Collecte. Dieselbe ist, wie schon beim Versikel bemerkt worden, die Weise des Leisetons (modus choraliter legendi) — ein „Recitativ-Gesang, der bekanntlich die Mitte hält zwischen der eigentlichen Rede und dem eigentlichen Gesang und nichts anderes ist und sein soll, als eine höhere, feierliche Erhebung des Tons und des Redegangs, in dem das wichtige Wort weithin schalle“. Auch sie hat die lutherische Kirche aus der römischen Messe mit herübergenommen und sich in der Beibehaltung derselben bis auf diesen Tag trotz alles Widerspruchs der Reformirten und geistreicher Schwärmer überhaupt aus guten Gründen nicht irre machen lassen; denn nicht nur wollte man
\quellseite{168}
„keine merkliche Veränderung“ machen, nicht nur ist es den Christen freigelassen, ein Gebet sprechweise oder singweise vor Gott zu bringen, und nicht nur ist dieses choraliter Lesen aus dem Psalmengesang hervorgegangen, sondern es hat sich auch diese Weise gerade für diese kurzen „Gemeingebete“ als sehr passend erwiesen. Es ist nämlich diese Weise des Lesens feierlich, und vor allem, sie macht auch allen Subjectivismus im Vortrag und in der Betonung unmöglich. Hier ist ja der Liturg Organ der betenden Gemeinde, hier soll er nur die helle, objective Stimme der Beter sein und darf darum von dem Seinen nichts hinzuthun. Eben darum muß auch beim Sprechen der Collecte und anderer Gebete im Gottesdienste alles Eigene, alles Declamationsmäßige wegfallen und ist insbesondere bei der Collecte nöthig, daß sie da, wo sie gesprochen wird, ohne sonderliche Betonung, aber mit Erhebung der Stimme vom Liturgen vorgetragen werde.

In der römischen Kirche ist die Gesangesweise der Collecten entweder festiv oder ferial. Festiv, feierlich werden die Collecten bei mehr oder weniger festlichen, feierlichen Handlungen, sowie in den Gottesdiensten des Sonn- und Festtages gesungen, in allem übrigen ist die feriale, einfache Weise vorgeschrieben. Diese besteht darin, daß die Collecte vom Anfang bis zum Ende nach dem laufenden Ton (tonus currens) gesungen wird, allein unter Beobachtung der Interpunction und Accentuation des Textes. Bei der festiven Weise jedoch kommt der Accentus ecclesiasticus zu seiner Verwendung. In seiner „Lehre vom römischen Choralgesang“, 1842, stellt Vilsecker die dabei geltenden Regeln und veranschaulicht sie durch Beispiele. Obwohl diese natürlich es nur mit der lateinischen Oration zu thun haben, so sind sie doch für das Singen deutscher Collecten instructiv, und wenn schon man in unserer Kirche sich nicht strenge an diese Weise band, so nahm man sich dieselbe doch mehr oder weniger zum Muster, ja, schloß sich mitunter, soweit es die deutsche Sprache zuließ, möglichst an dieselbe an, wie wir dieses an seinem Orte auch insonderheit an den Präfationsweisen sehen werden.

Die in Rede stehenden Regeln sind diese:

Der Anfang der Oration geschieht hier mit erhobener Stimme, den angefangenen Ton c oder f festhaltend,\footnote{c für eine höhere, f für eine tiefere Mannesstimme. In beiden Tonhöhen pflegte man auch in der lutherischen Kirche in der Regel zu singen. In diesem Werke ist, wie auch in der St. Louiser Agende, g als Durchschnittston angenommen.} bis ungefähr in der Mitte der Oration
\quellseite{169}
die erste Veränderung eintritt, welche man den Hauptpunkt (punctum principale) nennt, und zwar bei dem Doppelpuncte, welcher den ersten Theil der Oration schließt.

Man singt nämlich bei der vierten oder fünften Silbe vor dem Doppelpuncte zwei Secunden abwärts, und kehrt wieder bei den letzten zwei oder drei Silben in den Hauptton zurück entweder durch c h a c oder f e d f. Ist nun die Oration lang genug, daß sie zwei Veränderungen ertragen kann, so tritt bei dem zweiten Schlußsatze der Halbpunkt (Semipunctum) ein und zwar vermittelst der Töne c h oder f e. Sollte jedoch die Oration zu kurz sein, so wird nur der Hauptpunkt gehalten. Sind sowohl der Haupt- als auch der Halbpunkt vorüber, so wird weder ein zweiter Haupt- noch ein Halbpunkt gehalten, sollten auch noch mehrere Schlüsse (clausulae) folgen. Bei beiden Punkten wird ein wenig ausgehalten, um Athem zu holen. — Kommen in einer Oration zwei Doppelpuncte vor, so trifft der Hauptpunkt auf den letzteren Doppelpunkt unmittelbar vor dem Halbpunkte.

Bei den Schlußsätzen des festlichen Tones einer Oration: Per Dominum und per eundem Dominum („durch den HEerrn“ und „durch denselben HEerrn“) fällt der Halbpunkt, und zwar auf das Wort tuum: c h oder f e, und der Hauptpunkt h a | c c oder e d | f f auf Sancti Deus. Bei den Schlußsätzen: Qui tecum oder qui vivis kommt nur der Hauptpunkt e d f f oder h a c c vor und der Halbpunkt fällt ganz weg. Die letzten Silben einer jeden Oration müssen immer etwas länger gehalten werden, als die übrigen Worte. (In den deutschen entsprechenden Schlußsätzen käme der Halbpunkt statt auf „deinen“ auf „Sohn“ c c h und statt auf „sancti Deus“ auf „wahrer Gott“ e d f oder h a c.)

Zur „Wiederholung der Regel“ faßt dann Vilsecker alles in folgendes Beispiel:

\notensatz{music/169-01}{Lateinische Formel des festlichen Collectentons: Sic incipit Oratio, sic canitur Punctum principale: sic autem Semipunctum; et sic finitur oratio. Dreißig Halbe mit abwärts gerichteten Hälsen und vier Viertel auf ci, ti, ni und ti.}

\quellseite{170}
Zu desto besserer Veranschaulichung ist jedoch auch die lateinische Collecte vom Tage Petri und Pauli in Noten beigegeben. Da diese Weise mehr oder weniger für manche der ältern deutschen Collecten nachgeahmt, ja, beispielsweise eine derselben ganz angepaßt sich findet (vgl. Nr. 3 der nachfolgenden Collectenweisen), so möge dies Beispiel zur Orientierung auch hier stehen, zugleich jedoch auch dem deutschen Text angepaßt, wie er sich in der dritten Auflage der Löhe’schen Agende findet.

\section{Beispiel des festlichen Tones}

\notensatz{music/170-01}{Festlicher Collectenton. Die lateinische Collecte vom Tage Petri und Pauli und ihre deutsche Fassung, Anfang.}

\quellseite{171}
\notensatz{music/171-01}{Festlicher Collectenton. Die lateinische Collecte und ihre deutsche Fassung, Fortsetzung und beide Amen.}

Für den Gesang der Collecte nach der kirchlichen Tradition gibt Dr. Herzog im musikalischen Anhang zur bayerischen Agende folgendes als Hauptregel an:

\begin{quote}
„Er beginnt in der Quinte; die Silbe oder das Wort mit besonderem Accent enthält die untere Terz; das Komma oder Semicolon geht herab zur Terz und von da wieder zurück zur Quinte; ein Punktum geht von der Terz mit durchgehender erhöhter Quarte zur Quinte; der Schluß steigt zur Terz und geht stufenweise herab bis zur Terz. Z. B.:“
\end{quote}

\notensatz{music/171-02}{Herzogs Collectenton: Anfang, Accent, Komma oder Punkt, Kadenz.}

\quellseite{172}
\begin{quote}
„Die Monotonie, welche bei dem Collectengesang, namentlich bei längeren Collecten, hervortritt, kann dadurch gemildert werden, daß auch die zwischen der Terz und Quinte liegende Quarte benützt wird.“
\end{quote}

Zur Erläuterung durch ein Beispiel diene Nr. 5 der nachfolgenden Collectenweisen. Die zur Milderung der Monotonie benützte Quarte zeigen die Sätze: „daß dein Wort unter uns laufe und wachse“, und: „auf daß wir mit beständigem Glauben dir dienen“.

Was nun schließlich die hier folgenden Weisen für die Collecten und für die Amen zu denselben betrifft, so sind sie sämtlich den älteren Sammlungen entnommen; namentlich denen des 16. Jahrhunderts. Es sind ganz einfache und es sind schwungvollere Collectenweisen, alle aber von gleichem Charakter. Die einfachsten unter ihnen haben nicht nur in solcher Einfachheit ihr Schönes, sondern dürften sich um derselben willen auch den Pastoren empfehlen, denen es schwer fällt, sich eine etwas schwungvollere Weise anzueignen.

Auch der Weisen für das Amen zu den Collecten gibt es vielerlei. Ob aber einfacher oder reicherer Form, immer nur ist Ein Amen der Collecte zugewiesen, nicht ein zweifaches oder dreifaches, wie beim Segen.

\section{A. Weisen zu den Collecten}

\subsection{Nr. 1}

Libellus antiphon. Hof 1605.\\
Libellus continens Ant. Onolzbach 1627.

\notensatz{music/172-01}{Collecte Nr. 1, Anfang: Laßt uns beten! O Gott, von dem alle gute Gabe kommt … der mit dir lebet und regieret in.}

\quellseite{173}
\notensatz{music/173-no1-close}{Collecte Nr. 1, Schluß: Einigkeit des Heiligen Geistes von nun an bis in Ewigkeit. Mit Medius-Variante.}

\subsection{Nr. 2}

Braunsch.-Lüneb. K.-O. 1531. 1569. 1657.\\
Oldenb. K.-O. 1573. Hoya K.-O. 1531.\\
Ludecus Missale 1589. Verden K.-O. 1606.

\notensatz{music/173-no2}{Collecte Nr. 2: Allmächtiger HErre Gott, weck uns auf.}

\subsection{Nr. 3}

Hommel, Lit. luth. Gem.-Gottesd. 1851.

\notensatz{music/173-no3-opening}{Collecte Nr. 3, Anfang: Allmächtiger, ewiger Gott … mit allen ihren Gliedern.}

\quellseite{174}
\notensatz{music/174-no3-normal}{Collecte Nr. 3, Schluß: in reinem Glauben … um Jesu Christi … willen.}

\notensatz{music/174-no3-alternative}{Collecte Nr. 3, anderer Schluß: durch unsern HErrn Jesum Christum … von Ewigkeit zu Ewigkeit.}

\subsection{Nr. 4}

Hommel 1851.

\notensatz{music/174-no4}{Collecte Nr. 4: O HErr, wir bitten, gib uns allzeit gnädiglich einen Geist, zu gedenken und zu thun, was recht ist.}

\subsection{Nr. 5}

Layriz, Kern des deutschen A. G., IV. Abtheil., lit. Weisen enth., 1855.

\notensatz{music/174-no5-opening}{Collecte Nr. 5, Anfang: Allmächtiger HErre Gott, wir bitten dich, gib dei-.}

\quellseite{175}
\notensatz{music/175-no5-middle}{Collecte Nr. 5, Fortsetzung: -ner Gemeine deinen Geist und göttliche Weisheit … dadurch gebess-.}

\notensatz{music/175-no5-close}{Collecte Nr. 5, Schluß: -sert werde … durch Jesum Christum, deinen lieben Sohn, unsern HErren.}

\subsection{Nr. 6}

Layriz, Kern v. 1855.

\notensatz{music/175-no6-opening}{Collecte Nr. 6, Anfang: Wir danken dir, HErr Gott, himmlischer Vater … solch selig Licht deines Worts uns gnä-.}

\quellseite{176}
\notensatz{music/176-no6-close}{Collecte Nr. 6, Schluß: -diglich erhalten … durch Jesum Christum, deinen Sohn, unsern HErren.}

\subsection{Nr. 7}

Layriz, Kern v. 1855.

\notensatz{music/176-no7-opening}{Collecte Nr. 7, Anfang: HErr Gott, himmlischer Vater, der du heiligen Muth, guten Rath und rechte Werke schaffest … auf daß un-.}

\notensatz{music/176-no7-middle}{Collecte Nr. 7, Fortsetzung: -sere Herzen an deinen Geboten hangen … vor Feinden leben.}

\notensatz{music/176-no7-close}{Collecte Nr. 7, Schluß: durch Jesum Christum, deinen Sohn, unsern HErren. Mit zweiter Schlußweise.}

\quellseite{177}
\section{B. Amen zu den Collecten}

\subsubsection{1}

\notensatz{music/177-01}{Amen zu den Collecten, Nr. 1. Vierstimmig.}

\subsubsection{2}

\notensatz{music/177-02}{Amen zu den Collecten, Nr. 2. Vierstimmig.}

\subsubsection{3}

\notensatz{music/177-03}{Amen zu den Collecten, Nr. 3. Vierstimmig.}

\subsubsection{4}

\notensatz{music/177-04}{Amen zu den Collecten, Nr. 4. Vierstimmig.}

\subsubsection{5. (Zu Nr. 6.)}

\notensatz{music/177-05}{Amen zu den Collecten, Nr. 5. Vierstimmig.}

\subsubsection{6}

\notensatz{music/177-06}{Amen zu den Collecten, Nr. 6. Vierstimmig.}

\subsubsection{7}

\notensatz{music/177-07}{Amen zu den Collecten, Nr. 7. Vierstimmig.}

\subsubsection{8}

\notensatz{music/177-08}{Amen zu den Collecten, Nr. 8. Vierstimmig.}

\subsubsection{9. (Zu Nr. 5.)}

\notensatz{music/177-09}{Amen zu den Collecten, Nr. 9. Vierstimmig.}

\subsubsection{10. (Zu Nr. 7.)}

\notensatz{music/177-10}{Amen zu den Collecten, Nr. 10. Vierstimmig.}

\subsubsection{11}

\notensatz{music/177-11}{Amen zu den Collecten, Nr. 11. Vierstimmig.}

\quellseite{178}
\chapter{VII. Epistel}


Vorbereitet durch das „kurz Gemeingebet“, die Collecte, vernimmt die Gemeinde die Epistel und nach dieser das Evangelium des Sonn- oder Festtages, in denen nun der HEerr HEerr durch den Mund des Liturgen zu ihr redet. Warum erst Epistel und dann Evangelium, soll bei Besprechung des letzteren zu zeigen versucht werden.

Auch in diesen Lesungen behielt die lutherische Kirche nur eine altkirchliche Einrichtung bei; denn die gottesdienstliche Sitte des Alten Bundes, Abschnitte aus dem Gesetz und den Propheten vorzulesen (Luc. 4, 16. 17. Apost. 13, 14. 15. 15, 21.), nahm die alte Kirche frühzeitig in ihren Cultus herüber, verband aber bald damit die Lesung theils von Erzählungen aus dem Leben des HEerrn, theils von Abschnitten aus Briefen der Apostel. So berichtet Justin der Märtyrer († 163), daß in den sonntäglichen Zusammenkünften der Christen „die Denkwürdigkeiten der Apostel, die man Evangelien nenne, und die Schriften der Propheten“ gelesen wurden und daß sich hieran Predigt und Gebet schloß. Desgleichen rühmt Tertullian († 220) von der Kirche, daß sie das Gesetz und die Propheten mit den evangelischen und apostolischen Schriften zugleich lese. Auf diesem Wege fortschreitend ordnen die apostolischen Constitutionen die Lesung des Gesetzes und der Propheten, und zwar vor der des Neuen Testamentes, der Apostelgeschichte, der paulinischen Briefe und der Evangelien an (II, 39. 57). Später wurden die alttestamentlichen Lectionen in die Nebengottesdienste, speciell in die Vesper verlegt. Bald aber ergab sich auch die Nothwendigkeit, für die Hauptfeste und dann auch für die übrigen Feiertage diejenigen Abschnitte aus der heiligen Schrift auszuwählen, welche die Geschichte des Tages enthielten oder zu ihr in ausdrücklicher innerer Beziehung standen. Es bildete sich in der römischen Kirche das Perikopensystem.

Die lutherische Kirche hat diese Perikopenordnung für den Hauptgottesdienst nicht nur im Allgemeinen aufgenommen, sondern dieselbe auch noch weiter ausgebildet, indem sie namentlich den letzten Trinitatissonntagen Schriftabschnitte zuwies, welche von den letzten Dingen handeln. Zwar versuchten einige reformatorische Kirchen fortlaufende Lesungen; allein es fand das keine weitere Verbreitung. Man hielt sich mehr an Luthers Weise, der die fortlaufenden Schriftlectionen für die Nebengottesdienste, die Perikopenlesung aber
\quellseite{179}
für den Hauptgottesdienst ordnete. Da im Hauptgottesdienst vornehmlich über die Evangelien gepredigt werden sollte, so verband man hie und da mit der Epistel Summarien oder die Lesung einer kurzen Zusammenfassung ihres Inhalts, so z. B. nach der Officia missae, Hof 1605, und der Hennebergischen Kirchenordnung, 1582, „eine Auslegung Pankratii“ — eine Anordnung, die sich auch noch in der Ausgabe des Nürnberger Agendbüchleins von 1755 findet. Die im Hannover’schen noch heute bestehende Einrichtung, an die Stelle der Lesung des Evangeliums besondere fortlaufende Schriftlectionen zu setzen und so das Evangelium allein auf die Kanzel zu verlegen, datirt sich erst aus dem Jahre 1769 und empfiehlt sich nicht in Rücksicht auf die Gliederung des Hauptgottesdienstes. (Es ist daher auch nicht zu loben, wenn mancher Prediger bei den Lesungen sich nicht an die Perikopen und an die Aufeinanderfolge von Epistel und Evangelium hält.)

Selbst auch da, wo die Lesung anfänglich noch in lateinischer Sprache um der Lateinschule willen geschah, unterließ man selbstverständlich nicht die deutsche Lesung um des Volks willen, und schärfte dabei ein, daß die Lesung „fein laut, deutlich und langsam, ohne Auslassung einiger Worte und Syllaben“ geschehe.

Nach Luthers Vorgang findet sich durchweg die Bestimmung, daß die Lesung „gegen dem Volk geschehe“, der Liturg sich „zum Volke“ kehre, während dies in der römischen Messe bei den liturgischen Lesungen eigentlich nicht geschieht. Meistentheils stand er dabei am Altar oder auf der Südseite desselben oder er hielt die Lesung von einem besonderen Pulte, der „kleinen Kanzel“ aus, wie man dies noch auf Abbildungen hin und wieder sieht.

Daß nach Vorgang der römischen Messe die Weise der Lesung der Epistel und des Evangeliums auch Gesangesvortrag war und Luther in seiner deutschen Messe dazu Regeln und Beispiele gab, ist bereits § 5 des ersten Theils erwähnt worden. Es finden sich dieselben auch als Musikbeilage zu Band 22 der Werke Luthers, Erlanger Ausgabe, wie in Band X der St. Louiser Ausgabe, wo die Noten dem Text beigefügt sind. Obwohl nun auch hier Luther conservativ verfuhr, so bewog und trieb ihn doch auch der an der deutschen Bibelübersetzung erwachte Geist der deutschen Kirchensprache, „an Stelle des traditionellen einfachen, und doch in jeder Diöcese verschieden verschnörkelten Episteltons und Evangelientons eine schwungreichere, ausdrucksfähigere Form oder vielmehr zwei reformatorische Formen zu setzen, welche nicht als örtliche Umbildung oder zeitliche Fortentwicklung eines nachweisbar vorhandenen, vorhermotorischen
\quellseite{180}
Schemas zu erklären sind, obschon sie das Gepräge bestimmter gregorianischer Tonarten fest und deutlich an der Stirne tragen. Diese Formen waren, so neu sie dazumal erscheinen mußten, allerdings geeignet, ihre Vorgänger zu verdrängen für alle Zeit, so lange der feierlich cantillirende Vortrag der Perikopen an den lutherischen Altären nicht verstummen sollte“ (Lyra). Wie auch hier wieder Luther seine Meisterschaft bewies, hat uns das im ersten Theil, § 5, beigebrachte Zeugniß Johann Walthers gezeigt.

Der liturgischen Vortragsweise der Perikopen folgte man durchweg in Sachsen, desgleichen in Pommern, der Mark und einem Theil Niedersachsens; in einem andern Theile Niedersachsens, wie auch in Oesterreich, ließ man Sprechen oder Singen der Perikopen frei, während man anderwärts nur dann das Sprechen gestattete, wenn der Pastor keine hinreichende Gabe zum Singen hatte. Und sie hielt sich lange. So bezeugt Joh. Schrader in seinem „Formular-Buch“, 1621, daß dieselbe noch „in großen Städten gebräuchlich“ sei und daß zu Magdeburg etliche Jahre vorher ein Buch in Folio gedruckt worden sei, „darin alle Episteln und Evangelien unter die Noten gesetzt sein“. Die dem Schreiber vorliegende Leipziger Agende, 1681, enthält gleichfalls Anweisung und Exempel zum liturgischen Perikopen-Singen. Ja, noch bis in das 18. Jahrhundert hatte sie sich vielfach erhalten, und in der norwegisch-lutherischen Kirche besteht sie zum Theil noch.

Da unter der Herrschaft des Rationalismus in Deutschland und der des Puritanismus und Methodismus allhier das liturgische Singen überhaupt gefallen ist und man aus Unverstand in so vielen Gemeinden noch in demselben etwas Katholisches sieht, so kann von einer Wiedereinführung des Perikopen-Singens auch an den Orten nicht wohl die Rede sein, an welchen das liturgische Singen der übrigen Theile der „Messe“ überhaupt noch bräuchlich ist oder wieder eingeführt werden konnte. Damit sich jedoch die Kenntniß auch dieser Vortragsweise behufs liturgischen Studiums unter uns erhalte, wie auch der Vollständigkeit wegen, sei die eine von den drei bis vier gebräuchlich gewordenen Singweisen hier für die Epistel und an seinem Ort für das Evangelium mitgetheilt. Es ist dies die Lutherische Weise in ihrer, jedoch nur ganz geringen Umbildung, wie sie meist vorkommt und auch in der genannten Leipziger Agende von 1681 sich findet.

Nur noch die Bemerkung, daß für das gesanglose Lesen der Perikope gleichwohl der Declamations-ton möglichst zu vermeiden und mehr ein Recitativton anzuwenden ist.

\quellseite{181}
\section{1. Regeln für die Modulation der Epistel}

Deutsche Messe 1526. Braunsch.-Lüneb. 1542. Magdeb. 1615. Leipzig 1611.

\subsubsection{Initium (Anfang). Comma}

\notensatz{music/181-01}{Initium der Epistel: „(So beginnt man die Epistel,“}

\notensatz{music/181-02}{Comma der Epistel: und so singt man das Komma,}

\subsubsection[Comma aliud. Colon. Periodus]{Comma aliud (anderes Komma). Colon. Periodus}

\notensatz{music/181-03}{Anderes Komma der Epistel: das zweite Komma so,}

\notensatz{music/181-04}{Kolon der Epistel: so aber ein Kolon,}

\notensatz{music/181-05}{Punkt der Epistel: in dieser Weise den Punkt,}

\subsubsection{Quaestio (Frage). Finale (Ende)}

\notensatz{music/181-06}{Fragezeichen im Epistelton: das Fragezeichen so,}

\notensatz{music/181-07}{Schluß im Epistelton: den Schluß der Epistel aber so.}

\section{2. Epistelton}

\notensatz{music/181-epistel-a}{Epistelton, 1. Corinther 4. So schreibt Sanct Paulus an die Corinther. Lieben Brüder, dafür halt uns Jedermann, nämlich für Christus Diener, und Haushalter über Gottes Geheimniß. Nun sucht man nicht mehr an den Haushaltern,}

\notensatz{music/181-epistel-b}{Epistelton, Fortsetzung. Denn daß sie treu erfunden werden. Mir aber ist’s ein Geringes, daß ich von euch gerichtet werde, oder von einem menschlichen Tage; auch richte ich mich selber nicht. Ich bin mir wohl nichts bewußt, aber dar-}

\quellseite{182}
\notensatz{music/182-epistel}{Epistelton, Fortsetzung. -um bin ich nicht gerechtfertiget. Der HErr aber ist’s, der mich richtet. Darum richtet nicht vor der Zeit, bis der HErr komme, welcher auch wird ans Licht bringen, was im Finstern verborgen ist, und den Rath der Herzen offenbaren. Alsdann wird einem Jeglichen von Gott Lob widerfahren.}

\chapter{VIII. Graduale}


Mit dem Ausdruck Graduale, das ist, Stufengesang, werden die zwischen Epistel und Evangelium stattfindenden Gesänge bezeichnet. Diesen Namen haben dieselben nämlich von der Weise der alten Kirche, nach welcher der Diakon, ehe er zu dem erhöhten Raum (Ambon) emporstieg, auf welchem er die vorgeschriebenen beiden Lectionen sowohl des Alten als auch des Neuen Testamentes vorzulesen hatte und auf welchem zugleich die Chorsänger standen, an der untersten Stufe einen Gesang anstimmte, den dann der Chor wiederholte. Da nun in der römischen Messe gewisse Gesänge angestimmt werden, wenn nach der Epistel, die auf der linken Seite des Altars vom Subdiakon gelesen wird, der Diakon die Altarstufen besteigt, um auf der rechten Seite das Evangelium zu verlesen, so ist auf diese der altkirchliche Name übertragen worden.

Im engeren Sinne des Wortes versteht man unter Graduale den Gesang einer sich insbesondere auf die Epistel beziehenden Schriftstelle, die entweder ganz oder theilweise von einer Einzelstimme gesungen und dann von einer andern oder dem vollen Chor wiederholt und wohl daher im besonderen Sinne das „Responsorium“ genannt wird. Auf dieses Responsorium folgte ein einfaches oder längeres Halleluja in Nachahmung des alttestamentlichen Hallel (Ps. 113–118, vgl. Matth. 26, 30.) und des apokalyptischen Halleluja
\quellseite{183}
(Offenb. 19, 4.), und sodann ein drei- oder sechzeiliger Hymnus, Sequenz oder Folgegesang genannt, wie z. B. Grates nunc omnes, verdeutscht: Danksagen wir alle Gott etc., das von Weihnachten bis Mariä Reinigung gesungen wurde. Seit Gregor dem Großen sang man auch an diesem Orte Prosen und den Tractus. Sofern nämlich die dem Halleluja folgenden Gesänge nicht eine poetische, sondern eine prosaische Form hatten, nannte man dieselben Prosen. Was nun aber den Tractus betrifft, so hat es damit folgende Bewandtniß. Wenn nämlich der Lector nach Beendigung des Graduale von dem Ambon herabstieg, so zog man in Procession nach dem erhöhten Standort, von dem aus das Evangelium gelesen wurde. Da nun dieser Hinzug unter Psalmengesang mit dem Halleluja der Gemeinde zwischen einzelnen Versen geschah, so nannte man diesen Gesang tractus (Fortzug). Später, wo die Procession wegfiel, beschränkte man ihn auf einen Vers, vor und nach welchem mehrere Hallelujas gesungen wurden, und beließ ihm seinen Namen wohl von den gehaltenen gezogenen Tönen, in denen man denselben sang, und zwar während der Fastenzeit anstatt des Halleluja. Alle diese Gesänge zusammen nannte man auch im weiteren Sinne Graduale.

Zwar ordnet Dr. Luther (Deutsche Messe etc., 1526): „Auf die Epistel singet man ein deutsch Lied: Nun bitten wir den Heiligen Geist etc. oder sonst eins, und das mit dem ganzen Chor.“ Weil aber „Gradual und Alleluja ..., sonderlich die man singet (de tempore) von der Zeit oder an Sonntagen“, auch in „die alte Lauterkeit des Gottesdienstes“ gehören, so nahm man eine Anzahl der Gradualgesänge mit herüber, wie dies schon die herrlichen „Kirchengesänge“ von Spangenberg, 1545, dann die von Wittenberg, 1551, und das Cantional von Lossius, 1553, zeigen.

Das Graduale wurde vom Chor gesungen, wobei einzelne Knabenstimmen, meist das Halleluja anstimmten, die daher die „Halleluja-Knaben“ hießen. Aber die Hauptsache war und blieb beim Graduale doch das deutsche Gemeindelied. Wie sollte es nicht Bedürfniß der ganzen Gemeinde sein, mit ihrem Munde und ihrem eigenen Liede, als ihrem mit dem Evangelio geborenen Charisma, den HEerrn für die Gnadengabe seines Wortes oder auch für die zu feiernde Heilsthatsache zu loben. Das Gemeindelied wird daher an dieser Stelle des Hauptgottesdienstes als der Hauptgesang betrachtet und bezeichnet, wie es denn auch jetzt noch das „Hauptlied“ heißt. Dasselbe wurde laut verschiedener Kirchenordnungen auf mancherlei Weise mit den Chorgesängen des Graduals verbunden. Das bei weitem Gewöhnlichste war, den Gemeindegesang
\quellseite{184}
unmittelbar auf den Chorgesang folgen zu lassen, wie z. B. in der Kirchenordnung von Braunschweig-Lüneburg, 1544, von Wittenberg, 1533 und 1559 etc., mitunter auch wohl nach beiden Kirchenordnungen Chor- und Gemeindegesang Vers um Vers abwechseln zu lassen. In späterer Zeit, zum Theil schon im 17. Jahrhundert, traten an die Stelle des einfachen oder des mit einem Vers verbundenen Halleluja und der herkömmlichen Sequenzen etc. „ein Stück figuraliter gesungen“, wobei „zugleich mit der Orgel darein geschlagen“ wurde. -- Alle älteren und späteren Kirchenordnungen zeigen jedoch, daß, da an dieser Stelle der Chor mit einer besonderen Lobpreisung neben dem Hauptliede der Gemeinde hervortrat, unser heutiger „Chorgesang“ seine Stelle nach der Epistel, also vor dem Hauptliede, haben sollte (wobei jedoch zu wünschen wäre, daß derselbe, wenn das betreffende Chorstück ohne Orgelbegleitung gesungen werden soll, doch durch einige Accorde der Orgel eingeleitet würde, damit nicht jene unschöne und unliebliche Pause entsteht, bei der man den Dirigenten die Stimmen angeben hört, die Gemeinde aber erwartungsvoll dasitzt, ob es „nicht bald angeht“).

Was nun die Weise der Ausführung hinsichtlich des Hauptliedes betrifft, so hatten die Väter freilich auch hier mehr Ausdauer bei der Abhaltung des öffentlichen Gottesdienstes, als ihre heutigen Nachkommen. Während nämlich diese auch im Betreff des Singens immer mehr Kürze auch an Orten begehren, wo die Verhältnisse keine Abkürzung des Gottesdienstes fordern, pflegen die Alten das Hauptlied in der Regel ganz zu singen.

Eine besondere Beachtung erfordert die Wahl des Hauptliedes. In den Kirchenordnungen lautet die gewöhnliche Forderung dahin, daß der Gesang „auf die Zeit (de tempore) oder auf das Evangelium“ sich schicke. Zu dem Ende war ehedem das Hauptlied agendarisch festgestellt. Einer späteren Zeit gehört die Bestimmung an, daß sich das Hauptlied auch „auf die Predigt schicken“ solle. Wir halten dafür, daß bei der Wahl beides zu berücksichtigen ist, vor allem aber die Zeit, und zwar nicht bloß der betreffende Festtag, sondern auch die mit demselben in Verbindung stehenden Sonntage, so daß für die Adventssonntage ein Adventslied, für die Epiphaniensonntage, wenigstens für etliche, ein Epiphaniaslied, und für die Sonntage von Ostern bis Pfingsten in den drei ersten Sonntagen ein Osterlied und in den übrigen ein Pfingstlied das Hauptlied sei, während, wenn keine Communion stattfindet, das Lied nach der Predigt alsdann vor allem den Inhalt der Predigt zu berücksichtigen hätte.

\quellseite{185}
Im Nachstehenden geben wir einige der unmittelbar auf die Epistel folgenden und dem Hauptlied vorhergehenden Halleluja-Weisen.

\section{1. Auf die gewöhnlichen Sonntage}

(mit Ausnahme der Fastenzeit, in welcher das Halleluja ganz wegfällt).

Hommels Lit.

\notensatz{music/185-01}{Halleluja für die gewöhnlichen Sonntage nach Hommels Liturgie, vierstimmig.}

\section{2. Auf Festtage überhaupt}

Mel.: L. Lossius 1553.

a.

\notensatz{music/185-02}{Festtags-Halleluja a nach Lossius 1553, vierstimmig.}

Mel.: Preußische K.-O. 1558.

b.

\notensatz{music/185-03}{Festtags-Halleluja b, Preußische Kirchenordnung 1558, vierstimmig.}

\quellseite{186}
\section{3. Auf Ostern insonderheit}

Herz. Elisabeth v. Braunschwg. K.-O. 1542.

Knaben. Unisono.

\notensatz{music/186-01}{Osterhalleluja der Knaben, unisono.}

Chor.

\notensatz{music/186-02}{Osterhalleluja, vierstimmiger Chor. Unser Osterlamm ist Christus, der am Kreuz für uns geopfert ist. Halleluja!}

\quellseite{187}
\chapter{IX. Evangelium}


Nach dem Vorgang des christlichen Alterthums und im ziemlichen Anschluß an die römische Messe hat ursprünglich die lutherische Kirche die Lesung des Evangeliums vor der Epistel sehr merklich ausgezeichnet. Dies geschieht schon durch die Vortragsweise. Während Luther nämlich für das choraliter Lesen der Epistel einen der ernsteren der acht Kirchentonweisen vorschreibt, will er das Evangelium in einer der freudigeren gelesen wissen und begründet dies in seiner Besprechung mit dem Kapellmeister Walther so: „Christus ist ein freundlicher HEerr und seine Reden sind lieblich, darum wollen wir sextum tonum zum Evangelio nehmen; und weil St. Paulus ein ernster Apostel ist, wollen wir octavum tonum zur Epistel ordnen.“ Bekanntlich hat sich Luther dann in der deutschen Messe für die fünfte Tonweise (quintum tonum) entschieden und nach dieser auch die Einsetzungsworte componirt.

Aber es ist diese Lesung auch noch durch anderweitigen liturgischen Schmuck ausgezeichnet. In seiner Formula missae wollte Luther noch Lichter und Weihrauch beibehalten sehen. Kam auch der Weihrauch bald außer Gebrauch, so doch nicht der brennenden Lichter während des Lesens des Evangeliums, welche man ohne Zweifel mit der alten Kirche als das Symbol des Lichtes der von dem Evangelio ausgehenden Gnade betrachtet. Wie in der römischen Messe geschah mancherorts die Lesung auf der rechten Seite des Altars, während die der Epistel auf der linken Seite stattfand, auch hörte die Gemeinde dieselbe stehend an, während sie bei der Epistel sitzen blieb. Voraus ging eine abermalige Salutation: „Der HEerr sei mit euch“ etc. und dazu in einer noch schwungvolleren Weise als die der Salutation vor der Collecte. Nach der Ankündigung: „So schreibt der heilige Evangelist St. N. N. im N. N. Kapitel“ oder: „Höret zu dem heiligen Evangelium“, respondirte erst der Chor: „Ehre sei dir, HEerre!“, worauf dann die Lesung selbst quinto tono erfolgte. Im Laufe des 17. Jahrhunderts begann man auch, nach dem Evangelio und vor dem Credo eine Motette („ein Figuralstück, wo man es haben kann“, wie es in der Schwarzb. Agende 1675 heißt) durch den Chor singen zu lassen, was Schöberlein als Preis der Heilsthatsache auffaßt, die in der Regel den Inhalt dieser Schriftlectionen aus den vier Evangelien bildet.\footnote{Auch die Liturgie jener nestorianischen Gemeinden, welche standhaft die Unterwerfung unter den römischen Antichrist zurückgewiesen haben, und bei denen, beiläufig bemerkt, heute noch die Messe nicht als unblutiges Sühnopfer, sondern nur als Communion zur Besiegelung der Sündenvergebung unter beiderlei Gestalt gehalten werden soll, zeigt in ihrer dreimaligen Lesung folgende Steigerung: Alttestamentliche Lection, Epistel und Evangelium, die erste vom Vorleser an der Südseite des Altarzugangs, die Epistel vom Subdiakon an der Nordseite desselben, das Evangelium aber vom Priester oder dem ordinirten Diakon vor der Altarschranke, wobei der Lesende immer das Angesicht gegen die Gemeinde kehrt. -- In der orientalischen Kirche ist überhaupt beim Lesen oder Singen des Evangeliums das Ceremoniell noch reichhaltiger und feierlicher, als in der römischen Kirche.}

\quellseite{188}
Weshalb nun diese liturgische Auszeichnung? Man hat in dem Stufengang von der Epistel zum Evangelium eine Wiederholung der Weise der jüdischen Kirche erblicken wollen, nach welcher erst ein Abschnitt aus dem Gesetze Mosis und dann einer aus der Weissagung der Propheten, diesem Evangelium des Alten Bundes, vernommen wurde. Da aber der Inhalt der Episteln doch nicht immer in Ermahnung zu einem heiligen Wandel besteht, sondern oft auch theilweise, manchmal auch ganz, Lehre, Heilslehre, ist, so möchten wir der Stufenfolge von der Epistel zum Evangelium im Hauptgottesdienst doch eine etwas andere Deutung geben und in dieser den Grund der liturgischen Auszeichnung suchen. In einem Gutachten über die altkirchlichen Perikopen bemerkt Rudelbach: „So wie das Alte Testament mit Geschichte anfängt, also auch das Neue; so wie das Geschichtliche dort die breite und feste Grundlage bildet, also auch hier. So wie im Alten Testament das Didaktische überall auf die historische Grundlage zurückweist und nimmer ohne dieselbe völlig verstanden werden mag, also auch im Neuen. So wie im Alten Testament das Prophetische das Geschichtliche und Didaktische beides in eine höhere Ordnung aufnimmt und den lebendigen Schlußstein der Offenbarung bildet, also auch im Neuen. Mit einem Worte: Alles in der heiligen Schrift zeigt uns das Werk des göttlichen Geistes an, in welchem That und Wort, Lehre und Geschichte aufs innigste verbunden und gleichsam zusammengewachsen sind.“ (Sacrament-Worte etc. S. 130.) Den Hauptinhalt von Epistel und Evangelium zusammenhaltend, möchten wir daher uns die Aufeinanderfolge von Epistel und Evangelium als aufsteigend\footnote{In der alten Kirche geschah die Lesung des Evangeliums sogar auch von einer höheren Stufe des Altarchors aus als der des Epistelpultes.} denken, und zwar in folgender Weise. Die Epistel ist entweder apostolische Verkündigung, die bald als Lehre „den vollen Reichthum Christi entfaltet“, bald als Vermahnung es mit den Früchten des Glaubens zu thun hat, oder sie ist mitunter auch prophetisches Wort, Heilsweissagung (Weihnachten Jes. 9, Epiphanias, den Marienfesten und den Johannistag); im
\quellseite{189}
Evangelium aber schauen wir das letztere in seiner Erfüllung, erblicken das fleischgewordene Wort in seiner Person, in seiner That zu unserm Heil von seiner Empfängniß bis zu seiner Himmelfahrt, hören aus seinem Munde ipsissima verba und sehen seine Zeichen und Wunder. Man könnte hierbei auch an das „Erbauet auf den Grund der Apostel und Propheten (Epistel), da JEsus Christus der Eckstein ist“ (Evangelium), denken. Daß auch Luthern ein solches Aufsteigen von der Epistel zum Evangelium vorschwebte, beweist das zu Anfang gebrachte Citat, nach welchem er für das Evangelium als Christi liebliche Stimme sextum tonum, für die Epistel aber als des ernsten Apostels Stimme octavum tonum nehmen will.

Die liturgische Auszeichnung der Lesung des Evangeliums ist im Allgemeinen nach und nach außer Gebrauch gekommen. Nur in etlichen uns bekannten Gemeinden dahier brennen die Altarlichter beim Lesen des Evangeliums und in etlichen andern wird, wenn auch kein Singen des Evangeliums stattfindet, doch das Lesen desselben nach altem Brauch mit der Salutation und der Lobpreisung eingeleitet. Obwohl ja an eine Wiederaufnahme des Singens der Perikopen nicht mehr zu denken ist, so könnte doch noch immer nebst dem Anbrennen der Lichter nach der erwähnten Weise die Lesung des Evangeliums ausgezeichnet werden. Diese Weise wäre dann folgende:

\begin{verse}
Pastor singt (oder spricht): Der HEerr sei mit euch!\\
Gemeinde singt: Und mit deinem Geiste.\\
Pastor spricht: Dies (oder das) heilige Evangelium am (ersten Sonntag des Advents etc.) schreibt (St. Matthäus etc.) im . . . Kapitel . . . Vers . . .\\
Gemeinde erhebt sich und singt: Ehre sei dir, HEerre.\\
Pastor verliest das Evangelium.
\end{verse}

Unmittelbar darauf intonirt der Pastor das Credo. Soll anstatt nach der Epistel nach dem Evangelium ein Chorgesang eintreten, so kann derselbe entweder auf den Glauben und vor der Predigt oder auch nach derselben folgen.

Wie auch aus dem X. Band der St. Louiser Ausgabe der Werke Luthers zu ersehen ist, hat Luther in seiner deutschen Messe nicht nur für das Evangelium Singregeln aufgestellt, sondern sogar zweierlei Melodien gegeben als „Exercitatio oder Uebung der Melodien, auf daß man sich wohl lerne schicken in Melodien und wohl gewöhne der Colen, Commaten u. dgl. Pausen“. Aus jenen Regeln ist auch zu ersehen, daß in damaliger Zeit nicht bloß von Einem, sondern auch von drei Priestern gesungen wurde, von denen der eine die erzähl-enden
\quellseite{190}
Worte des Evangelisten vortrug (vox evangelista), der andere die Worte anderer Personen (außer Christus, vox personarum) und der dritte die Worte Christi (vox Christi). Mit Weglassung dieser Regeln geben wir beispielsweise nur die eine der Melodien.

\subsection{Einleitung}

Mel.: Bopelius 1682.

\notensatz{music/190-01}{Einleitung zum Evangelium, erste Salutation des Pastors: Der HErr sei mit euch! G–G–Fis–G–E–G, mit Melisma auf euch.}

\notensatz{music/190-04}{Einleitung zum Evangelium, Antwort nach Bopelius 1682. Und mit deinem Geiste.}

\subsubsection{Oder:}

Hoya K.-O. 1581.

\notensatz{music/190-02}{Alternative Salutation des Pastors: Der HErr sei mit euch! G–A–H–G–A–H–A–A, abschließend als Brevis; Melismen auf Der, sei und mit.}

\notensatz{music/190-05}{Einleitung zum Evangelium, Antwort nach der Hoyaer Kirchenordnung 1581. Und mit deinem Geiste.}

Mel.: Bopelius 1682.

\notensatz{music/190-03}{Ankündigung des Evangeliums, Bopelius 1682, Pastor: Dies heilige Evangelium beschreibt der Evangelist St. Matthäus. Zwanzig ganze Noten, überwiegend G, mit A auf ge in Evangelium und St., dann H–A–G auf Matthäus.}

\notensatz{music/190-06}{Vierstimmige Antwort. Ehre sei dir, o HErre. Alternative Textfassung: HErr.}

Oder:

\notensatz{music/190-07}{Zweite vierstimmige Antwort. Ehre sei dir, HErre. Alternative Textfassung: HErr.}

\quellseite{191}
\emph{Braunsch.-Lüneburg. 1542.}\footnote{Eine andere Weise für das Evangelium siehe Luthers Werke, St. Louiser Ausgabe, Bd. X, Col. 243 ff. und Col. 253 ff.}

\notensatz{music/191-01}{Evangelienton des Pastors, Braunschweig-Lüneburg 1542: Matthäus 21, von „Die nachfolgenden Worte des heiligen Evangeliums“ bis „durch den Propheten, der“.}

\quellseite{192}
\notensatz{music/192-01}{Fortsetzung des Evangelientons, Matthäus 21: von „da spricht: Saget der Tochter Zion“ bis „Hosianna in der Höhe!“; Schlussmelisma auf „an“ in Hosianna.}

Hierauf folgt die Ueberleitung zur Intonation des Credo.

\quellseite{193}
\chapter{X. Credo}


In Folge des arianischen Streits entstand nach dem apostolischen Symbolum das nicänische Symbolum, das zweite ökumenische Bekenntniß. Im 5. Jahrhundert fand es bereits in der orientalischen Kirche seine Verwendung im Hauptgottesdienst, erst aber zu Ende des 8. Jahrhunderts folgte die occidentalische Kirche nach, und zwar in derselben zuerst die gallicanische Kirche, während erst auf Andringen Kaiser Heinrichs Papst Benedict VIII. dasselbe der Messe einverleibte.

In der Mitte des Altars stehend spricht entweder der Priester das ganze Nicänum nach der Lection des Evangeliums oder, wenn die Messe eine feierliche ist, intonirt er: „Credo in unum Deum“, worauf der Chor vom „Patrem omnipotentem“ an das Nicänum hinaussingt (daher auch der Name „Patrem“). Wie auch Luther in einer Predigt darauf hinweist, war es (und ist noch) Sitte, daß bei den Worten: „Et incarnatus est de Spiritu Sancto, ex Maria virgine: et homo factus est“ Priester und Volk zum Zeichen der Anbetung des kündlich großen Geheimnisses der Menschwerdung des Sohnes die Kniee beugen. Diesem Brauch folgte man auch zum Theil in der lutherischen Kirche darin, daß zur Reformationszeit schon in den österreichischen Landen die Sitte aufkam und sich dann auch nachher theilweise nach Sachsen verbreitete, daß bei diesen Worten, wie auch bei den Worten des Luther’schen Patrem: „Von Maria der Jungfrauen ist ein wahrer Mensch geboren durch den Heiligen Geist im Glauben für uns, die wir waren verloren“ die ganze Gemeinde sich erhob.

In seiner „Deutschen Messe“ 1526 ordnet Luther: „Nach dem Evangelio singt die ganze Kirche den Glauben zu deutsch: Wir glauben all an einen Gott. Darnach gehet die Predigt vom Evangelio des Sonntags oder Festes.“ Hiernach schließt sich Luther der herkömmlichen Weise darin an, daß er das Credo auf die Lection des Evangeliums folgen läßt; er macht aber darin eine Aenderung, daß er die auf das Credo folgende Predigt zum organischen Glied der Messe macht, während in der römischen Messe die Predigt auch hier wegfallen kann und auch oft wegfällt, so daß gleich auf das Credo die Präfation folgt, der Priester schon während des Credo die Secreta spricht, und somit die Predigt mehr wie ein Einschiebsel erscheint. Ursprünglich gehört auch das Credo in die missa fidelium, und steht so an der Spitze des Communionactes.

\quellseite{194}
Da nur Gläubige sich zum Tisch des HEerrn nahen sollen, so soll sich auch jeder durch das Credo als einen Gläubigen kund geben. Das meint die ursprüngliche Stellung des Credo. In seiner „Weise, christlich Meß zu halten“ sagt daher Luther: „Wir halten von der deutschen Predigt, daß es nicht daran gelegen ist, sie werde nach dem Patrem oder vor dem Introit gehalten, wiewohl Ursache könnte angezeigt werden, warum sie besser vor der Messe geschehe, weil das Evangelium eine rufende Stimme ist in der Wüste, welche der Ungläubigen Herzen zum Glauben laden soll, die Messe aber soll sein der Brauch oder Nutz des Evangelii und Austheilung des Sacraments, so allein den Gläubigen gebührt, und sollte billig an einem andern Ort, außerhalb der Ungläubigen, gehalten werden. Doch weil wir frei sind, soll uns diese Weise nicht binden noch fangen.“ Da nun Luther mit Recht in christlicher Freiheit von dieser ursprünglichen Weise abwich und die Predigt nicht vor dem Introitus, sondern nach dem Credo stellte, so fragt sich’s, welche Bedeutung in dieser Stellung das letztere habe? Es dürfte dieselbe wohl diese sein: Dem Evangelio folgend ist es als Bekenntniß des Mundes dessen aufzufassen, was die Kirche nach dem Evangelio und durch dasselbe von Gottes Wesen und Willen im Herzen glaubt; der Predigt aber vorausgehend steht es zugleich als regula fidei da, nach welcher der Prediger nun das gelesene Wort auslegen soll, gemäß der Weisung des Apostels: „Hat jemand Weissagung, so sei sie dem Glauben ähnlich“ (Röm. 12, 7.).

Zwar sang man in der lutherischen Kirche Anfangs nicht allein das Luther’sche Patrem, sondern auch das Nicänum, und zwar an manchen Orten lateinisch, an manchen Orten auch deutsch, doch immer entweder so, daß Luthers Lied von der Gemeinde gesungen folgte oder daß man den Gebrauch beider im Wechsel empfahl, „auf daß Beides geübt werde“ (Veit Dietrich, Agendbüchlein, 1548). Häufig aber beschränkte man sich auf nur Ein Credo, und zwar am liebsten auf das Luther’sche, vom Chor und der Gemeinde zugleich gesungen, bis in der Folgezeit man überall in der Regel sich desselben bediente. Um 1680 dichtete Tobias Clausnitzer den sogenannten „kleinen“ Glauben (S. St. Louiser Gesangb. Nr. 184), der dann neben dem „großen“ Glauben in Gottesdiensten gebraucht wurde und dermalen gebraucht wird, in welchen ein längeres Hauptlied gesungen wurde oder die Zahl der Communicanten groß war oder die Umstände sonst einmal eine Abkürzung erheischten. HEerr im Hause blieb jedoch der „große“ Glaube, das in deutsche Reime gebrachte Symbolum Nicänum.

\quellseite{195}
Und er blieb es zugleich mit seiner grandiosen Melodie, mit der er zuerst in Walthers Chorgesangbüchlein vom Jahre 1524 erschien und dann auch im Straßburger „Kirchenamt 1525“. Da erfand im Jahre 1790 Joh. Christ. Kittel, Bachs letzter Schüler, seine C-dur-Melodie: c c h a g a d g e (S. auch Nr. 297 des St. Louiser Choralb. 1886.), und so geschah es, daß dies Kind einer Zeit, die am wenigsten mehr weder ein Kirchenlied, noch eine Kirchenmelodie hervorzubringen vermochte, die Weise Luthers,\footnote{Joh. Walther, der Kapellmeister, schreibt diese Melodie, wie die zum deutschen Sanctus („Jesaja dem Propheten“ etc.) Luthern zu, desgleichen auch v. Winterfeld. Weil aber eine gereimte Verdeutschung des Nicänischen Symbolums vor Luther schon vorhanden war und auf der Breslauer Bibliothek handschriftlich eine Melodie von Nik. v. Kosel 1417 zu derselben sich befinden soll, so dürfte wohl auch die Melodie, gleichwie der Text, eine und zwar gleichfalls meisterliche Bearbeitung Luthers sein.} diese der großartigsten Weisen eine im lutherischen Melodienschatz, bald fast überall verdrängt hatte. Im Hinblick auf diese und manch andere der verdrängten reformatorischen Weisen stößt Palmer in seinem Wort zum „neuen Choralbuch“ (Süddeutscher Schulbote 1845, Nr. 5) den Seufzer aus: „Wollte Gott, wir könnten solche grandiose Choräle mit unsern Gemeinden singen! Vielleicht wenn es wieder buchstäblicher wahr wird, daß wir alle an Einen Gott glauben, lernen wir das auch wieder besser singen.“ In unsern heutigen Kirchen und Schulen sind jene grandiosen Choräle der Reformationszeit, Gott Lob! wieder heimisch geworden. Man singt daher sonn- und festtäglich allüberall Luthers Credo auch in Luthers Weise, und zwar da und dort theils von Anfang an, theils sie mit der zuerst gebrauchten Melodie Kittels vertauschend. Leider hat man aber schon da und dort einen umgekehrten Tausch wieder eingegangen — nicht weil das Erlernen der Lutherweise zu schwer wurde, sondern weil die mehr arienmäßige Melodie Kittels dem verdorbenen Geschmack mehr zusagt, als der von aller Sentimentalität freie und erhabene Kirchenton. Möchte dieser unvortheilhafte Tausch doch nicht weiter einreißen, sondern die Luther’sche Weise das Regiment im Hauptgottesdienst allüberall führen!

\section[In Betreff der Ausführung]{In Betreff der Ausführung nun noch einige Erinnerungen}

1. Wie in der römischen Messe, so wurde auch im lutherischen Hauptgottesdienste ehedem das Credo „gegen den Altar“ von dem in der Mitte desselben stehenden Liturgen lateinisch oder deutsch intonirt. Nicht nur entspricht solches Intoniren des Credo dem des Gloria: „Ehre sei Gott in der Höhe!“, sondern es kommt auch damit das unpassend lange Präludiren
\quellseite{196}
zwischen der Evangelium-Section und dem Credo der Gemeinde von selbst in Wegfall, indem entweder auf das Evangelium das Responsorium mit der Credo-Intonation oder letztere allein folgt. Selbst wo dieselbe nicht üblich ist, sollten nur ein paar passende Accorde auf der Orgel dieselbe einleiten. Ein längeres Orgelpräludium zwischen dem Lesen des Evangeliums und dem: „Wir glauben all“ etc. ist durchaus unstatthaft.

In der Hoffnung, daß die Credo-Intonation, die in etlichen unserer Gemeinden längst üblich ist, in weitere Aufnahme komme, sind hier etliche der Intonationen mitgetheilt. Die sub A zu dem Credo Luthers findet sich durchgängig mit der zum Gloria in den alten Agenden. Die eine sub B ist dem Nürnberger „Agendbüchlein“ entnommen, wo sie zum Luther’schen Credo sich findet, während sie hier für den „kleinen“ Glauben transponirt worden ist; die andere wird in ein paar unserer hiesigen Gemeinden gebraucht. Die sub C ist die festive Weise zu einem deutschen, sehr schönen Nicänum in D-dur bei L. Lossius, Psalmodia sacra 1553 und 1579, in der Pomm. K.-O. 1563, bei Keuchenthal, Wittenberg 1573, Ludeccus Missale 1589, im Nürnb. Agendbüchl. 1639 und aus diesen bei Hommel.

2. Da das Bekenntniß des Glaubens überhaupt frisch, tapfer und deutlich geschehen soll, so darf insbesondere die mächtige Luthermelodie durchaus nicht schleppend gesungen und müssen dabei auch alle Worte recht deutlich ausgesprochen werden, wobei der Organist viel des Seinigen thun kann, wenn er beim Spiel gerade dieser Melodie mehr ein staccato einhält. Hierbei die Bemerkung, daß in manchen älteren Werken die auf die Anfangsnote folgenden vier Schleifnoten dem „Wir“ gehören, und daß sich daran die Sage knüpft, Luther habe damit andeuten wollen, daß die nach den vier Himmelsgegenden zerstreuten Völker in Einigkeit des Glaubens versammelt seien. Der Anfang des Luther’schen Credo lautete hiernach so:

\notensatz{music/196-01}{Anfang des Luther’schen Credo: Wir glau:ben. Wir auf D und den vier folgenden Schleifnoten A–G–A–E; glau auf F, ben auf E–G. Die D- und F-Noten sind Halbe, die übrigen Viertel.}

3. Bei den Worten: „Die ganz Christenheit auf Erden“ pflegt man die Kanzel zu betreten.\footnote{Ein älterer schlesischer Pastor, den sein treuherziger Kirchendiener mit den Worten: „Herr Pastor, nun kommt die ganze Christenheit“ an das Besteigen der Kanzel zu erinnern pflegte, bemerkt: „Dies Wort bringt mir dann allemal wundersam ins Herz; erinnert mich’s doch an mein mächtig Amt in der ganzen großen Christenheit. Mit meiner Gemeinde kommt die ganze Christenheit, auch die nun schon triumphirende am Stuhl des Lammes, mir entgegen, mir armen Menschenkinde mit blödem Herzen und matten Lippen. Und dann seufze ich wohl aus Herzensgrund: „Ach, daß sie nicht vergebens käme, die ganze Christenheit! HEerr, hilf, HEerr, sei mit mir, daß sie Dich bei mir durch mich finde!“}

\quellseite{197}
4. Da man zur Abwechselung, besonders aber an Festen, nicht nur das Nicänum, sondern mitunter auch an dessen Statt das Apostolicum zu singen pflegte, so ist hier die schöne, schon im Straßburger Kirchenamt 1525 vorkommende, dann auch in einige Kirchenordnungen des 16. Jahrhunderts und in neuerer Zeit auch von Hommel und Layriz aufgenommene Melodie des letzteren mitgetheilt. Doch kann auch der Pastor den Anfang der Melodie unter Orgelbegleitung zur Intonation machen, indem er singt: „Ich glaub an Gott Vater den Allmächtigen“, und die Gemeinde dann fortfährt: „Schöpfer Himmels und der Erden.“ Wie schnell bei der Bekanntschaft des Textes eine Gemeinde dem Chor diese Melodie nachsingen lernt, hat der Verfasser in seiner früheren Gemeinde erfahren.

5. Da die Gemeinde bei der Lesung des Evangeliums steht, so ist es passend, noch während der Intonation des Credo und der ersten Zeile desselben stehen zu bleiben.

\section{1. Credo-Intonationen}

\subsubsection{A}

\notensatz{music/197-01}{Credo-Intonation A, Orgel und Liturg. Ich glaub(e) an einen Gott.}

\notensatz{music/197-02}{Credo A, vierstimmiger Einsatz der Gemeinde. Wir glauben etc.}

\subsubsection{B}

a.

\notensatz{music/197-03}{Credo-Intonation B a, Orgel und Liturg. Ich glaub(e) an einen Gott.}

b.

\notensatz{music/197-04}{Credo-Intonation B b, Orgel und Liturg. Ich glaube}

\quellseite{198}
\notensatz{music/198-01}{Credo-Intonation B b, Fortsetzung. An einen Gott.}

\notensatz{music/198-02}{Credo B b, vierstimmiger Einsatz der Gemeinde. Wir glauben all an einen Gott etc.}

\subsubsection{C}

\notensatz{music/198-03}{Credo-Intonation C, Orgel und Liturg. Ich glaube an Einen Gott.}

\subsubsection[Apostolisches Glaubensbekenntniß]{Gemeinde. (Apostolisches Glaubensbekenntniß.)}

\notensatz{music/198-04}{Credo C, vierstimmiger Einsatz der Gemeinde. Ich glaub an Gott Vater, den Allmächtigen.}

\section{2. a. Der apostolische Glaube}

Mel.: Braunschweig-Lüneb. K.-O. 1542. Harm. v. F. Jubitz.

\notensatz{music/198-05}{Apostolischer Glaube: Ich glaub an Gott Vater, den Allmächtigen, Schöpfer.}

\quellseite{199}
\notensatz{music/199-01}{Apostolischer Glaube: Himmels und der Erden. Und an Jesum Christum, seinen.}

\notensatz{music/199-02}{Apostolischer Glaube: einigen Sohn, unsern HErren, der empfangen ist vom.}

\notensatz{music/199-03}{Apostolischer Glaube: heiligen Geist, geboren aus Maria der Jungfrau-.}

\notensatz{music/199-04}{Apostolischer Glaube: -en, gelitten unter Pontio Pilato.}

\quellseite{200}
\notensatz{music/200-01}{Apostolischer Glaube: gekreuziget, gestorben und begraben, niederge-.}

\notensatz{music/200-02}{Apostolischer Glaube: -fahren zu der Höllen, am dritten Tage aufer-.}

\notensatz{music/200-03}{Apostolischer Glaube: -standen von den Todten, aufgefahren gen.}

\notensatz{music/200-04}{Apostolischer Glaube: Himmel, sitzend zu der Rechten Gottes.}

\quellseite{201}
\notensatz{music/201-01}{Apostolischer Glaube: des allmächtigen Vaters, von dannen er kommen.}

\notensatz{music/201-02}{Apostolischer Glaube: wird, zu richten die Lebendigen und die Todten.}

\notensatz{music/201-03}{Apostolischer Glaube: Ich glaube an den Heiligen Geist, Eine heilige christliche.}

\notensatz{music/201-04}{Apostolischer Glaube: Kirch, die Gemeine der Heiligen, Vergebung der Sünden, Auferste-.}

\quellseite{202}
\notensatz{music/202-01}{Apostolischer Glaube: -hung des Fleisches und ein ewiges Leben. Amen.}

\section{b. Das Nicänische Bekenntniß}

Mel.: L. Lossius 1553. Pomm. K.-O. 1563. Wittenb. Kirchengef. 1573. Ludecus Missale 1589 etc. Harm.: Dr. Herzog.

\subsubsection{Pastor}

\notensatz{music/202-02}{Nicänisches Bekenntniß, Intonation des Pastors: Ich glaube an Einen Gott.}

\subsubsection{Gemeinde}

\notensatz{music/202-03}{Nicänisches Bekenntniß: Allmächtigen Vater, Schöpfer Himmels und der Erden.}

\notensatz{music/202-04}{Nicänisches Bekenntniß: aller sichtbaren und unsichtbaren Ding. Ich glaube an den.}

\quellseite{203}
\notensatz{music/203-01}{Nicänisches Bekenntniß: einigen HErrn Jesum Christum, eingebornen Sohn Gottes.}

\notensatz{music/203-02}{Nicänisches Bekenntniß: aus dem Vater vor aller Zeit geboren, Gott von.}

\notensatz{music/203-03}{Nicänisches Bekenntniß: Gotte, Licht vom Lichte, wahrer Gott vom wahren Gott, ge-.}

\notensatz{music/203-04}{Nicänisches Bekenntniß: -boren, nicht geschaffen, Eines Wesens mit dem Vater, durch welchen.}

\quellseite{204}
\notensatz{music/204-01}{Nicänisches Bekenntniß: alle Ding gemacht sind, der um uns Menschen und um unsrer}

\notensatz{music/204-02}{Nicänisches Bekenntniß: Seligkeit willen vom Himmel herabgestiegen, und}

\notensatz{music/204-03}{Nicänisches Bekenntniß: Fleisch geworden ist vom Heiligen Geiste aus Maria der}

\notensatz{music/204-04}{Nicänisches Bekenntniß: Jungfrau, und ist Mensch geworden, auch gekreuzigt für}

\quellseite{205}
\notensatz{music/205-01}{Nicänisches Bekenntniß: uns unter Pontio Pilato, gestorben und begraben,}

\notensatz{music/205-02}{Nicänisches Bekenntniß: auferstanden am dritten Tage nach der Schrift, aufge-}

\notensatz{music/205-03}{Nicänisches Bekenntniß: -fahren gen Himmel, sitzet zur Rechten des Vaters,}

\notensatz{music/205-04}{Nicänisches Bekenntniß: und wird wiederkommen mit Herrlichkeit, zu richten die Leben-}

\quellseite{206}
\notensatz{music/206-01}{Nicänisches Bekenntniß: -digen und die Todten, so daß seines Reichs kein Ende sein wird.}

\notensatz{music/206-02}{Nicänisches Bekenntniß: Ich glaube auch an den Heiligen Geist, der da ist HErr und}

\notensatz{music/206-03}{Nicänisches Bekenntniß: macht lebendig, der von dem Vater und dem Sohne ausgeht,}

\notensatz{music/206-04}{Nicänisches Bekenntniß: der mit dem Vater und dem Sohne zugleich angebetet und geehret}

\quellseite{207}
\notensatz{music/207-01}{Nicänisches Bekenntniß: wird, der durch die Propheten geredet hat. Ich glaube Eine hei-}

\notensatz{music/207-02}{Nicänisches Bekenntniß: -lige, allgemeine und apostolische Kirche.}

\notensatz{music/207-03}{Nicänisches Bekenntniß: Ich bekenne Eine Taufe zur Vergebung der Sünden,}

\notensatz{music/207-04}{Nicänisches Bekenntniß: und warte auf die Auferstehung der Todten, und das}

\quellseite{208}
\notensatz{music/208-01}{Nicänisches Bekenntniß: Leben der zukünftigen Welt. Amen.}

\chapter[XI. Die Predigt und die Acte]{XI. Die Predigt und die mit ihr zusammenhängenden und angeschlossenen Acte}


\section{A}

\section{1. Die Predigt}

Unter den drei großen Mißbräuchen des papistischen Gottesdienstes überhaupt bezeichnet Luther in seiner Schrift: „Ordnung des Gottesdienstes in der Gemeinde 1523“ als ersten und ärgsten Mißbrauch den, „daß man Gottes Wort geschwiegen hat und alleine gelesen und gesungen in den Kirchen“. Er will daher, „daß die christliche Gemeine nimmer soll zusammenkommen, es werde denn daselbst Gottes Wort gepredigt und gebetet, es sei auch aufs Kürzeste“. Ja, es „zeigt die ganze Schrift an und Christus auch selbst sagt, Luc. 10: Eins ist vonnöthen, daß Maria zu Christi Füßen sitze und höre sein Wort täglich, das ist das beste Theil, das zu erwählen ist und nimmer weggenommen wird. Es ist ein ewig Wort, das Ander muß Alles vergehen, wie viel es auch der Martha zu schaffen gibt. Dazu helf uns Gott. Amen“. (Luther, Erl. Ausg. 22, 154. 156. St. L. Ausg. X, 220. 225.)

Und so ordnete er denn die Predigt als einen Haupttheil der gereinigten Messe, und obwohl er erst in seiner „Weise, christlich Meß zu halten“ geneigter war, um ihres missionirenden Charakters willen die Predigt vor die ganze Messe zu stellen, so reihte er sie doch später in seiner deutschen Messe 1526 da
\quellseite{209}
ein, wo sie jetzt noch allüberall im lutherischen Hauptgottesdienst steht. An dieser Stelle aber wird sie zum Ruhepunkt im Gang der Liturgie, da Maria zu des HEerrn Füßen sitzt, und zu einem Höhepunkt, von dem aus die feiernde Gemeinde zum andern Höhepunkt wandelt, zum Sacrament als dem Siegel des Worts.

Der Text für die Predigt des Hauptgottesdienstes war das Evangelium des Tages oder Festes, während in einem der dem Hauptgottesdienste vorhergehenden oder nachfolgenden Nebengottesdienste, der Sonntagsmette oder Sonntagvesper, über die Epistel gepredigt wurde. Erst später ließ man einen Wechsel zwischen Evangelium und Epistel für die Predigt des Hauptgottesdienstes zu. Obgleich schon am Altar gelesen, wurde das Evangelium doch nochmals auf der Kanzel verlesen, um so mehr, da es am Altar choraliter gelesen war.

Verschieden nun ist die Einleitung zur Predigt. Nach Luthers Vorgang in der Hauspostille bediente man sich eines stehenden Eingangs zur Textverlesung und zur Predigt. Sie war Vermahnung zum Gebet um fruchtbare Handlung des Wortes, welches Gebet in einem Vaterunser bestand, das an manchen Orten der Prediger laut, an manchen aber mit ihm jeder für sich im Stillen betete, eine Anordnung, die wohl im Hinblick auf Offenb. 8, 1. Hab. 2, 20. Jes. 14, 7. gemacht wurde. Später fügte man zwischen dem lauten oder stillen Vaterunser und der Verlesung des Textes oder schon vor dem Vaterunser ein kurzes Lied de tempore ein, während andere gleichzeitige Kirchenordnungen für den Eingang mit Vaterunser ein stehendes Eingangsgebet vorschreiben, wie aus den hier nachstehenden Formularen zu ersehen. Doch ordnen andere Kirchenordnungen der nachreformatorischen Zeit auch schon einen freien Eingang zur Predigt an, wiewohl derselbe nicht Einleitung zum Thema der Predigt, sondern immer noch Vermahnung zum Gebete ist. So heißt es z. B. in der Braunsch.-Lüneburg. K.-O. Herzog Augusts 1657: „Für der Predigt soll der Prediger den gewöhnlichen Wunsch sprechen: Gnade, Friede und Barmherzigkeit von Gott, dem himmlischen Vater, seinem eingebornen Sohn Jesu Christo, sammt dem Heiligen Geist sei und bleibe bei uns von nun an bis in Ewigkeit. Amen. Darauf das Volk zum christlichen Gebet vermahnen, daß Gott ihnen seinen Heiligen Geist geben wolle, sein heiliges Wort nützlich zu handeln und in das Herz zu fassen, daß sie ihr sündlich Leben bessern, Trost und Stärkung ihres Glaubens erlangen und hinführo christlich leben und endlich selig sterben mögen. Alsbald darauf bete er mit der Gemeinde in der
\quellseite{210}
Stille, und zwar die ganze Gemeinde mit ihm knieend,\footnote{Von einem Niederknien des Predigers bei Betreten der Kanzel, wie es jetzt hie und da üblich ist, findet sich in den älteren Kirchenordnungen unseres Wissens nichts.} das Vaterunser. Vorhero aber singe man unterweilen den Gesang: Nun bitten wir den Heiligen Geist etc. Ingleichen um Weihnachten: Ein Kindelein so löbelich etc. Um die Ostern: Also heilig ist der Tag, oder: Christ ist erstanden etc. Auf Pfingsten: Komm, Heiliger Geist, HEerre Gott etc. Nach geendigtem Gebet soll der Text des Sonn- oder Festtagsevangelii nach der Version des H. Lutherus sel. abgelesen und darauf die Predigt angefangen werden.“ Die hier genannten Lieder sind als Kanzellieder auch in der Schwarzburgischen Agenda 1675 angegeben.

Betreffs der Dauer der Predigt im Hauptgottesdienst bestimmt die Braunschw.-Lüneb. K.-O. 1657, „daß alle Prediger in Städten und auf dem Lande der langen Predigten sich abthun und enthalten und über drei Viertel- oder aufs höchste über eine geringe Stunde nicht predigen sollen“. Eine dreiviertelstündige Predigt gehörte also damals nicht zu den langen Predigten.

\section[2. Die angeschlossenen Acte]{2. Die mit der Predigt zusammenhängenden und angeschlossenen Acte}

\section[a. Das Confiteor]{a. Das Confiteor oder die allgemeine Beichte}

Im engen Anschluß an die römische Messe (s. die Tabelle Theil I.) stellen eine Anzahl alter Kirchenordnungen und Agenden das Confiteor oder die allgemeine Beichte des Hauptgottesdienstes demselben voran. Unter diesen finden wir in Löhes „Sammlung liturgischer Formulare 1839“ verzeichnet: Andreas Döbers ev. Meß im neuen Spital zu Nürnberg 1525; Straßburg. deutsches Kirchenamt 1525; Brandenb. (Ansbach) und Nürnberg 1533; Wittenberg 1559. 1563; Mecklenburg 1552. 1650; Frankfurt a. M. 1565; Lüneburg 1564. 1569; Lippe 1571; Oldenburg 1573; Liegnitz 1594 etc.

Die Form und Weise des Confiteor an dieser Stelle ist folgende: An den Stufen des Altars stehend spricht der Liturg: „Im Namen des Vaters und des Sohnes und des Heiligen Geistes. Amen. Meine Allerliebsten in Gott, eröffnet eure Herzen, lasset uns Gott unsere Sünden bekennen und um Vergebung im Namen unseres HEerrn Jesu Christi bitten. Sprecht mir nach mit herzlicher Begierde zu Gott, im Glauben an den HEerrn Jesum Christum, durch den Heiligen Geist.“ Darauf kniet der Pastor mit der Gemeinde nieder, spricht mit derselben antiphonatim: „Unsere Hülfe stehet etc.
\quellseite{211}
Ich sprach: „Ich will dem HEerrn meine Uebertretung“ etc. und fährt dann allein fort: „Ich armer sündiger Mensch bekenne Gott, dem Allmächtigen, meinem Schöpfer und Erlöser, daß ich nicht allein gesündigt habe mit Gedanken, Worten und Werken, sondern auch in Sünden empfangen und geboren bin, also daß alle meine Natur und Wesen vor seiner Gerechtigkeit sträflich und verdammlich ist. Darum fliehe ich zu seiner grundlosen Barmherzigkeit, such und bitt um Gnade. HEerr, sei gnädig mir armen Sünder.“ Die Gemeinde spricht: „Der barmherzige Gott wolle sich unser erbarmen, und uns unsere Sünden verzeihen und den Heiligen Geist geben, daß wir durch ihn seinen göttlichen Willen vollbringen und das ewige Leben empfangen. Amen.“ Der Liturg besteigt hierauf den Altar und spricht zu der noch knieenden Gemeinde: „Neiget eure Herzen zu Gott. Der allmächtige und barmherzige Gott hat sich unser erbarmt und seinen einigen Sohn für unsere Sünde in den Tod gegeben und um seinetwillen uns verziehen, auch allen denen, die an seinen heiligen Namen glauben, Gewalt gegeben, Gottes Kinder zu werden, und den Heiligen Geist verheißen. Wer glaubt und getauft wird, soll selig werden. Das verleihe uns Gott allen. Amen.“ (Manche Kirchenordnungen bedienen sich jedoch an dieser Stelle auch der Form der Beichte und Absolution, wie sie nach der Predigt anderwärts und allgemein üblich geworden ist.)

Nach Luthers Vorgang jedoch weisen viele andere Kirchenordnungen und Agenden dem Confiteor seine Stelle nach der Predigt an. Der Paraphrase des Vaterunsers und der Vermahnung an die Communicanten zwischen Predigt und Consecration (s. St. L. Agende S. 93 ff.) fügt nämlich Luther in seiner deutschen Messe 1526 die Bemerkung bei: „Ob man aber solche Paraphrasin und Vermahnung wolle auf der Kanzel flugs auf die Predigt thun oder für den Altar, lasse ich frei, einem jeglichen sein Willküre. Es siehet, als haben’s die Alten bisher auf der Kanzel gethan; daher noch blieben ist, daß man auf der Kanzel gemein Gebet thut, oder das Vaterunser fürspricht, aber die Vermahnung zu einer öffentlichen Beicht worden ist. Denn damit bliebe das Vaterunser mit einer kurzen Auslegung im Volk, und würde des HEerrn gedacht, wie er befohlen hat am Abendessen.“ Uebereinstimmend damit ließ Bugenhagen als „Pfarrherr zu Wittenberg“ bereits im Jahre 1524 „ein Ordnung christlicher Meß“ ausgehen, in welcher es heißt: „Darnach (nämlich nach gehaltener Predigt) mag er (der Pastor) die Beichtväter, so er wieder zum Altar gangen ist, oder auf der Kanzel mit solchen oder andern Worten anreden: Allerliebsten in Christo! Auf daß ihr mit fröhlichem,
\quellseite{212}
begierigem Herzen empfahet den zarten Fronleichnam und rosenfarbnes Blut unsers Seligmachers Jesu Christi: So demüthiget eure Herzen gegen Gott, bekennet ihm eure Sünd und Gebrechen mit herzlicher Begier und Sehnen nach seiner göttlichen Gnad und Hülff, so will ich jetzt, euer erwählter Priester, euch mittheilen die heilige Absolution.“ Eine Beichtformel ist dann zwar nicht angegeben; aber es folgt nicht nur eine Form der Absolution, sondern auch nach derselben „die Bekenntniß des Priesters“. Es heißt nämlich weiter: „Der allmächtig barmherzig Gott vergebe euch eure Sünd, und ich aus Befehl unsers HEerrn Jesu Christi, anstatt der heiligen Kirche, sage euch frei, ledig und los aller eurer Sünden im Namen des Vaters und des Sohnes und des Heiligen Geistes. Amen. Dieweil ich aber auch beschweret bin mit mancherlei Gebrechen, so bittet Gott treulich für mich, daß ich diesen seinen Dienst ihm zu Lobe, euch und mir zu Trost und Heil möge fleißig ausrichten und vollenden, wie geschrieben stehet am 123. Psalm: Unser Hülff soll sein im Namen des HEerrn, der Himmel und Erden gemacht hat.“ Im engen Anschluß an Luther und Bugenhagen findet sich daher auch Beichte und Absolution als Schluß der schönen Abendmahlvermahnung des Nürnberger Augustiner-Priors Joh. Volprecht gleichfalls aus dem Jahre 1524. Diese Abendmahlsvermahnung ist in viele Agenden übergegangen und so auch in die St. Louiser Agende S. 233–235 aufgenommen. Nach den Worten: „Denn wir alle sind Ein Brod und Ein Leib, dieweil wir alle Eines Brodes theilhaftig sind und aus Einem Kelche trinken“ heißt es nämlich weiter: „Und dieweil wir alle gesündiget haben und der Gnade Gottes bedürfen, so demüthiget eure Herzen vor Gott, dem Allmächtigen, bekennet eure Sünden und Gebrechen und sprecht mit dem öffentlichen Sünder: Gott, sei mir armen Sünder gnädig! mit herzlicher Liebe und Begierde nach seiner göttlichen Gnade und Hülfe, mit festem Glauben und Vertrauen auf sein gnädiges Zusagen, und vergebet von Herzen ein jeder seinem Nächsten, auf daß euch euer himmlischer Vater eure Sünden und Missethaten auch vergebe. So ihr solches thut, entbinde ich euch nochmals, anstatt der heiligen christlichen Kirche und aus Befehl und nach Zusage unsers HEerrn Jesu Christi, da er sprach: Wem ihr die Sünden vergebet, dem sind sie vergeben, von allen euren Sünden im Namen des Vaters und des Sohnes und des Heiligen Geistes. Amen.“\footnote{War etwa dies jene allgemeine Beichte vor dem Abendmahl, welche Osiander nicht neben der Privatbeichte geübt wissen wollte, und über welche sich dann der Nürnberger Beichtstreit erhob?} Wie oft hörte der Verfasser in seiner Jugendzeit
\quellseite{213}
die Vermahnung mit diesem Schluß in den Kirchen Nürnbergs! Ebenso enthält eine Wittenberger Gottesdienstordnung vom Jahre 1559 eine köstliche Form, in welcher Beichte, Absolution, Vermahnung zum Gebet und allgemeines Gebet zu einem so schönen Ganzen verbunden sind, daß man um so lieber der Beichte nach der Predigt den Vorzug gibt, da jene Wittenberger Form den Zusammenhang des Confiteor mit der vorausgehenden Predigt und dem nachfolgenden Abendmahl vortheilhaft erkennen läßt. Nur Mangel an Raum hält den Verfasser ab, diese Form mitzutheilen.

Indeß schon in der Kirchenordnung Herzog Erichs II., 1544, finden wir eine andere Weise des Confiteor nach der Predigt. Voraus geht hier eine Vermahnung zum Gebet, die im Vaterunser ausläuft. Dann folgen mit einleitender kurzer Vermahnung das apostolische Symbolum und die zehn Gebote und auf diese hinweisend Beichte und Absolution. Von etwa 1569 aber an wurde, soweit es der Verfasser erkunden kann, die Form des Confiteor bräuchlich, wie wir sie jetzt haben und immer mehr überall da aufgenommen, wo dasselbe seine Stelle nicht zum Anfang des Gottesdienstes hatte, nur mit dem Unterschiede, daß die Absolution in einem Theil der Kirchen noch eine Retention hatte. So schon jene Form der K.-O. Erichs, in welcher es nach der Absolution heißt: „Ueber denen aber, die ihre Sünden nicht von Herzen bekennen, auch die Zuversicht ihres Herzens nicht setzen auf das Verdienst Christi, bleibt der Zorn Gottes.“ Absolution mit Retention finden wir in den Kirchenordnungen von Braunschweig-Lüneburg 1657, Magdeburg 1663, desgleichen derer von Pommern, Schlesien u. a. O., in welchen gewöhnlich die Retention in dieser Fassung vorkommt: „Dagegen aber sage ich allen Unbußfertigen, Ungläubigen, Sichern, Gottes und seines Wortes und der hochheiligen Sacramente Verächtern aus Gottes Wort und im Namen Jesu Christi, daß ihnen Gott ihre Sünden vorbehalten hat, und gewißlich zeitlich und auch ewig strafen wird, wofern sie nicht in der Gnadenzeit ablassen und Buße thun, welche Buße wir ihnen von Herzen wünschen und sie dazu wollen vermahnt haben.“ Absolution ohne diese besondere Retention enthalten namentlich die verschiedenen sächsischen Agenden.

Folgt das Confiteor der Predigt, so erscheint es im Zusammenhang mit dieser als Uebung und Applicatio des in der Predigt verkündigten Gesetzes und Evangeliums, im Zusammenhang aber mit dem nachfolgenden allgemeinen Kirchengebet als Anleitung, in desto größerer Versicherung der Gnade Gottes in Vergebung der Sünden Gott im allgemeinen
\quellseite{214}
Gebete getrost und mit aller Zuversicht zu nahen. Dies drücken auch die einleitenden Worte zu einem jeden aus, wie sie in den betreffenden Agenden vorkommen. Die Einleitung zur Beichte lautet: „Nachdem wir das Wort Gottes haben angehöret, so wollen wir uns auch vor der hohen Majestät Gottes jetzo demüthigen und erstlich also beichten und sprechen:“. Und der Uebergang zum Gebet, wie z. B. in der Leipziger Agende 1681: „Weil uns der gnädige und barmherzige Gott unsere Sünde und Missethat vergeben, so wollen wir ihm auch ferner die Noth der ganzen Christenheit vortragen und miteinander also beten.“ Dieser Zusammenhang des Confiteor mit der Predigt und besonders mit dem „gemeinen Gebet“ steht auch Luther vor der Seele, wenn er schon 1520 in seinem „Sermon von den guten Werken“ schreibt: „Dieses gemeinen Gebetes ist noch von alter Gewohnheit blieben eine Anzeigung, wenn man am Ende der Predigt die Beichte erzählet und für alle Christenheit auf der Kanzel bittet. Aber es sollte nicht damit ausgerichtet sein, wie nun der Brauch und Weise ist; sondern sollte es lassen eine Vermahnung sein, durch die ganze Messe für solche Nothdurft zu bitten, zu welchem der Prediger uns reget; und auf daß wir würdiglich bitten, uns unserer Sünde zuvor ermahnet und dadurch demüthiget: welches aufs kürzeste soll geschehen, daß darnach das Volk im Haufen sämmtlich Gott seine Sünde selbst klage und für jedermann bitte mit Ernst und Glauben. O wann Gott wollte, daß irgend ein Haufe dieser Weise noch Mess hörte und betete, daß ein gemein, ein ernst Herzensgeschrei des ganzen Volkes zu Gott aufginge, wie unermeßliche Tugend und Hilfe sollte aus dem Gebet folgen? Was möchte schrecklicher allen bösen Geistern begegnen? Was möchte größer Werk auf Erden geschehen? Dadurch so viele Fromme erhalten, so viel Sünder bekehrt würden. Denn fürwahr, die christliche Kirche auf Erden nicht größere Macht noch Werk hat, denn solch gemein Gebet wider alles, was sie anstoßen mag. Das weiß der böse Geist wohl, darum thut er auch alles, was er vermag, dies Gebet zu verhindern.“ (Luther, Erl. Ausg. 20, 241. 242.) Ebenso sagt auch die zu Marburg 1566 gedruckte Kirchenordnung von Hessen vom Confiteor in seiner Beziehung zum Gebete: „Was können wir Größeres und Besseres von Gott begehren, denn daß uns unsere Sünden um Christi willen vergeben werden? Für dies allein sollen wir von Herzen sorgfältig und bekümmert sein. Denn wie wollen wir für Andere bitten oder etwas Fruchtbarliches in der christlichen Versammlung verrichten, so wir nicht ernstlich mit Gott versöhnt sind? Sollen wir aber von Gott
\quellseite{215}
Vergebung der Sünden begehren, so ist vonnöthen, daß wir uns ernstlich vor ihm demüthigen und für arme Sünder erkennen und bekennen. Derhalben sollen diese drei Stücke, nämlich die wahre Demuth, die Beicht und Bekenntniß mit Reue und Leid über die Sünde und das gläubige Gebet um Vergebung der Sünde um Christi willen nach dem Gesang in der christlichen Gemeinde geschehen.“


Wie vieles sich auch aus dem zuletzt Gesagten für die Stellung des Confiteor am Anfang des Hauptgottesdienstes anführen läßt — alles wohl erwogen dürfte die Stellung des Confiteor nach der Predigt und vor dem Kirchengebet doch den Vorzug verdienen.

Daß heute noch in den Gemeinden der separirten Freikirche Deutschlands und auch in mancher der hiesigen, namentlich aus der separirt lutherischen Kirche Preußens eingewanderten Gemeinden die ganze Versammlung bei dieser Beichte kniet, beruht auf altem, kirchenordnungsmäßigem Brauch. Sitzenbleiben ist jedenfalls unpassend.

\section[b. Das Kirchengebet]{b. Das Kirchengebet mit den Fürbitten etc}

Da das gemeinsame Gebet der im Namen Jesu versammelten Gläubigen eine so herrliche Verheißung hat, Matth. 18, 19., und Paulus in Bezug auf die gottesdienstlichen Versammlungen 1 Tim. 2, 1–4. ermahnt, „daß man vor allem Dingen zuerst thue Bitte, Gebet, Fürbitte und Danksagung für alle Menschen, für die Könige und für alle Obrigkeit“, so bekam schon im 2. Jahrhundert das „gemeine Gebet für allerlei Nothdurft“ seine liturgische Stelle. Es stand hier meist als Ueberleitung von der Katechumenen- zur Gläubigenmesse. Im 3. Jahrhundert ist es in ausgeführterer Form mit der Communion in Verbindung gebracht, theils vor, theils nach der Consecration stehend. In der römischen Messe schließt es sich im Zusammenhang mit der ganzen Meßopferidee um die Consecration.

Für den lutherischen Hauptgottesdienst reihen nun zwar einige Kirchenordnungen der Reformationszeit das allgemeine Kirchengebet schon bei den Lectionen, also vor der Predigt, ein; die meisten Kirchenordnungen jedoch haben es an unserer Stelle.

Verschieden aber ist die Form. Nach Luthers Vorgang in seiner deutschen Messe, wie in seiner Hauspostille („Eine gemeine Form, wie zum Beschluß der
\quellseite{216}
Predigt das Volk zum gemeinen Gebet soll vermahnet werden“), bestand nach einem Theil der Kirchenordnungen das Gebet im Vaterunser mit vorausgehender Vermahnung, in welcher die Dinge bezeichnet sind, um die man im Vaterunser vereiniget bitten will, und zugleich auch die speciell begehrten Fürbitten und Danksagungen namhaft gemacht werden. So die Agenden von Breslau 1526, Lüneburg 1569 und 1598, Oldenburg 1573, von Nürnberg etc. Häufig folgte auf das Vaterunser als stehendes Lied: „Erhalt uns, HErr, bei deinem Wort“ mit dem „Verleih uns Frieden gnädiglich“, welche als zum Kirchengebet gehörig betrachtet wurden. Die directe Gebetsform, deren sich schon frühzeitig andere Agenden bedienen, wurde jedoch die herrschende. Die verbreitetste Form des gemeinen Gebetes dieser Art stammt aus der Mitte des 16. Jahrhunderts und ist schon in der Baden’schen Agende 1556 und in der von Herzog Wolfgang 1557 aufgenommen. Mit verschiedenen Eingängen existirt sie in mancherlei Erweiterungen und Variationen. Außerdem findet man seine Spuren in vielen andern selbstständigen Formen des gemeinen Gebetes und unvermuthet begegnet man bald diesem, bald jenem Stück dieser Erstlingsform. So ist z. B. in dem „Kirchengebet“ der St. Louiser Agende der erste Absatz, die Bitte für den Lehrstand, wörtlich wiedergegeben, desgleichen auch die Bitte für die Angefochtenen und Verfolgten.

Auch der Gebrauch besonderer Festgebete im Hauptgottesdienst gehört nicht einer späteren Zeit an, sondern datirt sich schon aus dem Ende des 16. Jahrhunderts.

Unter den auf das allgemeine Kirchengebet folgenden speciellen Fürbitten und Danksagungen steht überall die Fürbitte für die Communicanten oben an.\footnote{Formulare für die speciellen Fürbitten und Danksagungen dieses Orts in directer Gebetsform vermochte der Verfasser nicht zu finden. Man hat wohl ehedem die indirecte Form in möglichster Kürze gebraucht. So heißt es z. B. in der Braunschw.-Lüneb. Kirchenordnung 1657: „Für etliche Kranken, Schwangere oder Andere etc. zu bitten oder Dank zu sagen, ist solches dieses Orts kurz zu thun.“ In der Schwarzb. Kirchenordnung 1657 findet sich allein folgende „Formul der Vorbitte für die communicanten“: „In unser Gebet wollen wir auch mit einschließen diejenigen Personen, welche jetzo zum heiligen Abendmahl des HEerrn gehen wollen, an der Zahl —. Der allmächtige Gott, der ihnen das Wollen gegeben hat, wirke in ihnen auch das Vollbringen und verhelfe gnädig, daß sie solche heilige Speise nicht zum Gericht, sondern zu ihrer Seelen Heil und Seligkeit empfangen mögen um Jesu Christi willen. Amen.“ Nach dieser Form sollten dann wohl auch die andern Fürbitten etc. geschehen.} Eine vorkommende öffentliche Abbitte oder eine Ercommunication ist am passenden Orte eingereiht.

\quellseite{217}
Nach der Braunschw.-Lüneb. Ordnung soll das Vaterunser gleich auf diese Fürbitten etc. folgen, also schon vor den übrigen kirchlichen Verkündigungen gebetet werden, „damit die ganze Gemein unverrückt in ihrer Andacht bis zu Ende des Gebets gelassen werde“. Die Aufforderung zum Beten desselben geschieht nach der Kirchenordnung von Dels 1592 etwa mit diesen Worten: „Diese und alle andere Noth der Christenheit, was auch ein jeder für eigen Noth und Anliegen hat und befindet, wollen wir Gott im Namen Christi seines Sohnes in unserm andächtigen, gläubigen und demüthigen Gebet und Vaterunser fürtragen.“

Nicht überflüssig für den Prediger ist auch heute noch die Bemerkung der erstgenannten Kirchenordnung: „Hierauf liest er der Gemeine sein langsam, laut und deutlich vor das gemeine Gebet.“ Bei der Schwachheit des Fleisches und dem sonntäglichen und nothwendigen Brauch einer und derselben Formel liegt die Gefahr eines eiligen und andachtslosen Ablesens für den Prediger nur zu nahe.\footnote{Es ist daher statthaft, daß mancher Prediger sich willkürlich allerlei Kürzungen im Kirchengebete erlaubt. Er ist der Mund der Gemeinde; und gerade diejenigen, welche andächtig mitbeten und denen daher die Worte des Gebets sich eingeprägt haben, empfinden solche Willkür des Predigers schmerzlich und ist ihnen dieselbe ärgerlich.} Dieselbe Kirchenordnung empfiehlt auch das Knien der ganzen Gemeinde beim Vaterunser. Jedenfalls sollte die Gemeinde nicht nur bei dem Vaterunser, sondern auch beim Kirchengebet stehen, wie es schon von Alters her allewege Brauch war, namentlich an Sonn- und Festtagen.

Die übrigen kirchlichen Verkündigungen, bei denen, wie es scheint, die Gemeinde saß, schlossen an vielen Orten mit einer „Vermahnung zum Opfer für die Armen“. So heißt es im Frankfurter Agendbüchlein 1688: „Arme und kranke Leute laßt euch in eurem Gebet und Almosen um des HEerrn Christi willen jederzeit treulich und fleißig befohlen sein.“ Ebenso in der Agende von Dels 1592: „Zum Beschluß lasset euch das Armuth befohlen sein, bedenket den Gotteskasten, Hospital und andere hausarme Leute, die ein gut Zeugniß haben.“ Somit wäre hier die geeignete Stelle zur Abkündigung der Collecten für besondere kirchliche Zwecke.

Das Schluß-Votum: „Der Friede Gottes, welcher höher ist“ etc. findet sich überall.

\quellseite{218}
\section{B}

\section{1. Eingangsformeln zur Predigt}

\subsection{a}

Auf daß wir unserm lieben HEerrn Gott heut seinen Dienst leisten und nach seinem Befehl den Feiertag heiligen, das ist, mit Gottes Wort, welches allein heilig ist und alles heilig macht, zubringen; so wollen wir jezund das heilige Evangelium hören, und Gott um Gnade bitten, daß wir’s also hören, auf daß seine Ehre gepreiset, unser Glaube gestärket und unser Leben gebessert werde. Betet ein Vaterunser. (Luther, Hauspostille: „Eine gemeine Vorrede auf die Hauspredigten.“)

\subsection{b}

Liebe Andächtige in Christo! Lasset uns Gott bitten um den Heiligen Geist zu seinem göttlichen Wort, der uns in alle Wahrheit leite, unsere Herzen in wahrer Erkenntniß und Furcht Gottes heilige, daß wir sein göttlich Wort zur gottseligen Lehre, Trost und Besserung hören mögen und durch desselben Kraft im Glauben JEsu Christi ewiglich leben. In dessen Namen sprechet: Vaterunser etc.

Erinnerung: „In solchem Gebet können die Prediger auch auf den Text sehen, den sie predigen wollen, und auf die Gelegenheit der Feste, aus denselbigen Ursache und Schluß nehmen, die christliche Gemeinde zum Gebet um die Gnade des Heiligen Geistes zum Wort Gottes zu vermahnen.“ (Pommern 1568. 1691.)

\subsection{c}

Gnad, Fried und Barmherzigkeit von Gott, unserm lieben himmlischen Vater, sei durch JEsum Christum, seinen geliebten Sohn, unserm HErrn, sammt dem Heiligen Geist jetzt und zu allen Zeiten mit uns allen. Amen.

Geliebte im HErrn! Wir wollen den allmächtigen, ewigen, gütigen Gott, unsern getreuen Vater im Himmel, von Herzen mit einander anrufen und bitten um die Gnade, Hülfe und den Beistand des Heiligen Geistes, damit wir sein heiliges göttliches Wort dermaßen hören und handeln mögen, daß dadurch sein großmächtiger Name gepriesen, unser schwacher Glaube gestärket und unser sündliches Leben wahrhaft gebessert werde. Solches und alles Gute von Gott, dem Allmächtigen, zu erlangen, thut euer Gebet zum HErrn und singet mit einander von Herzen: „Nun bitten wir den Heiligen Geist“ etc.\footnote{Nach der Schwarzb. Agende 1650 soll das sogenannte Kanzellied je nach Gelegenheit des Festes wechseln.}

Sprecht auch darauf ein heiliges Vaterunser. (Frankf. a. M. 1688.)

\quellseite{219}
\section{2. Eingangsgebete zur Predigt}

\subsection{a}

Allmächtiger Gott, lieber himmlischer Vater, dieweil all unser Heil daran liegt, daß wir deines heiligen Wortes einen rechten Verstand haben; so verleih uns allen, die im Namen deines Sohnes versammelt sind, daß unsere Herzen, gefreiet von allen weltlichen Geschäften, mit allem Fleiß und Ernst dein heiliges Wort anhören und fassen mögen, damit wir daraus deinen väterlichen Willen recht erkennen und demselben nachleben, dir zu Lob und Preis, unserm Nächsten zu fruchtbarer Besserung und uns allen zum ewigen Heil. Durch deinen lieben Sohn, unsern HEerrn und Heiland Jesum Christum. Amen. (Hanau 1659.)

\subsection{b}

O allmächtiger Gott, himmlischer Vater, auf deinen Befehl sind wir in dein Haus gekommen zu beten und dein seligmachendes Wort zu hören, Lehre, Trost und Erquickung für unsere Seelen zu holen. Weil aber leider unsere Herzen und Verstand verfinstert, wie auch von Natur untüchtig, träg und schläfrig zu hören, und unser Gedächtnis zu merken und zu behalten schwach ist; so bitten wir dich demüthiglich, du wollest unsern Mund, Ohren, Herz und Gedächtnis durch deinen Heiligen Geist eröffnen, damit wir dein heiliges Wort recht lehren und mit herzlicher Begierde, Eifer und Andacht hören, dich daraus erkennen lernen und im Glauben und Gottseligkeit gestärket werden, auf daß wir uns kraft desselben in allen Anfechtungen und Nöthen bis an unser letztes Seufzen wider Sünde, Tod und Teufel kräftiglich trösten und das ewige Leben dadurch erlangen mögen. Durch Jesum Christum, unsern HEerrn. Amen. (Limburg 1666.)

\subsection{c}

\begin{verse}
O Christe, Wahrheit und Leben, wir bitten, du wollest uns geben\\
Deinen Geist von oben mit seinen heiligen Gaben,\\
Daß dein Wort rein hie auf Erden mög verkündigt werden.\\
O gib, daß die Red vorhanden, recht erklärt werd und verstanden,\\
Laß es ihr gelingen und sie unser Herz durchdringen,\\
Zu Lob und Ehr deinem Namen. Darauf sprechen wir Amen. Amen.
\end{verse}

(Dieses aus dem Jahre 1550 zum Singen vorhandene Lied dürfte sich auch zum Eingangsgebet eignen.)

\quellseite{220}
\chapter{XII. Offertorium}


Bekanntlich ging im apostolischen Zeitalter der Feier des heiligen Abendmahles ein Liebesmahl (Agape) voraus, zu welch letzterem die wohlhabenderen Christen Speise und Trank, namentlich Brod und Wein, mitbrachten, und diese mit den Aermeren theilten. Von dem Brod und Wein wurde dann das Nöthige für das darauf folgende Abendmahl ausgesondert. Als man nun in der ersten Hälfte des zweiten Jahrhunderts aus triftigen Gründen die Agapen einstellte und die Abendmahlsfeier mit dem vormittägigen Predigtgottesdienste vereinigte, nahm man aus der Verbindung mit den Agapen in die Abendmahlsfeier die Sitte mit herüber, daß die dazu nöthigen Elemente, Brod und Wein, als freiwillige Gaben der Gemeindeglieder dargebracht wurden, und galt die Annahme dieser Gaben von Seiten des Clerus als eine förmliche Anerkennung der Kirchengemeinschaft (wie denn auch überhaupt von groben, unbußfertigen Sündern, von Gebannten u. dgl. keine Gaben für kirchliche Zwecke angenommen wurden). Was von diesen Gaben nach dem Abendmahl noch übrig blieb, wurde nebst andern dabei dargebrachten Gaben zum Unterhalt der Diener der Kirche und zur Unterstützung der Armen verwendet.

Diese Spenden an Brod und Wein, wie auch die andern Liebesgaben, nannte man Opfer, die Darbringung derselben Opferung (oblatio) und den Tisch oder den Ort, auf welchen sie niedergelegt wurden, Altar. Insofern jedoch diese dargebrachten Gaben der Liebe als primitiae creaturarum (Erstlinge der Creaturen) galten, sollten sie zugleich der symbolische Ausdruck für die innere Selbsthingabe der versammelten Gläubigen an den HEerrn sein, als die sich selbst, nebst allem, was sie haben, mit Freuden in den Dienst Gottes stellen. Insbesondere aber wurden die für die Abendmahlsfeier bereit liegenden Elemente als primitiae creaturarum nicht nur mit einem erhebenden Dankgebete, der Präfation, dargebracht, sondern auch mit der hinzugefügten Bitte, Gott wolle diese Gaben des Brodes und Weines zu ihrem sacramentalen Zwecke annehmen und durch die Herabsendung des Heiligen Geistes Brod und Wein zur Gemeinschaft des Leibes und Blutes Christi und die Empfänger zum würdigen Genuß des Sacramentes heiligen.

Es war also dies alles ursprünglich Gemeinde-Oblation und erhielt sich diese Sitte Anfangs auch noch in den Kirchen Roms, „indem die Gläubigen, zuerst die Männer, dann die Weiber, ihre Gaben an Brod und Wein in weißen
\quellseite{221}
Tüchern, dann die Priester und Diakonen bloß in Brod und zuletzt die Sänger in Wasser\footnote{Im Morgenlande pflegte man überhaupt dem Wein wegen seiner Stärke beim Genuß Wasser beizumischen. Lediglich aus diesem Grunde bediente man sich auch beim heiligen Abendmahl ursprünglich des Wassers.} darbrachten, während der Chor das Offertorium sang, bestehend aus antiphonischen Psalm- und andern Stellen (versus offertorii), die mit jedem Gottesdienste wechselten. Zu Anfang des Mittelalters begann man jene Gaben in Geld zu verwandeln, womit diese Oblationen ihre eigentliche Bedeutung verloren und ganz verschwinden mußten.“ (Schöberlein.)

Seit der Mitte des dritten Jahrhunderts jedoch hatte allmählich die unevangelische Vorstellung von einem clericalen christlichen Priesterstande Eingang gefunden. Folgerichtig bildete sich nun auch die Vorstellung von dem Abendmahl als einer Opferhandlung des christlichen Priesters mehr und mehr aus, und nach und nach gewöhnte man sich, den mit dem Abendmahl verbundenen, jetzt unmittelbar an die vollzogene Consecration sich anschließenden Gebeten, insbesondere denen für die Verstorbenen, eine so hohe Bedeutung beizulegen, daß bereits jetzt das Abendmahl als eine oblatio pro mortuis zu gelten begann. Und so kam es denn endlich zum Greuel des papistischen Meßopfers. Was erst bloßes Darbringen des Brodes und Weines war, im guten Sinne Opfer hieß und im Gegensatz zu dem blutigen Opfer des alten Testamentes als „unblutig“ bezeichnet wurde, das wurde nun ein sogenanntes Opfer des Leibes und Blutes Christi selbst, zum Unterschied von seinem blutigen Opfer am Kreuz als „unblutig“ bezeichnet und als Sühnopfer für die Lebendigen und die Todten fortwährend wiederholt, ja die angeblich sühnende Kraft schon auf die noch ungesegneten Elemente in den Opfergebeten ausgedehnt.

Und so traten denn an die Stelle der alten Gemeinde-Oblationen in der Liturgie diejenigen Handlungen des Priesters, welche die wirkliche Darbringung des Sühnopfers vorbereiten. Zunächst spricht der Priester den nach Zeit und Umständen verschiedenen versus offertorii selbst, und zwar im leisen Flüsterton. Dann ergreift er die Hostie, legt sie auf der Patene wieder und ruft, gleichfalls im Flüsterton, nun Gott an, er wolle dies „unbefleckte Opfer“, das er für seine Sünden, für die Sünden aller Anwesenden und aller Gläubigen, für die Lebendigen und die Todten darbringe, gnädig annehmen und denselben zum Heile gedeihen lassen.\footnote{Dieses erste der Opfergebete lautet: „Suscipe, sancte Pater, omnipotens, aeterne Deus, hanc immaculatam hostiam, quam ego indignus famulus tuus offero tibi. Deo meo vivo et vero, pro innumeris peccatis et offensionibus et negligentiis meis, et pro omnibus circumstantibus, sed et pro omnibus fidelibus christianis, vivis atque defunctis: ut mihi et illis proficiat ad salutem in vitam aeternam. Amen.“ In seiner Schrift „Von dem Greuel der Stillmesse, so man den Canon nennt“, in der er die einzelnen Theile des Canons glossirt, bewertet Luther zu diesem und dem hernach sich auf den Kelch beziehenden Opfergebet: „Siehe, da stehet der Pfaff und hat ein Hostien oder Oblat, von lauter Brod, und ein Kelch mit Wein für sich und redet mit Gott, daß er das ansehen solle, und solle so köstlich sein und so viel gelten, daß es für die ganze Christenheit ein Opfer sei für Gott. Ist das nicht Christus Blut gelästert, daß ihm Brod und Wein gleich geacht wird? Es merckt und siehet niemand, was es für Lästerung ist.“ (Erl. A. 29, 118.) — An einem andern Ort bemerkt Luther: „Was böse in der Messe ist, ist vom Opfer und Werk, das hat Gott wunderlich geschickt, daß alles der Priester heimlich liest, und heißt die Stillmesse.“} Diesem Gebete schließen sich dann noch vier andere Opfergebete
\quellseite{222}
an, die gleichfalls leise gesprochen werden, während die darauf folgende Präfation der Priester laut spricht oder singt und die daher Secreta heißen, mit dem ersten zusammen aber das Offertorium der römischen Messe bilden.

Wie Luther zu diesem Stück der vorgefundenen Messe stand, sowie zu den der Präfation folgenden, die Testamentsworte umgebenden Opfergebeten, dem gleichfalls mit den Testamentsworten leise gesprochenen Canon missae, das spricht er unter anderm in seiner ersten liturgischen Schrift, der Formula missae oder „Weise, christlich Meß zu halten 1523“ also aus: „Hernach, was dem folgt“ — nämlich dem eben erwähnten „Opfergesang“, dem versus offertorii, der Psalmodie — „klingt und stinkt allzumal eitel Opfer und sind die heiligen Worte des Lebens der Seligkeit“ — die Testamentsworte — „mitten unter das greuliche Opferwesen dermaßen hineingesetzt, wie vor Zeiten die Arche des HEerrn im abgöttischen Tempel neben dem Götzen Dagon. . . . Darum wollen wir alles unterlassen, so nach Opfer klingt, sammt dem ganzen Canon, und allein behalten, das rein und heilig ist.“

Demgemäß behielt man als „rein und heilig“ zunächst die altkirchliche Psalmodie, den versus offertorii, bei. „In protestantium ecclesiis restituta est loco offertorii Psalmodia“ sagt Calvör in seinem Rituale etc. I. c. I. p. 565. Zwar war dies mancherorts erst ein von der Predigt zum Abendmahl passendes Uebergangslied, anderorts aber ein specifischer, mit den Festtagen wechselnder Offertorien-Gesang des Chors, „der Gesang, den man hiervor pro offertorio gehabt hat“, wie es in der Märkischen Kirchenordnung heißt, vielfach auch nahm man zur Psalmodie den 111. Psalm, für welchen zu dem Ende Spangenberg 1545 die Melodie nach einem Psalmenton componirte, oder man gebrauchte auch als Preis der Freundlichkeit des HEerrn Ps. 34, 9., oder als Ausdruck der Brudergemeinschaft Ps. 133 u. dgl. m. Seit der Mitte des 17. Jahrhunderts jedoch fand das aus Ps. 51 genommene „Schaffe in mir, Gott, ein
\quellseite{223}
reines Herze“ immer mehr Verbreitung und wurde zuletzt fast allgemein die Psalmodie. Und fürwahr, die Wahl derselben dürfte wohl als die sinnigste und darum auch als die passendste erscheinen, denn diese Psalmodie ist der schönste Ausdruck für die Darbringung der Opfer, die Gott am besten gefallen und die zum würdigen Genuß des heiligen Abendmahls gehören, denn von diesen heißt es V. 19.: „Die Opfer, die Gott gefallen, sind ein geängstigeter Geist; ein geängstiget und zerschlagen Herz wirst du, Gott, nicht verachten.“ Die in alle Choralbücher übergegangene schöne Melodie findet sich zuerst im Gothaischen Cantional 1646 und wird G. Winter, „pastor at Henrici Henneberg“, zugeschrieben. Es gibt zu derselben auch eine nicht minder schöne Variante, weshalb sie gleichfalls in das von H. F. Hölter herausgegebene St. Louiser Choralbuch unter Nr. 335 aufgenommen ist. Dieser Psalmodie vorausgehend fügt Hommels musikalische Liturgie Vers 20. und 21. des 51. Psalms hinzu: „HEerr, thu wohl an Zion nach deiner Gnade“ etc. Da die mitgetheilte Melodie so schön ist, so folgt sie hier mit. Es könnte diese Psalmodie etwa vom Chor gesungen werden und die andere: „Schaffe in mir, Gott“ etc. von der Gemeinde unmittelbar folgen.

Während nun aber Luther mit Recht alles unterlassen haben wollte, was nach dem papistischen Opfer auch nur roch oder wie dasselbe klang, so war er doch von Anfang geneigt, nach Brauch der alten Kirche die gottesdienstliche Darbringung der Liebesopfer mit der Abendmahlsfeier zu verbinden. Indem er in seiner „deutschen Messe und Ordnung des Gottesdienstes 1526“ einer missa fidelium das Wort redet, „wenn man die Leute und Personen hätte, die mit Ernst Christen zu sein begehrten“, fährt er fort: „Hier könnt man auch ein gemein Almosen den Christen auflegen, das man williglich gäbe und austheilete unter die Armen nach dem Exempel St. Pauli 2 Cor. 9, 1. 2. 12.“ (St. L. Ausg. X, 229. Erl. Ausg. 22, 231.) Und im „Sermon vom Neuen Testament, das ist, der heiligen Messe 1520“ schreibt er: „Von diesem Brauch (da man nämlich im Alten Testament die Speise „webete“) sein überblieben noch drei Wahrzeichen. Das erst, daß man das erst und letzt Gebet der Meß heißet Collecten, das ist, Versammlung; damit bezeuget wird, daß dieselben Gebet seien gesprochen als ein Benedictie und Gratias über dieselben zusammengetragen Speis, dieselben zu segnen und Gott zu danken, und nach der Lehre St. Pauli (1 Cor. 10, 30. 31.). Das ander, daß die Leut nach dem Evangelio noch zum Opfer gehen, davon der Gesang, den man daselbst singt, heißt Offertorium, das ist, ein Opfer.

\quellseite{224}
\begin{quote}
• Das dritte, daß der Priester mit der Patene aufhebt und opfert Gott die ungesegnete Hostie, eben dieselbe Zeit, wenn man singt das Offertorium und die Leut opfern; damit angezeigt wird, daß nicht das Sacrament Gott geopfert wird von uns, sondern allein dieselben Collecten und Opfer der versammelten (gesammelten) Speis und Güter, daß da Gott für gedankt wird und sie gesegnet, auszutheilen allen Dürftigen.“ (Erl. Ausg. 27, 157. 158.)
\end{quote}

Diese Darbringung oder Einsammlung der Liebesopfer für die Bedürfnisse des Gottesdienstes, der Armen oder sonstige wohltätige oder kirchliche Zwecke ist denn auch allgemeiner Brauch geworden und geblieben. Während es nun aber bei den Reformirten herrschende Sitte wurde, diese Einsammlung nach dem Schluß des Gottesdienstes durch an den Kirchstühlen ausgestellte Becken vornehmen zu lassen, zog man es allgemein bei den Lutheranern vor, ihr eine Stelle im Gottesdienste anzuweisen und dieselbe durch Umhertragen des sogenannten Klingelbeutels\footnote{Von dem Umtragen des Klingelbeutels, das heutzutage mancher auch nach Beseitigung des Glöckleins für unpassend hält, schreibt Dr. Slevogt: „Dieser modus collectandi kömmt sonder Zweifel von den alten Oblationen her, wie denn die Art, die Oblationes in einem Sacke in der Kirche zu colligiren, schon in Ordine Rom. vorgeschrieben ist: Et accedentes subdiaconi sequentes cum acolythis, qui sacculum portant a dextris et sinistris altaris“ etc. (Rechte der Altäre, Taufsteine etc. 1732. S. 458.)} oder einer Büchse oder eines andern Gefäßes geschehen zu lassen. Meist erfolgte die Einsammlung während der Psalmodie, oder wo die Zeit während derselben nicht ausreichte, schon unter dem Credo-Gesang. Mit welchem bestimmten Bewußtsein eines innern Zusammenhangs des Lob- und Liebesopfers der Gemeinde diese Einrichtung getroffen wurde, zeigt folgende Stelle der „Cölnischen Reformation“ 1544: „Also sollen die Gläubigen, dieweil man den Glauben singt, ihre freiwilligen Opfer bringen, ein jeder nach dem Segen, der ihm von der milden Hand Gottes verliehen ist; dazu sie die Prediger allemal fleißig vermahnen sollen und anzeigen, wie solche Opfer der wahren Bekenntniß des Glaubens und dem gläubigen Gebet alleweg anhangen sollen.“ Wie sehr wird man hierbei an den Brauch des Alten Testaments erinnert, nach welchem die Darbringung des Brandopfers unter dem Gesang der Sänger und dem Schall der Cymbeln, Psalter und Harfen der Leviten und den Trommeten der Priester geschah (2 Chron. 29, 25–28.), womit ohne Zweifel schon damals angezeigt werden sollte, daß man freiwillig und mit Freuden dem HEerrn seine Gabe darbringe und worüber ein David weissagend hinwies auf die Zeit des Messias, dem sein Volk nach seinem Siege williglich opfern werde im heiligen Schmuck (Ps. 110, 3.).

\quellseite{225}
Wie sehr überhaupt die Sammlung von Liebesgaben im Gottesdienst als Stück des geistlichen Priesterthums aufgefaßt wurde, beweist der Name Opfer, mit dem noch an vielen Orten jegliche kirchliche Sammlung, betreffe sie die Armen, die Gemeindebedürfnisse, Schulen, Mission etc. bezeichnet wird, die Sammler heißen Opfermänner, die in der Kirche angebrachten Stätten zum sonstigen Einlegen von Almosen Opferstock, Opferbüchse. In manchen der hiesigen Gemeinden der aus Preußen eingewanderten Lutheraner sah der Verfasser bei der Communion eine besondere Liebesgabe auf den Altar legen, und bei einer Kirchweih eingewanderter norwegischer Lutheraner nach der Predigt die ganze Gemeinde singend um den Altar ziehen und auf der rechten Seite desselben die Gabe niederlegen. Im Common Prayer der Episcopalkirche ist folgende schöne Weise vorgeschrieben: „Nach der Predigt kehrt der Priester zum Altar zurück und beginnt mit einem oder nach Belieben auch mehreren der nachstehenden Sprüche das Offertorium (Offertory).“ Diese von dem „Priester“ vorzulesenden Sprüche beziehen sich sämmtlich aufs Geben, wie z. B. Matth. 5, 16. 6, 19. 20. Luc. 19, 8. 1 Cor. 9, 7. 11. 13. 14. 2 Cor. 9, 6. 7. Gal. 6, 6. 7. 10. 1 Tim. 6, 17. 18. Hebr. 6, 10. 13, 16. u. dgl. Darauf heißt es nach einer Bemerkung, daß während des Sammelns des Opfers auch ein „Opfer- und Danklied“ gesungen werden könne, weiter: „Während der Verlesung der Sprüche sollen die Diakonen oder andere dafür bestimmte Männer die Almosen für die Armen und andere Opfer der Gemeinde in einem passenden vom Kirchspiel dafür beschaffenen Becken einsammeln und dasselbe dem Geistlichen ehrerbietig überreichen, welch letzterer es dem HEerrn demüthig darbringen und auf den Altar stellen soll.“

Das Offertorium des lutherischen Hauptgottesdienstes oder wie man es sonst nennen mag, ist hiernach etwas ganz anderes als das der papistischen Messe. Es ist der kirchliche Ausdruck für die gemeinsame Darbringung von Opfern des geistlichen Priesterthums — des bußfertigen Herzens durch die Psalmodie: „Schaffe in mir, Gott“ etc., und der dankbaren Liebe durch die damit verbundene oder unter Umständen ihr vorhergehende oder nachfolgende Einsammlung der Liebesgaben. In der Vorrede zur zweiten Auflage seiner Agende, 1853, findet Löhe jedoch in dieser Form des Offertoriums noch nichts Vollendetes. Für den „Höhepunkt des Offertoriums“ hält er die altkirchliche „Darbringung der primitiae creaturarum, das ist, die betende Darbringung der Elemente, damit der HEerr seinen heiligen Leib und sein theures Blut damit vereinige“. Gleichwohl trägt er Bedenken, sie in eine lutherische Agende so
\quellseite{226}
ohne Weiteres aufzunehmen, lediglich jedoch nur wegen „der Geneigtheit zu Mißverständnissen und friedensstörenden Auslegungen unschuldiger Reden und der Handlungen, welche gegenwärtig an der Tagesordnung ist“. Wenn er nun dabei bemerkt: „Sie ist auch dem lutherischen Alterthume nicht unbekannt, wenngleich sie nicht durchgeschlagen hat“, so dachte er ohne Zweifel insbesondere an eine Stelle in der von ihm in seiner 1839 herausgegebenen „Sammlung liturgischer Formulare“, Heft 4, mitgetheilten deutschen Messe von Bugenhagen,\footnote{„Ain ordnung Christlicher Meßien wie gehalten wirdt, von dem Ehrwürdigen herren Johann Bugenhagen auß Pommern, Pfarrherr zu Wittenberg 1524.“} indem daselbst nach Zurichtung der Elemente „der Priester mit dem Volk“ also betet: „O allergütigster Vater, barmherziger Gott, hilf, daß dieses Brod und Wein uns werde der wahrhaftig Leib und das unschuldig Blut deines allerliebsten Sohns unsers HEerrn Jesu Christi. Amen.“ Daß nun diese Darbringung trotz ihres unverfänglichen Wortlauts „nicht durchgeschlagen“ hat und allgemein geworden ist, wird erklärlich, wenn man bedenkt, daß Luther ganz entschieden alles unterlassen haben will, was nach Opfer auch nur riecht oder klingt. Es ist auch gut, daß dieselbe nicht „durchgeschlagen“ hat und in die lutherischen Gottesdienstordnungen des 16. und 17. Jahrhunderts übergegangen ist, denn es möchte, wie ehedem, aus der an sich unverfänglichen Darbringung der primitiae creaturarum die lästerliche Darbringung eines „unbefleckten Opfers“ für die Lebendigen und Todten wieder werden, sonderlich in unserer Zeit! Beiläufig sei hier noch bemerkt, daß eine solche Darbringung der Elemente auch im Book of Common Prayer der Episkopalen sich findet, und zwar als eine “Oblation”, verbunden mit einer “Invocation”, dabei charakteristisch ist, daß dieselbe der Consecration nicht voransteht, sondern erst unmittelbar derselben folgt.\footnote{The Oblation: Wherefore, O Lord and heavenly Father, according to the institution of thy dearly beloved Son our Saviour Jesus Christ, we, thy humble servants, do celebrate and make here before thy divine Majesty, with these thy holy gifts, which we now offer unto thee, the memorial thy Son hath commanded us to make; having in remembrance his blessed passion, etc. — The Invocation: And we most humbly beseech thee, O merciful Father, to hear us, and, of thy almighty goodness, vouchsafe to bless and sanctify, with thy Word and Holy Spirit, these thy gifts and creatures of bread and wine; that we, receiving them according to thy Son our Saviour Jesus Christ’s holy institution, in remembrance of his Death and Passion, may be partakers of his most blessed Body and Blood, etc.}

Während des Offertoriums pflegte nicht allein der Pastor an den Altar zu treten, um die heiligen Gefäße zu enthüllen und sich besser zur Handlung des Sacraments vorzubereiten,\footnote{Wird ihm assistirt, so steht während der Consecration der Assistent vor den Stufen des Altars entweder in der Mitte oder zur Seite.} sondern auch die Communicanten sich nach dem
\quellseite{227}
Altar-Chor zu begeben und, falls er groß genug war, sich in demselben aufzustellen, die männlichen Communicanten zur rechten, die weiblichen zur linken Seite des Altars. Schöberlein stellt in Bezug auf das letztere aus verschiedenen Kirchenordnungen folgende Bestimmungen zusammen. Es sollten die Communicanten „nicht mit üppiger und stolzer Kleidung noch leichter Geberde, sondern in ehrbarer, reinlicher Kirchenkleidung, und zwar in schwarzem Habit etc. fein züchtiglich, mit großer Reverenz und Andacht in den Chor treten, auch nicht aus dem Chor gehen bis zum Segen“, indem „die Chör in den Kirchen der Ursachen gemacht, daß die Communion darin gehalten werde und allein die Communicanten, wenn man die Communion hält, von Anfang bis zum Ende darin stehen, sitzen und knieen sollen, damit sie alle Ding, was bei solchem heiligen Gottesdienst mit Lesen, Singen und Beten geschieht, desto andächtiger und besser vernehmen und anhören mögen“. Daß die Communicanten an einem Platz besonders stehen, findet auch Luther für passend. In seiner „Weise, christlich Meß zu halten“, 1523, schreibt er: „Wenn die Communion gehalten wird, schicket sich’s fein, daß die, so zum hochwürdigen Sacrament gehen wollen, sich zusammenhalten und an einem sonderen Ort alleine stehen; denn auch dazu beide, Altar und Chor, gebauet sind. Nicht, daß es vor Gott etwas sei, man stehe hier oder dort, oder daß es etwas zum Glauben thäte, sondern darum vonnöthen, daß die Personen öffentlich gesehen und erkannt werden, sowohl von denen, die das Sacrament empfangen, als von denen, die nicht hinzu gehen, damit hernach ihr Leben auch desto baß gesehen, geprüft und geoffenbaret könnte werden. Denn die Nießung dieses Sacraments in der Gemeinde ist ein Stück christlicher Bekenntniß, dadurch die, so hinzugehen, vor Gott, Engeln und Menschen bekennen, daß sie Christen seien. Um deswillen ist fleißig wahrzunehmen, daß nicht Etliche das Sacrament heimlich abstehlen und nachmals, unter andern Christen vermenget, nicht können erkannt werden, ob sie wohl oder übel leben. Wiewohl ich hier auch kein Gesetz stellen will, sondern dies allein anzeigen, daß von Christen, so allerdings frei sind, frei ohne Zwang gethan und gehalten werde alles, so ehrlich und ordentlich ist, 1 Cor. 14, 40.“ Solche schöne Weise kann natürlich nur da stattfinden, wo der Altarchor groß genug ist. Ist das nicht der Fall oder ist die Gemeinde so groß, daß ein die Andacht störendes Gedränge um den Altar entsteht, so ist es besser, die Communicanten bleiben auf ihren Sitzen und erheben sich dafür beim Beginn der Consecration.

\quellseite{228}
\section{Zum Offertorium}

Chor. Hommels Liturgie.

\notensatz{music/228-01}{Zum Offertorium, Hommels Liturgie, vierstimmig. HErr, thue wohl an Zion nach deiner Gnade, baue die Mauern zu Jerusalem. Dann werden dir gefallen die}

\notensatz{music/228-02}{Offertorium, Fortsetzung. Opfer der Gerechtigkeit, die Brandopfer und ganzen Opfer.}

\notensatz{music/228-03}{Offertorium, Schluß. Dann wird man Farren auf deinem Altar opfern.}


\quellseite{229}
\chapter{XIII. Präfation mit Sanctus}


\section{1. Die Präfation}

Dem Acte der Consecration geht ein feierliches Dank- und Preisgebet voraus, das hernach der ganzen nun folgenden Communionhandlung den Namen Eucharistie oder Danksagung gegeben hat. Es ist dies die Präfation (Vorrede, Vorgang sc. Dankeinleitung). Zwar verstand man in der alten Kirche unter Präfation sämmtliche der Consecration vorhergehende Gebete. Der Mittelpunkt aber war das Lob- und Dankgebet, mit dem erwecklichen „Sursum corda!“ (die Herzen in die Höhe!) eingeleitet, mit dem „Trisagion“, dem Dreimalheilig oder dem seraphischen Hymnus schließend.

Das im I. Theil, § 2, Anmerk. 1 beigebrachte Zeugniß aus Justin und den „apostolischen Constitutionen“ hat bereits auf das hohe Alter der Präfation hingewiesen. Aus einer der ältesten aufgezeichneten Liturgien, genannt die „Liturgie des heiligen Jacobus“, nachweislich aber noch im 4. Jahrhundert in Jerusalem gebraucht, stehe sie hier in einer ihrer frühesten morgenländischen Formen, ursprünglich griechisch verfaßt:

\begin{quote}
„Hierauf spricht er (der Priester) laut: Die Liebe des HEerrn und Vaters, die Gnade des HEerrn und Sohnes und die Gemeinschaft und Gabe des Heiligen Geistes sei mit uns allen.

Das Volk: Und mit deinem Geiste.

Der Priester: Empor laßt uns Sinn und Herzen heben.

Das Volk: Würdig und recht ist es.

Hierauf betet der Priester: Wahrhaft würdig und recht ist es, geziemend und schuldig, dich zu loben, dir zu singen, dich zu preisen, dich anzubeten, dich zu verherrlichen, dir zu danken, dem Hervorbringer der ganzen Schöpfung, der sichtbaren sowohl wie der unsichtbaren, dem Schatze der ewigen Güter, dem Quell des Lebens und der Unsterblichkeit, dem Gott und Herrscher des Alls, den besingen die Himmel der Himmel und all ihr Heer, Sonne und Mond und der ganze Chor der Gestirne, Erde, Meer und Alles in ihnen, Jerusalems himmlische Versammlung, die Gemeinde der Erstgebornen, die in den Himmeln angeschrieben sind, die Geister der Gerechten und Propheten, die Seelen der Märtyrer und Apostel, Engel, Erzengel, Thronen, Herrschaften, Obrigkeiten und Vollmächten und furchtbare Kräfte und die vieläugigen Cherubim und sechsflügeligen Seraphim, welche mit zwei Schwingen ihre Angesichter
\quellseite{230}
bedecken, mit zweien aber die Füße, und mit zweien fliegend einer zum andern rufen mit unaufhörlichem Munde, in nie schweigenden Verherrlichungen.

(Laut gesprochen:) Den Siegesgesang deines herrlichen Ruhmes mit erhabener Stimme singend, rufend, lobpreisend, laut tönend und sprechend.

Das Volk: „Heiliger, heiliger, heiliger HEerr Zebaoth, voll ist der Himmel und die Erde deiner Herrlichkeit. Hosianna (du) in den höchsten Höhen, gepriesen ist, der da kommt im Namen des HEerrn; Hosianna (du) in den höchsten Höhen.“\footnote{Die Hauptliturgien der alten Kirche in wortgetreuer Uebersetzung nebst Einleitung von J. Ludwig König. 1865. S. 146. 147.}
\end{quote}

Aus der morgenländischen Kirche ging schon frühzeitig der Gebrauch der Präfation in die abendländische Kirche über. Dort erhielt sie jedoch nicht nur die einfachere Form, in der sie jetzt noch in der römischen Messe sich findet, sondern im 4. Jahrhundert fing man daselbst auch an, Fest-Präfationen zu formuliren, indem man der allgemeinen Form einen auf das jeweilige Fest bezüglichen Passus einfügte. Bei der Zunahme der Feste im Papstthum wuchs dann freilich auch die Zahl dieser besonderen Präfationen, propriae genannt, doch schon der Vorgänger Gregors I. reducirte dieselbe.

Auch in der lutherischen Kirche fand die Präfation alsbald Aufnahme im Hauptgottesdienst, und zwar nicht bloß die allgemeine (quotidiana), sondern auch die besondere, die für die wichtigsten Feste (propria). Zwar setzte Luther in seiner deutschen Messe 1526 an die Stelle der altkirchlichen Präfation die Vermahnung an die Communicanten mit vorausgehender Vaterunser-Paraphrase, und behielt allein das Sanctus bei, indem er das von ihm gedichtete und in Musik gesetzte „Jesaja dem Propheten“ etc. als das deutsche Sanctus anordnete und ihm seine Stelle nach der Consecration anwies. Dagegen aber hatte er sie in seiner „Weise, christlich Meß zu halten“ 1523 wenigstens theilweise, nämlich bis zum „per Christum Dominum nostrum“ (durch Christum unsern HEerrn) aufgenommen und nach älterem Vorgang daran die Consecration gefügt. Ein Theil der Kirchenordnungen richtete sich nach der Weise von 1526, ein anderer Theil nach der von 1523, doch also, daß man die ganze Präfation aufnahm, mit dieser aber auch die Abendmahlsvermahnung verband, wie hernach folgt.

Da man im Allgemeinen die Präfation der römischen Messe unverändert in den lutherischen Hauptgottesdienst aufnahm, so bediente man sich zum großen Theil hierbei auch noch der lateinischen Sprache, oder bestimmte die lateinische
\quellseite{231}
Präfation für die Städte und die deutsche für die Dörfer. Anderorts jedoch sang oder sprach man die Präfation gleich von vornherein deutsch. So nach den K.-O. A. Döbers zu Nürnberg 1525, Riga 1530, Braunschw.-Lüneburg 1542, Cöllnische Reformat. 1544 u. a. m.

Im 18. Jahrhundert fingen die Präfationen allmählich an, aus den Kirchenbüchern zu verschwinden; am längsten noch hielten sie sich in Sachsen. Kläglich war der Ersatz, den der Rationalismus vorschlug, wie z. B.: „Erhebe ein jeder von uns nach seiner Fassungskraft nun Sinn, Herz und Geist zur unsichtbaren Welt“ etc. Wo immer man jedoch in neuerer Zeit zur alten Gottesdienstordnung zurückkehrt, läßt man die Präfation nicht beiseite. So die anfänglich hier vielfach gebrauchte Agende von Löhe, so die Agende der Missouri-synode, so auch das Kirchenbuch des General-Councils. Bildet sie doch einen wesentlichen Bestandteil aller Liturgien der ersten Jahrhunderte und ist in der That ein Kleinod unter den gottesdienstlichen Gebeten, dessen Stellung zugleich so entsprechend ist. Wie es in Israel Gewohnheit war, Brod und Wein beim Mahle mit besonderer Danksagung zu genießen,\footnote{Das Brod mit beiden Händen bedeckend, spricht der jüdische Hausvater: „Gelobet seist du, HErr, unser Gott, König der Welt, welcher uns das Brod aus der Erde hervorgebracht hat.“ Und ähnlich segnet er den Wein.} und wir Christen Dank sagen sollen für alle Gabe dem Vater des Lichts und der Barmherzigkeit, so ist es wahrhaft würdig oder billig und recht, wie die alte Kirche sich ausdrückte, daß wir für die unaussprechliche, unter Brod und Wein uns dargereichte Gabe des Leibes und Blutes unsers HEerrn Jesu Christi Gott vorher schon danken und seinen Namen mit allen Engeln und Auserwählten loben und preisen.

In der Regel wurde und wird die Präfation vom Liturgen nicht gesprochen, sondern gesungen. Die ursprünglichen Weisen sind aber mit Recht schwungvoller als die der Collecten. Die hier nach Hommel für sämmtliche Präfationen wieder gegebenen Melodien finden sich im Directorium chori von Guidetti 1589 (Freund und Schüler von Palestrina). Die dem Verfasser vorliegende Melodie für eine Sonntagspräfation in der Ausgabe der Braunschw.-Lüneb. K.-O. 1657 ist der Guidetti’schen Weise des lateinischen Textes offenbar nachgebildet. Da die Melodie von Guidetti nicht nur schön, sondern auch leicht singbar ist, so geschieht es wohl, daß man bei der Rückkehr zur alten Liturgie sie jener gleichfalls schwungvollen, aber für manche etwas schwereren Weise vorzieht, die in einer Anzahl älterer Agenden mit unbedeutender Variation sich findet, und deren Quelle Spangenbergs Cantional: „Kirchengesänge lateinisch
\quellseite{232}
und deutsch. Magdeburg 1545“ ist. Aus der Schwarzburgischen Agenda von 1675 folge nachstehend als Beispiel die Singweise eine der Präfationen für die drei hohen Feste, nur daß dabei der Text der St. Louiser Agenda vom Herausgeber untergelegt worden ist, was sich sehr leicht thun ließ.

\section{2. Das Sanctus}

Unter der Hinweisung auf die Vereinigung mit den Lobgesängen der himmlischen Heerschaaren leitet die Präfation zu dem Sanctus über, dem Hymnus seraphicus aus Jesaia 6, 3. Seine Einführung in den Gottesdienst scheint schon der frühesten Zeit anzugehören, denn man schreibt dasselbe dem Ignatius († 116) zu. Bestandteil der römischen Messe soll es bereits durch Sixtus I. († 130) geworden sein.

Dem Sanctus wurde im 4. Jahrhundert noch das Benedictus („Gebenedeiet sei“ etc.) nebst dem Hosianna zugefügt. Eine sinnige Verbindung! Wie einst den im Fleisch erschienenen Gottessohn bei seinem Einzug in Jerusalem Israel mit dem Benedictus und Hosianna als den Davidssohn begrüßte, so begrüßt hier die feiernde Gemeinde ihn als den Mariensohn und das Passahlamm, der mit seinem Leib und Blut unter dem Brod und Wein wahrhaftig zu ihr, der Tochter Zion, kommt, um alle im Wort verkündigte und dargereichte Gnade zu besiegeln und sich mit ihr auf das innigste zu vereinigen.

Noch bis zur Zeit Gregors I. und Karls des Großen, ja wohl drüber hinaus, sang das Volk das ganze (lateinische) Sanctus mit den Clerikern. Allmählich jedoch hörte es auf, Gemeindegesang zu sein. Die Reformation fand es allein als Chorgesang.

Wo immer man durch Aufnahme der Präfation in den Hauptgottesdienst der römischen Messe sich anschloß, wie z. B. in den Kirchenordnungen von Straßburg 1525, Erfurt 1526, Wittenberg 1559, Mecklenburg 1552, Frankfurt a. M. 1565 etc., pflegte man dieselbe mit dem Sanctus abzuschließen. Wo man aber nach Vorgang von Luthers deutscher Messe 1526 die Vermahnung an die Stelle der Präfation treten ließ, folgte das Sanctus unmittelbar auf die Consecration, namentlich das deutsche Sanctus „Jesaja dem Propheten“ etc. Wie verschieden hierbei der Vortrag war, und welche besondere Auszeichnung das dreimalige Heilig meist erfuhr, zeigt Schöberlein aus folgender Zusammenstellung aus verschiedenen Kirchenordnungen und Cantionalen. So heißt es: „Den Vers, Heilig, soll die Gemeinde dreimal langsam, mit
\quellseite{233}
Andacht und großer Ehrerbietung singen“\footnote{Vom Herausgeber unterstrichen, da er das Sanctus von so manchem Organisten hier zu Lande schon in so raschem Tempo spielen hörte, daß es einem nicht zu Muthe sein konnte, als sollte man jetzt mit den Seraphinen den dreimal Heiligen anbeten.} (Pommern 1563). Desgleichen: „In Städten aber, da Schulen sind, sollen drei Knaben den Vers langsam, vor dem Altar knieend, singen, und das Chor mit der Gemeine denselbigen antworten (oder: „es soll von den Knaben zweimal mit lauter Stimme gesungen werden“) (Ludecus, Missale 1589). Wenn die Knaben den Vers zum dritten Mal gesungen haben, dann beschließt das Chor mit dem folgenden Vers: „Sein Ehr die ganze Welt“ etc. (Pomm.) Wohl suchte man auch in das dreimalige „Heilig ist unser Gott, der HEerre Zebaoth“, dadurch noch eine Steigerung zu legen, daß man dasselbe zuerst zweistimmig, dann dreistimmig und das dritte Mal vierstimmig von einzelnen Stimmen singen ließ, während der übrige Theil des Sanctus vom ganzen Chor oder der Gemeinde gesungen wurde (Mich. Prätorius Missodia Sionia 1606); oder daß man das „Heilig ist“ etc. zuerst in Choralnoten, dann vier-, fünf- und sechsstimmig und hierauf das „Sein Ehr“ etc. siebenstimmig sang, mit der letzten Strophe aber („Von dem Geschrei“ etc.) zum einfacheren Tonsatz wieder zurückkehrte. (M. Vulpius, Kirchengesänge 1604.)

\section{Einleitung zu allen Präfationen}

\notensatz{music/233-01}{Pastor: Der HErr sei mit euch!}
\notensatz{music/233-02}{Gemeinde: Und mit deinem Geiste.}

\notensatz{music/233-03}{Pastor: Die Herzen in die Höhe!}
\notensatz{music/233-04}{Gemeinde: Erheben wir zum HErrn.}

\quellseite{234}
\notensatz{music/234-01}{Pastor: Laßset uns Dank sagen dem HErren, unserm Gotte!}
\notensatz{music/234-02}{Gemeinde: Das ist würdig und recht.}

\section[A. Die gemeine Präfation]{A. Die gemeine Präfation\footnote{Zum Vergleich stehe hier ein Bruchstück der quodlitiana aus dem Directorium chori von Guidetti 1589:
\notensatz{music/234-04}{Lateinische Präfation nach Guidetti 1589: Vere dignum et justum est … aeterne Deus etc.}
}}

\subsubsection{An gewöhnlichen Sonntagen zu gebrauchen}

\notensatz{music/234-03}{Gemeine Präfation, Anfang: Wahrhaft würdig und recht … fürchten.}

\quellseite{235}
\notensatz{music/235-01}{Gemeine Präfation, Schluß: die Mächte … und anbetend zu dir sprechen.}

\section{B. Die besonderen Präfationen}

\subsection{Am Christtage}

\notensatz{music/235-02}{Präfation am Christtage, Anfang: Wahrhaft würdig und recht … mit den Thronen und.}

\quellseite{236}
\notensatz{music/236-01}{Präfation am Christtage, Schluß: Herrschaften … und sprechen ohn Ende.}

\subsection[Am Epiphaniasfest]{Am Epiphaniasfest oder in der Epiphaniaszeit}

\notensatz{music/236-02}{Präfation am Epiphaniasfest: Wahrhaft würdig und recht … und sprechen ohn Ende.}

\quellseite{237}
\subsection{In der Passionszeit}

\notensatz{music/237-01}{Präfation in der Passionszeit: Wahrhaft würdig und recht … und anbetend zu dir sprechen.}

\quellseite{238}
\subsection{Am Osterfeste}

\notensatz{music/238-01}{Präfation am Osterfeste: Wahrhaft würdig und recht … und sprechen ohn Ende.}

\subsection{Am Himmelfahrtsfeste}

\notensatz{music/238-02}{Präfation am Himmelfahrtsfeste, Anfang: Wahrhaft würdig und recht … daß wir dir.}

\quellseite{239}
\notensatz{music/239-01}{Präfation am Himmelfahrtsfeste, Schluß: HErr, heiliger, allmächtiger Vater … und sprechen ohn Ende.}

\subsection{Am Pfingstfeste}

\notensatz{music/239-02}{Präfation am Pfingstfeste, Anfang: Wahrhaft würdig und recht … allezeit und.}

\quellseite{240}
\notensatz{music/240-01}{Präfation am Pfingstfeste, Schluß: allenthalben danksagen … und sprechen ohn Ende.}

\subsection{Am Trinitatisfeste}

\notensatz{music/240-02}{Präfation am Trinitatisfeste, Anfang: Wahrhaft würdig und recht … und dem Heiligen Geist Ein.}

\quellseite{241}
\notensatz{music/241-01}{Präfation am Trinitatisfeste, Schluß: Gott und Ein HErr bist … und anbetend zu dir sprechen.}

\section[Eine andere Weise]{Eine andere Weise, die Präfation zu singen}
\quellseite{242}
\subsection{An Weihnachten}

Spangenberg 1545. Mansfeld 1580. Schwarzburg 1675.

\notensatz{music/242-01}{Andere Weise der Präfation an Weihnachten: Wahrhaft würdig und recht … und sprechen ohn Ende.}

\section{Sanctus}
\quellseite{243}
\subsection{I. Sanctus}
Mel.: L. Lossius' Psalmodia 1851. Tonsatz: F. Jubitz u. G. v. Tucher.

\notensatz{music/243-01}{Sanctus I: Heilig! Heilig! Heilig ist Gott.}

\notensatz{music/243-02}{Sanctus I: der HErr Zebaoth. Voll sind Him-.}

\notensatz{music/243-03}{Sanctus I: -mel und Erdreich deiner Ehren. Hosi-.}

\notensatz{music/243-04}{Sanctus I: -anna in der Höhe! Gebenedeit.}

\quellseite{244}
\notensatz{music/244-01}{Sanctus I: sei Marien Sohn, der da kommt im Namen des HErrn. Abweichende Unterlage: das Passahlamm, das.}

\notensatz{music/244-02}{Sanctus I: Hosianna in der Höhe.}

\subsection{II. Sanctus minus summum}
\subsubsection{Chor}
Mel.: L. Lossius 1553. Wolffg. Pfalzg. K.-O. 1557. Ludecus 1589. Harm.: F. Kiegel 1865.

\notensatz{music/244-03}{Sanctus II, Chor: Heilig, heilig, heilig bist du. Mit der ersten alternativen Unterlage.}

\notensatz{music/244-04}{Sanctus II, Chor und Gemeinde: HErr Gott Zebaoth! Voll sind die Himmel und Erde von.}

\quellseite{245}
\notensatz{music/245-01}{Sanctus II, Gemeinde und Chor: deiner Ehre! Hosianna.}

\notensatz{music/245-02}{Sanctus II, Chor und Gemeinde: in der Höhe! Gelobet sei, der da kommt, im. Zweite Unterlage: Gebenedeit.}

\notensatz{music/245-03}{Sanctus II, Gemeinde, dann Chor und Gemeinde: Namen des HErren! Hosi-.}

\notensatz{music/245-04}{Sanctus II, Chor und Gemeinde: -anna in der Höhe!}

Die im vorstehenden Sanctus unter die Baßnoten von dem Herausgeber gesetzten Worte sollen zeigen, wie man die unter uns gebräuchliche Fassung des Sanctus anpassen kann.

\quellseite{246}
\subsection{III. Sanctus}
\subsubsection{Knaben}
Mel.: Luther. Aus: Jesaia, dem Propheten etc.

\notensatz{music/246-iii-knaben}{Sanctus III, Knaben: Heilig ist Gott, der HErre Zebaoth! Heilig ist.}

\subsubsection{Gemeinde}

\notensatz{music/246-iii-gemeinde-a}{Sanctus III, Gemeinde: Gott, der HErre Zebaoth! Heilig ist Gott, der.}

\notensatz{music/246-iii-gemeinde-b}{Sanctus III, Gemeinde: HErre Zebaoth. Sein Ehr die ganze Welt erfüllet hat.}

\subsection{IV. Sanctus}
Mel.: Braunschw. Wolfenb. A.-C. 1543. Pomm. K.-O. 1563. Harm.: M. Prätorius 1606.

\notensatz{music/246-01}{Sanctus IV: Heilig ist Gott, der Vater, heilig ist Gott, der Sohn.}

\quellseite{247}
\notensatz{music/247-01}{Sanctus IV: heilig ist Gott, der Heilige Geist. Er ist der HErre.}

\notensatz{music/247-02}{Sanctus IV: Zebaoth. Alle Welt ist seiner Ehre voll. Hosianna.}

\notensatz{music/247-03}{Sanctus IV: in der Höhe! Gelobet sei, der da kommt im Na-.}

\notensatz{music/247-04}{Sanctus IV: -men des HErren! Hosianna in der Höhe!}

\quellseite{248}
\subsection{V. Sanctus}
Unbekannt.

\notensatz{music/248-01}{Sanctus V, Feierlich. Gemeinde: Heilig, heilig, heilig ist Gott, der.}

\notensatz{music/248-02}{Sanctus V, Gemeinde: HErre Zebaoth! Voll sind Himmel und Erdreich deiner.}

\notensatz{music/248-03}{Sanctus V: Ehre. Drei Hosianna für Knaben oder Einzelne, Chor und Gemeinde.}

\notensatz{music/248-04}{Sanctus V: in der Höhe! Gebenedeit, Solo und Gemeinde.}

\quellseite{249}
\notensatz{music/249-01}{Sanctus V: sei Mariens Sohn, der da kommt im Namen des HErrn. Mit der Unterlage sei das Passahlamm, das, und Hosianna der Knaben.}

\notensatz{music/249-02}{Sanctus V, Chor und Gemeinde: Hosianna in der Höhe!}

\chapter{XIV. Abendmahlsvermahnung}

Die Abendmahlsvermahnung hat seinen Vorgang in der römischen Messe. Sie ist ein eigentliches Erzeugniß der Reformation. Heißt es doch auch im Artikel von der Messe A. C. Art. 24: „So werden auch die Leute mit höchstem Fleiß zum öfternmal unterrichtet vom heiligen Sacrament, wozu es eingesetzt und wie es zu gebrauchen.“ Dazu sollte nun auch eine besondere Vermahnung vor dem Abendmahl dienen.

Zur Einfügung dieser Vermahnung in die Abendmahlsliturgie machte Luther seinerseits den Anfang eigentlich schon im Jahre 1525. In einem Briefe an Nikolaus Hausmann vom 26. März, mit welchem er einige der ihm mitgetheilten Kirchengesänge zurücksendet, legt er eine von ihm verfaßte Abendmahlsvermahnung bei, mit der Ueberschrift: „Die Präfaz.“ Im Jahre 1524 erschien auch in Nürnberg in der deutschen Meßordnung der Pröbste zu St. Sebald und St. Lorenz, Preßler und Bömer, und 1525 in der des Kaplans am neuen Spital, Andreas Döber, eine Abendmahlsvermahnung.

\quellseite{250}
(Es ist dies die bekannte W. Volprecht’sche\footnote{Wolfgang Volprecht, Prior des Augustinerklosters zu Nürnberg, ein Freund Staupitzens und Luthers, hatte schon 1518 einen Tractat Luthers gegen den Ablaß in Nürnberg drucken lassen und gehört überhaupt zu den redlichsten und rückhaltlosesten Freunden der Reformation. Er war auch der erste, welcher, dem Verlangen der Nürnberger in einem gewissen Maße nachgebend, bereits im Jahre 1523 seinen bekehrten Mönchen und wenigen Bürgern das Abendmahl unter beider Gestalt austheilte. Er starb 1528, von Luther sehr betrauert.} Vermahnung in ihrer, wie es scheint, ursprünglichen Gestalt. Dieselbe beginnt bekanntlich mit den Worten: „Ihr Allerliebsten in Gott! Dieweil wir das heilige Abendmahl unsers HEerrn Jesu Christi wollen bedenken und halten“ etc. Bekanntlich setzte dann auch Luther in seiner deutschen Messe 1526 an die Stelle der Präfation die Vaterunser-Paraphrase mit ausgehender Abendmahlsvermahnung, auf welche dann unmittelbar die Consecration folgt. Wegen ihrer Schönheit und Kürze, vielfach jedoch mit Weglassung der Vaterunser-Paraphrase, hat sie nebst der ebenso köstlichen, aber längeren Volprecht’schen Vermahnung am meisten Aufnahme gefunden.

Wie viel man auf eine derartige Abendmahlsvermahnung hielt, geht daraus hervor, daß es u. A. in der Pommerschen Kirchenordnung heißt: „Solche Vermahnungen sollen nicht leichtlich, so es die Zeit duldet, unterlassen, sondern viel lieber die Gesänge gekürzt oder eingestellt werden.“ Wie in der Braunschw.-Lüneb. Kirchenordnung und anderwärts waren mehrere Formulare zum beliebigen Gebrauch dargeboten. Ueberhaupt finden sich in den Kirchenordnungen des 16., 17. und zum Theil 18. Jahrhunderts viele und der Form nach bei aller Uebereinstimmung der Gedanken sehr verschiedene Abendmahlsvermahnungen.

Verschieden jedoch ist die Stellung, welche diese Vermahnung in der Liturgie einnimmt. Einige Kirchenordnungen, z. B. die Brandenb.-Nürnb., setzen sie an die Stelle der Präfation, oder vor die Präfation, z. B. Erfurt 1526 und Oesterreich 1571, oder nach der Präfation, nämlich zwischen dieser und der Consecration, z. B. Mecklenb. 1552, Wittenberg 1559, Pommern 1563, Braunschw.-Lüneb. 1594, Schwarzb. 1675, Leipzig 1681 etc., etliche auch nach der Consecration vor der Austheilung, z. B. Döber 1525, Mkt. Brandenb. 1540. Am häufigsten hat sie ihre Stellung zwischen Präfation und Consecration. Es dürfte wohl schwer zu entscheiden sein, welche Stellung den Vorzug verdiene — die Stellung unmittelbar vor oder nach der Präfation. Stellt man die Vermahnung der Präfation voran, so wird allerdings der schwungvolle Gang der Liturgie nicht durch die Prosa der Vermahnung unterbrochen. In Absicht auf Inhalt und Länge dürfte z. B. die Volprecht’sche Vermahnung, desgleichen auch in Absicht auf Inhalt und Form jene ältere
\quellseite{251}
Luther’sche Vermahnung nach unserm Dafürhalten ihre geeignete Stellung vor der Präfation finden. Wenn aber die meisten der Kirchenordnungen, welche die Präfation beibehalten haben, gerade die kurze Luther’sche Vermahnung nach der Präfation, das ist, zwischen dieser und der Consecration stellen, so geschieht das gewiß nicht nur, weil sie so kurz und dabei doch auch mehr schwungvoll gehalten ist, sondern auch, weil sie sonderlich hinweist auf „die Worte, darinnen uns Christus seinen Leib und Blut zu Vergebung schenket“, und durch ihren Schluß in directer Verbindung mit der Consecration steht. Der Herausgeber, der in seiner vorigen Gemeinde in dieser Stellung die Luther’sche Vermahnung vorstand und so sie lange Jahre brauchte, muß auch bekennen, daß er hierbei nichts von einer prosaischen Unterbrechung inne wurde. Im Gegentheil war es ihm und Andern immer ein ergreifender Moment, wenn nach dem Sanctus diese Vermahnung gelesen wurde und ein paar Accorde der Orgel darauf die liturgische Recitation des Vaterunsers einleitete.

Von einer freien Vermahnung in der Abendmahlsliturgie wollte man nichts wissen. Man ließ sich gesagt sein, was Luther von seiner Paraphrase und Vermahnung bemerkt: „Ich will aber gebeten haben, daß man dieselbige Paraphrasis und Vermahnung conceptis, seu praescriptis verbis, oder auf eine sonderliche Weise stelle, um des Volks willen, daß nicht heute einer also, der andere morgen anders stelle; und ein jeglicher seine Kunst beweise, das Volk irre zu machen, daß es nichts lernen noch behalten kann. Denn es ist ja um das Volk zu lehren und zu führen zu thun; darum ist es noth, daß man die Freiheit hie breche und einerlei Weise führe in solcher Paraphrasie und Vermahnung, sonderlich in einerlei Kirchen oder Gemeine für sich, ob sie einer andern nicht folgen wollen um ihrer Freiheit willen.“

Zur Kenntnißnahme und zu entsprechendem Gebrauch folgen dieses Orts einige der alten Vermahnungen. Es sind dies:

\section[Nr. 1. Luthers Vermahnung]{Nr. 1. Luthers Vermahnung in der deutschen Messe 1526}
Wo dieselbe an die Stelle der Präfation tritt, steht sie immer in Verbindung mit der Vaterunser-Paraphrase, wie dies denn auch sehr passend in dem Frühcommunion-Formular der St. Louiser Agende geschieht. Wo sie mit der Präfation verbunden ist, erscheint sie immer getrennt von der Paraphrase. In dieser Gestalt ist sie hier gegeben.

\section[Nr. 2. Luthers andere Vermahnung]{Nr. 2. Luthers andere Vermahnung aus seinem Briefe an Nicol. Hausmann 1525}
Wiewohl dieselbe unseres Wissens nicht weiter verbreitet ist, sondern sich nur in der Kirchenordnung von Coburg 1626 und Gotha findet,
\quellseite{252}
so stehe sie doch hier als ein deutsch-lutherischer Wiederhall des altkirchlichen „das Heilige den Heiligen“. Wollte man sich ihrer mitunter bedienen, so könnte man entweder das im Einschluß beigegebene oder das mit der Volprecht’schen Vermahnung verbundene Gebet folgen lassen.

\section[Nr. 3. Wolfgang Volprechts Vermahnung]{Nr. 3. Wolfgang Volprechts Vermahnung 1524}
Ursprünglich reicht dieselbe jedoch bis zu dem von uns gemachten Zeichen. Die nach diesem folgende Stelle ist nach Löhe in der Kirchenordnung von Baden 1556 angefügt. Wegen seiner Schönheit fand der Zusatz denn auch vielfach Aufnahme. Ob das der Vermahnung auch in der St. Louiser Agenda folgende schöne und passende Gebet jener baden’schen Kirchenordnung gleichfalls entstammt, vermag der Herausgeber zur Zeit nicht zu sagen. Wiewohl in der St. Louiser Agenda, dem „Kirchenbuch des General Council“ und andern Agenden befindlich, ist sie doch hier der Vollständigkeit wegen aufgenommen.

\section[Nr. 4. Vermahnung aus der Hamb. Agende]{Nr. 4. Vermahnung aus der Hamb. Agende 1538 und 1788}
Dieselbe war viel in Norddeutschland und zum Theil in Mitteldeutschland verbreitet. Bemerkenswerth ist, daß hier auf die Vermahnung der Gesang des „Heilig ist Gott, der HEerre Zebaoth“, aus dem Luther-Sanctus: „Jesaja, dem Propheten“ etc. und nach diesem die Consecration folgt. Es könnte dies als ein Wink betrachtet werden, da, wo die Präfation überhaupt wegfällt und die Volprecht’sche oder eine andere Vermahnung anstatt der Luther’schen von 1526 gebraucht wird, nach dieser und vor der Consecration den Gesang des ganzen deutschen Sanctus (Jesaja, dem Propheten) folgen zu lassen. Es wäre dies dann für denselben wohl die geeignetste Stelle, falls man nicht es während der Austheilung singen will, wie es Luther in der deutschen Messe anordnet.

\section[5. Vermahnung und Gebet]{5. Vermahnung und Gebet aus der Braunschw.-Wolfenb. A.-O. 1657}
Zu derselben findet sich in dieser Kirchenordnung folgende Bemerkung: „Weil an vielen Orten, löblich hergebracht, daß entweder vor oder nach der Präfation die Gemeinde und insonderheit diejenigen, so zum Sacrament gehen wollen, durch eine deutsche Ermahnung zur Andacht aufgemuntert werden; wir auch dafür halten, daß zu mehrerer Andacht diene, daß dasjenige, wozu die Gemeinde vermahnt wird, im Werk und effectu an ihm selbst vielmehr geschehe, als daß man der Vermahnung zu solchem Werk noch lange inhärire; so haben wir zu dem Behuf eine Formulam eines Gebets, so der Gemeinde öffentlich anstatt der Vermahnung vorgelesen werden soll, abfassen lassen.“ Man wird auf den ersten Blick erkennen, daß diese besondere Form Nachbildung der Volprecht’schen Vermahnung ist.

\quellseite{253}
\section{Verschiedene Formulare}
\subsection{1}
Lieben Freunde Christi, weil wir hie versammelt sind in dem Namen des HErrn, sein heiliges Testament zu empfahen, so vermahne ich euch in Christo, daß ihr mit rechtem Glauben des Testaments Christi wahrnehmet und allermeist die Worte, darinnen uns Christus seinen Leib und Blut zur Vergebung schenket, im Herzen feste fasset, daß ihr gedenket und danket der grundlosen Liebe, die er uns bewiesen hat, da er uns durch sein Blut von Gottes Zorn, Sünde, Tod und Hölle erlöset hat, und darauf das Brod und Wein, das ist, seinen Leib und Blut, zur Sicherung und Pfand zu euch nehmet. Demnach wollen wir in seinem Namen und auf seinen Befehl durch seine eigene Worte das Testament also handeln und gebrauchen:

\subsection{2}
Allerliebsten Freunde in Christo. Ihr wisset, daß unser HErr Jesus Christus aus unaussprechlicher Liebe dies sein Abendmahl zur Letze hat eingesetzt zum Gedächtniß und Verkündigung seines Todes, für unsere Sünden erlitten, zu welchem Gedächtniß gehöret ein fester Glaube, der eines jeglichen Gewissen und Herz, der sein brauchen und genießen will, sicher und gewisser machen will, daß also der Tod für alle seine Sünde von Christo erlitten sei.

Wo aber Jemand daran zweifelt und solchen Glauben nicht etlichermaßen bei ihm fühlet, der soll wissen, daß ihm dies Abendmahl nicht nütze, sondern schädlich sei, und er soll davon bleiben. Welchen Glauben, weil wir ihn nicht sehen, und alleine Gott bewußt ist, wollen wir einem jeglichen, so herzu geht, auf sein Gewissen gestellet haben, und auf sein Bitten und Begehren zulassen.

Welche aber noch in öffentlichen Sünden stecken, als: Geiz, Haß, Zorn, Neid, Wucher, Unkeuschheit und dergleichen, und nicht abzulassen gedenken, denen sei hiermit abgesagt, und warnen sie treulich, daß sie nicht ein Gericht und Schaden über ihre Seele holen, wie St. Paulus saget. Jedoch, so jemand gefallen, aus Gebrechlichkeit, und sich zu bessern ernstlich beweiset mit der That, soll ihm solche Gnade und Gemeinschaft des Leibes und Blutes Christi nicht versagt sein. Darnach sich habe und wisse ein jeglicher zu richten und sehe sich vor; denn unser Gott lässet sein nicht spotten, so will er auch nicht das Heiligthum den Hunden geben, noch die Perlen vor die Säue werfen lassen. Amen.

(Derhalben laßt uns beten: Allmächtiger Gott, himmlischer Vater, sintemal wir dir nicht, denn allein in deinem geliebten Sohn, unserm HErrn, wohlgefallen
\quellseite{254}
mögen, so heilige unsern Leib und unsere Seele und gib uns seine selige Gemeinschaft in dem heiligen Abendmahl mit rechtgläubiger Begierde und Dankbarkeit zu empfahen, daß wir, deiner ewigen Güte und Liebe gegen uns abermals getröstet und in dem neuen Leben gestärket, dir zum Preis deines heiligen Namens und zur Besserung deines Volkes mit mehr Fleiß und Furcht leben und dienen mögen. Durch denselben unsern HErrn JEsum Christum. Amen.)


\section{3}

Ihr Allerliebsten in Gott! Dieweil wir das heilige Abendmahl unsers HErrn JEsu Christi wollen bedenken und halten, darin Er uns sein Fleisch zu einer Speise und sein Blut zu einem Trank, den Glauben damit zu stärken, gegeben hat, sollen wir billig mit großem Fleiß und inbrünstiger Andacht ein jeder sich prüfen, wie St. Paulus sagt, und alsdann von diesem Brode essen und von diesem Kelche trinken. Denn dies heilige Sacrament ist zu einem besonderen Trost und Stärkung gegeben den armen betrübten Gewissen, die ihre Sünden bekennen, Gottes Zorn und den Tod fürchten und nach der Gerechtigkeit hungrig und durstig sind. So wir aber uns selbst prüfen und ein jeder in sein eigen Gewissen geht, wie uns der heilige Paulus lehrt, werden wir gewiß nichts anders finden, denn allerlei greuliche Sünde und den Tod, den wir mit der Sünde verschuldet haben, und können uns doch selbst in keinem Wege daraus helfen. — Darum hat unser lieber HErr JEsus Christus sich über uns erbarmt und ist um unsrer Sünde willen Mensch worden, auf daß Er das Gesetz und allen Willen Gottes für uns und uns zu gut erfüllete und den Tod und alles, was wir mit unsren Sünden verschuldet haben, für uns und zu unsrer Erledigung auf sich nehme und erlitte. — Und daß wir das ja festiglich glauben und durch den Glauben fröhlich in seinem Willen möchten leben, nahm Er nach dem Abendmahl das Brod, sagte Dank, brach’s und sprach: „Nehmet hin und esset, das ist mein Leib, der für euch dargegeben wird“, das ist, daß ich Mensch geworden bin und alles, was ich thue und leide, ist alles euer eigen, für euch und euch zu gute geschehen. Deß zu einem gewissen Anzeichen und Zeugniß gebe ich euch meinen Leib zur Speise. — Desgleichen nahm er auch den Kelch und sprach: „Nehmet hin und trinket aus diesem alle, das ist der Kelch des Neuen Testaments in meinem Blute, das für euch und für viele vergossen wird zur Vergebung der Sünde; so oft ihr das thut, sollt ihr mein dabei gedenken“, das ist, dieweil ich mich euer angenommen und eure Sünden auf mich geladen habe, will ich mich selbst für die Sünden in den Tod opfern, mein
\quellseite{255}
Blut vergießen, Gnade und Vergebung der Sünden erwerben und also ein neu Testament aufrichten, darin die Sünde vergeben und ihrer ewig nicht mehr gedacht werden soll. Deß zu einem gewissen Anzeichen und Zeugniß gebe ich euch mein Blut zu trinken.

Wer nun also von diesem Brod isset und von diesem Kelch trinket, das ist, wer diesen Worten, die er von Christo höret, auch diesen Zeichen, die er von Christo empfängt, festiglich glaubt, der bleibt in dem HEerrn Christo und Christus in ihm, und wird ewiglich leben. Dabei sollen wir nun sein gedenken und seinen Tod verkündigen, nämlich daß Er für unsre Sünde sei gestorben und zu unsrer Rechtfertigung wieder auferstanden, und Ihm darum Dank sagen, ein jeder sein Kreuz auf sich nehmen und Ihm nachfolge, und nach seinem Gebot einander lieben, wie er uns geliebt hat.\footnote{So weit die aus dem Jahre 1524 stammende und weitverbreitete Vermahnung des Nürnberger Abtes Volprecht. Das von dem Zeichen an Folgende fügte die K.-O. von Baden 1556 hinzu, während die Brandenb.-Nürnb. K.-O. 1592 und 1639 folgenden Zusatz hat: „Und dieweil wir Alle gesündiget haben, so demüthiget eure Herzen vor Gott dem Allmächtigen, bekennet eure Sünde und Gebrechen mit herzlicher Liebe und Begierde seiner göttlichen Gnade und Hülfe, mit festem Vertrauen auf sein gnädiges Zusagen, und vergebet von Herzen ein Jeder seinem Nächsten, auf daß euch euer himmlischer Vater eure Sünden und Mißethaten auch vergebe. So ihr solches thut, so entbinde ich euch nochmals anstatt der h. christlichen Kirche und aus Befehl und Zusagen unsers HEerrn Jesu Christi, da er sprach: Wenn ihr die Sünden vergebet, dem sind sie vergeben, von allen euren Sünden im Namen Gottes des Vaters und des Sohnes und des Heiligen Geistes. Amen.“} Denn wir alle sind ein Brod und ein Leib, dieweil wir alle eines Brodes theilhaftig sind und aus einem Kelche trinken. Denn gleichwie aus vielen Beerlein zusammen gekeltert ein Wein und ein Trank fleußt und sich in einander mengt, und aus vielen Körnlein ein Mehl gemahlen, ein Brod und Kuchen gebacken wird: also sollen wir alle, die Christo durch den Glauben eingeleibt sind, durch brüderliche Liebe um Christi, unsers Heilandes, willen alle ein Leib und Brod und Trank werden und solches gegen einander nicht allein mit Worten, sondern mit der That und Wahrheit, wie St. Johannes lehret (1. Br. 3), ohne allen Trug treulich gegen einander beweisen. Das helfe uns der allmächtige, barmherzige Gott und Vater unsers lieben HEerrn Jesu Christi durch seinen Heiligen Geist! Amen.

Solches Alles zu erlangen, lasset uns beten:

O mein HEerr Jesu Christe, du ewiger und allmächtiger Sohn Gottes, der du nach Willen und beschlossenem Rath deines himmlischen Vaters durch ein einiges, heiliges, vollkommenes Opfer deines Leibes und Blutes in Ewigkeit
\quellseite{256}
vollendet hast, die geheiligt werden, auch solchen deinen wahren kostbaren Leib und Blut in dem heiligen Abendmahl uns zu essen und zu trinken verordnet hast: mit was ernster und herzlicher Reu und Leid über die Sünde, mit was starkgläubiger Begierde und Ehrerbietung soll man billig dein heilig Abendmahl begehen und genießen! Nun sind aber wir Menschen aus uns selbst anklebender Sünde halben viel zu gering, schwach und unwürdig, solches deinem Willen und Befehl nach mit Nutzen im heiligen Schmuck zu begehen. Heiliger demnach, getreuer HEerr und Heiland, du selbst unsre Seel und Leib durch deinen Heiligen Geist, und mache uns also bereit und geschickt, würdiglich zu dieser himmlischen Wohlthat zu kommen. Auch was uns an rechtschaffener Buße, Reu und Leid über die Sünde, an festem, kindlichem Vertrauen und Glauben auf dein Verdienst und ernstlichem wahrhaftigem Vorsatz, das Leben zu bessern, in dieser Schwachheit mangelt, das erstatte und erfülle du gnädiglich mit dem Reichthum und Verdienst deiner bittern Schmerzen, deiner Noth und deines Todes, die du am Stamm des Kreuzes für uns alle erlitten und ausgestanden hast, auf daß also wir, die wir jetzt noch in der Welt auf dem Wege unsrer Pilgerfahrt dich, unsern einigen Trost und Seligmacher, verkleidet in dem heiligen Sacrament, begehren zu genießen, dich dermaleinst dort auch in dem ewigen, rechten Vaterland von Angesicht zu Angesicht klärlich anschauen und vor dir stehen mögen, der du mit Gott dem Vater und Heiligen Geist lebst und regierest, wahrer Gott, hochgelobet in Ewigkeit! Amen.

\section{4}

Meine Allerliebsten! Uns wird stets durch die Predigt des Evangelii vorgehalten, daß wir von uns selbst unwissend, arme Sünder und verloren sind. Und dieweil wir nichts mehr von uns selbst sind, denn Fleisch und Blut, derowegen wir uns auch mit unserm Verstande und Vermögen nicht können losmachen aus dem gestrengen Gerichte Gottes und von der Gewalt des Teufels, darein wir gefallen sind durch die Uebertretung der Gebote und des Willens Gottes: so hat Gott unser Unvermögen besser erkannt denn wir, und hat für uns gegeben, als ein gnädiger Vater, seinen eingebornen Sohn, Jesum Christum, daß wir durch seinen Tod erlöset würden von unseren Sünden, und durch ihn Kinder Gottes würden, ewig selig, so wir das glauben. Solches läßt er uns stets predigen; wer das glaubet, der hat gewiß das ewige Leben. Auf solchen Glauben und zu solcher Seligkeit werden wir auch getauft; da sollen wir stets in bleiben, so bleiben wir in Christo, und Christus in uns; so essen
\quellseite{257}
wir stets ohn Unterlaß geistlich mit dem Glauben den Leib Christi, und trinken sein Blut. Das ist: Wir werden Christo einverleibt, daß wir eins mit ihm werden; damit, daß wir glauben, daß er seinen Leib für uns in den Tod gegeben hat, und sein Blut für uns am Kreuze vergossen. Darauf verlassen wir uns zur Seligkeit, wider alle falsche Lehre, alle Sünde, Anfechtung und Noth. Aus welcher Wohlthat Christi wir auch lernen, welche Liebe und Geduld wir üben sollen gegen unsern Nächsten, auch gegen unsere Feinde; was wollen wir mehr? Doch daß wir nicht vergessen, oder träge würden, als wir leider werden, zu solchem Glauben der Menschwerdung und des Todes Christi, hat er uns ein besonderes Gedächtniß oder Verkündigung seines Todes, so oft wir wollen, befohlen; daß wir auch im auswendigen Sacrament, der Vernunft verborgen, allein dem Glauben aus den Worten Christi bekannt, essen sollen seinen Leib, und trinken sein Blut, daß wir ja nicht zweifeln sollen, sein Tod und Blutvergießen am Kreuze sei unsre gewisse Seligkeit. Davon sollen wir singen, lesen, predigen hören, gleichwie wir (in dieser unserer Gemeinde) jetzund thun, und nachmals auch davon reden und unter einander verkündigen, uns zum Troste und vielen zur Seligkeit, nach dem Befehl Christi: Solches thut zu meinem Gedächtniß! Wer nun würdig will essen und trinken dies Sacrament, der soll zwei Dinge thun: Er soll gläuben, was Christus saget, und thun, was er gebeut. Er sagt: Das ist mein Leib, der für euch gegeben wird! Das ist mein Blut, das für euch vergossen wird zur Vergebung der Sünden. Solches sollet ihr gläuben. Er gebeut aber: Nehmet hin und esset, trinket alle daraus, und gedenket mein! Solches sollt ihr thun, nach seinem Gnadenworte und Befehle, Amen!

\section{5}

(Wenn solches — der Gesang der Präfation und des Sanctus — geschehen, kehre sich der Priester gegen die Communicanten und die ganze Gemeine mit folgenden Worten:)

\begin{quote}
Meine Allerliebsten in Gott, dieweil wir das Abendmahl unsers HErrn JEsu Christi wollen halten, darin uns sein Fleisch zu einer Speise und sein Blut zu einem Trank gegeben wird, daß wir dabei seines Todes gedenken und ihm von Herzen Dank sagen sollen, so wollen wir Gott, den Vater, im Namen JEsu Christi anrufen und von Grund des Herzens also beten:
\end{quote}

(Die bete der Priester mit heller Stimme, verständlich und langsam, daß es die ganze Gemeine hören und ein jeder bei ihm selbst mit Andacht nachsprechen möge.)

\begin{quote}
HErr Gott, himmlischer Vater, wir danken dir von Herzen, daß du deinen einigen lieben Sohn, unsern HErrn und Heiland JEsum Christum, für unsere
\quellseite{258}
Sünde zu unserer Seligkeit hat Mensch werden, so viele unaussprechliche Marter leiden und am Kreuz eines bittern Todes sterben lassen, welcher uns auch befohlen hat, damit wir dessen, was er für uns ausgestanden, nicht vergessen möchten, diese gegenwärtige Gedächtniß seiner Marter und Todes im heiligen Abendmahl zu halten, in welchem er uns aus unaussprechlicher Liebe zu uns eben denselben Leib, welcher am Kreuz für uns gestorben, und eben dasselbe Blut, welches am Kreuz für uns vergossen, in und mit gegenwärtigem Brod und Wein wahrhaftig zu essen und zu trinken gibt, damit wir ja nicht zweifeln, sondern desto fester glauben möchten, sein Tod und Blutvergießen am Kreuz sei unsere gewisse Seligkeit.

Durch denselbigen deinen lieben Sohn als den einzigen Mittler zwischen dir, dem gerechten Gott, und uns armen Sündern, und dessen blutiges Verdienst bitten wir dich, o himmlischer Vater, du wollest dich unser aller erbarmen. Siehe an das Angesicht solches unsers getreuen Mittlers, deines lieben Sohnes Jesu Christi, welcher zu deiner Rechten für uns bittet und uns vertritt; siehe an sein heiliges unschuldiges Blutvergießen, Leiden und Sterben, welches er für unsere Missethat aus unaussprechlicher Liebe gegen uns verlorne Menschen mit höchster Geduld ausgestanden: Vergib uns unsere Sünde und gib uns deinen Heiligen Geist, daß wir verleugnen das ungöttliche Wesen und die weltlichen Lüste, hingegen aber mäßig gegen uns selbst, gerecht gegen unsern Nächsten und in steter kindlicher Furcht gegen dich, Gott, unsern himmlischen Vater, leben in dieser Welt und warten auf die selige Hoffnung und Erscheinung der Herrlichkeit des großen Gottes, deines Sohnes, unsers Heilandes Jesu Christi. Laß auch einen jeglichen unter uns sein Kreuz auf sich nehmen, deinem lieben Sohn nachfolgen, einer den andern lieben, gleichwie dein lieber Sohn uns geliebt hat. Denn es ist dies heilige Sacrament ein Brod, wir alle aber sind ein Leib, weil wir eines Brodes theilhaftig werden und alle aus einem Kelche trinken. Gib auch uns und allen gläubigen Christen, was uns sonst nützlich und dienlich ist, dieses Leben in deiner Furcht zu führen, selig zu beschließen und die ewige Seligkeit zu erlangen. Darauf sprechen wir von Herzen Amen.\footnote{Anmerkung des Herausgebers: Ein darauf folgendes gesungenes Amen der Gemeinde würde nicht nur am besten dieses gemeinsame Amen ausdrücken, sondern auch zugleich den Uebergang zum Singen des Vaterunsers etc. bilden.}
\end{quote}

\quellseite{259}
\chapter{XV. Consecration}


Da nach dem Bericht der Evangelisten Matthäus, Marcus und Lucas, sowie des Apostels Paulus unser HEerr Jesus Christus, „nachdem er die Elemente genommen, dankte und dann sie mit bestimmten Darreichungsworten austheilte, so ist es klar, daß Consecration und Distribution, und zwar beides in bestimmter Form, das Wesentliche beim heiligen Abendmahl ausmachen, dasjenige, ohne welches kein Abendmahl ist . . ., während das Uebrige, namentlich die vorhergehende Präfation und allenfalls, wie in der Griechischen Kirche, die „Anrufung des Heiligen Geistes“, zu dem, was die kirchliche Vollständigkeit ausmacht, zu zählen ist“. (Rudelbach, Sacramentsworte S. 56. 57.) Wenn nämlich der Apostel 1 Cor. 10, 16. das Brod und den Kelch, über welchen man nach des HEerrn Thun und Befehl „dankete“, das „gesegnete“ Brod, den „gesegneten“ Kelch nennt und diese nun als die „Gemeinschaft“ des Leibes und Blutes Christi bezeichnet, so spricht er damit wie die sacramentliche Verbindung des himmlischen Gutes mit den Elementen, so auch die Aussonderung der letzteren zu solcher Verbindung und ihrem entsprechenden Gebrauch, der im Essen und Trinken besteht, aus.

Die Handlung der Consecration besteht darin, daß über die Elemente die Einsetzungsworte (verba testamenti) gesprochen werden. Mit ihnen hat die Kirche zu allen Zeiten die Consecration vollzogen. Sie hat damit die Ueberzeugung ausgesprochen, „daß sie nicht vollständiger, nicht ergreifender des HEerrn Befehl ausführen könne, als indem sie das Lebenswort mit der urkundlich bezeugten That verbindet, und also in jeder Art thut, was Christus that, oder vielmehr Christum selbst als noch fort und fort wirkend, schaffend, heiligend durch sein Wort darstellt“. (R.) Durch des Dieners Mund wiederholt derselbe die einmal ausgesprochenen Stiftungs-Worte, sondert dieselben ab vom gemeinen Gebrauch, und wie das Schöpfungswort: „Seid fruchtbar und mehret euch und füllet die Erde“, einmal nur ausgesprochen, allezeit kräftig ist in der Natur, so ist auch dies Testamentwort: „Das ist mein Leib, das ist mein Blut“ kräftig und wirksam, so oft man seine Stiftung hält, also daß die Elemente, wenn sie genossen werden, nicht mehr schlechtweg Brod und Wein, sondern die Organe, Träger und Mittel zur mündlichen Genießung des wahren Leibes und Blutes Christi sind.

\quellseite{260}
Mit den Einsetzungsworten verbindet die Kirche jedoch das heilige Vaterunser (oratio dominica). Bekanntlich besteht hierin zwischen der morgenländischen und abendländischen Kirche seit dem 5. Jahrhundert eine liturgische Differenz. Jene entschied sich für ein besonderes Weihegebet, die vorhin erwähnte „Anrufung des Heiligen Geistes“, wie sich dieselbe in verschiedener Version in den unter dem Namen des Clemens, Jacobus, Marcus, Basilius, Chrysostomus etc. bekannt gewordenen Liturgien der Griechen sich findet. Die abendländische Kirche entschied sich für das Gebet des HEerrn und die lutherische Kirche folgte ihr darin. Wie in dem Stiftungsworte soll man auch hier ipsissima verba vernehmen. Sehr bezeichnend sagt Gregor der Große: „Mir will es sehr unziemlich vorkommen, daß mir ein Gebet, von irgend einem Kirchenlehrer (Scholasticus) abgefaßt, über die Oblation sprechen und hingegen des Gebetes gleichwiegen sollten über den Leib und das Blut des Erlösers, welches er selbst verfaßt.“\footnote{Beispielsweise stehe hier eine solche Form der „Anrufung des Heiligen Geistes“. Sie findet sich in der im 8. Buch der „Apostol. Constitutionen“ enthaltenen „Liturgie des Clemens“, wohl der ältesten der hernach schriftlich fixirten Liturgien, und lautet dieselbe, nachdem die Sacramentsworte, jedoch interpolirt, wie in allen jenen Liturgien, vorausgegangen sind, also: „Eingedenkt nun seines Leidens und Todes und seiner Auferstehung etc. opfern wir dir, dem König und Gotte, nach seiner Anordnung dieses Brod und diesen Kelch, indem wir dir durch ihn dafür danken, daß du uns gewürdigt hast, vor dir zu stehen und dir zu dienen; und bitten dich, daß du wohlwollend blickest auf diese vor dir daliegenden Gaben, du bedürfnißloser Gott, und Wohlgefallen an ihnen habest zur Ehre deines Christus, und herabsendest auf dieses Opfer deinen Heiligen Geist, den Zeugen der Leiden des HEerrn Jesu, daß er dieses Brod mache zu dem Leibe deines Christus und diesen Kelch zum Blute deines Christus, damit die daran Theilnehmenden befestigt werden in der Gottseligkeit, die Vergebung der Sünden erlangen, von dem Teufel und seinem Betrug befreit werden, vom Heiligen Geiste erfüllt, deines Christus würdig werden, das ewige Leben erlangen, nachdem du mit ihnen versöhnt bist, allmächtiger Herrscher.“ (König: „Die Haupt-Liturgien der alten Kirche“, S. 114.) Dieser Formel sind augenscheinlich die Formeln in jenen andern Liturgien nachgebildet. Uebrigens haben auch etliche der lutherischen Kirchenordnungen neben dem Vaterunser ein ähnliches Weihegebet. So lautet ein solches in der Kirchenordnung von Oettheinrich 1543: „HEerr Jesu Christe, du einiger wahrer Sohn des lebendigen Gottes etc. — wir bringen vor deine göttliche Majestät diese deine Gaben, Brod und Wein, und bitten, du wollest dieselben durch deine göttliche Gnade, Güte und Kraft heiligen, segnen, und schaffen, daß dieses Brod dein Leib und dieser Wein dein Blut sei, und allen denen, die davon essen und trinken, zum ewigen Leben lassen gedeihen“ etc.} Daß jedoch gleichwohl von den ältesten Zeiten her das Vaterunser bei der Consecration gebraucht worden ist, erhellt aus den Aussprüchen des Justin, Cyrill, Hieronymus, Augustin und Gregor des Großen, welche Rudelbach anführt (Sacraments-Worte, S. 64. 67. 70.).

\quellseite{261}
Verschieden nun aber ist in lutherischen Abendmahlsliturgien die Stellung des Vaterunsers. Im Anschluß an die römische Messe lassen die einen es nach den Einsetzungsworten folgen, während die andern es ihnen unmittelbar voransetzen. Selbst in Luthers beiden Gottesdienstordnungen findet sich diese Verschiedenheit, denn in der „Weise, christlich Meß zu halten (Formula Missae), 1523“ folgt das Vaterunser den Einsetzungsworten, und in der „deutschen Messe und Ordnung des Gottesdienstes, 1526“ geht es ihnen voran, wenn auch in Gestalt der Paraphrasis. Nach Vorgang der Formula Missae beließen es an seiner Stelle u. a. die Märkische, die Brandenb.-Nürnb. K.-O. und die von Ottheinrich. Die übrigen Kirchenordnungen aber, die in Nord-, Mittel- und Südwestdeutschland, speciell die in Sachsen, Niedersachsen, Mecklenburg, Schleswig-Holstein, Pommern, Schlesien, Württemberg etc., folgen der andern Weise Luthers und reihen das Gebet des HEerrn vor den Einsetzungsworten ein, jedoch so, daß sie dasselbe nicht nur an den Altar verlegen, da es von Luther selbst freigegeben wurde, „ob man solche Paraphrasis oder Vermahnung wolle auf der Kanzel flugs auf die Predigt thun, oder vor dem Altar“, sondern auch es ohne Paraphrase zu brauchen vorschreiben, mit der eigentlichen Luther’schen kurzen Abendmahlsvermahnung die mit dem Vaterunser verbundene Consecration einleiten, und zwar auch hier mit den Luther’schen Schlußworten: „Demnach wollen wir in seinem Namen und aus seinem Befehl, durch seine eigenen Worte, das Testament also handeln und gebrauchen.“ Mit Recht, denn auch im Vaterunser vernimmt man, wie schon bemerkt, ipsissima verba.\footnote{Sehr passend enthält das Formular für eine nur aus Beichte und Abendmahl bestehende „Frühcommunion“ in der St. Louiser Agende gleichwohl die vielsagende, alle Noth der ganzen Christenheit umfassende schöne Vaterunser-Paraphrase, wodurch dieselbe Vermahnung zum Gebet des HEerrn wird.}

Welches ist nun die Bedeutung des Vaterunsers im Communionact? Weil dasselbe auf Grund seiner Anfangsworte und als Gebet des HEerrn von Anfang an als „das Gebet der Gläubigen“ gegolten hat, durch welches wir „als Kinder Gottes vor das Angesicht Gottes treten und ihn kraft unsers Kinderrechtes um jegliche Güter anflehen dürfen“, im Communionact aber ursprünglich seine Stellung nach den Einsetzungsworten hat, so weist Schöberlein die Annahme, daß es die Bedeutung eines Consecrationsgebetes habe, als irrthümlich ab und faßt es in solcher Stellung allein als Gebet der Gläubigen „vor dem Hinzutritt zum heiligen Mahle“ auf, das von ihnen, und zwar
\quellseite{262}
ehedem unisono, hier gesprochen wurde, „um sich dadurch im Bewußtsein ihrer Kindschaft zu bekräftigen, welche ihnen das Recht dazu gewährte“. Hierbei ist jedoch nicht zu übersehen, daß aus den Abendmahlsliturgien der alten Kirche ein Weihegebet für die Elemente gleichwohl in die Liturgie des Abendlandes mit übergegangen ist und, wie schon angegeben, selbst etliche der lutherischen Kirchenordnungen, welche sich nach der Formula missae richteten, ein solches Gebet in rechtgläubigem Sinne und in rechtgläubiger Form enthalten. Anders jedoch Rudelbach in seiner Schrift: „Die Sacrament-Worte“. In derselben heißt es S. 67: „Mit der Consecration, antworteten wir, hat sie (die alte Kirche) das Gebet des HErrn verbunden, wiederum als die Summa alles desjenigen, was wir bitten können, und darum auch als die rechte Eulogie. Und so lehrt denn unsere Kirche, das Ganze zusammenfassend: durch das Gebet des HErrn werden die Elemente zum heiligen Gebrauche definirt; dieser Gebrauch selbst, die Consecration, wird vollzogen mit den Einsetzungsworten; und mit der Distribution verbunden, ist das Mysterium vollbracht, wird uns Christi Leib und Blut in, mit und unter dem Brod und Wein gereicht.“ Dieser Auffassung, die wir unsererseits als die genügende halten, entspricht dann um so mehr die Stellung des Vaterunsers vor den Einsetzungsworten. Thatsache ist auch, daß man hernach in der lutherischen Kirche vorwiegend für diese Stellung sich entschied.

Wir haben nun die liturgische Ausführung der Consecration ins Auge zu fassen.

Zunächst den Vortrag der Einsetzungsworte und des Vaterunsers.

In dem römischen Meßcanon sind die Einsetzungsworte theils aus den Liturgien des Clemens und Jacobus interpolirt, theils ist von denselben das: „der für euch gegeben wird“ geradezu weggelassen. Luther hat die ipsissima verba Domini unvermischt und unverkürzt für die Consecration wieder hergestellt und an allen lutherischen Altären vernimmt man sie seitdem unverändert.

Im christlichen Alterthum consecrirte man mit lauter, vernehmlicher Stimme. Dadurch, daß man hernach die Messe zu einem sogenannten unblutigen Opfer für die Lebendigen und die Todten machte, ist man von diesem Gebrauch im Papstthum abgeweichen und es ist zum „Greuel der Stillmesse“ gekommen. Unmittelbar nach der Präfation nämlich beginnt der Canon missae
\quellseite{263}
im engeren Sinne. Man versteht darunter jene zunächst auf die Consecration sich beziehenden fünf bis sechs Opfergebete. Luther nennt ihn in der Formula Missae „den zufließenden, göttlichen, greulichen Canon oder Stillmesse, der aus vielen stinkenden Pfützen zusammengesetzet und geflickt ist“. Diese Gebete werden vom Priester „heimlich“, mit leiser Stimme gesprochen, während der Chor das Sanctus singt — zum Glück für die arme Kirche freilich. Aber wenn nun diese Gebete gemurmelt sind und die Einsetzungsworte kommen, so schweigt zwar der Gesang und es erfolgt eine tiefe Stille, deren Eintritt das Glöcklein des Ministranten kündet, allein mit dem Qui pridie („Welcher in der Nacht, da er verrathen ward“ etc.) wird nun die Stimme des Priesters vollends zum Flüsterton und betet derselbe nach der mit der Consecration verbundenen Elevation die noch übrigen paar Opfergebete mit leiser Stimme, während der Chor weiter singt, und erst mit dem Paternoster wird seine Stimme laut und wird dasselbe von ihm gesungen, wenn überhaupt die betreffende Messe eine musikalische ist. Indem nun Luther den ganzen Opfergreuel und was nach demselben auch nur schmeckte, schonungslos hinausfegte, that er auch den Mißbrauch des Leisesprechens insbesondere der Consecrationsworte ab. Gleich in seiner ersten liturgischen Hauptschrift, der Formula Missae, ordnet er das Singen der Einsetzungsworte, und zwar in deutscher Sprache, während er die Präfation und manch anderes Stück lateinisch noch beläßt, und obwohl er da aus Rücksicht den lauten oder leisen Vortrag in die Freiheit des Einzelnen stellt, so kann er doch zehn Jahre später in seiner Schrift „Von der Winkelmesse und Pfaffenweihe“ freudig auszurufen: „Da tritt vor den Altar unser Pfarrherr ...., der singt fröhlich und deutlich die Ordnung Christi im Abendmahl eingesetzt.“ Und so finden sich denn auch in den verschiedenen alten Kirchenordnungen Vorschriften, wie: daß der Liturg die Einsetzungsworte „mit sonderer Tapferkeit und langsamem gutem Verstande“ — „mit lauten, verständlichen Worten“ — „mit hoher deutlicher Stimme und besonderer Tapferkeit“ singen solle.

Es ist damit schon angedeutet, daß die Consecration gesangsweise geschah. Zwar forderten etliche der rechtgläubigen Kirchenordnungen der Reformationszeit ausdrücklich, daß die Einsetzungsworte und mit ihnen das Vaterunser „gesprochen und nit gesungen“ werden sollen. So die Kirchenordnung von Augsburg, Ulm, Schwäbisch-Hall und die von Brenz verfaßte Württembergische Kirchenordnung; ebenso auch wurde im Elsaß, in der Pfalz, in Frankfurt a. M., in Nassau und Hessen nicht gesungen — kurz, im südwestlichen
\quellseite{264}
Deutschland, wo sich zwar nicht in der Lehre, wohl aber in der Liturgie der Einfluß der reformirten Nachbarschaft namentlich in Folge der Bucer’schen Unionsversuche geltend machen konnte. Sonst aber verbreitete sich die gesangsweise Behandlung dieser Stücke über ganz Nord- und Mitteldeutschland und verschaffte sich selbst in einzelnen Theilen Süddeutschlands noch Eingang. Warum auch nicht? Mit Recht bemerkt Schöberlein: „Und auch wenn man dem liturgischen Gesang des Geistlichen nicht hold ist, kann man ihm an dieser Stelle geneigt sein. Denn es ist hier nicht ein Gebetsact, der vollzogen wird, sondern Worte des HErrn, welche die heilige Grundlage der ganzen Handlung bilden, sollen feierlich recitirt werden. Und wenn irgendwo, so ist hier der liturgische Altargesang am Orte.“ Wo daher trotzdem es nicht möglich ist, eine Gemeinde für das liturgische Singen auch nur der Consecrationsworte zu gewinnen — und ein Adiaphoron bleibt ja doch auch diese entsprechende Weise der rechtgläubigen Kirche — da befleißige sich der Liturg, diese Worte immer mit gehobener Stimme zu sprechen!

Da in der römischen Messe die Einsetzungsworte flüsternd gesprochen werden, so fand Luther wohl eine Melodie für das Vaterunser vor, keine aber für jene. In der Formula Missae äußert er nun zwar: „Gerne wollte ich, daß diese Worte Christi über eine kleine Weile nach der Präfation im selben Ton gesungen würden, darin man sonst das Pater noster in canone zu singen pflegt“; in der deutschen Messe jedoch legt er selbst eine Melodie vor, zu deren Schönheit u. a. gehört, daß sie beim Eintreten der Worte des HErrn („Nehmet hin“ etc.) kühn nach F herunterschreitet. Sie fand Eingang mit etlichen Veränderungen namentlich in Sachsen, Thüringen und zum Theil in Niedersachsen. Zu der Luther’schen Weise fanden sich dann noch einige andere Weisen. Dahin gehört namentlich die Bugenhagen’sche von 1528 und die Brandenburgisch-Nürnbergische Melodie von 1533. Erstere verbreitete sich mit etlichen, doch unbedeutenden Variationen zunächst über Norddeutschland und Pommern, drang aber auch nach Sachsen und in die Lausitz und wurde nach 1662 selbst in Hessen aufgenommen. Im Lauf der Zeit ist Bugenhagens Weise wohl am meisten gebraucht und findet sich dieselbe daher auch in der St. Louis’er Agende. Die letztere, die Brandenburg-Nürnbergische Weise, die mehr „in nota praefationis“ gehalten ist, fand ihre Aufnahme in Franken, wo man sich ihrer heute noch fast allüberall bedient, und wurde auch nach Oesterreich und selbst in die lutherische Pfalz verpflanzt. Gleichfalls eine herrliche Melodie. Obwohl nun von diesen Sangesweisen man je eine in eine Landeskirche einzuführen
\quellseite{265}
pflegte, so bieten doch häufig die älteren Kirchenordnungen jener und der nachfolgenden Zeit mehrere Weisen zur Auswahl dar, wie z. B. die uns hier vorliegende Leipziger Agende 1681, damit „ein jeder Priester eine nehme, die ihm und seiner Stimme am bequemsten“, wie es in einer andern Agende heißt.

Theils zum Studium, theils zum beliebigen Gebrauch folgen die genannten drei liturgischen Weisen der Einsetzungsworte und des Vaterunsers denn auch hier. „Weil sich in diesen Melodien der liturgische Vortrag mehr als in andern dem wirklichen Gesange nähert und deshalb die Begleitung der Orgel eher verträgt“ (Schöberlein), so wird unter Hinweisung auf das in der Einleitung zu diesem Theil § 3 Gesagte hier zu beliebigem Gebrauch eine Orgelbegleitung beigegeben, wobei man sich mit nur wenigen Ausnahmen der Riegel’schen Harmonisirung in Schöberleins Werk bediente, eine Begleitung in continuo aber vermied, und zwar auf Grund langjähriger Erfahrung und durch dieselbe befestigter Ueberzeugung, daß, wenn eine Orgelbegleitung hier stattfinden soll, sie so anzuwenden sei, daß die Begleitung nur am Schluß, oder beim Wechsel der Harmonie hie und da eintrete, dagegen, wo mehrere Töne der Melodie denselben Accord zur Grundlage haben, ganz schweige und allein der Liturg vernommen werde.

Noch ist zu bemerken, daß die der Luther’schen Weise hier beigegebene Vaterunser-Melodie aus Thüringen stammt. Da nämlich in der Deutschen Messe Luthers wegen der paraphrastischen Form des Vaterunsers sich nur für die Einsetzungsworte die Melodie findet, so ist in Thüringen, wo die Luther’sche Weise vorherrschend im Gebrauch war und wohl noch ist, eine entsprechende Vaterunser-Melodie entstanden. Der Verfasser, der derselben über ein Vierteljahrhundert in seiner vorigen Gemeinde sich bedient hat, theilt dieselbe in der ihm bekannt gewordenen Gestalt und Harmonisirung mit.

\section[das Decorum Betreffende]{Schließlich einiges das Decorum Betreffende}

Selbstverständlich wird die Consecration dem Altar zugekehrt verrichtet.

Bei den Worten: „Nahm er das Brod“ und: „Nahm er den Kelch“ ergreift oder berührt man mit der Hand Patene und Kelch, und bei den Worten: „Das ist mein Leib“ und „Das ist mein Blut“ oder: „Dieser Kelch ist das neue Testament in meinem Blut“ macht man über die Elemente ein Kreuz. Manche Agenden, wie z. B. die Schwarzburgische Agende von 1675, setzen das Kreuzeszeichen zu dem Wörtlein
\quellseite{266}
„ist“, andere wieder, wie z. B. die St. Louiser Agende, zu den Worten: „Leib“ und „Blut“.

Altkirchlich ist’s, daß entweder die ganze Gemeinde oder der Haufe der Communicanten während der Consecration knieet. Luther: „Da tritt vor den Altar unser Pfarrherr etc. . . . . Und wir, sonderlich so das Sacrament nehmen wollen, knien neben, hinter und um ihn her“ etc. (In der früheren Gemeinde des Verfassers knieten beim Beginn der Consecration die Communicanten, die unter dem „Schaffe in mir, Gott“, theils in den Altarraum, theils in die Gänge getreten waren, während die übrigen Zuhörer standen.)

\section{Vaterunser und Testamentworte}

\section{A. Luthers Weise}

\subsection{a. Vaterunser. (Oratio dominica.)}

(Aus Thüringen.)

\subsubsection{Pastor}

\notensatz{music/266-01}{Luthers Vaterunser, Pastor und Orgel: Laßt uns beten! Vater unser, der du bist im Himmel,}

\subsubsection{Orgel}

Ohne Pedal.

\notensatz{music/266-02}{Luthers Vaterunser: Geheiliget werde dein Name. Zukomme dein Reich. Dein Wille}

\quellseite{267}
\notensatz{music/267-01}{Luthers Vaterunser: geschehe, wie im Himmel, also auch auf Erden.}

\notensatz{music/267-02}{Luthers Vaterunser: Unser täglich Brod gib uns heut. Und vergib uns unsre Schuld,}

\notensatz{music/267-03}{Luthers Vaterunser: wie wir vergeben unsern Schuldigern. Und führe uns nicht in Ver-}

\notensatz{music/267-04}{Luthers Vaterunser: -suchung. Sondern erlöse uns von dem Uebel.}

\quellseite{268}
\subsubsection{Gemeinde mit Orgel}

\notensatz{music/268-01}{Luthers Vaterunser, erste Doxologie: Denn dein ist das Reich, und die Kraft, und die}

\notensatz{music/268-02}{Luthers Vaterunser, erste Doxologie: Herrlichkeit in Ewigkeit. Amen.}

\subsubsection{Oder:}

\notensatz{music/268-03}{Luthers Vaterunser, zweite Doxologie: Denn dein ist das Reich, und die Kraft, und die Herrlichkeit in}

\notensatz{music/268-04}{Luthers Vaterunser, zweite Doxologie: Ewigkeit. Amen.}

Anmerkung. Statt dieser beiden Doxologien kann auch die nächstfolgende gebraucht werden.

\quellseite{269}
\subsection{b. Testamentworte. (Verba testamenti.)}

(Luther 1526.)

\subsubsection{Orgel}

\notensatz{music/269-01}{Luthers Testamentworte, Intonation und Orgel: Unser HErr JEsus Christ(us), in der Nacht, da er verrathen ward,}

\notensatz{music/269-02}{Luthers Testamentworte: nahm er das Brod, dankt und brach’s, und gab’s seinen Jüngern und sprach:}

\notensatz{music/269-03}{Luthers Testamentworte: Nehmet hin und esset. Das ist mein † Leib, der für euch gegeben wird.}

\notensatz{music/269-04}{Luthers Testamentworte: Solches thut (, so oft ihr’s thut,) zu meinem Gedächtniß. Desselben}

(, so oft ihr’s thut,)\footnote{Dieser eingeklammerte Satz des Originals kann auch weggelassen werden.}

\quellseite{270}
\notensatz{music/270-01}{Luthers Testamentworte: gleichen auch den Kelch, nach dem Abendmahl, danket, gab ihnen den}

\notensatz{music/270-02}{Luthers Testamentworte: und sprach: Nehmet hin und trinket alle daraus. Dieser Kelch}

und sprach:\footnote{Deutsche Messe: „nach dem Abendmahl und sprach“ — a a a a c h a.}
Dieser Kelch\footnote{Ebendaselbst: „Das ist der Kelch, ein neu Testament“ etc. f f e f$'$ i f f f f f i etc.}

\notensatz{music/270-03}{Luthers Testamentworte: ist das neue Testament in meinem † Blut, das für euch vergos-}

\notensatz{music/270-04}{Luthers Testamentworte: -sen wird zur Vergebung der Sünden. Solches thut, so oft ihr’s trinkt,}

\quellseite{271}
\notensatz{music/271-01}{Luthers Testamentworte, Schluß: zu meinem Gedächtniß.}

\subsubsection{Anderorts auch:}

\notensatz{music/271-02}{Luthers Testamentworte, alternative Formeln: Das ist mein Leib etc. Dieser Kelch etc.}

\section{B. Bugenhagen’sche Weise}

\subsection{a. Vaterunser}

1528.

\notensatz{music/271-03}{Bugenhagen’sches Vaterunser, Einleitung: Laßt uns beten! Vater unser,}

\notensatz{music/271-04}{Bugenhagen’sches Vaterunser: der du bist im Himmel. Geheiliget werde dein Name.}

\quellseite{272}
\notensatz{music/272-01}{Bugenhagen’sches Vaterunser: Zukomm dein Reich. Dein Wille geschehe, wie im Himmel, also}

\notensatz{music/272-02}{Bugenhagen’sches Vaterunser: auch auf Erden. Unser täglich Brod gib uns heut. Und vergib}

\notensatz{music/272-03}{Bugenhagen’sches Vaterunser: uns unsre Schuld, wie wir vergeben unsern Schuldigern. Und führe}

\notensatz{music/272-04}{Bugenhagen’sches Vaterunser: uns nicht in Versuchung. Sondern erlöse uns von dem Uebel.}

\quellseite{273}
\subsubsection{Gemeinde}

\notensatz{music/273-01}{Bugenhagen’sches Vaterunser, Gemeinde: Denn dein ist das Reich, und die Kraft, und die}

\notensatz{music/273-02}{Bugenhagen’sches Vaterunser, Gemeinde: Herrlichkeit in Ewigkeit. Amen.}

\subsection{b. Testamentsworte}

\notensatz{music/273-03}{Bugenhagen’sche Testamentsworte: Unser HErr JEsus Christus, in der Nacht, da er verrathen ward,}

\notensatz{music/273-04}{Bugenhagen’sche Testamentsworte: nahm das Brod, dankt und brach’s, und gab’s seinen Jüngern und sprach:}

\quellseite{274}
\notensatz{music/274-01}{Bugenhagen’sche Testamentsworte: Nehmet hin und esset, das ist mein Leib†, der für euch gege-}

\notensatz{music/274-02}{Bugenhagen’sche Testamentsworte: -ben wird. Solches thut zu meinem Gedächtniß. Desselben glei-}

\notensatz{music/274-03}{Bugenhagen’sche Testamentsworte: -chen nahm Er auch den Kelch, nach dem Abendmahl, danket,}

\notensatz{music/274-04}{Bugenhagen’sche Testamentsworte: gab ihnen den und sprach: Nehmet hin und trinket alle daraus.}

\quellseite{275}
\notensatz{music/275-01}{Bugenhagen’sche Testamentsworte: Dieser Kelch ist das neue Testament in meinem Blut†,}

\notensatz{music/275-02}{Bugenhagen’sche Testamentsworte: das für euch vergossen wird zur Vergebung der Sünden.}

\notensatz{music/275-03}{Bugenhagen’sche Testamentsworte: Solches thut, so oft ihr’s trinkt, zu meinem Gedächtniß.}

\section{C. Brandenburg-Nürnbergische Weise}

\subsection{a. Vaterunser}

1583.

\subsubsection{Intonation}

\subsubsection{Orgel}

\notensatz{music/275-04}{Brandenburg-Nürnbergisches Vaterunser, Intonation und Orgel: Laßt uns beten! Vater unser,}

\quellseite{276}
\notensatz{music/276-01}{Brandenburg-Nürnbergisches Vaterunser: der du bist im Himmel. Geheiligt werd dein Name. Zukomm}

\notensatz{music/276-02}{Brandenburg-Nürnbergisches Vaterunser: dein Reich. Dein Will geschehe, wie im Himmel, auch auf Erden.}

\notensatz{music/276-03}{Brandenburg-Nürnbergisches Vaterunser: Unser täglich Brod gib uns heut. Und vergib uns unsre Schulden,}

\notensatz{music/276-04}{Brandenburg-Nürnbergisches Vaterunser: als wir vergeben unsern Schuldigern. Und führ uns nicht}

\quellseite{277}
\notensatz{music/277-01}{Brandenburg-Nürnbergisches Vaterunser: in Versuchung. Sondern erlös uns vom Uebel.}

\subsubsection{Gemeine}

\notensatz{music/277-02}{Brandenburg-Nürnbergisches Vaterunser, Gemeine: Denn dein ist das Reich, und die Kraft, und die Herrlichkeit}

\notensatz{music/277-03}{Brandenburg-Nürnbergisches Vaterunser, Gemeine: in Ewigkeit. Amen.}

\subsubsection{Amen beim Wegfall der Doxologie}

\notensatz{music/277-04}{Erstes Amen beim Wegfall der Doxologie.}

\subsubsection{Oder:}

\notensatz{music/277-05}{Zweites Amen beim Wegfall der Doxologie.}

\subsection{b. Testamentworte}

\notensatz{music/277-06}{Brandenburg-Nürnbergische Testamentworte: Unser HErr JEsus, in der Nacht, da er verrathen ward, nahm er}

\quellseite{278}
\notensatz{music/278-01}{Brandenburg-Nürnbergische Testamentworte: das Brod, dancket und brach’s, und gab’s seinen Jüngern}

\notensatz{music/278-02}{Brandenburg-Nürnbergische Testamentworte: und sprach: Nehmet hin und esset, das ist mein Leib†,}

\notensatz{music/278-03}{Brandenburg-Nürnbergische Testamentworte: der für euch gegeben wird. Das thut zu meinem Gedächtniß.}

\notensatz{music/278-04}{Brandenburg-Nürnbergische Testamentworte: Desselben gleichen nahm er auch den Kelch, nach dem Abendmahl,}

\quellseite{279}
\notensatz{music/279-01}{Brandenburg-Nürnbergische Testamentworte: und danket, und gab ihnen den und sprach:}

\notensatz{music/279-02}{Brandenburg-Nürnbergische Testamentworte: Trinkt alle daraus. Das ist mein Blut† des neuen Te-}

\notensatz{music/279-03}{Brandenburg-Nürnbergische Testamentworte: -stamentes, das für euch und für viel vergossen wird zur Ver-}

\notensatz{music/279-04}{Brandenburg-Nürnbergische Testamentworte: -gebung der Sünden. Solchs thut, so}

\quellseite{280}
\notensatz{music/280-01}{Brandenburg-Nürnbergische Testamentworte: oft ihr’s trinkt, zu meinem Gedächtniß.}

\chapter{XVI. Agnus Dei. Pax Domini}


In seiner „Weise, christlich Meß zu halten“ (Formula missae) 1523 schreibt Luther: „Nach dem soll das Pater noster gelesen oder gesungen werden. Doch daß das folgende Gebet: Wir bitten, erlöse uns etc. ausgelassen werde mit allen Schirmschlagen, die man bisher gemacht hat über die Hostien, und mit der Hostien über den Kelch;\footnote{Dieses Schirm- oder Kreuzschlagen besteht darin, daß der Priester ein Stück von der geeigneten Hostie abbricht, mit diesem Stückchen den geeigneten Kelch dreimal betrauft und dann singt: „Pax Domini sit semper vobiscum“, welches Stückchen er dann in den Kelch wirft.} auch soll die Hostien nicht gebrochen werden, auch kein Theil derselben (unter den Wein) im Kelch vermischt werden, sondern bald aufs Pater noster sollen diese Worte: Pax Domini etc. gelesen werden, welche sind eine öffentliche Absolution von Sünden aller, die zum Sacrament gehen, und ist fürwahr ein recht evangelisch Wort, das Vergebung der Sünden verkündiget, und die einige und allerwürdigste Bereitung zum Tisch des HEerrn, so es mit dem Glauben gefaßt wird, nicht anders, als hätte Christus aus seinem eigenen Munde geredet. Darum wollte ich, daß sich der Priester mit dem Angesicht zum Volk kehrete, wenn er dies Wort spricht, wie die Bischöfe zu thun pflegen, in welchem Stück allein die jetzigen Bischöfe den alten nachfolgen. Hernach reiche er das Sacrament beide ihm selbst und dem Volke, indes singe man das Agnus Dei. So aber der Priester dies Gebet: HEerr Jesu Christe, du Sohn des lebendigen Gottes, der du nach dem Willen des Vaters etc. vor der Communion sprechen wollte,
\quellseite{281}
thut er nicht unrecht daran; allein daß er für die Worte meine und mich, unsere und uns lese. Mag auch dies Gebet sprechen: Der Leichnam unsers HEerrn etc. bewahre meine oder deine Seele zum ewigen Leben, und das Blut unsers HEerrn bewahre meine oder deine Seele zum ewigen Leben, und das Blut unsers HEerrn bewahre“ etc.

Was nun das Joh. 1, 29. entnommene Agnus Dei betrifft, so reicht dessen Gebrauch überhaupt weit in das christliche Alterthum zurück, und war es in der orientalischen Kirche meist Morgengesang. Im 6. Jahrhundert verpflanzte es dann Gregor der Große durch Uebersetzung in die lateinische Kirche. Seine Einfügung in die Messe aber und seine Stellung zwischen Consecration und Distribution erhielt es erst durch Papst Sergius I. (687), und zwar mit der Bestimmung, daß es „von der Priesterschaft und vom Volke gesungen werden“ solle, was jedoch Papst Hadrian I. ein Jahrhundert später dahin abänderte, daß es, wie es nun auch bis heute geschieht, allein der Chor singen dürfe. Die dreimalige Wiederholung kam erst um 1120 auf, wobei für die dritte Wiederholung das „miserere nobis“ (erbarm dich unser) mit dem „dona nobis pacem“ (gib uns deinen Frieden) vertauscht wurde, und zwar in Bezug auf den Friedenskuß (osculum pacis), womit sich alsdann die Gemeinde zu begrüßen pflegte. Sowohl in seiner „Geschichte des Kirchenlieds und Kirchengesangs“ (IV, S. 191) als auch in Herzogs „Real-Encycl.“ (I, 179) gibt Koch in einer älteren Uebersetzung einer Stelle aus dem Werke des Bischofs Durandus († 1270): „Rationali officiorum divinorum“ folgendes als Grund der dreimaligen Wiederholung an: „Niemand hat größere Geduld in den allergrößten Leiden, Versuchungen und Anfechtungen von sich leuchten lassen, als der Erlöser, deshalb sich die christliche Kirche darüber verwundert und wohl bedächtlich dreimal singet: ‚Agnus Dei, qui tollis‘ etc. Bei der dreimaligen Wiederholung sollen wir bedenken, wie der HEerr Jesus unsere Sünden 1. weggetragen, 2. die Strafe selber getragen, 3. uns in der Predigt des Evangelii und im hochheiligen Abendmahl sein Verdienst ins Herz getragen habe.“

Da Christus bei der Einsetzung des heiligen Abendmahls spricht: „Solches thut zu meinem Gedächtniß“ und Paulus demgemäß schreibt: „So oft ihr von diesem Brode esset und von dem Kelch des HEerrn trinket, sollt ihr des HEerrn Tod verkündigen, bis daß er kommt“, so fand Luther den Gebrauch des Agnus bei der Communion nur um so passender
\quellseite{282}
und sinniger. Er schreibt: „Sonderlich dient das Agnus über alle Gesänge aus der Massen wohl zum Sacrament, denn es klärlich daher singt und lobt Christum, daß er unsere Sünden getragen habe, und mit schönen kurzen Worten das Gedächtnis Christi gewaltiglich und lieblich treibt.“ (Erl. Ausg. 23, 191, St. L. Ausg. X, 2195.) Während er nun aber in seiner Formula Missae auch hier enger an die Gliederung der römischen Messe sich anschließt und deshalb ordnet: Consecratio, Pater noster, Pax Domini, Agnus mit dem Gebet: Domine Jesu Christe etc., läßt er in seiner „deutschen Messe“ das Vaterunser den Einsetzungsworten voraus gehen, wie wir bereits gesehen haben. Weil nun hier Luther die Consecration und Distribution des Brodes und Weines theilt, so ordnet er für die Austheilung des ersteren das Singen von Liedern, wie „Jesaia, dem Propheten“, oder: „Gott sei gelobet“, oder: „JEsus Christus, unser Heiland“, nach der Segnung und bei Austheilung des Kelches aber „was übrig ist von obgenannten Liedern oder das deutsche Agnus Dei“. (Erl. Ausg. 22, 241. St. L. Ausg. X, 249.) Diejenigen Kirchenordnungen nun, welche nach Luthers Vorgang das Vaterunser den Einsetzungsworten voranstellen, ohne jedoch die von Luther vorgeschlagene Theilung betreffs der Consecration und Distribution zu acceptiren, stellen vielfach das Agnus an den Schluß der Austheilung, andere jedoch an den Anfang derselben, resp. gleich hinter die Testamentworte. So heißt es z. B. in der Braunsch.-Lüneb. K.-O.: „Sobald die Communion angefangen wird, soll die ganze Gemeinde getracks ohne einiges vorhergegangenes Spiel auf der Orgel oder einigen andern Gesang langsam und mit Andacht singen: ‚O Lamm Gottes unschuldig.‘ Wenn dasselbe zu Ende und noch Communicanten vorhanden, soll gesungen werden der Gesang vom Leiden Christi: ‚Christus, der uns selig macht‘ etc. Und nach demselben, darnach viel Communicanten sein: ‚HEerr Jesu Christ, du höchstes Gut.‘ ‚Jesus Christus, unser Heiland.‘“ Also unmittelbar nach den Worten Christi: „Solches thut zu meinem Gedächtnis“ die besondere Verkündigung des Todes Christi durch das Agnus Dei! Wo man heutzutage wieder zu der alten Weise des Gottesdienstes zurückkehrt, wie sie die meisten sächsischen Agenden vorschreiben, zieht man es vor, das Agnus an die Testamentsworte zu reihen. So z. B. die Agende der Missourisynode. Wie aus dem Mitgetheilten zu ersehen, bediente man sich mancherorts auch der liedförmigen Bearbeitung des Agnus von Decius: „O Lamm Gottes unschuldig.“

Für das Agnus Dei gibt es verschiedene Melodien. Die fast überall nun gebräuchliche ist die bekannte einfache schöne Melodie, wie sie sich in allen
\quellseite{283}
guten Choralbüchern und so auch im St. Louiser Choralbuch findet. Um deswillen ist eine Mittheilung derselben hier überflüssig. Nur für den Organisten die Bemerkung, daß, wenn in der römischen Messe das Agnus „in sanften Klängen, öfters mit leisen Flötentönen und zartem Saitenspiel begleitet gesungen wird“, bei der Orgelbegleitung Rohrflöte mit Gamba oder Salicional die entsprechenden Register sein dürften.

Pax Domini, den Friedensgruß (Joh. 20, 21.) betreffend, so verhält es sich mit demselben in der römischen Messe also: Wie schon erwähnt, bricht nach dem Vaterunser der Priester die geweihte Hostie in zwei Hälften und die eine davon wieder in zwei Theile. Mit dem kleinsten derselben macht er über den Kelch ein Kreuz und singt dabei: „Pax Domini sit semper vobiscum“ (der Friede des HEerrn sei allezeit mit euch), worauf die Antwort: „Et cum Spiritu tuo“ (und mit deinem Geiste). Hierauf läßt er dies Theilchen der Hostie in den Kelch fallen mit den Worten: „Haec commixtio et consecratio Corporis et Sanguinis Domini nostri Jesu Christi fiat accipientibus nobis in vitam aeternam. Amen. (Diese Vermengung und Weihung des Leibes und Blutes unsers HEerrn Jesu Christi gereiche uns, den Empfängern, zum ewigen Leben. Amen.)\footnote{In seiner Schrift: „Vom Greuel der Stillmesse“ 1523 macht Luther hierzu folgende Glosse: „Da redet er wieder, der Narr, von vielen, die da nehmen, und nimmt es alleine. Und ist abermal ein Lästerwort, daß er sagt: Er vermenge den Leib und Blut Christi untereinander. Ich wollte gern wissen, wie man den Leib und das Blut mengen könnte, daß doch Ein Christus ist? Brod und Wein vermenget er wohl. Das ist aber viel ärger, daß er’s heißet eine Heiligung des Leibes und Blutes; gerade als wäre etwas da an Christo zu heiligen. Und sagt dazu, die Vermengung solle gedeihen denen, so sie nehmen, zum ewigen Leben, das ist, das Werk soll es thun. Wo bleiben die Wort, darin das Leben stehet? Ist’s doch alles nur auf uns und aufs Opfer gestellet!“ (Erl. Ausg. 29, 128 ff.)}

So unerbittlich Luther nun diese letztere Ceremonie, obwohl gleichfalls aus dem Alterthum herrührend,\footnote{Vgl. „Die Haupt-Liturgien der alten Kirche in wortgetreuer Uebersetzung“ etc. von Joh. Ludw. König.} weil aber lästerlich, abthat, so schüttete er doch nicht das Kind mit dem Bade aus, indem er in seiner „Formula Missae“ den Friedensgruß aufnahm, ihn so schön deutend als „eine öffentliche Absolution von Sünden aller, die zum Sacrament gehen“, und als „ein recht evangelisch Wort, das Vergebung der Sünden verkündigt, und die einige und allerwürdigste Bereitung zum Tisch des HEerrn, so
\quellseite{284}
es mit dem Glauben gefaßt wird, nicht anders, als hätte Christus aus seinem eigenen Munde geredet“. Eben darum will er, daß dieser Gruß nicht gegen den Altar gesungen oder gesprochen wird, wie in der römischen Messe, sondern gegen das Volk.

Der Friedensgruß findet sich nur in den lutherischen Liturgien, welche sich mehr an die Formula missae anschließen, wie z. B. die von Straßburg, Erfurt und Brandenburg-Nürnberg. Neuere gute Liturgien haben um seiner schönen Deutung willen den Friedensgruß meist aufgenommen und ihm mit Recht seine Stelle nach dem Agnus und unmittelbar vor der Austheilung angewiesen. So z. B. das „Kirchenbuch“ des General Council, das, obwohl an die Brandenb.-Nürnb. K.-O. nach dem Vorbild Löhe’s sich eng anschließend, doch das Vaterunser den Einsetzungsworten vorangehen läßt. Für diejenigen, welche sich vor der Austheilung des Friedensgrußes bedienen können oder wollen, sind hier ein paar Melodien beigegeben. —

Nach dem Agnus und vor dem Genuß folgen in der römischen Messe noch ein paar Vorbereitungsgebete. Das eine derselben: „O HEerr Jesu Christe, du Sohn des lebendigen Gottes, der du nach dem Willen des Vaters“ etc. hält Luther in der Formula Missae für gut, stellt jedoch den Gebrauch in des Einzelnen Willkür. Offenbar ist das schöne, auch in der St. Louiser Agende S. 235 ff. der Vollbrecht’schen Abendmahlsvermahnung später angehängte Vorbereitungsgebet eine Erweiterung und daher auch an dieser Stelle zu gebrauchen. In der neuen Ausgabe der Löhe’schen Agende (1884) steht es daher auch an dieser Stelle erst verkürzt, dann wörtlich von den Worten an: „Heilige darum du selbst uns“ etc. Die Breslauer Agende enthält folgendes, so viel mir wissen, in den meisten schlesisch-lutherischen Kirchen gebrauchtes Vorbereitungsgebet:

Andächtige und Auserwählte, damit aber auch Gott euch Gnade gebe, das heilige Abendmahl selig zu genießen, so laßt uns ihn mit Mund und Herzen anrufen und also beten:

\begin{quote}
O HEerr, ob ich zwar nicht würdig bin, daß du unter mein Dach eingehest, so bin ich doch nothdürftig deiner Hülfe und begierig deiner Gnade, daß ich möge fromm und selig werden. Nun komme ich in keiner andern Zuversicht, denn auf dein Wort, da du selbst mich zu diesem Tische läßt und sagst mir Unwürdigen zu, ich solle Vergebung meiner Sünden haben durch dein Leib und Blut, so ich esse und trinke in diesem Sacrament. O lieber HEerr, ich weiß, daß deine göttliche Zusage und deine Worte wahrhaftig sind, daran zweifle ich
\quellseite{285}
nicht und darauf esse und trinke ich: mir geschehe nach deinem Worte. Amen.
\end{quote}

(Luther.)

O HEerr Jesu, vereinige dich mit mir, auf daß ich bleibe in dir und du in mir und ich von dir ungeschieden sei, hie zeitlich und dort ewiglich.

Dein heiliger Leib, HEerr Jesu Christe, speise mich; dein theures Blut tränke mich; dein bitter Leiden und Sterben stärke mich; in deine heiligen Wunden verberg ich mich, laß mich von dir nimmermehr geschieden werden. Vom bösen Feind errette mich. Im wahren Glauben erhalte mich, auf daß ich dich sammt allen Auserwählten lobe und preise, hie zeitlich und dort ewiglich. Amen.

\section{Pax Domini}

\subsubsection{(Zu den Testamentsworten sub A und B.)}

Braunschweig 1531.

\notensatz{music/285-01}{Pax Domini, Braunschweig 1531, Versikel. Der Fried des HErren sei mit euch allen.}

\notensatz{music/285-02}{Pax Domini, Braunschweig 1531, Gemeinde, vierstimmig. Amen.}

\subsubsection{(Zu den Testamentsworten sub C.)}

Brandenb.-Nürnb. 1533.

\notensatz{music/285-03}{Pax Domini, Brandenburg-Nürnberg 1533, Versikel. Der Fried des HErren sei mit euch allen. Dazu eine zweite Melodie für allen.}

\notensatz{music/285-04}{Pax Domini, Brandenburg-Nürnberg 1533, Gemeinde. Zwei vierstimmige Amen.}

\quellseite{286}

\chapter{XVII. Distribution}


„Mit der Consecration aufs innigste verbunden ist die Distribution des heiligen Abendmahls nebst der Spendungsformel, so daß unsere Kirche beides eigentlich als ein Continuum betrachtet, wovon das Letztere ebenso wenig als das Erstere hinweg gedacht werden kann, wenn von der Handlung des Sacraments die Rede ist“ (Rudelbach). Anders in der römischen Messe; denn ob in derselben nun auch auf die Consecration die Sumtio (Selbstcommunion) des Priesters folgt und an diese die Communion der Laien unter Einer Gestalt sich schließt, so betrachtet doch speciell die römische Kirche die Distribution als einen bloßen, nicht wesentlichen Appendix der Messe und redet mehr von einem geistigen Mitgenuß der Gemeinde, während sie die Sumtio des Priesters als wesentlich zur Messe, resp. zum Meßopfer gehörig hält und dabei vorgibt, dem zuschauenden Volke erwüchse aus der Sumtio des opfernden Priesters ein Nutzen.

Daß im lutherischen Hauptgottesdienst die Sumtio nicht schlechtweg abgethan ist, zeigen die im vorigen Abschnitt mit angeführten Worte aus Luthers Formula Missae: „Hernach reiche er das Sacrament beide ihm selbst und dem Volke, indes singe man das Agnus Dei.“ Ihm nach findet sich die Sumtion mit der Distribution in der Kirchenordnung von Wittenberg durch Bugenhagen, von Straßburg und des Andreas Döber in Nürnberg 1525. Daß nun aber die Schmalkaldischen Artikel in den bekannten Worten nicht gegen diese Selbstcommunion des Pastors bei der Communion der Gemeinde reden, sondern allein gegen die Selbstcommunion des Priesters, die mit Ausschluß der Gemeinde in der sogenannten Still- oder Opfermesse geschieht, ist längst sattsam nachgewiesen. Nur als die nicht normale Weise der Dispensatio wird sie angesehen, indem sie nur für den Nothfall als zulässig erklärt wird, „wenn nämlich, wie hier nicht selten, Prediger so einsam und entfernt von Amtsbrüdern stehen, daß sie, wenn sie sich das heilige Abendmahl nicht selbst reichen wollten, desselben oft über Jahr und Tag entbehren müßten“. (Für das Nähere überhaupt s. „Americ.-luth. Pastoraltheologie“ von Dr. Walther, § 18. Anmerk. 8.)

Mit der Austheilung des heiligen Abendmahls unter beiderlei Gestalt ist die lutherische Kirche zur Einsetzung des HErrn und damit zum Brauch der alten Kirche zurückgekehrt. Ob man nun aber nach alter Weise die vereinigte
\quellseite{287}
Consecration der beiden Gestalten der Distribution wolle vorausgegeben lassen, oder ob man das Brod gleich nach seiner Segnung austheile und so auch den Kelch, stellt Luther sowohl in seiner „Formula Missae“ als auch in seiner „deutschen Messe“ etc. dem einzelnen Pastor völlig anheim. Er schreibt in ersterer: „Es soll auch hier dem Bischof oder Pfarrherren frei sein, mit was für Ordnung er beide Gestalt wolle nehmen oder reichen. Denn er mag sie beide zumal, Brod und Wein, segnen, ehe er das Brod genießet oder zwischen der Segnung des Brods und Weins sich und andere, so viel ihr begehren, mit dem Brode speisen, nochmals auch den Wein segnen und alsdann allen zu trinken geben, welcher Weise Christus, achte ich, gebraucht habe, wie die Worte des Evangelii lauten, da er die Jünger hieß essen das Brod, ehe er den Kelch segnete. Nochmals spricht der Evangelist deutlicher: Dieselben gleichen nahm er auch den Kelch nach dem Abendmahl etc., daß man wohl vernehmen kann, daß der Kelch erst nach der Nießung des Brods gesegnet sei. Aber diese allzu neue Weise wird nicht leiden, daß dies alles geschehe, das wir bisher nach der Segnung gesagt haben, man wollte sie denn auch ändern.“ (St. L. Ausg. X, 2245.) Und in letzteren, in welcher er zugleich auch die entsprechende Aenderung der Form und Weise vornimmt, schreibt er: „Es dünkt mich aber, daß es dem Abendmahl gemäß sei, so man flugs auf die Consecration des Brods das Sacrament reiche und gebe, ehe man den Kelch segnet. Denn so reden beide Lucas und Paulus: Dieselben gleichen den Kelch, nachdem sie gegessen etc., und dieweil singe das deutsche Sanctus, oder das Lied: Gott sei gelobet; oder Johanns Hussens Lied: JEsus Christus, unser Heiland. Darnach segne man den Kelch und gebe denselbigen auch, und singe, was übrig ist von obgenannten Liedern, oder das deutsche Agnus Dei.“ (Erl. Ausg. 22, 241. St. L. Ausg. X, 249.) Gleichwohl aber wurde die vereinigte Consecration der beiden Gestalten der herrschende Brauch und die andere Weise nur für die Krankencommunion da und dort, wie z. B. noch im Nürnberger Agendbüchlein von 1755, verwendet, für welche sie auch am anwendbarsten sein dürfte.

Nach Vorgang der alten Kirche hat die lutherische Kirche auch eine Distributionsformel. In der Formula Missae behielt Luther die der römischen Kirche bei: „Der Leichnam unsers HEerrn Jesu Christi bewahre deine Seele zum ewigen Leben.“ In der Agende Herzog Heinrichs zu Sachsen 1539 heißt es etwas ausführlicher: „Der Leib unsers HEerrn Jesu Christi, für dich in den Tod gegeben, stärke und bewahre dich im rechten Glauben zum ewigen Leben.“ Seit Dietrichs
\quellseite{288}
Agendbüchlein 1543 hat die Formel: „Nimm hin und iß, das ist der Leib Christi, der für dich gegeben ist. Nimm hin und trink“ etc. Gegenüber den Sacramentirern aber hielt man es in der Folgezeit für nöthig, das Bekenntniß der wahrhaftigen Gegenwart des Leibes und Blutes Christi unter Brod und Wein noch bestimmter zu fassen, und entstand so die noch überall in der lutherischen Kirche gebräuchliche Spendeformel: „Nehmet hin und esset, das ist der wahre Leib“ etc. Mit Recht, denn wenn „die Bewahrung der Reinheit des Wortes die Hauptsache bei der Consecration war, so wird sie es bei der Distribution nicht weniger sein, und eine jede Willkür um so mehr zu verwerfen, als in den wenigen Worten der Aneignung die Summe des ganzen Vorhergehenden befaßt sein muß“. (Rudelbach.) Wo nun aber der Pastor sich selbst communicirt, nimmt er das Brod allein mit dem Segenswunsch: „Der Leib unsers HEerrn Jesu Christi bewahre meine Seele zum ewigen Leben. Amen“; desgleichen den Kelch mit den Worten: „Das Blut unsers HEerrn Jesu Christi bewahre meine Seele zum ewigen Leben. Amen.“

Die Ordnung des Empfangs betreffend, so wird, wie in der alten Kirche, das heilige Abendmahl erst den Männern und dann den Weibern gereicht. An manchen Orten wurde und wird es so gehalten, „daß die Männer und Junggesellen erstlich, hernach die Jungfrauen mit den Weibern und Müttern hinzugehen“ (Brandenb.-Magdeb.), während an andern Orten die Jünglinge die Reihe der Männer und die Jungfrauen die Reihe der Weiber eröffnen. Communicirt der Prediger, so empfängt derselbe nach altkirchlichem Brauch zuerst das Sacrament. Doch kommt, wiewohl mehr vereinzelt, die Bestimmung vor: „Wenn jedermann das heilige Sacrament empfangen, so soll es dann der Priester auch empfangen.“ (Oetheinrich.) Ueberhaupt setzte man im Allgemeinen voraus, daß der Prediger der fleißigste Genosse am Tisch des HEerrn sei und daher womöglich jedesmal mitcommunicire. Beherzigenswerth ist demnach u. a. die Vermahnung der Agende von Oesterreich: „Die Kirchendiener und Pastores sollen sich auch gewöhnen, wie zu der Apostel Zeit und in der Kirche allermalen gebräuchlich gewesen, daß sie mit communiciren, und das zu gutem Exempel der Kirchen und der Gemein.“

Zwar empfing man schon früher an diesem und jenem Orte das Sacrament stehend. Doch gehörte das mehr zu den Ausnahmen. In der Regel empfing und empfängt man es knieend, vorher ehrerbietig sich neigend. Es sollte daher am Altar „eine kurze Bank“ angebracht werden, „darauf die
\quellseite{289}
Communicanten fein knieen können, wenn sie das Nachtmahl empfangen\footnote{In der römischen Kirche empfänget nur der Meßhalter am Altar das Sacrament, allen andern wird es gewöhnlich an der Schranke gereicht, welche den Altarraum vom Schiff der Kirche trennt.} (Höna 1571 etc.). Das Knieen beim Empfang des Sacraments ist auch wegen der Darreichung des Kelchs passender. Ist der Umgang um den Altar bräuchlich und möglich, wie er auch meistentheils geschah, so treten zum Empfang des gesegneten Brodes je drei an die linke Seite, die nördliche, wenn der Altar gen Osten steht. Knieen oder stehen sie um den Altar, so daß die Austheilung schichtenweise stattfindet, so kann auch bei der Empfangung des heiligen Leibes unter dem gesegneten Brode bei der Spendeformel der Segenswunsch noch weggelassen und erst wenn die Schicht das heilige Blut empfangen hat, von dem Consecrator der Segenswunsch mit den Worten gesprochen werden: „Der Leib unsers HEerrn Jesu Christi und sein theures Blut stärke und bewahre euch im wahren Glauben zum ewigen Leben. Amen.“ Doch kann auch die Spendeformel in Verbindung mit dem Segenswunsch wie beim Umgang um den Altar gebraucht werden.

Will das gesegnete Brod oder der gesegnete Wein nicht ausreichen, so schweigt Orgel und Gesang und wird von neuem consecrirt, sei es das Brod oder der Wein oder beides. Nach der Ulmer Kirchenordnung soll man, um auszureichen, „keineswegs einige Hostien entzweybrechen“, sondern es sollen „die abgängigen Hostien hergebracht“ werden. Diese Kirchenordnung will auch bei der Nach-Consecration nicht nur die diesbezüglichen Testamentesworte, sondern auch das Vaterunser wiederholt haben, während anderwärts nur jene gefordert werden.

Wenn es nun in Luthers „deutscher Messe“ in Bezug auf den Empfang noch heißt, „daß man fein ordentlich und züchtig zugehe“, so finden sich in den Kirchenordnungen hin und her hierzu einzelne Bestimmungen. So lesen wir z. B. in der Magdeburger Kirchenordnung von 1685: „Und sollen alle Communicanten mit schuldiger Ehrerbietung, Andacht, Zucht und Demuth ohne alles Abziehen auf einige praecedenz zum Tische des HEerrn gehen, auch mit der Kleidung, Geberden und allem sich also bezeigen, daß keine Leichtfertigkeit, Pracht oder Hoffarth, sondern vielmehr der rechte christliche Schmuck an ihnen gespüret werde, welcher nicht stehet im Haarflechten, Goldumhängen etc. ... Insonderheit sollen die Weibes-Personen sich, wie sonst allezeit, also auch insonderheit bei Empfahung des heiligen Nachtmahls nicht entblößen und alle
\quellseite{290}
die Trachten, so in Austheilung Hinderung bringen\footnote{Es ist noch nicht lange her, daß eine weibliche Huttracht bei der Austheilung, sonderlich des Kelchs, geradezu hinderlich wurde.} oder sonsten Aergerniß geben, bei der Communion ablegen.“ Und in Corp. jur. eccl. Saxon. 1735 heißt es nach Augustis Denkmal.: „Sonderlich aber (sollen) die Männer mit ihren Zwick-Bärten sich nicht ärgerlich in den Kelch legen und damit nachmals Aergerniß geben. Desgleichen auch die Weiber die Schleier vor dem Mund und die Parete vor denen Augen, desgleichen den Mund aufgethan, auf daß ihnen das Sacrament in den Mund gegeben, und nicht, wie etwa geschehen, die Communicanten gar nichts aus dem Kelche trinken.“ (Auch in manchen unserer hiesigen Gemeinden bestand ehedem die feine Sitte, daß die Frauen nicht nur in dunklem einfachen Kleide zum Tisch des HEerrn gingen, sondern auch die Jungfrauen entblößten Hauptes, die Frauen aber im Hernahen ein weißes Häubchen trugen. Leider ist diese Sitte mehr und mehr schon geschwunden.)

So lange die Austheilung währt, erbaut sich die übrige Gemeinde, deren Gegenwart dabei vorausgesetzt wird, — nach der Braunschw.-Lüneb. Kirchenordnung soll deshalb „der Priester öfter im Jahr die Gemeine ernstlich vermahnen, bei der Communion zu bleiben“ — durch den Gesang von Communionliedern, namentlich der Lieder „Jesus Christus, unser Heiland“ und „Gott sei gelobet und gebenedeiet“, auch des deutschen Sanctus: „Jesaja, dem Propheten“ etc. Bei großen Communionen bediente man sich dann zur Erhaltung und Vermehrung der Andacht gern des Wechselgesangs. Es sang z. B. „der Chor und das Volk einen Vers um den andern“ (Pomm.). An hohen Festtagen verband sich damit hie und da noch Instrumentalmusik.\footnote{So heißt es z. B. in dem Communionlied von S. Scherertz: „Mein Seel, dich freu“ B. 21–23:
\begin{verse}
Zur Freude mein Violen rein\quad Die Orgel auch, nach altem Brauch,\\
Beweglich laßt erklingen,\quad Im Tempel Gottes schöne,\\
Von Herzen thut mit gutem Muth\quad Posaunenklang ziert den Gesang\\
In Harmonie bald singen.\quad Mit lieblichen Getöne,\\[1ex]
Auf daß die Gäst zu diesem Fest\\
Hierüber jubiliren\\
Und sich mög heut die Traurigkeit\\
Aus dem Gemüth verlieren.
\end{verse}} Ausdrücklich aber wird für den Gebrauch der Instrumentalmusik die Beschränkung geordnet, „daß der Choralgesang nicht gar hintangesetzt werde“, und für den
\quellseite{291}
Organisten wird bemerkt, daß das Zwischenspiel auf der Orgel „mit Maß“ geschehe. Eine nöthige Bemerkung auch für heutige Zeit. Wie nicht selten kommt es vor, daß ein Organist glaubt, zwischen den einzelnen Versen des Abendmahlsliedes recht lang orgeln zu müssen, entweder weil er fürchtet, es möchte das Communionlied nicht ausreichen, oder noch mehr, es möchte die Gemeinde durch zu vieles Singen ermüden, während er nicht merkt, daß gerade das maßlose Zwischenorgeln die Gemeinde ermüdet und zum Weiterfingen verdrossen macht, zumal wenn der Spieler kein Meister auf der Orgel ist. Man halte also auch hier mit dem Spiel zwischen den einzelnen Versen Maß und bedenke, daß die Gemeinde nicht um des Organisten, sondern der Organist um der Gemeinde willen da ist. Reicht das gewählte Communionlied nicht aus, so stecke man einfach eine andere Liednummer noch an, oder nehme von dem „Gott sei gelobet“ die ersten beiden Verse und mache den dritten zum Schlußvers.\footnote{Das non plus ultra erlebte ich vor Jahren bei Gelegenheit einer Versammlung von Vertretern mehrerer Synoden, wo man während der Austheilung des heiligen Abendmahls es bis zum Singen von nur zwei Versen des Liedes: „Schmücke dich, o liebe Seele“ brachte, während man das ganze Lied hätte hinaussingen können, wenn der HEerr Organist in seinen Zwischenspielen, die noch dazu recht opernmäßig klangen, bescheiden Maß gehalten hätte. Noch sehe ich den vor mir sitzenden sel. W., wie er einmal uns andere unwillig sein graues Haupt schüttelte und zornige Blicke nach dem Mann auf der Orgelbank schoß.}

\chapter{XVIII. Postcommunio}


In der Formula Missae heißt es: „Will man das Commun singen, so singe man’s. Anstatt aber der Complend oder letzten Collecte, weil sie gemeiniglich aufs Opfer stimmen, lese man im selben Ton dies Gebet: HEerr, das wir mit dem Munde empfangen haben etc. Auch könnte man das Gebet lesen: Deinen Leib, HEerr, den wir empfangen haben etc., und schließe es mit diesen Worten: Der du lebest und regierest etc. Dann soll der Priester singen: Dominus vobiscum“ etc. (St. L. Ausg. X, 2244.) Nach der Genießung singt oder spricht der römische Meßpriester auf der Epistelseite einen Spruch, gewöhnlich Tob. 12, 7.: „Lobet und danket ihr Gott vom Himmel bei jedermann, daß er euch solche Gnade erzeiget hat“, oder sonst nach Gelegenheit einen biblischen Spruch. Das ist das „Commun“, der römischen Messe. Dann läßt sich
\quellseite{292}
derselbe von dem Ministranten etwas nicht consecrirten Wein in den Kelch gießen, denselben auszuspülen,\footnote{Zuweilen wird dieser Spülwein auch Laien dargereicht und meinen dann die armen Leute, sie hätten das Sacrament unter beiderlei Gestalt empfangen.} wobei er das erste von Luther angeführte Gebet spricht: „HEerr, das wir\footnote{„Wir“ — welch ein Selbstzeugniß des Meßhalters gegen seine Winkelmesse!} mit dem Munde“ etc. Hierauf wischt er den Kelch mit einem Tüchlein aus, hält Daumen und Zeigefinger über den Kelch, läßt sich über dieselben Wasser zur Reinigung gießen und spricht jenes zweite Gebet: „Deinen Leib, HEerr“ etc.

Was nun nach diesem zunächst folgt, heißt die Postcommunio. Dieselbe besteht in einer nochmaligen Bitte, daß der geschehene Sacramentsgenuß für die Theilnehmer ein geeigneter gewesen sein möge, welcher dann nach dem Abschiedeswort an die Gemeinde: „Ite, missa est“ noch die „Complend“ oder die Complenda sich anschließt und mit den Worten beginnt: „Placeat tibi, sancta Trinitas“ etc.,\footnote{In seiner Schrift „Vom Greuel der Stillmesse“ verdeutscht Luther diese „Complend“ also: „Laß dir gefallen, du heilige Dreifaltigkeit, unsern Dienst, und gib, daß das Opfer, so ich für den Augen deiner Majestät unwürdig geopfert habe, sei dir angenehm, und mir sammt allen, für welche ich das geopfert habe, durch deine Erbarmung versöhnlich, durch Christum, unsern HEerrn. Amen.“ Dazu macht L. die Glosse: „Wie er’s angefangen hat, so führet er es auch hinaus, opfert immerdar und bittet, daß es angenehme sei“ etc.} das letzte Gebet, das, weil es noch einmal auf das Opfer recurrirt und darum bittet, daß dasselbe „Deo acceptabile“, dem Priester aber und denen, für die er es dargebracht hat „propitiabile“ sei, Luther streicht und dafür jene zwei Gebete zur Postcommunion der Formula Missae macht. In seiner deutschen Messe jedoch setzt er für diese Gebete eine Dank-collecte, die mit dem Aufruf beginnt: „Lasset uns dem HEerrn danken und beten.“ Sie fand nebst einer ähnlichen überall Aufnahme und bildet seitdem im lutherischen Hauptgottesdienst die Postcommunio.

Und mit Recht! Sie ist das Gratias nach der himmlischen Mahlzeit. Wie die Abendmahlshandlung durch die Präfation mit einer Danksagung beginnt, so schließt sie durch die Postcommunion mit einer Danksagung. Sehr schön gibt Chrysostomus den Grund dazu aus der Schrift an, indem er über Matth. 26, 30.: „Und da sie den Lobgesang gesprochen hatten, gingen sie hinaus an den Oelberg“ diejenigen, welche vor dem Schluß des Gottesdienstes die Kirche zu verlassen pflegen, also anredet: „Höret ihr, die ihr das letzte Gebet des Sacraments nicht erwartet, was die Bedeutung des
\quellseite{293}
desselben ist. Der HErr dankte, bevor er es den Jüngern gegeben hatte, damit auch wir von dieser Dankstagung anheben sollten. Er dankte, nachdem er es ihnen gegeben hatte, und sang einen Lobgesang, damit auch wir gleicherweise thun sollen.“ (Rudelbach, Sacrament-Worte, S. 81.)

\section{Eine Weise der Dankcollecte}

\subsubsection{Latria}

\notensatz{music/293-01}{Dankcollecte, erster Versikel des Pastors. Danket dem HErrn, denn er ist freundlich. Halleluja!}

\notensatz{music/293-02}{Dankcollecte, vierstimmige Antwort. Und seine Güte währet ewiglich. Halleluja!}

\subsubsection{Oder:}

\notensatz{music/293-03}{Dankcollecte, zweiter Versikel des Pastors. So oft ihr von diesem Brode esset und von diesem Kelche trinket,}

\notensatz{music/293-04}{Dankcollecte, vierstimmige Antwort. Sollt ihr des HErrn Tod verkündigen, bis daß er kommt.}

\notensatz{music/293-05}{Dankcollecte des Pastors. Laßt uns dem HErrn danken und beten! Wir danken dir, allmächtiger HErre Gott, daß du uns durch diese heilsame Gabe}

\quellseite{294}
\notensatz{music/294-01}{Dankcollecte, Fortsetzung. Des Leibes und Blutes deines lieben Sohnes JEsu Christi hast erquicket, und bitten deine Barmherzigkeit, daß du uns solches gedeihen lassest zu starkem Glauben gegen dich und zu brünstiger Liebe unter uns allen.}

\notensatz{music/294-02}{Dankcollecte, Schluß. Durch unsern HErrn JEsum Christum, deinen Sohn, der mit dir in Einigkeit des Heiligen Geistes, wahrer Gott, lebet und herrschet immer und ewiglich.}

\notensatz{music/294-03}{Dankcollecte, Gemeinde, erstes vierstimmiges Amen.}

Oder:

\notensatz{music/294-04}{Dankcollecte, Gemeinde, zweites vierstimmiges Amen.}

Oder:

\notensatz{music/294-05}{Dankcollecte, Gemeinde, drittes vierstimmiges Amen.}



\quellseite{295}
\chapter{XIX. Schlußacte}


Den Schluß des Hauptgottesdienstes gestaltet Luther in der Formula Missae folgendermaßen: „Denn (nämlich nach der Abendmahlscollecte) soll der Priester singen: Dominus vobiscum etc. Und für das Ite, missa est, singe man Benedicamus Domino, mit dem Alleluja, in gewöhnlichem Ton, wo und wenn’s gefällt, oder nehme die Benedicamus, so man in der Vesper braucht. Zuletzt spreche der Priester den gewöhnlichen Segen übers Volk, und nehme den aus dem 6. Cap. des 4. Buchs Mosis V. 24–26., so der HEerr selbst geordnet hat: Der HEerr segne dich etc. Oder nehme den aus dem 67. Psalm, V. 8.: Es segne uns Gott, unser Gott etc. Ich glaube auch, daß Christus dergleichen Segen gesprochen habe über seine Jünger, da er gen Himmel aufgefahren.“ Den Schluß bilden hiernach zwei Acte: das mit der Salutation verbundene Benedicamus und der Aaronische Segen.

1. Was nun zunächst das Benedicamus Domino (Laßt uns benedeien dem HEerrn) betrifft, so hat es Luther zwar aus der „deutschen Messe“ hernach weggelassen (gleichwie auch die Präfation), wohl aber hat er für die täglichen Metten und Vespern es geordnet (Erl. Ausg. 22, 236. St. L. Ausg. X, 235.), denen es ursprünglich angehört und denen er es für den Hauptgottesdienst entnommen und in der Formula Missae es an die Stelle des „Ite, missa est“ und dessen Responsorium: „Deo gratias“ gesetzt hat. (Gleichwohl hat diese Verwerthung für den Hauptgottesdienst damals Beifall gefunden. Es findet sich in den Kirchenordnungen von Bugenhagen 1524, Andr. Döber 1525, Straßburg 1525, Erfurt 1526 und Brandenb.-Nürnberg 1533. Weil es nun ein antiphonischer Lobspruch von eigenthümlicher Schönheit ist, so wird es auch von neueren Liturgikern meist in den Hauptgottesdienst wieder aufgenommen.

Welches ist nun hier seine Bedeutung? Da es Luther an die Stelle des die römische Messe abschließenden Ite, missa est (Geht, die Messe ist aus!) setzt, so bildet es den eigentlichen Schluß des Hauptgottesdienstes, gleichwie es den Schluß der Metten und der Vesper bildet, und da die Postcommunion-Collecte der Dank insonderheit der Communicanten für die Gabe des Leibes und Blutes Christi ist, so soll nach unserm Dafürhalten das Benedicamus der Dank nicht allein für die Gabe des Leibes und Blutes Christi, sondern auch für die Gabe des gehörten Wortes und darum nicht allein der Dank der
\quellseite{296}
jedesmaligen Communicanten, sondern auch der Dank der ganzen versammelten Gemeinde sein, unter welcher der HEerr laut seiner Verheißung Matth. 18, 20. durch Wort und Sacrament seine Gnadengegenwart abermals geoffenbart hat.

Wie reich in musikalischer Beziehung das Benedicamus auch in den Agenden und Ritualien der lutherischen Kirche ausgestattet ist, zeigt ein Blick in dieselben. Da findet sich je eines mit einem dreifachen Halleluja für die einzelnen Festtage, wie auch für die Sonntage insgemein, ja auch für etliche Sonntage insonderheit und wiederum ohne Halleluja für Messe und Vesper an den Wochentagen. Da jedoch die Ausführung dieser mancherlei Weisen, sonderlich der Hallelujas, für ungeübte Sänger zu viel Schwierigkeit bieten, auch den Ohren zu ungewohnt sein dürfte, so wählen wir mit Hommel, Löhe, Layriz etc. eine der einfachen Weisen.

2. Dem Benedicamus nun folgt der Aaronische Segen. Es ist dies der zweite Schlußact. Mit ihm wird die versammelte Gemeinde entlassen. In der römischen Messe lautet die Segensformel zur Entlassung: „Benedicat vos omnipotens Deus, Pater et Filius et Spiritus sanctus. Amen.“ (Es segne euch der allmächtige Gott, der Vater, der Sohn und der Heilige Geist. Amen.) Sie findet sich nebst der von Luther mit vorgeschlagenen aus Ps. 67 und noch zwei andern in der Brandenb.-Nürnb. Agende, wie in dem Agendenbüchlein von Veit Dietrich. Der Aaronische Segen 4 Mos. 6 hat indes die allgemeinste Verbreitung gefunden und ist die Segensformel zur Entlassung der Gemeinde bei allen gottesdienstlichen Veranlassungen, wie auch für einzelne Handlungen geworden.\footnote{Bekanntlich hat Luther auch eine herrliche Auslegung des Segens, „so man nach der Messe spricht über das Volk“, geschrieben. S. Erl. Ausg. 36, 156 ff.} (Er ist die passendste Segensform zur „Bekräftigung und Besiegelung aller durch Wort und Sacrament empfangenen Gnadengaben durch eine letzte Gnadenmittheilung“.

Die Segensformel, gesprochen oder gesungen, wird mit symbolischer Bewegung der Hände begleitet, indem man während desselben die Hände oder wenigstens die rechte Hand erhebt, am Schlusse aber bei dem Worte „Friede“ das Kreuzzeichen macht. Das Erheben der Hände über die Gemeinde ist die biblische Gebärde des Segnens, das Kreuzeszeichen aber über die Gemeinde am Schlusse desselben „prägt dem Symbole seinen specifisch christlichen Character auf und erinnert hier, wie überall, daß alle Gnade und aller
\quellseite{297}
Segen aus dem Verdienste des Gekreuzigten fließt“. Das Zeichen des Kreuzes ist nach alter Vorschrift in dieser Form und Folge zu machen:

\begin{center}
1\\
4\quad 2\quad 5\\
3
\end{center}

und zwar, es mag dasselbe mit der ganzen Hand oder mit dem Daumen und den beiden ersten Fingern formirt werden, so daß der kleine Finger dem zu segnenden Gegenstand zunächst zugekehrt ist, hier darum mit senkrecht emporgerichteter Hand und langsamer, mäßiger Bewegung des Armes — also nicht mit der der Gemeinde zugekehrten Handfläche und in rascher Bewegung, was eher wie ein der beim Segen stehenden Gemeinde gegebenes Zeichen zum Niedersitzen aussieht, denn als ein Zeichen zur Erinnerung an das Kreuz Christi, oder als sollte, wie einmal ein schlichter Landmann nicht so ganz unrichtig meinte, „das Kreuz gar wieder weggewischt werden“.

Wiewohl nun aber wie bei den Collecten so auch bei dem Segen in den Kirchenordnungen meist nur ein einmaliges Amen vorgeschrieben ist, so hat man doch in der nachfolgenden Zeit sich vorzugsweise eines zweimaligen, mitunter auch dreimaligen Amen’s bedient.

3. Wir haben hier jedoch noch eines dritten Schlußactes zu gedenken. Dies ist nämlich der sogenannte Schlußvers. Während nämlich in der römischen Messe auf das vorausgehende Ite, missa est die Versammlung nach der Benediction die Kirche verläßt, bleibt der Priester noch am Altar stehen, um nach vorausgehender letztmaliger Salutation auf der Evangelienseite das Evangelium vom fleischgewordenen Wort, Joh. 1, 1–14., zu lesen oder zu singen. Das ist nun im lutherischen Hauptgottesdienst von Anfang an mit Recht allerwege abgethan. Während nun aber bei Luther in seiner Formula Missae, wie in seiner deutschen Messe, bei Bugenhagen, Döber, in der Straßburger-, Erfurter- und Brandenb.-Nürnberger Kirchenordnung auf das Amen nach dem Segen nichts weiter folgt, setzen andere, etwas spätere Kirchenordnungen, an die Stelle des Johannisevangeliums der römischen Messe noch einen Schlußgesang, eben den sogenannten Schlußvers, meist ein stehender Gesang. So ordnet Wittenberg 1559 und 1565, Mecklenb. 1562 und 1650, Frankfurt a. M. 1565, Lüneburg 1564 und 1569, Liegnitz 1594 u. A. das Lied: „Erhalt uns, HEerr, bei deinem Wort“, Westphalen 1585: „Verleih uns Frieden gnädiglich“ etc. und Magdeburg 1692: „Gott sei gelobet und gebenedeiet.“ In
\quellseite{298}
Augustis Braunschw.-Lüneb. Kirchenordnung 1657 aber heißt es nach dem Segen: „Darauf singe die ganze Gemeinde choraliter: HEerr, nun lässest du deinen Diener in Frieden fahren etc. Und darauf: Gott sei gelobet etc. ganz aus.“ Beides würde nun freilich für unsere heutigen Verhältnisse den Gottesdienst zu sehr verlängern, da man gewohnt ist, nur einen Vers oder höchstens ein paar kleine Verse nach dem Segen noch zu singen. Aber wie leicht ließe sich von den beiden Schlußgesängen Gebrauch machen, wenn man sich ihrer abwechselnd bediente, wenn man also einmal während der Austheilung das ganze Lied „Gott sei gelobet und gebenedeiet“ singen ließe und in diesem Fall den Simeonischen Lobgesang mit seiner schönen Choralweise zum Schlußvers machte, wie dies auch von Layriz geschehen ist. Auch Löhe, der in den Ausgaben seiner Agende das Nunc dimittis der Danksagungscollecte voransetzt, findet es auch als Schlußgesang sehr passend. Er bemerkt: „Gewiß recht schön, vor dem Heimgehen, nach Schluß des Ganzen, den Lobgesang Simeons zu singen.“ Calvör theilt I, 715 über diesen Lobgesang eine schöne Stelle von Casaubonus mit, die in Uebersetzung auch hier stehen möge. Sie lautet:

\begin{quote}
„Wenn die heilige Handlung vollendet und die Feier des Abendmahls geschlossen ist, so wird nach dem vom Diakon gesprochenen Schlußgebet das Volk entlassen, was in der Liturgie des Chrysostomus unter den Worten verstanden wird: ‚der Priester vollzieht die Entlassung, die missa, das heißt, er entläßt das Volk.‘ Die Schlußhandlung ist eben missa oder missio, ein Wort, das dem Gesang Simeons: ‚HEerr, nun lässest du deinen Diener in Friede fahren‘ entspricht, welch letzterer daselbst auch angeführt wird. In vielen Kirchen der Protestanten wird die ganze Handlung der Feier des heiligen Abendmahls des HEerrn mit diesem Hymnus geschlossen, den das Volk mit gebeugten Knien singt, was eine sehr herrliche und sehr heilige Einrichtung ist.“
\end{quote}

Löhe setzt hinzu: „Das plerisque (in vielen) dürfte vielleicht bezweifelt werden, aber das Urtheil über den Hymnus ist gewiß recht.“\footnote{Unter der Ueberschrift: „Ite, benedicti et electi“, findet sich bei Layriz ein aus den Gesängen der böhmischen Brüder 1566 mitgetheilter Schlußgesang, der offenbar anstatt des „Ite, missa est“ in deren Gemeinden gebraucht wurde. Derselbe lautet:
\begin{verse}
Geht hin, die ihr gebenedeit und in Christo auserwählt seid,\\
Geht hin mit Freuden in Fried, Gott richt all euer Schritt.\\
Gesegnet ist euer Ausgang, gesegnet ist euer Eingang,\\
Gesegnet all euer Thun durch Christum, Gottes Sohn.\\
Derselb unser HEerr und Heiland führ uns ein ins recht Vaterland\\
Zu Lob und Ehr seim Namen in Ewigkeit. Amen.
\end{verse}
}

\quellseite{299}
Sei es nun aber, daß mit diesem oder einem andern Hymnus der letzte Schlußact geschehen, so dürfte darauf aufmerksam gemacht werden, daß während desselben der Consecrator am Altar noch verweilt, um denselben zu ordnen, falls dasselbe nicht schon nach der Austheilung und vor der Dankcollecte geschehen konnte. Es sieht nämlich nicht fein aus, wenn nach dem Segen der Pastor alsbald den Altar verläßt, auf dem die heiligen Gefäße noch offen und durcheinander stehen. Es werden unter dem Schlußgesang daher von ihm Kannen, Kelch und Hostienschachtel geschlossen und an ihren vorigen Platz gestellt, der Kelch mit der Patene, diese mit der Palla (viereckiges, das Herabwehen der Hostien von der Patene verhütendes Deckelchen von weißer Leinwand) bedeckt und über das Ganze das altkirchliche und altlutherische Velum (Tüchlein von weißer oder farbiger Seide) gebreitet. Nach einem stillen Gebet verläßt der Pastor den Altar. Noch sei bemerkt, daß um die schöne Sitte der betenden Sammlung nach Schluß des Gottesdienstes und vor dem Auseinandergehen zu erhalten und dem geräuschvollen Weggehen unmittelbar nach dem Gesang vorzubeugen, es jetzt in manchen unserer hiesigen lutherischen Gemeinden Gebrauch geworden ist, daß der Pastor während des ganzen Schlußgesangs am Altar stehen bleibt, um nach demselben noch ein lautes Vaterunser zu beten.

\section{1. Benedicamus}

\subsubsection{Formel}

\notensatz{music/299-01}{Benedicamus, Salutation des Pastors. Der HErr sei mit euch!}

Hommel.

\notensatz{music/299-02}{Benedicamus, vierstimmige Antwort der Gemeinde. Und mit deinem Geiste.}

\notensatz{music/299-03}{Benedicamus, zweiter Versikel des Pastors. Laßt uns benedeien den HErren!}

\quellseite{300}
\notensatz{music/300-01}{Benedicamus, zwei vierstimmige Fassungen der Antwort. Gott sei ewiglich Dank.}

\section{2. Der Segen. (Benedictio.)}

Layriz.

\notensatz{music/300-02}{Aaronischer Segen nach Layriz. Der HErr segne dich und behüte dich. Der HErr erleuchte sein Angesicht über dich und sei dir gnädig. Der HErr erhebe sein Angesicht auf dich, und gebe dir Friede.}

\subsubsection{Gemeinde:}

\notensatz{music/300-03}{Segen, Gemeinde, erstes vierstimmiges Amen.}

Oder:

\notensatz{music/300-04}{Segen, Gemeinde, zweites vierstimmiges Amen.}

Oder:

\notensatz{music/300-05}{Segen, Gemeinde, vierstimmiges doppeltes Amen.}

\notensatz{music/300-06}{Segen, Gemeinde, vierstimmiges dreifaches Amen.}

\quellseite{301}
\section{3. Schlußgesang. (Nunc dimittis.)}

(Abwechselnd mit: „Gott sei gelobet“ etc. zu gebrauchen.)

\subsubsection{Sopran}

\notensatz{music/301-01}{Schlußgesang: HErre, nun läßt du deinen Diener in Frieden fahren, wie}

\notensatz{music/301-02}{Schlußgesang: du gesagt hast. Denn meine Augen haben deinen Heiland ge-}

\notensatz{music/301-03}{Schlußgesang: -sehn, den du bereitet hast für allen Völkern, ein Licht zu er-}

\notensatz{music/301-04}{Schlußgesang: -leuchten die Heiden und zum Preis deines Volks}

\quellseite{302}

\notensatz{music/302-01}{Schlußgesang: Israel. Ehr sei dem Vater und dem}

\notensatz{music/302-02}{Schlußgesang: Sohn und dem Heiligen Geist, wie es war im Anfang, jetzt und}

\notensatz{music/302-03}{Schlußgesang: immerdar, und von Ewigkeit zu Ewigkeit.}

\notensatz{music/302-04}{Schlußgesang: Amen.}

Israel.\footnote{Abkürzung.}

\notensatz{music/302-05}{Schlußgesang: Abkürzung. Israel.}

\quellseite{303}

\quellseite{304}
\chapter{Anhang}


\section{Die Stätten für die Liturgie}

\subsection{Einleitende Bemerkungen}

Bekanntlich hat im Gegensatz zum Mißbrauch der Kunst in der römischen Kirche die reformirte Kirche in Bezug auch auf die Stätten des Gotteshauses für die Handlung des Worts, der Sacramente und des Gebets Anfangs den Dienst der Kunst fast ganz abgewiesen oder hernach auf ein ganz Geringes beschränkt.\footnote{Die hiesigen fashionabelen Puritanergemeinden gerathen, immer weltförmiger werdend, in ein anderes Extrem. Obwohl nach Außen und Innen mehr oder weniger im gothischen Stil erbaut, selbst mit einem Kreuz auf der Thurmspitze, erscheinen deren Kirchen im Innern doch mehr wie ein bemalter Concertsaal, namentlich da über dem Stand des Redners in einer Nische die Orgel anstatt eines Altars angebracht ist und vor derselben die Chorsänger, Männer und Frauen, sich der ganzen Versammlung präsentieren. Dazu kommt anstatt einer Nebenkapelle außer einem lecture-room vielleicht auch noch ein besonderer Saal zur Abhaltung von Sociables, Fairs etc., und daneben eine Küche nebst Zubehör!} Anders die lutherische Kirche auch hier gemäß dem Ausspruch Luthers, der nicht der Meinung war, „daß durch das Evangelium sollten alle Künste zu Boden geschlagen werden und vergehen, wie etliche Abergläubische fürgeben“. Demnach hat sie auch in dieser Beziehung „keine merkliche Aenderung“ vorgenommen, am allerwenigsten eine auffällige, sondern allein den Mißbrauch und sonstigen Unrath abgethan. Wie den Taufstein und die Kanzel, so hat sie auch den Altar an seinem Ort und in seiner Gestalt gelassen, und das nicht nur in den aus dem Papstthum vorgefundenen, sondern auch in den von ihr erbauten Gotteshäusern. Und ob man auch in der sogenannten Aufklärungsperiode puritanisch entleerte und weltförmig schmückte, so ist dies doch nur vorübergehend gewesen.

So wenig nun aber die lutherische Kirche, frei von aller unevangelischen Befangenheit und gesetzlichen Aengstlichkeit, das Kirchengebäude geflissentlich auf
\quellseite{305}
der Stufe eines schmucklosen Betsaales zu halten sucht, so tritt doch auch hier ihr Unterschied von der päpstlichen Kirche hervor. Während sie für das Gotteshaus die profane Helligkeit einer Werkstatt oder einer Schule verschmäht, vermeidet sie gleichfalls das mystische Halbdunkel und erwählt ein leicht abgedämpftes Licht. Obwohl ihr durch Wort und Sacrament der Versammlungsort allemal zur Stätte der Offenbarung der Gnadengegenwart Gottes wird, so ist er ihr doch nicht die Stätte des im römischen Sinne continuirlich gegenwärtigen Christus. Eben darum hat der lutherische Altar eine ganz andere Bedeutung, einen ganz anderen Gebrauch. Wohl handelt sie an demselben das heilige Abendmahl mit dem deutlichen und bestimmten Bekenntniß der wahrhaftigen Gegenwart des Leibes und Blutes Christi unter den Gestalten des Brodes und Weines; aber weil sie weder von einer Transsubstantiation, noch auch von einer sacramentlichen Gegenwart außer dem Genuß des Abendmahls etwas weiß, so wird ihr auch nicht der Altar zum heiligen Schrein, in welchem die Monstranz mit der geweihten Hostie eingeschlossen ist und vor dem, als einem irdischen Thron des Gottmenschen, Priester und Laie beim Vorübergehen anbetend sich beugen oder neigen.\footnote{Wie weit römwärts auch in Betreff des Altars die Ritualisten der anglicanischen Kirche schon gekommen sind, zeigt “The Parish Annual of St. Paul’s Church, Springfield, Ills., 1884.” Dasselbe bezeichnet den Altar ausdrücklich als Gottes “earthly” Thron und gibt den Zuhörern für ihr Verhalten folgende Anweisung: “I. On entering the Church do not stop to chat at the door, but make your obeisance toward the Altar. It is the throne of the King of kings, Who sanctifies it by His Presence, especially in the Holy Eucharist. ‘Ye shall reverence my Sanctuary. I am the Lord.’ — IX. Leave the Church with the same reverence with which you entered it. It is always God’s house: the shelter of His Altar, whether service is in progress or not.”} Wohl macht in christlicher Freiheit auch die lutherische Kirche den Altar zu einer Opferstätte, indem sie sonderlich an ihm durch den berufenen Diener das Rauchopfer des Gebets und der Danksagung als der geeignetsten Stätte darbringen und feierliche Gelübde vor demselben ablegen läßt: das Confirmations-, Ehe- und Ordinationsgelübde; sie wendet sich aber mit Grauen von einem mit allerlei abergläubischen Ceremonien pomphaft geweihten Meßopferaltar, zu dessen Requisiten auch noch ein paar Heiligenknochen gehören, die der Altar als sepulchrum birgt und durch die er auch noch wunderthätig werden kann. Eben darum baut oder braucht man zu den angedeuteten Zwecken in der lutherischen Kirche nur Einen Altar, während in der römischen Kirche der Hochaltar mindestens zwei Seitenaltäre hat, davon der eine Maria, der andere Joseph geweiht ist.

\quellseite{306}
Wir sehen, auch in Bezug auf das Gotteshaus und seine innere Ausstattung hält die lutherische Kirche die rechte Mitte zwischen papistischem Aberglauben und schwärmerischer Abergeistlichkeit.

Da gemäß dem öffentlichen Character und dem Gemeinschaftsprincip des lutherischen Gottesdienstes überhaupt die heiligen Handlungen vor den Augen der Versammlung allewege zu verrichten sind, so ist dadurch eine solche bauliche Disposition gefordert, durch welche Kanzel, Altar und Taufstein von der Mehrzahl der Versammelten gesehen werden können. Desgleichen, da die lutherische Kirche die Kirche des Worts ist, das Wort den Mittelpunkt des öffentlichen Gottesdienstes bildet, „alles Gottesdienstes das größte und fürnehmste Stück ist, Gottes Wort predigen“, ja „auch das Sacrament selbst durch Gottes Wort gemacht und geheiliget wird“ (Luther), so muß das Kirchengebäude nicht nur überhaupt akustisch gebaut sein, sondern es müssen auch in demselben Kanzel, Altar und Taufstein so stehen, daß das Wort allüberall gehört und verstanden wird.

Nach diesen einleitenden Bemerkungen fassen wir nun Altar und Kanzel, wie auch den Taufstein in Bezug auf Stellung, Form und Ausstattung etwas näher ins Auge.

\subsection{I. Der Altar}

1. Dieses nach 1 Mos. 8, 20. und 12, 7. ff. älteste Monument des öffentlichen Gottesdienstes, hat von altus und ara den Namen Altar durch die abendländische Kirche erhalten und bedeutet derselbe im Allgemeinen eine erhöhte Weihestätte. Da nun die abendländischen Christen schon in ältester Zeit das Angesicht gen Osten gewendet zu beten pflegten, sich erinnernd, daß von daher uns das Heil gekommen ist, so entwickelte sich daraus um 420 n. Chr. die Sitte, den Altar am östlichen Ende des Kirchengebäudes aufzurichten. Um dieser sinnigen Symbolik willen, welche das alte Wort: ex oriente lux so schön zum Ausdruck bringt, indem die anbetende und lobpreisende Gemeinde, mit ihrem Wortführer voran, sich dem Altar zukehrt, hat, wo immer die örtlichen Verhältnisse es gestatten, auch der lutherische Altar seine Stellung im Osten.

Gewöhnlich befindet sich zwei bis drei Stufen erhöht der Altar in einem besonderen Raum, genannt „Chor“ oder auch „hohe Chor“. Auch zu diesem steigt man vom Schiff der Kirche, dem Raum für die Zuhörer, auf zwei bis drei Stufen wieder hinauf, wie denn überhaupt dafür zu sorgen ist, daß der Altar
\quellseite{307}
nicht zu tief zu stehen kommt. Ebenso darf auch der Chor nicht zu klein sein, damit für Antretende zur Communion und Confirmation, sowie für Assistenten und Zeugen bei der Ordination, Trauung und andern am Altar zu verrichtenden Handlungen der nöthige Raum vorhanden ist.

Bezeichnend richtet man in den katholischen Kirchen zwischen Altarchor und Schiff eine Schranke auf. Vor ihr empfangen außer dem Messe haltenden Priester alle übrigen Communicanten das Abendmahl (unter Einer Gestalt); der Altarchor gehört ausschließlich dem Clerus. Entsprechend der Lehre vom Amte und dem geistlichen Priesterthum aller Christen pflegt man in der lutherischen Kirche keine solche Schranke aufzurichten oder doch die aus andern Gründen angebrachte Schranke zu durchbrechen und in der Mitte einen weiten offenen Zugang vom Schiff in den Altarraum zu lassen.

Da man nun in der lutherischen Kirche nicht vor dem Chor, sondern im Chor das Abendmahl empfängt, so kniet man entweder im Halbkreis um den Altar oder man wandelt um ihn, auf der Nordseite das Brod, auf der Südseite den Wein empfangend. Bei der letzteren, der herrschenden, Weise des Empfangs, muß dann behufs des Umgangs der Altar von der Chorwand weit genug abstehen.

2. Während die reformirte Kirche für die Abhaltung ihres Abendmahls ausschließlich der nackten Tischform sich bedient, bleibt die lutherische Kirche bei der altkirchlichen Form. Es ist der längliche Würfel — mit oder ohne Aufsatz.

\begin{quote}
„Da bei dem Altar die Hauptsache der eigentliche Altarkörper ist, so hat man über die Zulässigkeit eines Altaraufsatzes gestritten. Sollte durch die Rückwand das Chorfenster verbaut und damit zugleich der wirkungsvolle östliche Lichtzugang durch die gemalten Scheiben abgedämmt werden, so ist selbstverständlich davon abzusehen. Sonst sollte dieser schöne Schmuck dem Altar nicht fehlen; er ist mehr als ein bloßes Zierstück, denn er gibt zugleich den Ort ab für symbolische und typologische Darstellungen und bietet eine willkommene Gelegenheit, die mittelalterliche sacramentale Typologie (Abel, Melchisedek, Abraham und Isaak, eherne Schlange u. a.) und Symbolik (Lamm Gottes, Weinstock, Rebe, Kelch u. a.) vor Augen zu stellen. Immer aber soll der Aufbau in organischer Verbindung mit dem Altarkörper stehen und nicht etwa durch ein Bild in Goldrahmen ersetzt werden.“ (Schulze.)
\end{quote}

Wo man nun eines Altaraufsatzes sich bedienen will, so paßt für die Rückwand des Chors anstatt eines Fensters eine Rosette, hoch genug angebracht, um über dem Aufsatz vom Schiff aus gesehen zu werden. Ist nun im Altaraufsatz ein Gemälde eingefügt —
\quellseite{308}
wenn auf dem Altar ein Crucifix steht, paßt am besten die Darstellung der Auferstehung oder der Himmelfahrt, oder auch der Einsetzung des heiligen Abendmahls — so ist um so mehr Sorge zu tragen, daß auf den Altar das nöthige Licht fällt, sei es rechts und links durch ein zwischen dem Altar und Chorbogen angebrachtes schmales Seitenfenster oder durch ein Oberlicht.

Wie schon bemerkt, ruht im erhöhten Chor der Altarkörper selbst wieder auf ein paar Stufen. An manchen Orten ist auf der obersten nördlichen und südlichen Stufe eine Brüstung angebracht, über die hin die Communicanten das Sacrament empfangen, die aber nicht hoch sein darf, damit der Kelch ungehindert dargereicht werden kann. Daß hierbei das Kniebänkchen nicht fehlen darf, braucht wohl kaum erwähnt zu werden. Da es für Trauungen, Confirmationen und Ordinationen einer Kniebank an der Vorderseite der Stufen bedarf, so eignet sich hierzu eine tragbare, deren Länge jedoch zum ungehinderten Besteigen des Altars Raum genug läßt.

3. „Ich wollte“, ruft Scriver in seinem Seelenschatz aus, „daß man die ganze Kirche und sonderlich den Altar wie einen kleinen Himmel könnte ausschmücken, alles dem hochwürdigen Abendmahl zu Ehren und denen gottseeligen Herzen zur Erinnerung.“ Um so mehr aber auch vermeide man beim Ausschmücken des Altars alles Ungehörige, Fremdartige, namentlich Weltförmige. Nachfolgendes wird zur Ausstattung eines angemessen geschmückten und wohlbestellten Altars der lutherischen Kirche im Ganzen erfordert:

a. Die Altarbekleidung, bestehend aus dem Altargewand, dem Altartuch und dem Fußteppich.

Ist das Material, bezw. das Aeußere des Altarkörpers, ein solches, das sich nicht unverhüllt zeigen oder das auch wohl verhüllt werden kann, so ist eine vollständige Bedeckung durch das Altargewand am Platz. Aber auch da, wo diese Rücksicht nicht zu nehmen ist, wo namentlich der Altarkörper monumental ist, soll das Altargewand in irgend einer Form oder Reduktion vorhanden sein. Es kann alsdann die Gewandung aus einem am oberen Rande hinauflaufenden, mit einem passenden Symbol oder einer ornamentalen Stickerei versehenen Antependienstreifen bestehen, der etwa ein Drittel der Altarbreite und die volle Länge der Altarhöhe hat.

Das Altargewand wird am besten ganz glatt um den Altarkörper gespannt; „faltig drapirt, erinnert zu sehr an weltliche Decoration“.

Zur Bezeichnung des Fortschrittes im Wechsel des Kirchenjahrs behielt die ältere Kirche auch mehr oder weniger die liturgischen Farben bei. Dieselben
\quellseite{309}
waren vor Alters: Violett in der Advents- und Fastenzeit von Septuagesimä an; Weiß an Weihnachten bis Epiphanias, an Ostern, sowie an Mariä Reinigung, Verkündigung und Heimsuchung; Grün in der Epiphaniaszeit bis Septuagesimä, sowie in der ganzen Trinitatiszeit; Roth an Pfingsten, Trinitatis, Johannisfest und an den Aposteltagen; Schwarz am Charfreitag, an Bußtagen und bei Trauergottesdiensten. Das Minimum für uns wäre eine schwarze Bekleidung für die Fastenzeit, die Bußtage (auch für die gewöhnlichen Trinitatissonntage) und für Trauergottesdienste, und eine farbige, etwa dunkelgrün, für die übrigen Tage. Kann und will man eine dritte Farbe beifügen, so wähle man zwischen Roth und Violett und vertheile entweder: Dunkelgrün für die Adventszeit und Epiphanias bis Septuagesimä, und Roth für die drei hohen Feste und andere besondere Feste; oder: Violett für die Advents-, Epiphanias- und Trinitatiszeit, und Dunkelgrün für die hohen Feste etc.

Was die Gold- oder Silberbordüren für die Gewandung betrifft, so nehme man nur ächte, statt unechte, da letztere bald schwarz werden. Sind jene nicht zu beschaffen, so thun es einfach wollene. Jedenfalls aber „fort mit allem leichten Posamentirkram oder gar aufgenähten Glasperlen, Schmelz und dergleichen, und statt dessen lieber eine einfache Stickerei in Wollfäden oder Schnur“ etc. (Meurer.)

Das Altartuch bedeckt die Altarplatte, ist stets von schneeweißer Leinwand und kann an seinen überhängenden Kanten einfach gestickt oder mit passenden Spitzen versehen sein.

Theils zur Zierde, theils zur Erinnerung an 2 Mos. 3, 5. werden die Altarstufen, desgleichen auch Podium und Stufen des Chors mit einem Fußteppich belegt.

\begin{quote}
„Alle unruhigen, grellen Muster sind hier strengstens zu vermeiden, insbesondere die sogenannten Blumenteppiche. Figürliche Darstellungen, sowie das Kreuz müssen ebenfalls wegfallen. ‚Es ist unziemlich, das Heilige mit Füßen treten.‘ Man begnüge sich mit Anwendung von symbolischen Thier- und Pflanzenornamenten in gedämpfter Farbe.“ (Sch.)
\end{quote}

Ein unentbehrliches Altarrequisit ist

\subsubsection{b. das Crucifix}

Da das Kreuz das älteste unter allen Denkzeichen ist, die der christliche Glaube zu seiner eigenen Bezeugung aus sich selbst hervorgebracht hat, so ziert mit Recht das Crucifix unsere Altäre, entsprechend der vorbildlichen ehernen Schlange, so lange man in Israel sie hernach allein zum Gedächtnis gebrauchte und keinen Abgott daraus machte. In der That erzieht
\quellseite{310}
auch kein anderes Sinnbild das Zeichen für die vollbrachte Erlösung und Versöhnung der Welt, wie für die Predigt von derselben und ihre Aneignung durch den Glauben. (Joh. 19, 30. 3, 14. 15. 2 Cor. 5.) Um so sorgfältiger sei man in der Wahl, namentlich der Gestalt des Crucifixes. „Der Christuskörper wird entweder aus versilberter Bronze oder aus naturfarbigem Holze hergestellt. Die Silberfarbe hat, weil sie der blassen Körperfarbe mehr entspricht, den Vorzug vor der Goldfarbe. Leicht zerbrechliches oder unmächtiges Material, wie Zink, Eisenguß oder gar Gyps und Porcellan, soll nicht gewählt werden. Das Kreuz selbst kann aus vergoldeter Messingbronze, aus Eichen-, Nußbaum- oder schwarzem Birnbaumholz sein. Ein nacktes Kreuz ohne den Gekreuzigten wird besser vermieden.“ (Sch.)

Was die Größe betrifft, so sollte es stets höher sein als die Altarleuchter, auch hat sich dieselbe nach der Größe der Kirche zu richten.

Entweder steht es frei oder es wird in den Altaraufsatz hineingebaut.

Links und rechts von dem Crucifix, das den eigentlichen Mittelpunkt des Altarmobiliars bildet, haben c. die Altarleuchter ihren Stand. Bei der Feier des heiligen Abendmahls brennend, dienen sie nicht nur zur Erinnerung an die Zeit der Einsetzung desselben und damit an des HEerrn Stiftung selbst, sondern sind zugleich Symbol Christi, als des Lichtes der Welt, das da scheinet im Evangelium und der durch dasselbe einen hellen Schein in unsere Herzen gegeben hat, daß durch uns entstünde die Erleuchtung von der Erkenntniß der Klarheit Gottes in dem Angesichte Jesu Christi (2 Cor. 4, 6.), weshalb man denn auch bei Verlesung des Evangeliums da und dort die Altarleuchter schon brennen läßt. „Die Frage, ob überhaupt Leuchter auf dem Altar zu stehen kommen sollen, ist so selbstverständlich zu bejahen, daß eine Erörterung derselben überflüssig zu sein erscheint. Die schöne Symbolik des Lichtes muß jedem von selbst klar sein.“ (Sch.)

Eine offene Frage ist dagegen die Zahl der Leuchter. Auf manchem Altar befinden sich vier Leuchter entweder von gleicher oder von ungleicher Größe, in letzterem Falle mit dem Crucifix in der Mitte eine Pyramide bildend. Vielfach jedoch gibt man der Zweizahl den Vorzug.

Selbstverständlich müssen Altarleuchter ein würdiges Aussehen haben. Die Form betreffend halte man sich an die guten alten Muster. Als Material eignet sich jedenfalls Metall, keineswegs aber Glas oder Porcellan.

Mancherorts in Deutschland, häufig aber in unserer Zeit hierzulande, stellt man neben oder zwischen Leuchter und Crucifix auch
\quellseite{311}
d. Blumenvasen. Scriver betrachtet sie nach Hohel. 4, 5. als Symbol der geistlichen Erquickung, welche der gläubigen Seele insonderheit beim heiligen Abendmahl von dem himmlischen Bräutigam widerfährt. Wo sie nur als besonderer Schmuck an Festen oder bei festlichen Anlässen gebraucht werden, bedient man sich der natürlichen Blumen, die auch wohl überhaupt am passendsten wären; wo sie ihre bleibende Stellung auf dem Altar haben, gebrauche man gemachte Blumen. In letzterem Fall lasse man jedoch nur etwas Geschmackvolles und nicht allzu Buntes zu. Vor allen Dingen aber wehre man den auch für die Kirche Mode werdenden gläsernen Ueberstürzen (Glasglocken), da sie zu sehr an den Parlor erinnern. Außerhalb des Gottesdienstes mag man sich ihrer ja zum Schutz gegen den Staub bedienen; vor Anfang desselben oder irgend einer Handlung sollten sie vom Kirchendiener abgenommen werden.

So oft das heilige Abendmahl und somit der volle Hauptgottesdienst gehalten wird, bilden ein Haupterforderniß und eine Hauptzierde des Altars.

e. Die heiligen Gefäße (vasa sacra), bestehend aus Patene sammt Hostienbüchse, und Kelch sammt Kannen, nebst Sieb-Löffelchen, alle von gleichem Material und gleichem Styl. In den ersten Jahrhunderten hatten die heiligen Gefäße die antike Form und waren gläsern, thönern, sogar hölzern. Als die Kirche zu äußerlichem Wohlstand kam, verwandte man dazu nur Metall, namentlich Edelmetall.

Die Patene hat die Form eines Tellers, dessen Vertiefung in den oberen Rand des Kelches paßt. Sie dient zur Austheilung des Brodes. Da um der Zahl der Communicanten willen die Patene oft nicht alle Hostien aufnehmen kann, so gehört zu derselben der Hostienbehälter, das Ciborium, von cibus, die Speise. In der katholischen Kirche geschieht vermittelst diesem die Austheilung und hat es deshalb die Form eines niedrigen mit einem gewölbten Deckel versehenen Kelches. In der lutherischen Kirche wird als Ciborium eine runde oder oblonge niedere Büchse gebraucht, deren Deckel häufig ein Lamm oder ein Kreuz schmückt.

Der Kelch, bestehend aus Fuß, Schaft mit Knauf und Becher (cuppa), ist möglichst kunstvoll gestaltet, manchmal auch sogar mit Edelsteinen besetzt. Bei den primitiven Zuständen früherer Jahre daher und dem Mangel an passenden Kelchen vertrat freilich gar oft ein etwas kelchartig geformtes Trinkglas die Stelle des Kelches und war dasselbe immerhin den vorhandenen metallenen schmalen und langcuppigen Bechern der Puritaner vorzuziehen. Im Lauf der Zeit kam es zu wirklichen zinnernen oder silberplattirten und innen vergoldeten Kelchen.

\quellseite{312}
Der lutherische Abendmahlskelch ist auch der Form nach ein anderer als der papistische Meßkelch. Während bei letzterem die Cuppa verhältnismäßig klein ist, da ja nur der Meßhalter trinkt, ist bei ersterem dieselbe größer und sonderlich oben weiter, da gemäß der Einsetzung Christi Alle aus demselben trinken. Es ist daher bei der Cuppa darauf zu achten, daß der Spendende den Inhalt des Kelches übersehen kann, um darnach mehr oder weniger kippen zu können. „Die Renaissance und der Zopf haben schließlich aus dem Kelch einen Pokal gemacht, welcher jedes kirchlichen Charakters entbehrt. Der Rand der Cuppa ist nach außen geschweift, wodurch der Kelch zum Spenden ungeeignet wird, weil der Inhalt eines solchen Bechers bei der geringsten Neigung das Bestreben hat, sehr schnell und breit auszufließen, so daß am Munde des Communicirenden der Wein leicht verschüttet wird.“ (Sch.)

In Folge des Kelchbrauches kennt die katholische Kirche auch keine eigentlichen Abendmahlskannen, sondern bedient sich für den Priester nur zweier Kännchen oder Meßpollen, das eine für den Meßwein mit V (vinum), das für das Wasser, mit welchem er den Kelch nachspülen muß, mit A (aqua) bezeichnet, beide auf einem Teller stehend. Somit fehlen uns gänzlich Vorbilder. Uebrigens sind die Abendmahlskannen, wie sie jetzt hier mit den besseren Kelchen zu haben sind, auch der Form nach ziemlich zweckentsprechend. Häufig ziert ein Kreuz den Deckel. Passend ist auch der Pelikan, der mit seinem Herzblute seine Jungen tränkt. Da nur Gold im Stande ist, der Säure des Weines Widerstand zu leisten und Grünspan-Ansetzung zu verhindern, so ist eine Vergoldung des Innern auch der Kannen sehr wünschenswerth. Wo dieselbe nicht stattfindet, ist sofortige Ausspülung und Austrocknung der Kanne nach dem Gebrauch nur um so nothwendiger.

Als Zuthat zu den heiligen Gefäßen sei das Corporale, die Palle und das Velum erwähnt. Das Corporale oder Leibtuch ist ein Leintuch, auf welches die heiligen Gefäße gestellt werden, in der lutherischen Kirche wohl allein zu dem Endzweck, daß das weiße Altartuch ganz rein bleibt. Die Palle ist ein mit dem Stoff des Altargewandes oder mit einem andern, doch gleichfarbigen, Stoff überzogenes quadratisches Pappstück, mit welchem während der Austheilung des Abendmahls, zumal wenn dieselbe durch Eine Person geschieht, Patene und Kelch bedeckt werden, jene, daß ein etwaiger Luftzug keine Hostie herabweht, dieser, daß während der Austheilung des Brods nichts in den Kelch fällt. Das Velum ist ein weißes linnenes oder seidenes Tuch, mit welchem vor und nach der Austheilung die auf der Südseite des Altars zusammengestellten
\quellseite{313}
Gefäße vielfach überdeckt zu werden pflegen. Das Sieb-Löffelchen, das gleichen Stoffes mit den Gefäßen ist, dient, in Kelch oder Kannen Gefallenes bequem herauszunehmen zu können.

Zum Altar gehörig rechnen wir auch

\subsubsection[e. das Pult]{e. das Pult, das aber ein doppeltes sein sollte — eines zum Beten, das andere zum Lesen}

Das erstere ist das eigentliche Altarpult. Weil es passender erscheint, daß beim Singen oder Lesen der Collecte, der Präfation, des Vaterunsers und der Einsetzungsworte die Agende vor dem Liturgen auf dem Altar liegt, statt daß er sie in der Hand hält — weshalb auch von jeher für die Agende möglichst große Lettern gebraucht werden, gleichwie in der römischen Kirche für die Meßbücher etc. —; weil aber derselbe eine zu gebückte, namentlich für den Gesangsvortrag nachtheilige Stellung einnehmen müßte, läge die Agende unmittelbar auf der horizontalen Altarfläche, so pflegt man sich vielfach eines auf dem Altar liegenden Pultes zu bedienen.

Das andere Pult, das Lesepult, ist für die kirchlichen Lectionen bestimmt, und zwar zur Verlesung derselben nicht nur in den Nebengottesdiensten, sondern auch im Hauptgottesdienst, indem es nach Vorgang der alten Kirche eigentlich kirchlicher Brauch ist, Epistel und Evangelium statt unmittelbar vom Altar aus an einem von demselben abseits stehenden Pult zu verlesen.\footnote{In einer Beschreibung der Kirche zu Alt-Treßden aus dem Jubeljahre 1717 wird ausdrücklich erwähnt, daß am Eingang zum Altarchor „ein messingnes vergoldetes pulpet“ zur Verlesung der Episteln und Evangelien etc. sich befinde, und auf dem Titelbild zur Leipziger Agende 1681, welches das Innere einer Kirche und zugleich die Abhaltung des Hauptgottesdienstes mit Predigt und Communion darstellt, ist auch das am Eingang zum Altarchor auf einigen Stufen sich befindliche pulpet zu sehen und vor demselben der als Liturg fungirende und mit der Alba bekleidete Diakonus.} Da dieser allgemeine Brauch früherer Zeit sich noch in manchen Kirchen Mitteldeutschlands erhalten hat, so dürfte er auch unter uns Aufnahme finden, zumal die Altarbibel, deren der Liturg auch für das Vorlesen der Perikopen sich da bedienen muß, wo man kein besonderes Perikopenbuch besitzt, nicht selten ziemlich schwer ist, auch mancher Liturg durch Anlehnen an den Altar eine dem Decorum nicht entsprechende Stellung beim Lesen annimmt. Wohl scheint durch das Schreiten des Liturgen vom Altar zum Lesepult und, wo die Intonation zum Credo noch üblich ist, zurück an den Altar eine störende Pause zu entstehen, allein die existirt wenigstens nicht für die Gemeinde, da sie bei der Epistel durch das Singen des Amen zur Collecte, und beim Evangelium durch das zur Intonation des Credo oder zu diesem selbst überleitende Orgelspiel ausgefüllt wird.

\quellseite{314}
„Das erste Erforderniß für ein solches Lesepult ist, daß es absolut sicher und fest steht, auch nicht zu hoch und zu niedrig sei. Die früheren Pulte wurden aus Holz, später vielfach auch aus Metall hergestellt. Häufig diente ein Adler, das Symbol des Evangelisten St. Johannes, als Träger des Pultbrettes, manchmal erscheinen auch die vier Evangelisten als Stützen des ganzen Pultes; so an einem sehr alten holzgeschnitzten Pulte in der Kirche zu Freudenstadt in Württemberg.“ (Sch.) Auch in einer Episcopalkirche sah der Verfasser einen vergoldeten Adler angebracht, auf dessen Rücken und ausgebreiteten Flügeln das Buch ruhte.

Wie schon bemerkt, befand sich das Lesepult vor dem Altarchor und auf einer besonderen Erhöhung. Nach den Thesen der Eisenacher Conferenz 1861 gehört der „Lehrstuhl“ oder das Lesepult „entweder vor den Altar auf eine der Stufen, die aus dem Schiff zum Chor emporführen, doch so, daß der Blick der Gemeinde nach dem Altar nicht verhindert werde, oder an einen Pfeiler des Chorbogens, um für den Zweck der Katechese, Bibelstunde etc. vor den Altar hingerückt zu werden“. Steht daher die Kanzel am Eingang des Chors, so nimmt ihr gegenüber das Lesepult Stellung.

Durch die Theilung von Altar- und Lesepult und den Gebrauch des letzteren zu allen kirchlichen Lesungen erhalten Agende und die sogenannte Altarbibel ihren gesonderten und geeigneten Ort: erstere liegt auf dem Altarpult, letztere auf dem Lesepult.

Schön gearbeitete Pulte bedürfen keiner besonderen Bekleidung. Bedient man sich für Altar- und Lesepult einer besonderen Decke, so müssen Stoff und Farbe die des Altargewandes sein.

Anmerkung. Durch das Vorhandensein eines Lesepultes wird bei Lesegottesdiensten „die ungehörige Sitte vermieden, daß der Küster von der Kanzel aus oder gar vom Altar aus vorliest“. (Sch.)

\subsection{II. Die Kanzel}


1. Wie ursprünglich im Alten Testament, so war auch ursprünglich im christlichen Alterthum der Altar die Stätte der Verkündigung des Worts. Hinter demselben saß nämlich der Bischof oder der Presbyter, der Gemeinde zugewendet, und redete von da aus zur Gemeinde, wobei er jedoch saß, während die Zuhörer standen. Außerdem bediente man sich hierzu des Ambons, des vor dem Altar und abseits desselben für den Lector und die Sänger angesichts der Gemeinde erhöhten Ortes; zugleich geschah die Predigt an dem Gitter oder
\quellseite{315}
den Schranken, welche den Altarraum von dem Schiff trennten und an welchen in größeren Kirchen eine besondere Erhöhung zum Predigen angebracht war, während in kleineren Kirchen der Ambon dem Gitter selbst nahe genug sich befand. Diese Gitter, diese Schranken hießen cancelli. Von ihnen hat dann der Predigtstuhl den bleibenden Namen Kanzel erhalten. War dieser nun schon hierdurch vom äußersten Punkte der Kirche der Mitte entgegengerückt, so trat er vollends in diese ein, als man zu Gunsten der Mönchshorden, die sich vom 13. Jahrhundert an der Predigt bemächtigten, die Kanzel in größeren, mehrschiffigen Kirchen an einen der Pfeiler, in kleineren an eine der Langseiten setzte.

Ueber die richtige Stellung der Kanzel, ob getrennt, ob verbunden mit dem Altar, hat man neuerdings viel gestritten. Die Kanzel in den Altar einschieben oder gar über dem Altar aufbauen, ist nach Schultze „unstatthaft“, nach den Bestimmungen der Dresdener Conferenz 1856 „völlig falsch und geradezu widersinnig“. Ebenso heißt es These 10 der Eisenacher Conferenz 1861: „Die Kanzel darf weder vor, noch hinter, noch über dem Altar, noch überhaupt im Chore stehen. Ihre richtige Stellung ist da, wo Chor und Schiff zusammenstoßen, an einem Pfeiler des Chorbogens nach außen (dem Schiffe zu); in mehrschiffigen großen Kirchen an einem der östlicheren Pfeiler des Mittelschiffes. Die Höhe der Kanzel hängt wesentlich von derjenigen der Emporen ab und ist überhaupt möglichst gering anzunehmen, um den Prediger auf und unter den Emporen sichtbar zu machen.“ Andere wieder behaupten das Gegenteil. So Hasse in seinem Büchlein: „Die Zeichensprache der evang.-luth. Kirche“ etc. 1877. Nachdem derselbe auf gewisse Uebelstände hingewiesen hat, welche für einen Theil der Zuhörer entstehen, wenn wegen der Größe der Kirche die Kanzel mehr in die Mitte derselben gerückt werden muß, so fährt er fort: „In nicht allzu großen Gotteshäusern lassen sich alle diese Uebelstände verhüten, wenn Altar und Kanzel in das Verhältnis von Wort und Sacrament zu einander gebracht, also zusammengefügt werden, indem die Kanzel im Altaraufsatz so eingefügt wird, daß in letzterem unter ihr hinlänglicher Raum zu einem Altarbilde bleibt und der Kanzelfuß nicht unmittelbar auf dem Altartische ruht. Dadurch erhalten alle Plätze der Kirchengänger gleichen Werth, und die gesammte Andacht der Gemeinde kommt, unter der selbstverständlichen Voraussetzung, daß wenigstens sämmtliches Mittelgestühl dem Altar zugekehrt ist, in unverändert gleiche Richtung. Diese Kanzelstellung entspricht auch der ursprünglichen Stellung des bischöflichen Lehrstuhls hinter dem Altartisch und der späteren auf dem Vorderrande des Altarplatzes“ (S. 81. 82). Gleichwohl
\quellseite{316}
könnten wir uns für eine Zusammenfügung von Kanzel und Altar nicht entscheiden. Nicht nur kommt bei derselben die Kanzel doch ziemlich hoch zu stehen, was für die ihr zunächst sitzenden Zuhörer das Hinaufsehen ziemlich erschwert, und nicht nur fällt alsdann der Altarchor weg, sondern es verliert der Altar selbst nicht wenig an seinem würdigen Aussehen. Kanzel — Altar — Taufstein, jedes für sich und an seinem Ort, so ist’s doch wohl am schönsten und dabei auch mehr an 1 Joh. 5, 8. erinnernd: „Drei sind, die da zeugen auf Erden: Der Geist und das Wasser und das Blut, und die drei sind beisammen.“ So die drei Stätten für diese drei Zeugen sondernd, und nicht zwei derselben zusammenfügend, sondern sie zusammenstellend vermittelst eines Altarchors wird in unsern hiesigen Kirchen die Kanzel ihre angemessenste Stellung doch wohl da haben, wo Chor und Schiff zusammenstoßen. Man wird jedoch, zumal wenn die Kirche nicht in Kreuzform erbaut ist, sehr darauf achten müssen, daß da, wo Empore angebracht werden, auch von der Empore, die sich auf der Kanzelseite befindet, der Prediger überall gesehen werden kann.

2. Wohl anspielend auf Ps. 116, 13. ist die gewöhnliche Form der Kanzel die eines Kelches, nur daß da und dort der Fuß fehlt und bloß die Cuppa nachgeahmt ist. Zu keinem Fall darf jedoch die letztere zu umfangreich sein. Mancherorts ist die Kanzeltreppe in der Sacristei angebracht und tritt man alsdann durch eine Thüre auf die Kanzel; mancherorts aber auch im Chor. Hat die Kanzel ihre Stellung an einer Säule, so windet sich um diese die Treppe mehr oder weniger. Bei Anbringung eines Kanzelpultes ist, wenn die Kanzel seitlich vom Altar steht, für die Stellung des Predigers die Diagonale zu beachten, damit derselbe den größten Theil der Zuhörer möglichst in Einen Blick fassen kann. Zum Hinlegen der Bücher ist ein Seitenbrettchen nicht zu vergessen, und wo die Kanzelbeichte üblich ist, nicht ein Kniebänkchen.

Auch da, wo die Kirche nicht so groß ist, daß sie einen Schalldeckel erheischt, pflegt man doch gern einen solchen anzubringen. Aber wie manche unserer hiesigen Kanzeln hat einen Schalldeckel, der sich wegen seiner Kleinheit wie der Deckel einer Zuckerdose ausnimmt und dabei noch ziemlich hoch über der Kanzel schwebt. Soll es in der That ein Schalldeckel sein, so muß sein Durchmesser etwas größer sein, als der des oberen Kanzelrandes, auch darf er nicht zu hoch gehängt werden.

\begin{quote}
3. „Wie auf die rechte Gestaltung des Altars, so ist auch auf die würdige Bildung und Ausschmückung der Kanzel, als der Stätte der Heilsverkündigung,
\quellseite{317}
besondere Aufmerksamkeit zu wenden. Die Kanzel soll weder ein weltliches Prachtstück, noch ein unscheinbarer Holzbau sein, sondern wie Wort und Sacramente neben einander stehen, so soll ihr die gleiche Sorgfalt wie dem Altar werden und ihr so wenig wie diesem heiliger Schmuck architektonischer oder figürlicher Art fehlen. Doch darf vor der künstlerischen Seite das praktische Bedürfniß nicht zurückgestellt werden.“ (Sch.)
\end{quote}

Folgendes ist die Beschreibung der Kanzel der Kirche in Alt-Dresden aus dem Jahre 1717: „Nicht weniger Ansehen gibt der Kirche die an dem andern Pfeiler von dem Altar her linker Hand aufgerichtete Kanzel, welche von weiß und vergoldeter Bildhauer-Arbeit ist, und von sechs schwebenden Engeln gleichsam getragen wird, die Decke aber eine geschlossene Krone vorstellt, worauf der Glaube abgebildet zu sehen, so in der einen Hand ein Crucifix, in der andern aber ein Buch hält, auf dessen Blätter mit güldenen Buchstaben geschrieben: ‚Fides ex auditu, auditus autem per verbum Dei, Rom. X, 17.‘, das ist: ‚Der Glaube kommt aus der Predigt, das Predigen aber durch das Wort Gottes.‘ Dabei denn als etwas Ungewöhnliches angemerket wird, daß man zu selbiger zwar nur eine Thüre, aber zwei unterschiedene und den Pfeiler in die Mitte fallende Treppen hat, welches dem Bau nicht nur eine feine Gestalt, sondern auch denen, so dieselbige betreten, eine gute Erinnerung gibt, daß, wie sie auf den Stufen der reinen Lehre hinauf, also auf den Stufen des heiligen Wandels wiederum heruntersteigen, und dasjenige selbsten durch einen christlichen Wandel beweisen sollen, was sie andere zu thun gelehret haben.“

Als figürliche Darstellung für die Kanzel eignen sich am passendsten Orte sonst noch z. B. der HEerr in sitzender Stellung („Predigtstuhl“), in der Linken ein aufgeschlagenes Buch haltend und die Rechte segnend erhebend; ferner die Evangelisten, entweder in wirklicher Darstellung oder in ihren Symbolen; ebenso Propheten und Apostel, doch jene nie ohne diese. Zur Erinnerung an das Amt des Neuen Testamentes aber, dem „Amt des Geistes, das den Geist gibt“, pflegt man im Schalldeckel das Taubenbild anzubringen, und zwar so, daß es über dem Prediger schwebend erscheint.

\subsection{III. Der Taufstein}


1. Nachdem im ersten Jahrhundert die Taufe an jedem beliebigen Orte stattgefunden hatte, erbaute man seit Anfang des 4. Jahrhunderts besondere, mit der Kirche etwa nur durch einen bedeckten Gang verbundene Baptisterien oder Taufkapellen. Man wollte durch diese örtliche Trennung symbolisch
\quellseite{318}
andeuten, daß die Katechumenen erst durch die Taufe die Aufnahme in den Gottesdienst der Christen und insbesondere in die Abendmahlsgemeinschaft empfingen. Weil aber damals der Bischof die Taufe der Katechumenen vorzunehmen pflegte, so besaßen nur die Kathedralkirchen solche Baptisterien, und weil an volkreichen Orten die Zahl der Taufcandidaten oft ziemlich groß war, so hatten diese Taufkapellen mitunter einen ziemlichen Umfang. In der Mitte des Baptisteriums befand sich der Taufbrunnen, ein manchmal überaus prächtiges Bassin, in das drei Stufen hinabführten, und um dasselbige herum ein großer Raum für die vielen Täuflinge und Zuschauer, daher auch das Baptisterium die Form der Rotunda, hernach des Achtecks hatte. Seit aber die Kindertaufe Regel, die Taufe der Erwachsenen Ausnahme wurde, trat an die Stelle des alten Baptisteriums der Taufstein und hatte derselbe meist links beim Eingang seine Stelle, und zwar in einem besonderen Gemach der Vorhalle.

Da die Taufe unserer Kinder in der Regel im Gotteshause und dabei Angesichts der Gemeinde vollzogen werden sollte, so nimmt ohne Zweifel am passendsten der Taufstein in oder unmittelbar vor dem Altarchor Stellung, in letzterem Falle etwas erhöht. Erinnert er doch durch diese Stellung auch, daß die Taufe „der Eingang zu allen göttlichen Gütern und aller Heiligen Gemeinschaft ist“ (Luther), und wiederum, daß jede würdige Theilnahme am Sacrament des Altars thatsächliche Erneuerung des Taufbundes ist und Erhaltung in der Taufgnade bezweckt.

2. Aus Stein, Metall oder Holz gearbeitet hat der Taufstein mehr oder weniger die Form eines niedrigen romanischen oder gothischen Kelchs, ist manchmal rund, meist aber achteckig, mit Kanzel und Altar stylgleich. Häufig hat derselbe einen kuppelförmigen oder auch thurmartigen Deckel. In einer Kirche Norddeutschlands hält ein schwebender Engel den Deckel in der Hand, der mit demselben bei dem Gebrauch durch ein Gegengewicht sich in die Höhe hinaufzieht. (In einer andern Kirche vertritt die Stelle des Taufsteins ein der Gemeinde zugekehrter Engel in knieender Stellung, in dessen Händen ein großes Taufbecken ruht.)

3. Da die lutherische Kirche kein geweihtes Taufwasser kennt, auch nach der Handlung stets das Wasser ausgießt, so pflegt man die Vertiefung des Taufsteins mit einem abnehmbaren Taufbecken zu versehen. Zu demselben gehört dann auch die Taufkanne, um das Taufwasser nach dem Taufstein zu bringen und es durch den Küster in das Taufbecken vor Beginn der Handlung oder, wie es mancherorts Brauch ist, beim Actus selbst, dreimal in die
\quellseite{319}
hohle Hand des Täufers gießen zu lassen. Damit der letztere die Hand sich wieder trocknen kann, gehört zum Taufstein ein Tüchlein, das der Küster bereit hält. Ein Requisit für den Taufstein war häufig früher auch das „Westerhemdlein“ und ist es jetzt noch, wiewohl selten, hie und da. Dasselbe ist ein weißer ärmelloser Ueberwurf, der nach der Taufe dem Kinde angezogen oder auch nur auf dasselbe gebreitet wird und wobei der Täufer unter Handauflegung (da und dort auch der Pathen) das votum baptismale („der allmächtige Gott und Vater unsers HEerrn Jesu Christi, der dich durchs Wasser“ etc.) spricht. Was aber die Taufsteinbekleidung betrifft, so ist dieselbe je überflüssiger, je kunstvoller der Taufstein gearbeitet ist.

Anmerkung. Als die Privatbeichte in der lutherischen Kirche noch die herrschende Weise des Beichtens war, hatte man auch noch einen besonderen Beichtstuhl und wies demselben am liebsten seine Stelle im Altarchor an, so daß der Beichtende und der Beichtiger konnten von jedermann zwar gesehen, von niemandem aber gehört werden. Daß nun ein Beichtstuhl für die lutherische Privatbeichte anders gestaltet war, als der für die katholische Ohrenbeichte, zeigen u. a. auch noch vorhandene Abbildungen, von denen zwei dem Verfasser vorliegen, eines im Kubach’schen Gebetbuch von 1688, das andere im Communionbuch des Nürnbergischen Predigers J. C. Marperger von 1776. Beide haben keine Beine, sondern massive Seitenlehnen und ruhen auf einer Stufe, auf der der Beichtende kniete; bei ersterem steigt die Rücklehne noch in die Höhe und ruht auf derselben ein kleiner einfacher Baldachin. Daneben erblickt man in einer Nische ein Bild, das den weinenden Petrus und den krähenden Hahn darstellt. (Wie es scheint, nahm denselben auch beim Hauptgottesdienste der als Liturg fungirende Kirchendiener ein, während der, welcher die Predigt zu halten hatte, bis zu seinem Auftritt in der Sacristei sich aufhielt. Nie aber legte der Liturg das Amtskleid ab, wenn er außerhalb der Sacristei der Predigt zuhören wollte, er mochte nun seinen Sitz im Altarchor oder vorn im Schiff haben.)

\quellseite{320}
\quellseite{321}
\quellseite{322}
\quellseite{323}
\quellseite{324}
\quellseite{325}
\quellseite{326}
\quellseite{327}
\quellseite{328}
\quellseite{329}
